% Vietnamese translation for OI Public Library

% Source-PDF: 模板 TEMPLATES/Macac_rio_Maratona_de_Programa__o.pdf

\documentclass[11pt]{article}

\usepackage{oipl-vn}

\usepackage{array}

\usepackage{booktabs}

\usepackage{longtable}

\usepackage{fvextra}

\DefineVerbatimEnvironment{vncode}{Verbatim}{fontsize=\scriptsize,breaklines=true,breakanywhere=true,tabsize=4}

\newcommand{\codepath}[1]{\par\noindent\textit{Mã nguồn mẫu:} \texttt{#1}\par}

\title{Macacário: Maratona lập trình v3.3}

\author{Viện Công nghệ Hàng không\\Lucas França de Oliveira (Mexicano T-18)\\\texttt{lucas.fra.oli18@gmail.com}}

\date{9 tháng 9 năm 2019}

\begin{document}

\maketitle

\origtitle{Macacário Maratona de Programação v3.3}

\tableofcontents

\clearpage

\section{Giới thiệu}

\subsection{Các lỗi thiên niên kỷ}

\paragraph{Lỗi lý thuyết}\begin{itemize}

\item Không đọc đề bài một cách bình tĩnh.

\item Vội vàng giả định một tính chất nào đó của lời giải.

\item Không đọc lại giới hạn trước khi nộp.

\item Khi sửa một thuật toán, chú ý mọi chi tiết của cấu trúc thuật toán, phần nào cần hoặc không cần thay đổi.

\item Bài có thể là NP, được ngụy trang hoặc không nêu rõ giới hạn; lúc đó lời giải sẽ khó thật sự, không phải lúc cố đoạt Nobel.

\end{itemize}

\paragraph{Lỗi về giá trị cực đại của biến}\begin{itemize}

\item Tính cẩn thận để xem vô cực có thật sự đủ vô cực không.

\item Kiểm tra các phép toán với vô cực có tràn 31 bit không.

\item Cẩn thận phép nhân int tràn 32 bit, ví dụ kiểm tra dấu bằng a*b > 0.

\end{itemize}

\paragraph{Lỗi trường hợp biên}\begin{itemize}

\item Đã thử n=0, n=1, n=MAXN chưa? Nhiều bài cần xử lý riêng.

\item Nghĩ qua mọi trường hợp có thể là biên hoặc cô lập.

\item Trường hợp biên có thể phá cả việc dựng cấu trúc, ví dụ danh sách kề.

\item Không quên self-loop hoặc đa cạnh trong đồ thị.

\item Bài đường đi Euler cần kiểm tra đồ thị liên thông.

\end{itemize}

\paragraph{Lỗi lơ đãng khi cài đặt}\begin{itemize}

\item Copy/paste sai mã, rất thường gặp.

\item Đặt phép gán trong if, ví dụ if (a = 0).

\item Quên khởi tạo biến.

\item Nhầm break và continue.

\item Không khởi tạo lại mỗi test: mảng đã xóa chưa, biến cộng dồn đã đưa về đúng giá trị chưa?

\item Output đã đúng định dạng chưa?

\item Mảng khai báo đủ lớn chưa?

\item Cẩn thận modulo số âm: dùng (x\%n+n)\%n.

\end{itemize}

\paragraph{Lỗi cài đặt khác}\begin{itemize}

\item Khai báo biến global và local cùng tên.

\item Khai báo sai kiểu biến, ví dụ int thay vì double hoặc char.

\item Không dùng tên biến max và min.

\item Nhớ rằng size() là unsigned.

\item Nhớ 1 là int: long long a = 1 << 40 không đúng, nên dùng 1LL << 40.

\end{itemize}

\paragraph{Lỗi giới hạn}\begin{itemize}

\item Độ phức tạp thời gian và bộ nhớ là bao nhiêu? 10\textasciicircum{}8 là mốc tham khảo cho thời gian; cũng cần kiểm nhanh bộ nhớ.

\item Có thể giảm hằng số bằng thuật toán tốt hơn hoặc phép toán nhanh hơn không?

\item Bài có phải trường hợp riêng cần tối ưu thay vì dùng thẳng thư viện không?

\end{itemize}

\paragraph{Lỗi với số thực}\begin{itemize}

\item Tránh float/double nếu không cần hoặc không đơn giản hơn.

\item Luôn dùng double, không dùng float, trừ khi đề yêu cầu rõ.

\item So sánh bằng với sai số, có thể cả tuyệt đối và tương đối.

\item Tránh trừ hai số gần bằng nhau.

\end{itemize}

\paragraph{Lỗi khác}\begin{itemize}

\item Tránh cấp phát động nếu không cần.

\item Không dùng STL khi không cần thiết, nhưng dùng khi nó làm lời giải rõ hoặc cần để đạt độ phức tạp.

\end{itemize}

\subsection{1010 điều răn}

Theo PUC-RJ.

\begin{enumerate}[start=0]

\item Không chia cho zero.

\item Không cấp phát động.

\item So sánh số thực bằng EPS.

\item Kiểm tra đồ thị có thể không liên thông.

\item Kiểm tra cạnh có thể có trọng số âm.

\item Kiểm tra có thể có nhiều cạnh nối cùng hai đỉnh.

\item Kiểm tra mọi chỉ số của quy hoạch động.

\item Giảm branching factor của DFS.

\item Cắt tỉa mọi thứ có thể trong DFS.

\item Cẩn thận điểm trùng và điểm thẳng hàng.

\end{enumerate}

\subsection{Mẹo bẩn nhưng hợp lệ}

\begin{itemize}

\item Phương pháp Steve Halim: nếu toàn bộ output khả dĩ vừa trong mã nguồn, brute force trên máy vài phút rồi ghi thẳng đáp án vào code để nộp. Kiểm tra giới hạn kích thước submission.

\item Giai thừa tới 10\textasciicircum{}9: chạy chương trình trên máy để brute force giai thừa, in mỗi 10\textasciicircum{}3 hoặc 10\textasciicircum{}6 bước, dán output vào code và dùng giá trị tiền xử lý.

\item Bài có hằng số: nếu một giá trị hữu ích không phụ thuộc input nhưng bạn không biết, brute force trên máy rồi dán vào code.

\item Debug bằng assert: có thể đặt assert trong code nộp để biến WA thành RTE. Chỉ dùng khi thật sự bí vì tốn submission.

\end{itemize}

\subsection{Giới hạn biểu diễn dữ liệu}

\begin{center}\begin{tabular}{llrrr}\toprule Kiểu & scanf & bit & nhỏ nhất & lớn nhất \\ \midrule
char & \%c & 8 & 0 & 255 \\
signed char & \%hhd & 8 & -128 & 127 \\
unsigned char & \%hhu & 8 & 0 & 255 \\
short & \%hd & 16 & -32.768 & 32.767 \\
unsigned short & \%hu & 16 & 0 & 65.535 \\
int & \%d & 32 & -2 x $10^9$ & 2 x $10^9$ \\
unsigned int & \%u & 32 & 0 & 4 x $10^9$ \\
long long & \%lld & 64 & -9 x $10^{18}$ & 9 x $10^{18}$ \\
unsigned long long & \%llu & 64 & 0 & 18 x $10^{18}$ \\
\bottomrule\end{tabular}\end{center}

\begin{center}\begin{tabular}{llrr}\toprule Kiểu & scanf & số mũ & chữ số thập phân chính xác \\ \midrule
float & \%f & 38 & 6 \\
double & \%lf & 308 & 15 \\
long double & \%Lf & 19.728 & 18 \\
\bottomrule\end{tabular}\end{center}

\subsection{\texorpdfstring{Số lượng số nguyên tố từ 1 đến $10^n$}{Số lượng số nguyên tố từ 1 đến lũy thừa 10 mũ n}}

Luôn đúng rằng $n/\ln(n) < \pi(n) < 1.26 n/\ln(n)$.

\begingroup\scriptsize
\begin{vncode}
pi(10^1) = 4              pi(10^2) = 25             pi(10^3) = 168
pi(10^4) = 1.229          pi(10^5) = 9.592          pi(10^6) = 78.498
pi(10^7) = 664.579        pi(10^8) = 5.761.455      pi(10^9) = 50.847.534
\end{vncode}
\endgroup


\subsection{Tam giác Pascal}

\begingroup\scriptsize
\begin{vncode}
n\p   0   1   2    3    4    5    6    7   8   9  10
0     1
1     1   1
2     1   2   1
3     1   3   3    1
4     1   4   6    4    1
5     1   5  10   10    5    1
6     1   6  15   20   15    6    1
7     1   7  21   35   35   21    7    1
8     1   8  28   56   70   56   28    8   1
9     1   9  36   84  126  126   84   36   9   1
10    1  10  45  120  210  252  210  120  45  10   1

C(33,16)       1.166.803.110             giới hạn int
C(34,17)       2.333.606.220             giới hạn unsigned int
C(66,33)       7.219.428.434.016.265.740 giới hạn long long
C(67,33)      14.226.520.737.620.288.370 giới hạn unsigned long long
\end{vncode}
\endgroup


\subsection{Giai thừa}

Giai thừa tới 20 cùng các mốc giới hạn kiểu.

\begingroup\scriptsize
\begin{vncode}
 0!                         1
 1!                         1
 2!                         2
 3!                         6
 4!                        24
 5!                       120
 6!                       720
 7!                     5.040
 8!                    40.320
 9!                   362.880
10!                 3.628.800
11!                39.916.800
12!               479.001.600  giới hạn unsigned int
13!             6.227.020.800
14!            87.178.291.200
15!         1.307.674.368.000
16!        20.922.789.888.000
17!       355.687.428.096.000
18!     6.402.373.705.728.000
19!   121.645.100.408.832.000
20! 2.432.902.008.176.640.000  giới hạn unsigned long long
\end{vncode}
\endgroup


\subsection{Bảng ASCII}

\begingroup\scriptsize
\begin{vncode}
  0 0x00 NUL      1 0x01 SOH      2 0x02 STX      3 0x03 ETX
  4 0x04 EOT      5 0x05 ENQ      6 0x06 ACK      7 0x07 BEL
  8 0x08  BS      9 0x09 TAB     10 0x0A  LF     11 0x0B  VT
 12 0x0C  FF     13 0x0D  CR     14 0x0E  SO     15 0x0F  SI
 16 0x10 DLE     17 0x11 DC1     18 0x12 DC2     19 0x13 DC3
 20 0x14 DC4     21 0x15 NAK     22 0x16 SYN     23 0x17 ETB
 24 0x18 CAN     25 0x19  EM     26 0x1A SUB     27 0x1B ESC
 28 0x1C  FS     29 0x1D  GS     30 0x1E  RS     31 0x1F  US
 32 0x20         33 0x21   !     34 0x22   "     35 0x23   #
 36 0x24   $     37 0x25   %     38 0x26   &     39 0x27   '
 40 0x28   (     41 0x29   )     42 0x2A   *     43 0x2B   +
 44 0x2C   ,     45 0x2D   -     46 0x2E   .     47 0x2F   /
 48 0x30   0     49 0x31   1     50 0x32   2     51 0x33   3
 52 0x34   4     53 0x35   5     54 0x36   6     55 0x37   7
 56 0x38   8     57 0x39   9     58 0x3A   :     59 0x3B   ;
 60 0x3C   <     61 0x3D   =     62 0x3E   >     63 0x3F   ?
 64 0x40   @     65 0x41   A     66 0x42   B     67 0x43   C
 68 0x44   D     69 0x45   E     70 0x46   F     71 0x47   G
 72 0x48   H     73 0x49   I     74 0x4A   J     75 0x4B   K
 76 0x4C   L     77 0x4D   M     78 0x4E   N     79 0x4F   O
 80 0x50   P     81 0x51   Q     82 0x52   R     83 0x53   S
 84 0x54   T     85 0x55   U     86 0x56   V     87 0x57   W
 88 0x58   X     89 0x59   Y     90 0x5A   Z     91 0x5B   [
 92 0x5C   \     93 0x5D   ]     94 0x5E   ^     95 0x5F   _
 96 0x60   `     97 0x61   a     98 0x62   b     99 0x63   c
100 0x64   d    101 0x65   e    102 0x66   f    103 0x67   g
104 0x68   h    105 0x69   i    106 0x6A   j    107 0x6B   k
108 0x6C   l    109 0x6D   m    110 0x6E   n    111 0x6F   o
112 0x70   p    113 0x71   q    114 0x72   r    115 0x73   s
116 0x74   t    117 0x75   u    118 0x76   v    119 0x77   w
120 0x78   x    121 0x79   y    122 0x7A   z    123 0x7B   {
124 0x7C   |    125 0x7D   }    126 0x7E   ~    127 0x7F DEL
\end{vncode}
\endgroup


\subsection{Số nguyên tố tới 10.000}

Có 1229 số nguyên tố tới 10.000.

\begingroup\scriptsize
\begin{vncode}
    2     3     5     7    11    13    17    19    23    29    31    37    41    43    47
   53    59    61    67    71    73    79    83    89    97   101   103   107   109   113
  127   131   137   139   149   151   157   163   167   173   179   181   191   193   197
  199   211   223   227   229   233   239   241   251   257   263   269   271   277   281
  283   293   307   311   313   317   331   337   347   349   353   359   367   373   379
  383   389   397   401   409   419   421   431   433   439   443   449   457   461   463
  467   479   487   491   499   503   509   521   523   541   547   557   563   569   571
  577   587   593   599   601   607   613   617   619   631   641   643   647   653   659
  661   673   677   683   691   701   709   719   727   733   739   743   751   757   761
  769   773   787   797   809   811   821   823   827   829   839   853   857   859   863
  877   881   883   887   907   911   919   929   937   941   947   953   967   971   977
  983   991   997  1009  1013  1019  1021  1031  1033  1039  1049  1051  1061  1063  1069
 1087  1091  1093  1097  1103  1109  1117  1123  1129  1151  1153  1163  1171  1181  1187
 1193  1201  1213  1217  1223  1229  1231  1237  1249  1259  1277  1279  1283  1289  1291
 1297  1301  1303  1307  1319  1321  1327  1361  1367  1373  1381  1399  1409  1423  1427
 1429  1433  1439  1447  1451  1453  1459  1471  1481  1483  1487  1489  1493  1499  1511
 1523  1531  1543  1549  1553  1559  1567  1571  1579  1583  1597  1601  1607  1609  1613
 1619  1621  1627  1637  1657  1663  1667  1669  1693  1697  1699  1709  1721  1723  1733
 1741  1747  1753  1759  1777  1783  1787  1789  1801  1811  1823  1831  1847  1861  1867
 1871  1873  1877  1879  1889  1901  1907  1913  1931  1933  1949  1951  1973  1979  1987
 1993  1997  1999  2003  2011  2017  2027  2029  2039  2053  2063  2069  2081  2083  2087
 2089  2099  2111  2113  2129  2131  2137  2141  2143  2153  2161  2179  2203  2207  2213
 2221  2237  2239  2243  2251  2267  2269  2273  2281  2287  2293  2297  2309  2311  2333
 2339  2341  2347  2351  2357  2371  2377  2381  2383  2389  2393  2399  2411  2417  2423
 2437  2441  2447  2459  2467  2473  2477  2503  2521  2531  2539  2543  2549  2551  2557
 2579  2591  2593  2609  2617  2621  2633  2647  2657  2659  2663  2671  2677  2683  2687
 2689  2693  2699  2707  2711  2713  2719  2729  2731  2741  2749  2753  2767  2777  2789
 2791  2797  2801  2803  2819  2833  2837  2843  2851  2857  2861  2879  2887  2897  2903
 2909  2917  2927  2939  2953  2957  2963  2969  2971  2999  3001  3011  3019  3023  3037
 3041  3049  3061  3067  3079  3083  3089  3109  3119  3121  3137  3163  3167  3169  3181
 3187  3191  3203  3209  3217  3221  3229  3251  3253  3257  3259  3271  3299  3301  3307
 3313  3319  3323  3329  3331  3343  3347  3359  3361  3371  3373  3389  3391  3407  3413
 3433  3449  3457  3461  3463  3467  3469  3491  3499  3511  3517  3527  3529  3533  3539
 3541  3547  3557  3559  3571  3581  3583  3593  3607  3613  3617  3623  3631  3637  3643
 3659  3671  3673  3677  3691  3697  3701  3709  3719  3727  3733  3739  3761  3767  3769
 3779  3793  3797  3803  3821  3823  3833  3847  3851  3853  3863  3877  3881  3889  3907
 3911  3917  3919  3923  3929  3931  3943  3947  3967  3989  4001  4003  4007  4013  4019
 4021  4027  4049  4051  4057  4073  4079  4091  4093  4099  4111  4127  4129  4133  4139
 4153  4157  4159  4177  4201  4211  4217  4219  4229  4231  4241  4243  4253  4259  4261
 4271  4273  4283  4289  4297  4327  4337  4339  4349  4357  4363  4373  4391  4397  4409
 4421  4423  4441  4447  4451  4457  4463  4481  4483  4493  4507  4513  4517  4519  4523
 4547  4549  4561  4567  4583  4591  4597  4603  4621  4637  4639  4643  4649  4651  4657
 4663  4673  4679  4691  4703  4721  4723  4729  4733  4751  4759  4783  4787  4789  4793
 4799  4801  4813  4817  4831  4861  4871  4877  4889  4903  4909  4919  4931  4933  4937
 4943  4951  4957  4967  4969  4973  4987  4993  4999  5003  5009  5011  5021  5023  5039
 5051  5059  5077  5081  5087  5099  5101  5107  5113  5119  5147  5153  5167  5171  5179
 5189  5197  5209  5227  5231  5233  5237  5261  5273  5279  5281  5297  5303  5309  5323
 5333  5347  5351  5381  5387  5393  5399  5407  5413  5417  5419  5431  5437  5441  5443
 5449  5471  5477  5479  5483  5501  5503  5507  5519  5521  5527  5531  5557  5563  5569
 5573  5581  5591  5623  5639  5641  5647  5651  5653  5657  5659  5669  5683  5689  5693
 5701  5711  5717  5737  5741  5743  5749  5779  5783  5791  5801  5807  5813  5821  5827
 5839  5843  5849  5851  5857  5861  5867  5869  5879  5881  5897  5903  5923  5927  5939
 5953  5981  5987  6007  6011  6029  6037  6043  6047  6053  6067  6073  6079  6089  6091
 6101  6113  6121  6131  6133  6143  6151  6163  6173  6197  6199  6203  6211  6217  6221
 6229  6247  6257  6263  6269  6271  6277  6287  6299  6301  6311  6317  6323  6329  6337
 6343  6353  6359  6361  6367  6373  6379  6389  6397  6421  6427  6449  6451  6469  6473
 6481  6491  6521  6529  6547  6551  6553  6563  6569  6571  6577  6581  6599  6607  6619
 6637  6653  6659  6661  6673  6679  6689  6691  6701  6703  6709  6719  6733  6737  6761
 6763  6779  6781  6791  6793  6803  6823  6827  6829  6833  6841  6857  6863  6869  6871
 6883  6899  6907  6911  6917  6947  6949  6959  6961  6967  6971  6977  6983  6991  6997
 7001  7013  7019  7027  7039  7043  7057  7069  7079  7103  7109  7121  7127  7129  7151
 7159  7177  7187  7193  7207  7211  7213  7219  7229  7237  7243  7247  7253  7283  7297
 7307  7309  7321  7331  7333  7349  7351  7369  7393  7411  7417  7433  7451  7457  7459
 7477  7481  7487  7489  7499  7507  7517  7523  7529  7537  7541  7547  7549  7559  7561
 7573  7577  7583  7589  7591  7603  7607  7621  7639  7643  7649  7669  7673  7681  7687
 7691  7699  7703  7717  7723  7727  7741  7753  7757  7759  7789  7793  7817  7823  7829
 7841  7853  7867  7873  7877  7879  7883  7901  7907  7919  7927  7933  7937  7949  7951
 7963  7993  8009  8011  8017  8039  8053  8059  8069  8081  8087  8089  8093  8101  8111
 8117  8123  8147  8161  8167  8171  8179  8191  8209  8219  8221  8231  8233  8237  8243
 8263  8269  8273  8287  8291  8293  8297  8311  8317  8329  8353  8363  8369  8377  8387
 8389  8419  8423  8429  8431  8443  8447  8461  8467  8501  8513  8521  8527  8537  8539
 8543  8563  8573  8581  8597  8599  8609  8623  8627  8629  8641  8647  8663  8669  8677
 8681  8689  8693  8699  8707  8713  8719  8731  8737  8741  8747  8753  8761  8779  8783
 8803  8807  8819  8821  8831  8837  8839  8849  8861  8863  8867  8887  8893  8923  8929
 8933  8941  8951  8963  8969  8971  8999  9001  9007  9011  9013  9029  9041  9043  9049
 9059  9067  9091  9103  9109  9127  9133  9137  9151  9157  9161  9173  9181  9187  9199
 9203  9209  9221  9227  9239  9241  9257  9277  9281  9283  9293  9311  9319  9323  9337
 9341  9343  9349  9371  9377  9391  9397  9403  9413  9419  9421  9431  9433  9437  9439
 9461  9463  9467  9473  9479  9491  9497  9511  9521  9533  9539  9547  9551  9587  9601
 9613  9619  9623  9629  9631  9643  9649  9661  9677  9679  9689  9697  9719  9721  9733
 9739  9743  9749  9767  9769  9781  9787  9791  9803  9811  9817  9829  9833  9839  9851
 9857  9859  9871  9883  9887  9901  9907  9923  9929  9931  9941  9949  9967  9973
\end{vncode}
\endgroup


\clearpage

\section{C++ và thư viện STD}

\subsection{Macros}

Các macro và alias thường dùng trong template. Phần mã được giữ nguyên.

\codepath{header.cpp}

\begingroup\scriptsize
\begin{vncode}
#include <bits/stdc++.h>
#define DEBUG false
#define debugf if (DEBUG) printf
#define MAXN 200309
#define MAXM 900009
#define MAXLOGN 20
#define ALFA 256
#define MOD 1000000007
#define INF 0x3f3f3f3f
#define INFLL 0x3f3f3f3f3f3f3f3f
#define EPS 1e-9
#define PI 3.141592653589793238462643383279502884
#define FOR(x,n) for(int x=0; (x)<int(n); (x)++)
#define FOR1(x,n) for(int x=1; (x)<=int(n); (x)++)
#define REP(x,n) for(int x=int(n)-1; (x)>=0; (x)--)
#define REP1(x,n) for(int x=(n); (x)>0; (x)--)
#define pb push_back
#define pf push_front
#define fi first
#define se second
#define mp make_pair
#define sz(x) int(x.size())
#define all(x) x.begin(), x.end()
#define mset(x,y) memset(&x, (y), sizeof(x))
using namespace std;
typedef long long ll;
typedef unsigned long long ull;
typedef long double ld;
typedef unsigned int uint;
typedef vector<int> vi;
typedef pair<int, int> ii;

int main() {}
\end{vncode}
\endgroup


\subsection{Trình biên dịch GNU}

Một số builtin của GNU dịch trực tiếp thành lệnh máy: \texttt{\_\_builtin\_ffs}, \texttt{\_\_builtin\_clz}, \texttt{\_\_builtin\_ctz}, \texttt{\_\_builtin\_popcount}, \texttt{\_\_builtin\_parity} và các biến thể \texttt{ll}. Không gọi các hàm đếm zero đầu/cuối với đối số 0.

\subsection{I/O nhanh}

Để tăng tốc \texttt{cin}/\texttt{cout}, dùng \texttt{ios\_base::sync\_with\_stdio(0)}. \texttt{gets} nhanh hơn \texttt{scanf}, \texttt{scanf} nhanh hơn \texttt{cin}; \texttt{puts} nhanh hơn \texttt{printf}, \texttt{printf} nhanh hơn \texttt{cout}. Không dùng \texttt{gets} ngoài lập trình thi đấu. Mẫu đọc nhanh giữ nguyên bên dưới.

\codepath{Miscelaneous/fastio.cpp}

\begingroup\scriptsize
\begin{vncode}
#include <cstdio>
#define MAXI 1000009

/*
 * Pre-read all input fast
 */

char input[MAXI];
void readInput() {
	input[fread(input, 1, sizeof(input)-4, stdin)] = 0;
}

/*
 * a bit faster than scanf("%d", &x)
 */

int readInt() {
	bool minus = false;
	int ans = 0;
	char c;
	while (c = getchar()) {
		if (c == '-') break;
		if (c >= '0' && c <= '9') break;
	}
	if (c == '-') minus = true;
	else ans = c-'0';
	while (c = getchar()) {
		if (c < '0' || c > '9') break;
		ans = ans*10 + (c - '0');
	}
	return minus ? -ans : ans;
}

/*
 * TEST MATRIX
 */

int main() {
	while(true) {
		int x = readInt();
		if (x < -10000) return 0;
		printf("x = %d\n", x);
	}
}
\end{vncode}
\endgroup


\subsection{Kiểm tra overflow}

Trước khi cộng, trừ, nhân với biên lớn, kiểm tra dấu và so với \texttt{INFLL}; ví dụ \texttt{b > 0 \&\& a > INFLL-b} báo \texttt{a+b} overflow, \texttt{b < 0 \&\& a < -INFLL-b} báo underflow.

\subsection{C++11}

\texttt{auto a = b} lấy kiểu của \texttt{b}; \texttt{auto a = b()} lấy kiểu trả về; \texttt{for (T a : b)} duyệt phần tử; \texttt{for (T \&a : b)} duyệt tham chiếu; lambda có dạng \texttt{[](params) -> type \{ body \}}.

\subsection{Complex}

Ví dụ: \texttt{\#include <complex>}, \texttt{complex<double> point}. Các hàm hữu ích: \texttt{real}, \texttt{imag}, \texttt{abs}, \texttt{arg}, \texttt{norm}, \texttt{conj}, \texttt{polar}.

\subsection{Pair}

\texttt{\#include <utility>}, \texttt{pair<tipo1,tipo2> P}, với \texttt{first} và \texttt{second}.

\subsection{List}

Các constructor và thao tác chính: \texttt{begin}, \texttt{end}, \texttt{rbegin}, \texttt{rend}, \texttt{size}, \texttt{empty}, \texttt{clear}, \texttt{swap}, \texttt{front}, \texttt{back}, \texttt{push\_back}, \texttt{pop\_back}, \texttt{push\_front}, \texttt{pop\_front}, \texttt{erase}, \texttt{insert}.

\subsection{Vector}

Các thao tác chính: \texttt{begin}, \texttt{end}, \texttt{size}, \texttt{empty}, \texttt{clear}, \texttt{reserve}, \texttt{front}, \texttt{back}, \texttt{erase}, \texttt{pop\_back}, \texttt{push\_back}, \texttt{swap}.

\subsection{Deque}

\texttt{\#include <queue>}, \texttt{deque<tipo> Q}, truy cập ngẫu nhiên được bằng \texttt{Q[i]}; hỗ trợ \texttt{front}, \texttt{back}, \texttt{pop\_front}, \texttt{pop\_back}, \texttt{push\_front}, \texttt{push\_back}.

\subsection{Queue}

\texttt{\#include <queue>}, \texttt{queue<tipo> Q}; các thao tác: \texttt{back}, \texttt{empty}, \texttt{front}, \texttt{pop}, \texttt{push}, \texttt{size}.

\subsection{Stack}

\texttt{\#include <stack>}, \texttt{stack<tipo> P}; các thao tác: \texttt{empty}, \texttt{pop}, \texttt{push}, \texttt{size}, \texttt{top}.

\subsection{Map}

\texttt{\#include <map>}, ví dụ \texttt{map<string,int> si}; dùng \texttt{find}, \texttt{lower\_bound}, \texttt{upper\_bound}, \texttt{erase}, \texttt{count}. Khi duyệt map, \texttt{it->first} là khóa và \texttt{it->second} là giá trị.

\subsection{Set}

\texttt{\#include <set>}, dùng \texttt{find}, \texttt{insert}, \texttt{erase}, \texttt{lower\_bound}, \texttt{upper\_bound}. Với comparator tùy chỉnh, định nghĩa \texttt{struct cmp} có \texttt{bool operator()(tipo,tipo) const}.

\subsection{Ordered set}

Dùng GNU PBDS \texttt{tree} để có \texttt{find\_by\_order(p)} và \texttt{order\_by\_key(v)}, tức set có thống kê thứ tự.

\subsection{Unordered set và map}

Tương tự \texttt{set} và \texttt{map}, nhưng dùng bảng băm, thường nhanh hơn; cần C++11: \texttt{unordered\_set<tipo>}, \texttt{unordered\_map<chave, valor>}.

\subsection{Priority Queue}

\texttt{\#include <queue>}, \texttt{priority\_queue<tipo> pq}; thao tác: \texttt{empty}, \texttt{size}, \texttt{top}, \texttt{push}, \texttt{pop}. Mặc định phần tử lớn hơn đứng trước.

\subsection{Bitset}

\texttt{\#include <bitset>}, \texttt{bitset<MAXN> bs}; hỗ trợ \texttt{count}, \texttt{to\_string}, \texttt{to\_ulong}, \texttt{to\_ullong}, \texttt{set}, \texttt{reset}, \texttt{flip} và các toán tử bit.

\subsection{String}

\texttt{\#include <string>}, hỗ trợ nối bằng \texttt{+}, \texttt{+=}, \texttt{append}, thêm/xóa bằng \texttt{push\_back}, \texttt{erase}, \texttt{replace}, tìm bằng \texttt{find}, lấy con bằng \texttt{substr}.

\subsection{Algorithm và numeric}

Trong các mô tả sau, \texttt{beg} và \texttt{end} là con trỏ hoặc iterator. Comparator là hàm \texttt{bool comp(T a,T b)}, evaluator là \texttt{bool eval(T a)}, accumulator là \texttt{T add(T a,T b)}. Có thể viết bằng lambda C++11.

\subsection{Algorithm: không sửa đổi}

\texttt{any\_of}, \texttt{all\_of}, \texttt{none\_of}, \texttt{for\_each}, \texttt{count}, \texttt{count\_if} duyệt đoạn \texttt{[beg,end)} mà không thay đổi nội dung.

\subsection{Algorithm: sửa đổi}

\texttt{copy}, \texttt{fill}, \texttt{generate}, \texttt{remove}, \texttt{replace}, \texttt{swap}, \texttt{reverse}, \texttt{rotate}, \texttt{random\_shuffle}, \texttt{unique} thao tác trực tiếp trên đoạn hoặc container.

\subsection{Algorithm: phân hoạch}

\texttt{partition} đưa phần tử thỏa evaluator lên trước; \texttt{stable\_partition} giữ thứ tự tương đối trong từng phần.

\subsection{Algorithm: sắp xếp}

\texttt{is\_sorted}, \texttt{sort}, \texttt{stable\_sort}, \texttt{nth\_element} dùng toán tử \texttt{<} hoặc comparator.

\subsection{Algorithm: tìm kiếm nhị phân}

\texttt{binary\_search}, \texttt{lower\_bound}, \texttt{upper\_bound}, \texttt{equal\_range} dùng trên đoạn đã sắp xếp.

\subsection{Algorithm: heap}

\texttt{make\_heap}, \texttt{push\_heap}, \texttt{pop\_heap}, \texttt{sort\_heap} thao tác trên heap trong một đoạn.

\subsection{Algorithm: cực đại và cực tiểu}

\texttt{min}, \texttt{max}, \texttt{minmax}, \texttt{min\_element}, \texttt{max\_element}, \texttt{minmax\_element}, \texttt{lexicographical\_compare}.

\subsection{Algorithm: hoán vị}

\texttt{next\_permutation} và \texttt{prev\_permutation} sinh hoán vị kế tiếp hoặc trước đó theo thứ tự từ điển.

\subsection{Numeric: tích lũy}

\texttt{accumulate}, \texttt{partial\_sum}, \texttt{adjacent\_difference}, \texttt{inner\_product}, \texttt{iota}.

\subsection{Functional}

Các function object chuẩn gồm \texttt{plus}, \texttt{minus}, \texttt{multiplies}, \texttt{divides}, \texttt{modulus}, \texttt{negate}, \texttt{equal\_to}, \texttt{less}, \texttt{greater}, \texttt{logical\_and}, \texttt{bit\_and} và các biến thể liên quan.

\clearpage

\section{Cấu trúc dữ liệu}

\subsection{Max-Queue}

Hàng đợi hỗ trợ lấy giá trị lớn nhất trong O(1) khấu hao.

\codepath{Data Structures/maxqueue.cpp}

\begingroup\scriptsize
\begin{vncode}
#include <queue>
#include <list>
using namespace std;
typedef pair<int, int> ii;

/*
 * MAX QUEUE
 */

class MaxQueue {
	list<ii> q, l;
	int cnt = 0;
public:
	MaxQueue() : cnt(0) {}
	void push(int x) {
		ii cur = ii(x, cnt++);
		while(!l.empty() && l.back() <= cur) l.pop_back();
		q.push_back(cur);
		l.push_back(cur);
	}
	int front() { return q.front().first; }
	void pop() {
		if (q.front().second == l.front().second) l.pop_front();
		q.pop_front();
	}
	int max() { return l.front().first; }
	int size() { return q.size(); }
};

/*
 * Codeforces 980F
 */

#include <bits/stdc++.h>
#define DEBUG false
#define debugf if (DEBUG) printf
#define MAXN 500309
#define MAXM 300009
#define MOD 1000000007
#define INF 0x3f3f3f3f
#define EPS 1e-9
#define FOR(x,n) for(int x=0; (x)<(n); (x)++)
#define FOR1(x,n) for(int x=1; (x)<=(n); (x)++)
#define ROF(x,n) for(int x=(n)-1; (x)>=0; (x)--)
#define ROF1(x,n) for(int x=(n); (x)>0; (x)--)
#define pb push_back
#define pf push_front
#define fi first
#define se second
#define all(x) x.begin(), x.end()
typedef vector<int> vi;
typedef long long ll;

int n, m;
vi g[MAXN], st;
int cycle[MAXN], num[MAXN], pos[MAXN], par[MAXN];
int cnt = 0;
vi cycles[MAXN];
int down[MAXN], up[MAXN];

void buildcycles(int u, int p) {
	num[u] = 1;
	pos[u] = st.size();
	st.pb(u);
	for(int v : g[u]) {
		if (v == p) continue;
		else if (num[v] == 1 && cycle[u] == -1) {
			cycles[cnt].clear();
			for(int i = pos[v]; i <= pos[u]; i++) {
				cycle[st[i]] = cnt;
				cycles[cnt].pb(st[i]);
			}
			cnt++;
		}
		else if (num[v] == 0) {
			buildcycles(v, u);
			if (cycle[u] == -1 || cycle[v] != cycle[u]) par[v] = u;
		}
	}
	num[u] = 2;
	st.pop_back();
}

int find(vi & v, int u) {
	FOR(i, int(v.size())) {
		if (v[i] == u) return i;
	}
	return -1;
}

int nodedown(int u);

int cycledown(int u) {
	if (down[u] != -1) return down[u];
	int c = cycle[u];
	if (c == -1) return nodedown(u);
	vi & arr = cycles[c];
	int sz = arr.size();
	nodedown(u);
	int j = find(arr, u);
	assert(j != -1);
	FOR(i, sz) {
		int d = abs(i-j);
		if (i == j) continue;
		down[u] = max(down[u], min(d, sz-d) + nodedown(arr[i]));
	}
	return down[u];
}

int nodedown(int u) {
	if (down[u] != -1) return down[u];
	down[u] = 0;
	for(int v : g[u]) {
		if (v == par[u]) continue;
		if (cycle[u] != -1 && cycle[u] == cycle[v]) continue;
		down[u] = max(down[u], 1 + cycledown(v));
	}
	return down[u];
}

void solve(vi & L, vi & Z) {
	int sz = L.size();
	int k = sz/2;
	FOR(i, sz) L.pb(L[i]);
	vi left = L, right = L;
	FOR(i, 2*sz) left[i] += i;
	FOR(i, 2*sz) right[i] += 2*sz-i-1;
	Z.assign(sz, 0);
	MaxQueue q;
	for(int i = 0; i < sz; i++) {
		q.push(right[i]);
	}
	while(q.size() > k) q.pop();
	for(int i = sz; i < 2*sz; i++) {
		Z[i-sz] = max(Z[i-sz], q.max() - (2*sz-i-1));
		q.push(right[i]);
		q.pop();
	}
	q = MaxQueue();
	for(int i = 2*sz-1; i >= sz; i--) {
		q.push(left[i]);
	}
	while(q.size() > k) q.pop();
	for(int i = sz-1; i >= 0; i--) {
		Z[i] = max(Z[i], q.max() - i);
		q.push(left[i]);
		q.pop();
	}
}

void nodeup(int u, int d);

void cycleup(int u, int d) {
	int c = cycle[u];
	if (c == -1) {
		nodeup(u, d);
		return;
	}
	cycledown(u);
	vi & arr = cycles[c];
	int sz = arr.size();
	up[u] = d;
	int j = find(arr, u);
	assert(j != -1);
	vi L(sz), Z;
	int daux = d;
	for(int v : g[u]) {
		if (v == par[u]) continue;
		if (cycle[u] != -1 && cycle[u] == cycle[v]) continue;
		daux = max(daux, 1 + nodedown(v));
	}
	FOR(i, sz) {
		if (i == j) L[i] = daux;
		else L[i] = nodedown(arr[i]);
	}
	solve(L, Z);
	FOR(i, sz) {
		if (i == j) nodeup(arr[i], max(Z[i], d));
		else nodeup(arr[i], Z[i]);
	}
}

void nodeup(int u, int d) {
	up[u] = d;
	nodedown(u);
	int sz = g[u].size();
	int x = d, y = -1;
	for(int v : g[u]) {
		if (v == par[u]) continue;
		if (cycle[u] != -1 && cycle[u] == cycle[v]) continue;
		if (x >= y && 1+nodedown(v) > y) y = 1+down[v];
		else if (x < y && 1+nodedown(v) > x) x = 1+down[v];
	}
	for(int v : g[u]) {
		if (v == par[u]) continue;
		if (cycle[u] != -1 && cycle[u] == cycle[v]) continue;
		if (1+nodedown(v) == max(x, y))
			cycleup(v, min(x, y)+1);
		else cycleup(v, max(x, y)+1);
	}
}

int main() {
	scanf("%d %d", &n, &m);
	FOR(j, m) {
		int u, v;
		scanf("%d %d", &u, &v);
		g[u].pb(v);
		g[v].pb(u);
	}
	memset(&cycle, -1, sizeof cycle);
	memset(&down, -1, sizeof down);
	memset(&up, -1, sizeof up);
	memset(&num, 0, sizeof num);
	buildcycles(1, -1);
	cycleup(1, 0);
	FOR1(u, n) {
		printf("%d ", max(up[u], down[u]));
	}
	printf("\n");
	return 0;
}
\end{vncode}
\endgroup


\subsection{Union-Find}

Disjoint Set Union với nén đường đi và union theo kích thước/hạng.

\codepath{Data Structures/unionfind.cpp}

\begingroup\scriptsize
\begin{vncode}
#include <cstdio>
#include <vector>
using namespace std;

class UnionFind {
private:
	vector<int> parent, rank;
public:
	UnionFind(int N) {
		rank.assign(N+1, 0);
		parent.assign(N+1, 0);
		for (int i = 0; i <= N; i++) parent[i] = i;
	}
	int find(int i) {
		while(i != parent[i]) i = parent[i];
		return i;
	}
	bool isSameSet(int i, int j) {
		return find(i) == find(j);
	}
	void unionSet (int i, int j) {
		if (isSameSet(i, j)) return;
		int x = find(i), y = find(j);
		if (rank[x] > rank[y]) parent[y] = x;
		else {
			parent[x] = y;
			if (rank[x] == rank[y]) rank[y]++;
		}
	}
};

int main() {
	int a, b;
	char c;
	UnionFind uf(100009);
	while(scanf(" %c %d %d", &c, &a, &b)!=EOF) {
		if (c=='U' || c=='u') {
			uf.unionSet(a,b);
		}
		else if (c=='i' || c=='I') {
			if (uf.isSameSet(a,b)) printf("yes\n");
			else printf("no\n");
		}
		else break;
	}
	return 0;
}
\end{vncode}
\endgroup


\subsection{Binary Indexed Tree / Fenwick Tree}

Fenwick tree cho cộng điểm và truy vấn tổng tiền tố.

\codepath{Trees/bit.cpp}

\begingroup\scriptsize
\begin{vncode}
#include <cstdio>
#include <vector>
using namespace std;

const int neutral = 0;
#define comp(a, b) ((a)+(b))

class FenwickTree {
private:
	vector<int> ft;
public:
	FenwickTree(int n) { ft.assign(n + 1, 0); }
	int rsq(int i) { // returns RSQ(1, i)
		int sum = neutral;
		for(; i; i -= (i & -i))
			sum = comp(sum, ft[i]);
		return sum;
	}
	int rsq(int i, int j) {
		return rsq(j) - rsq(i - 1);
	}
	void update(int i, int v) {
		for(; i > 0 && i < (int)ft.size(); i += (i & -i))
			ft[i] = comp(v, ft[i]);
	}
};

int main() {}
\end{vncode}
\endgroup


\subsection{Binary Indexed Tree / Fenwick Tree với cập nhật và truy vấn đoạn}

Biến thể Fenwick hỗ trợ cập nhật đoạn và truy vấn đoạn.

\codepath{Trees/rangebit.cpp}

\begingroup\scriptsize
\begin{vncode}
#include <cstdio>
#include <vector>
#include <ctime>
#include <algorithm>
using namespace std;

/*
 * 2D BIT
 */

class FenwickTree {
private:
	vector<int> ft1, ft2;
	int rsq(vector<int> & ft, int i) {
		int sum = 0;
		for(; i; i -= (i & -i)) sum += ft[i];
		return sum;
	}
	void update(vector<int> & ft, int i, int v) {
		for(; i < (int)ft.size(); i += (i & -i))
			ft[i] += v;
	}
public:
	FenwickTree(int n) {
		ft1.assign(n + 1, 0); //1-indexed
		ft2.assign(n + 1, 0); //1-indexed
	}
	void update(int i, int j, int v) {
		update(ft1, i, v);
		update(ft1, j+1, -v);
		update(ft2, i, v*(i-1));
		update(ft2, j+1, -v*j);
	}
	int rsq(int i) {
		return rsq(ft1, i)*i - rsq(ft2, i);
	}
	int rsq(int i, int j) {
		return rsq(j) - rsq(i-1);
	}
};

/*
 * TEST MATRIX
 */

int vet[10009];
int sum(int a, int b) {
	int ans = vet[a];
	for(int i=a+1; i<=b; i++) ans += vet[i];
	return ans;
}

void update(int a, int b, int k) {
	for(int i=a; i<=b; i++) vet[i] += k;
}


int main() {
	int N = 10000, a, b, k;
	srand(time(NULL));
	FenwickTree ft(N);
	for(int i=1; i<=N; i++) {
		vet[i] = rand()%N;
		ft.update(i, i, vet[i]);
	}
	for(int q=0; q<100; q++) {

		printf("test #%d:", q+1);
		a = (rand()%N) + 1;
		b = (rand()%N) + 1;
		if (a>b) swap(a, b);
		k = rand()%100;
		printf(" %d to %d plus %d\n", a, b, k);

		update(a, b, k);
		ft.update(a, b, k);

		a = (rand()%N) + 1;
		b = (rand()%N) + 1;
		if (a>b) swap(a, b);

		if (sum(a, b) != ft.rsq(a, b)) printf("test failed\n");
	}
	return 0;
}
\end{vncode}
\endgroup


\subsection{2D Binary Indexed Tree / Fenwick Tree}

Fenwick hai chiều.

\codepath{Trees/2dbit.cpp}

\begingroup\scriptsize
\begin{vncode}
#include <vector>
using namespace std;

const int neutral = 0;
#define comp(a, b) ((a)+(b))

class FenwickTree2D {
private:
	vector< vector<int> > ft;
public:
	FenwickTree2D(int n, int m) {
		ft.assign(n + 1, vector<int>(m + 1, 0));	//1-indexed
	}
	int rsq(int i, int j) { // returns RSQ((1,1), (i,j))
		int sum = 0, _j = j;
		while(i > 0) {
			j = _j;
			while(j > 0) {
				sum = comp(sum, ft[i][j]);
				j -= (j & -j);
			}
			i -= (i & -i);
		}
		return sum;
	}
	void update(int i, int j, int v) {
		int _j = j;
		while(i < (int)ft.size()) {
			j = _j;
			while(j < (int)ft[i].size()) {
				ft[i][j] = comp(v, ft[i][j]);
				j += (j & -j);
			}
			i += (i & -i);
		}
	}
};

/*
 *	TEST MATRIX
 */
#include <cstdio>
#include <algorithm>

vector< vector<int> > A;

int rsq(int x, int y) {
	int ans = neutral;
	for(int i=1; i<=x; i++) {
		for(int j=1; j<=y; j++) {
			ans = comp(ans, A[i][j]);
		}
	}
	return ans;
}

void update(int x, int y, int v) {
	A[x][y] += v;
}

bool test() {
	int N = 500, nTests = 10000;
	int x, y, v;
	A.resize(N+1);
	for(int i=1; i<=N; i++) A[i].resize(N+1);
	FenwickTree2D ft(N, N);
	for(int i=1; i<=N; i++) {
		for(int j=1; j<=N; j++) {
			A[i][j] = rand() % N;
			ft.update(i, j, A[i][j]);
		}
	}
	printf("starting tests...\n");
	for(int q=1; q<=nTests; q++) {
		x = rand()%N + 1;
		y = rand()%N + 1;
		v = rand()%N;
		ft.update(x, y, v);
		update(x, y, v);
		x = rand()%N + 1;
		y = rand()%N + 1;
		int q1 = rsq(x, y);
		int q2 = ft.rsq(x, y);
		if (q1 != q2) {
			printf("Failed test %d, q1 = %d, q2 = %d\n", q, q1, q2);
			return false;
		}
	}
	printf("all tests passed\n");
	return true;
}

int main() {
	test();
	return 0;
}
\end{vncode}
\endgroup


\subsection{Segment Tree}

Cây đoạn cơ bản.

\codepath{Trees/segmenttree.cpp}

\begingroup\scriptsize
\begin{vncode}
#include <cstdio>
#include <vector>
#include <algorithm>
#define INF (1<<30)
using namespace std;

/*
 * Segment Tree
 */

const int neutral = 0;
#define comp(a, b) ((a)+(b))

class SegmentTree {
	vector<int> a;
	int n;
public:
	SegmentTree(int* st, int* en) {
		int sz = int(en-st);
		for (n = 1; n < sz; n <<= 1);
		a.assign(n << 1, neutral);
		for(int i=0; i<sz; i++) a[i+n] = st[i];
		for(int i=n+sz-1; i; i--)
			a[i>>1] = comp(a[i>>1], a[i]);
	}
	void update(int i, int x) {
		a[i += n] = x; //substitui
		for (i >>= 1; i; i >>= 1)
			a[i] = comp(a[i<<1], a[1+(i<<1)]);
	}
	int query(int l, int r) {
		int ans = neutral;
		for (l += n, r += n+1; l < r; l >>= 1, r >>= 1) {
			if (l & 1) ans = comp(ans, a[l++]);
			if (r & 1) ans = comp(ans, a[--r]);
		}
		return ans;
	}
};

/*
 * TEST MATRIX
 */

int vet[16];

int main() {
	int n;
	scanf("%d", &n);
	for(int i=0; i<n; i++) scanf("%d", &vet[i]);
	char c;
	int a, b;
	SegmentTree st(vet, vet+n);
	while(scanf(" %c", &c), c=='u' || c=='s' ) {
		if (c == 'u') {
			scanf("%d %d", &a, &b);
			st.update(a, b);
		}
		else {
			scanf("%d %d", &a, &b);
			printf("soma(%d,%d)=%d\n", a, b, st.query(a, b));
		}
	}
	return 0;
}
\end{vncode}
\endgroup


\subsection{Segment Tree với Lazy Propagation}

Cây đoạn có lazy propagation.

\codepath{Trees/lazysegmenttree.cpp}

\begingroup\scriptsize
\begin{vncode}
#include <cstdio>
#include <vector>
#include <algorithm>
#define INF (1<<30)
using namespace std;

/*
 * Lazy Propagation Segment Tree
 */

const int neutral = 0; //comp(x, neutral) = x
#define comp(a, b) ((a)+(b))

class SegmentTree {
private:
	vector<int> st, lazy;
	int size;
#define left(p) (p << 1)
#define right(p) ((p << 1) + 1)
	void build(int p, int l, int r, int* A) {
		if (l == r) { st[p] = A[l]; return; }
		int m = (l + r) / 2;
		build(left(p), l, m, A);
		build(right(p), m+1, r, A);
		st[p] = comp(st[left(p)], st[right(p)]);
	}
	void push(int p, int l, int r) {
		st[p] += (r - l + 1)*lazy[p];	//Caso RSQ
		//st[p] += lazy[p]; 		    //Caso RMQ
		if (l != r) {
			lazy[right(p)] += lazy[p];
			lazy[left(p)] += lazy[p];
		}
		lazy[p] = 0;
	}
	void update(int p, int l, int r, int a, int b, int k) {
		push(p, l, r);
		if (a > r || b < l) return;
		else if (l >= a && r <= b) {
			lazy[p] = k; push(p, l, r); return;
		}
		int m = (l + r) / 2;
		update(left(p), l, m, a, b, k);
		update(right(p), m+1, r, a, b, k);
		st[p] = comp(st[left(p)], st[right(p)]);
	}
	int query(int p, int l, int r, int a, int b) {
		push(p, l, r);
		if (a > r || b < l) return neutral;
		if (l >= a && r <= b) return st[p];
		int m = (l + r) / 2;
		int p1 = query(left(p), l, m, a, b);
		int p2 = query(right(p), m+1, r, a, b);
		return comp(p1, p2);
	}
public:
	SegmentTree(int* bg, int* en) {
		size = (int)(en - bg);
		st.assign(4 * size, neutral);
		lazy.assign(4 * size, 0);
		build(1, 0, size - 1, bg);
	}
	int query(int a, int b) { return query(1, 0, size - 1, a, b); }
	void update(int a, int b, int k) { update(1, 0, size - 1, a, b, k); }
};

/*
 * TEST MATRIX
 */

int vet[16];

int main() {
	int n, k;
	scanf("%d", &n);
	for(int i=0; i<n; i++) scanf("%d", &vet[i]);
	char c;
	int a, b;
	SegmentTree st(vet, vet+n);
	while(scanf(" %c", &c), c=='u' || c=='s' ) {
		if (c == 'u') {
			scanf("%d %d %d", &a, &b, &k);
			st.update(a, b, k);
		}
		else {
			scanf("%d %d", &a, &b);
			printf("soma(%d,%d)=%d\n", a, b, st.query(a, b));
		}
	}
	return 0;
}
\end{vncode}
\endgroup


\subsection{2D Segment Tree}

Cây đoạn hai chiều.

\codepath{Trees/2dsegmenttree.cpp}

\begingroup\scriptsize
\begin{vncode}
#include <algorithm>
#include <vector>
using namespace std;
#define MAXM 100000009 //1e8+9
#define INF 0x3f3f3f3f

/*
 * 2D Segment Tree
 */

int rs[MAXM], ls[MAXM], st[MAXM], cnt = 0;

class SegmentTree2D {
	int sizex, sizey, v, root;
	int x, y, ix, jx, iy, jy;
	void updatex(int p, int lx, int rx) {
		if (x < lx || rx < x) return;
		st[p] += v;
		if (lx == rx) return;
		if (!rs[p]) rs[p] = ++cnt, ls[p] = ++cnt;
		int mx = (lx + rx) / 2;
		updatex(ls[p], lx, mx);
		updatex(rs[p], mx + 1, rx);
	}
	void updatey(int p, int ly, int ry) {
		if (y < ly || ry < y) return;
		if (!st[p]) st[p] = ++cnt;
		updatex(st[p], 0, sizex);
		if (ly == ry) return;
		if (!rs[p]) rs[p] = ++cnt, ls[p] = ++cnt;
		int my = (ly + ry) / 2;
		updatey(ls[p], ly, my);
		updatey(rs[p], my + 1, ry);
	}
	int queryx(int p, int lx, int rx) {
		if (!p || jx < lx || ix > rx) return 0;
		if (ix <= lx && rx <= jx) return st[p];
		int mx = (lx + rx) / 2;
		return queryx(ls[p], lx, mx) +
			queryx(rs[p], mx + 1, rx);
	}
	int queryy(int p, int ly, int ry) {
		if (!p || jy < ly || iy > ry) return 0;
		if (iy <= ly && ry <= jy) return queryx(st[p], 0, sizex);
		int my = (ly + ry) / 2;
		return queryy(ls[p], ly, my) +
			queryy(rs[p], my + 1, ry);
	}
public:
	SegmentTree2D(int nx, int ny) : sizex(nx), sizey(ny) {
		root = ++cnt;
	}
	void update(int _x, int _y, int _v) {
		x = _x; y = _y; v = _v;
		updatey(root, 0, sizey);
	}
	int query(int _ix, int _jx, int _iy, int _jy) {
		ix = _ix; jx = _jx; iy = _iy; jy = _jy;
		return queryy(root, 0, sizey);
	}
};

/*
 *	TEST MATRIX
 */
#include <cstdio>

vector< vector<int> > A;

int query(int lx, int rx, int ly, int ry) {
	int ans = 0;
	for(int i=lx; i<=rx; i++) {
		for(int j=ly; j<=ry; j++) {
			ans += A[i][j];
		}
	}
	return ans;
}

void update(int x, int y, int v) {
	A[x][y] += v;
}

bool test() {
	int N = 500, nTests = 100000;
	int lx, rx, ly, ry, x, y, v;
	A.resize(N);
	for(int i=0; i<N; i++) {
		A[i].assign(N, 0);
	}
	SegmentTree2D st(N, N);
	printf("starting tests...\n");
	for(int q=1; q<=nTests; q++) {
		x = rand()%N;
		y = rand()%N;
		v = rand()%N;
		st.update(x, y, v);
		update(x, y, v);
		lx = rand()%N;
		rx = rand()%N;
		ly = rand()%N;
		ry = rand()%N;
		int q1 = query(lx, rx, ly, ry);
		int q2 = st.query(lx, rx, ly, ry);
		if (q1 != q2) {
			printf("Failed test %d, q1 = %d, q2 = %d\n", q, q1, q2);
			return false;
		}
	}
	printf("all tests passed\n");
	return true;
}

int main() {
	test();
	return 0;
}
\end{vncode}
\endgroup


\subsection{Persistent Segment Tree}

Cây đoạn persistent; \texttt{update} trả về phiên bản mới, phiên bản 0 là cây ban đầu.

\codepath{Trees/persistentsegtree.cpp}

\begingroup\scriptsize
\begin{vncode}
#include <vector>
#include <algorithm>
#define INF (1<<30)
using namespace std;
#define MAXS 2000009
#define MAXN 200009

/*
 * Persistent Segment Tree
 */

const int neutral = 0; //comp(x, neutral) = x
#define comp(a, b) ((a)+(b))

int nds, st[MAXS], ls[MAXS], rs[MAXS];

class PersistentSegmentTree {
private:
	int vroot[MAXN];
	int size, nds, nv;
	int newnode() {
		ls[nds] = rs[nds] = -1;
		st[nds++] = neutral;
		return nds-1;
	}
	void build(int p, int l, int r, int* A) {
		if (l == r) {
			st[p] = A ? A[l] : neutral;
			return;
		}
		int m = (l + r) / 2;
		build(ls[p] = newnode(), l, m, A);
		build(rs[p] = newnode(), m+1, r, A);
		st[p] = comp(st[ls[p]], st[rs[p]]);
	}
	void update(int prv, int p, int l, int r, int i, int k) {
		if (l == r) {
			st[p] = k; //substitui
			//st[p] = st[prv] + k; //agrega
			return;
		}
		int m = (l + r) / 2;
		if (i <= m) {
			rs[p] = rs[prv];
			update(ls[prv], ls[p] = newnode(), l, m, i, k);
		}
		else {
			ls[p] = ls[prv];
			update(rs[prv], rs[p] = newnode(), m+1, r, i, k);
		}
		st[p] = comp(st[ls[p]], st[rs[p]]);
	}
	int query(int p, int l, int r, int a, int b) {
		if (a > r || b < l) return neutral;
		if (l >= a && r <= b) return st[p];
		int m = (l + r) / 2;
		int p1 = query(ls[p], l, m, a, b);
		int p2 = query(rs[p], m + 1, r, a, b);
		return comp(p1, p2);
	}
public:
	PersistentSegmentTree(int* begin, int* end) {
		nds = nv = 0; size = (int)(end-begin);
		vroot[nv++] = newnode();
		build(vroot[0], 0, size-1, begin);
	}
	PersistentSegmentTree(int _size = 1) {
		nds = nv = 0; size = _size;
		vroot[nv++] = newnode();
		build(vroot[0], 0, size-1, NULL);
	}
	int query(int a, int b, int v) {
		return query(vroot[v], 0, size-1, a, b);
	}
	int update(int i, int v, int k) {
		vroot[nv++] = newnode();
		update(vroot[v], vroot[nv-1], 0, size-1, i, k);
		return nv-1;
	}
	int nver() { return nv; }
};

/*
 * TEST MATRIX
 */

#include <ctime>
#include <cstdio>
#define DEBUG false
const int N = 500;
int arr[N+9][N+9];
int curver;
PersistentSegmentTree pst;
bool test(int NTests) {
	srand(time(NULL));
	for(int nt=0; nt<NTests; nt++) {
		for(int i=0; i<N; i++) arr[0][i] = rand()%N;
		
		if(DEBUG) printf("building seg tree...\n");
		pst = PersistentSegmentTree(&arr[0][0], &arr[0][0] + N);
		curver = 1;
		if(DEBUG) printf("arr[%d]:", curver-1);
		for(int j=0; j<N; j++) {
			if(DEBUG) printf(" %2d", arr[curver-1][j]);
		}
		if(DEBUG) printf("\n");
		
		for(int i=0; i<N; i++) {
			if(DEBUG) printf("test %d:\n", i+1);
			
			// update
			int id = rand()%N;
			int k = rand()%N;
			int v = rand()%curver;
			if(DEBUG) printf("update over version %d: arr[%d] = %d\n", v, id, k);
			for(int j=0; j<N; j++) arr[curver][j] = arr[v][j];
			curver++;
			arr[curver-1][id] = k;
			if(DEBUG) printf("arr[%d]:", curver-1);
			for(int j=0; j<N; j++) {
				if(DEBUG) printf(" %2d", arr[curver-1][j]);
			}
			if(DEBUG) printf("\n");
			
			if (pst.update(id, v, k) != curver-1) {
				printf("failed to update at test %d\n", i+1);
				return false;
			}
			
			//queries over recent version
			for(int t=0; t<N; t++) {
				int a = rand()%N;
				int b = rand()%N;
				if (a > b) swap(a, b);
				int val1 = 0;
				for(int j=a; j<=b; j++) val1 = comp(val1, arr[curver-1][j]);
				int val2 = pst.query(a, b, curver-1);
				if(DEBUG) printf("sum[%d,%d] = %d (ver %d)\n", a, b, val1, curver-1);
				if (val1 != val2) {
					printf("failed at test %d, case %d (%d==%d), ver %d\n", i+1, t+1, val1, val2, curver-1);
					printf("arr[%d]:", curver-1);
					for(int j=0; j<N; j++) {
						printf(" %2d", arr[curver-1][j]);
					}
					printf("\nquery [%d,%d]\n", a, b);
					return false;
				}
			}
			
			//queries over all versions
			for(int t=0; t<N; t++) {
				int a = rand()%N;
				int b = rand()%N;
				int v = rand()%curver;
				if (a > b) swap(a, b);
				int val1 = 0;
				for(int j=a; j<=b; j++) val1 = comp(val1, arr[v][j]);
				int val2 = pst.query(a, b, v);
				if(DEBUG) printf("sum[%d,%d] = %d (ver %d)\n", a, b, val1, v);
				if (val1 != val2) {
					printf("failed at test %d, case %d (%d==%d), ver %d\n", i+1, t+1, val1, val2, v);
					printf("arr[%d]:", v);
					for(int j=0; j<N; j++) {
						printf(" %2d", arr[v][j]);
					}
					printf("\nquery [%d,%d]\n", a, b);
					return false;
				}
			}		
		}
	}
	printf("all tests passed\n");
	return true;
}

int main() {
	test(10);
	return 0;
}

/*
 * SPOJ DQUERY
 */

/*#define MAXN 100009
#include <cstdio>
#include <map>

int a[MAXN], id[MAXN], ver[MAXN];

bool cmp(int i, int j) {
	return a[i] < a[j];
}

PersistentSegmentTree pst(MAXN);

int main() {
	int N, Q, u, v, x;
	scanf("%d", &N);
	map<int, int> last;
	for(int i=0; i<N; i++) {
		scanf("%d", &x);
		if (last.count(x)) a[i] = last[x];
		else a[i] = -1;
		last[x] = i;
		id[i] = i;
	}
	sort(id, id+N, cmp);
	pst = PersistentSegmentTree(N);
	for(int it=0, i; it<N; it++) {
		i = id[it];
		ver[it] = pst.update(i, pst.nver()-1, pst.query(i, i, pst.nver()-1) + 1);
	}
	scanf("%d", &Q);
	while(Q-->0) {
		scanf("%d %d", &u, &v);
		int lo = -1;
		int hi = N;
		while(hi > lo+1) {
			int m = (hi+lo)/2;
			if (a[id[m]] < u-1) lo = m;
			else hi = m;
		}
		if (lo < 0) printf("0\n");
		else printf("%d\n", pst.query(u-1, v-1, ver[lo]));
	}
	return 0;
}*/
\end{vncode}
\endgroup


\subsection{Li Chao Segment Tree}

Cho mảng điểm \texttt{x[]} đã sắp, hỗ trợ thêm đường thẳng và truy vấn max tại điểm.

\codepath{Trees/lichaosegtree.cpp}

\begingroup\scriptsize
\begin{vncode}
#include <cstdio>
#include <vector>
#define INF 0x3f3f3f3f3f3f3f3f
#define MAXN 1009
using namespace std;

typedef long long ll;

/*
 * LiChao Segment Tree
 */

class LiChao {
	vector<ll> m, b;
	int n, sz; ll *x;
#define gx(i) (i < sz ? x[i] : x[sz-1])
	void update(int t, int l, int r, ll nm, ll nb) {
		ll xl = nm * gx(l) + nb, xr = nm * gx(r) + nb;
		ll yl = m[t] * gx(l) + b[t], yr = m[t] * gx(r) + b[t];
        if (yl >= xl && yr >= xr) return;
		if (yl <= xl && yr <= xr) {
			m[t] = nm, b[t] = nb; return;
		}
		int mid = (l + r) / 2;
		update(t<<1, l, mid, nm, nb);
		update(1+(t<<1), mid+1, r, nm, nb);
	}
public:
	LiChao(ll *st, ll *en) : x(st) {
		sz = int(en - st);
		for(n = 1; n < sz; n <<= 1);
		m.assign(2*n, 0); b.assign(2*n, -INF);
	}
	void insert_line(ll nm, ll nb) {
		update(1, 0, n-1, nm, nb);
	}
	ll query(int i) {
		ll ans = -INF;
		for(int t = i+n; t; t >>= 1)
			ans = max(ans, m[t] * x[i] + b[t]);
		return ans;
	}
};

/*
 * UVa 12524
 */

ll w[MAXN], x[MAXN], A[MAXN], B[MAXN], dp[MAXN][MAXN];

int main(){
	int N, K;
	while(scanf("%d %d", &N, &K)!=EOF) {
		for(int i=0; i<N; i++){
			scanf("%lld %lld", x+i, w+i);
			A[i] = w[i] + (i>0 ? A[i-1] : 0);
			B[i] = w[i]*x[i] + (i>0 ? B[i-1] : 0);
			dp[i][1] = x[i]*A[i] - B[i];
		}
		for(int k=2; k<=K; k++){
			dp[0][k] = 0;
            LiChao lc(x, x+N);
			for(int i=1; i<N; i++){
				lc.insert_line(A[i-1], -dp[i-1][k-1]-B[i-1]);
				dp[i][k] = x[i]*A[i] - B[i] - lc.query(i);
			}
		}
		printf("%lld\n", dp[N-1][K]);
	}
	return 0;
}
\end{vncode}
\endgroup


\subsection{Split-Merge Segment Tree}

Tập các multiset hỗ trợ tạo, hợp nhất, tách, đếm trong đoạn, truy cập phần tử thứ i và xuất ra vector.

\codepath{Trees/splitmergesegtree.cpp}

\begingroup\scriptsize
\begin{vncode}
#include <vector>
using namespace std;
#define MAXN 200009
#define MAXS 5000009

/*
 * Split-merge Segment Tree
 */

const int neutral = 0;
#define comp(a, b) ((a)+(b))

class SegmentTree {
	int cnt, n, root[MAXN], nroot;
	int sz[MAXS], ls[MAXS], rs[MAXS];

	int newnode(int _s) {
		ls[cnt] = rs[cnt] = -1;
		sz[cnt++] = _s;
		return cnt-1;
	}
	int build(int l, int r, int i) {
		int t = newnode(1);
		if (l == r) return t;
		int m = (l + r) / 2;
		if (i <= m) ls[t] = build(l, m, i);
		else rs[t] = build(m + 1, r, i);
		return t;
	}
	int split(int t1, int k) {
		if (t1 == -1 || sz[t1] <= k) return -1;
		int t2 = newnode(sz[t1] - k);
		sz[t1] = k;
		int sl = ls[t1] == -1 ? 0 : sz[ls[t1]];
		if (k > sl) rs[t2] = split(rs[t1], k - sl);
		else swap(rs[t1], rs[t2]);
		if (k < sl) ls[t2] = split(ls[t1], k);
		return t2;
	}
	int merge(int t1, int t2) {
		if(t1 == -1 || t2 == -1) return t1+t2+1;
		ls[t1] = merge(ls[t1], ls[t2]);
		rs[t1] = merge(rs[t1], rs[t2]);
		sz[t1] = comp(sz[t1], sz[t2]);
		return t1;
	}
	int query(int t, int l, int r, int a, int b) {
		if (t == -1 || l > b || r < a) return neutral;
		if (l >= a && r <= b) return sz[t];
		int m = (l + r) / 2;
		int p1 = query(ls[t], l, m, a, b);
		int p2 = query(rs[t], m+1, r, a, b);
		return comp(p1, p2);
	}
	int at(int t, int l, int r, int i) {
		if (l == r) return l;
		int sl = ls[t] == -1 ? 0 : sz[ls[t]];
		int m = (l + r) / 2;
		if (i < sl) return at(ls[t], l, m, i);
		else return at(rs[t], m+1, r, i - sl);
	}
	void to_vector(int t, int l, int r, vector<int> & v) {
		if (t == -1) return;
		int m = (l + r) / 2;
		to_vector(ls[t], l, m, v);
		to_vector(rs[t], m+1, r, v);
		for(int i = 0; i < sz[t] && l == r; i++)
			v.push_back(l);
	}
public:
	SegmentTree() { }
	SegmentTree(int _n) : n(_n), cnt(0), nroot(0) { }
	int newSet(int i) {
		root[nroot++] = build(0, n, i);
		return nroot-1;
	}
	void unionSet(int r1, int r2) {
		root[r1] = merge(root[r1], root[r2]);
		root[r2] = -1;
	}
	int query(int r, int a, int b) { return query(root[r], 0, n, a, b); }
	int size(int r) { return root[r] == -1 ? 0 : sz[root[r]]; }
	int splitSet(int r, int k) {
		if (k >= size(r)) return -1;
		if (k <= 0) {
			root[nroot++] = root[r];
			root[r] = -1;
		}
		else root[nroot++] = split(root[r], k);
		return nroot-1;
	}
	int at(int r, int i) { return at(root[r], 0, n, i); }
	void to_vector(int r, vector<int> & v) { to_vector(root[r], 0, n, v); }
};

/*
 * TEST MATRIX
 */
/*
#include <cstdlib>
#include <algorithm>
#include <ctime>
#include <cstdio>
#define NTESTS 1000009
#define all(x) x.begin(), x.end()

vector<int> a[NTESTS];
int cnt;
SegmentTree st;

void print(int r) {
	printf("%d: {", r);
	bool ok = false;
	for (int v : a[r]) {
		if (ok) printf(", ");
		ok = true;
		printf("%d", v);
	}
	printf("} ");
	printf("\n");
}

void merge(int r1, int r2) {
	a[r1].insert(a[r1].end(), all(a[r2]));
	sort(all(a[r1]));
	a[r2].clear();
}

int query(int r, int i, int j) {
	int ans = 0;
	for(int v : a[r]) {
		if (v >= i && v <= j) ans++;
	}
	return ans;
}

int split(int r1, int k) {
	int sz = (r1 == -1 ? 0 : int(a[r1].size()));
	if (k >= sz) {
		return -1;
	}
	else if (k <= 0) {
		a[cnt].clear();
		a[r1].swap(a[cnt]);
		cnt++;
		return cnt-1;
	}
	else {
		a[cnt] = vector<int>(a[r1].begin()+k, a[r1].end());
		a[r1].resize(k);
		cnt++;
		return cnt-1;
	}
}

bool test(int ntests) {
	srand(time(NULL));
	vector<int> alive;
	st = SegmentTree(ntests);
	cnt = 0;
	for(int i = 0; i <= ntests; i++) {
		a[cnt].clear();
		a[cnt].push_back(i);
		st.newSet(i);
		alive.push_back(cnt);
		cnt++;
	}
	for(int test = 0; test < ntests; test++) {
		//printf("\ntest %d\n", test);
		int r1, r2, i1, i2;
		//merge
		do {
			i1 = rand()%alive.size();
			i2 = rand()%alive.size();
		} while(i1 == i2);
		r1 = alive[i1];
		r2 = alive[i2];
		merge(r1, r2);
		st.unionSet(r1, r2);
		swap(alive[i2], alive.back());
		alive.pop_back();
		//split
		i1 = rand()%alive.size();
		r1 = alive[i1];
		int k = rand()%(a[r1].size()+1);
		r2 = split(r1, k);
		int r2l = st.splitSet(r1, k);
		if (r2 != r2l) {
			printf("failed on test %d, split r2test=%d r2st=%d\n", test, r2, r2l);
			return false;
		}
		if (r2 != -1) alive.push_back(r2);
		//test element by element
		for(int j = 0; j <= ntests; j++) {
			int ans1 = query(r1, j, j);
			int ans2 = st.query(r1, j, j);
			if (ans1 != ans2) {
				printf("failed on test %d, element %d ans1=%d ans2=%d\n", test, j, ans1, ans2);
				return false;
			}
		}
		//test size
		if (a[r1].size() != st.size(r1)) {
			printf("failed on test %d, size1=%d size2=%d\n", test, a[r1].size(), st.size(r1));
			return false;
		}
		//test element by index
		for(int i = 0; i < (int)a[r1].size(); i++) {
			int ans1 = a[r1][i];
			int ans2 = st.at(r1, i);
			if (ans1 != ans2) {
				print(r1);
				printf("failed on test %d, pos %d ans1=%d ans2=%d\n", test, i, ans1, ans2);
				return false;
			}
		}
		//test random queries
		for(int j = 0; j <= ntests; j++) {
			int a = rand()%(ntests+1);
			int b = rand()%(ntests+1);
			int ans1 = query(r1, a, b);
			int ans2 = st.query(r1, a, b);
			if (ans1 != ans2) {
				printf("failed on test %d, segment [%d,%d] ans1=%d ans2=%d\n", test, a, b, ans1, ans2);
				return false;
			}
		}
	}
	printf("all tests passed\n");
	return true;
}

int main() {
	test(1000);
	return 0;
}*/

/*
 * Szkoput Tree Rotations
 */

#include <cstdio>
typedef long long ll;
typedef pair<ll, int> ii;
typedef vector<int> vi;

SegmentTree st;

int n;
void dfs(int & t, ll & ans) {
	int a;
	scanf("%d", &a);
	if (a == 0) {
		int t1, t2;
		dfs(t1, ans);
		dfs(t2, ans);
		ll sum1 = 0, sum2 = 0;
		vi arr;
		if (st.size(t1) > st.size(t2)) {
			t = t1;
			st.to_vector(t2, arr);
			for (int b : arr) {
				sum1 += st.query(t1, 0, b-1);
				sum2 += st.query(t1, b+1, n);
			}
			st.unionSet(t1, t2);
		}
		else {
			t = t2;
			st.to_vector(t1, arr);
			for (int b : arr) {
				sum1 += st.query(t2, 0, b-1);
				sum2 += st.query(t2, b+1, n);
			}
			st.unionSet(t2, t1);
		}
		ans += min(sum1, sum2);
	}
	else {
		t = st.newSet(a);
	}
}

int main() {
	scanf("%d", &n);
	st = SegmentTree(n);
	ll ans = 0;
	int t;
	dfs(t, ans);
	printf("%lld\n", ans);
	return 0;
}
\end{vncode}
\endgroup


\subsection{Segment Tree Beats}

Cập nhật \texttt{a[i]=min(a[i],k)} trên đoạn và truy vấn tổng.

\codepath{Trees/segmenttreebeats.cpp}

\begingroup\scriptsize
\begin{vncode}
#include <vector>
#include <map>
using namespace std;
#define INF 0x3f3f3f3f

/*
 * Segment Tree Beats
 */

class SegmentTreeBeats {
	vector<int> st, lazy;
	vector<int> val[2], cnt[2];
	int size;
#define left(p) (p << 1)
#define right(p) ((p << 1) + 1)
	void update(int u) {
		map<int, int> aux;
		for(int k = 0; k < 2; k++) {
			aux[val[k][left(u)]] += cnt[k][left(u)];
			aux[val[k][right(u)]] += cnt[k][right(u)];
		}
		for(int k = 0; k < 2; k++) {
			val[k][u] = aux.rbegin()->first;
			cnt[k][u] = aux.rbegin()->second;
			aux.erase(aux.rbegin()->first);
		}
		st[u] = st[left(u)] + st[right(u)];
	}
	void push(int u, int l, int r) {
		if (lazy[u] < val[0][u]) {
			st[u] -= cnt[0][u]*val[0][u];
			val[0][u] = lazy[u];
			st[u] += cnt[0][u]*val[0][u];
		}
		if (l < r) {
			lazy[left(u)] = min(lazy[left(u)], lazy[u]);
			lazy[right(u)] = min(lazy[right(u)], lazy[u]);
		}
	}
	void build(int u, int l, int r, int A[]) {
		if (l == r) {
			val[0][u] = A ? A[r] : INF; cnt[0][u] = 1;
			val[1][u] = -INF; cnt[1][u] = 0;
			st[u] = A[r]; return;
		}
		int m = (l + r) / 2;
		build(left(u), l, m, A);
		build(right(u), m+1, r, A);
		update(u);
	}
	void update(int u, int l, int r, int a, int b, int k) {
		push(u, l, r);
		if (b < l || r < a || k >= val[0][u]) return;
		if (a <= l && r <= b && val[1][u] < k) {
			lazy[u] = k;
			push(u, l, r); return;
		}
		int m = (l + r) / 2;
		update(left(u), l, m, a, b, k);
		update(right(u), m+1, r, a, b, k);
		update(u);
	}
	int query(int u, int l, int r, int a, int b) {
		push(u, l, r);
		if (b < l || r < a) return 0;
		if (a <= l && r <= b) return st[u];
		int m = (l + r) / 2;
		int p1 = query(left(u), l, m, a, b);
		int p2 = query(right(u), m+1, r, a, b);
		return p1 + p2;
	}
public:
	SegmentTreeBeats(int *bg, int *en) {
		size = int(en - bg);
		st.resize(4 * size);
		for(int k = 0; k < 2; k++) {
			val[k].resize(4 * size);
			cnt[k].resize(4 * size);
		}
		lazy.assign(4 * size, INF);
		build(1, 0, size-1, bg);
	}
	void update(int a, int b, int k) {
		update(1, 0, size-1, a, b, k);
	}
	int query(int a, int b) {
		return query(1, 0, size-1, a, b);
	}
};

/*
 * TEST MATRIX
 */

#include <cstdio>
#include <ctime>
#include <cstdlib>
#define MAXN 1000000
#define RANGE 1000
#define NQUERY 10

int arr[MAXN];

void update(int a, int b, int k) {
	for(int i = a; i <= b; i++) arr[i] = min(arr[i], k);
}

int query(int a, int b) {
	int sum = 0;
	for(int i = a; i <= b; i++) sum += arr[i];
	return sum;
}

bool test(int suits, int n) {
	srand(time(NULL));
	clock_t t = clock();
	for(int suit = 0; suit < suits; suit++) {
		for(int i = 0; i < n; i++) arr[i] = 2*RANGE;
		SegmentTreeBeats stb(arr, arr+n);
		for(int test = 0; test < n; test++) {
			int i = rand()%n;
			int j = rand()%n;
			int k = (rand()%100) + n-test;
			if (i > j) swap(i, j);
			update(i, j, k);
			stb.update(i, j, k);
			for(int q = 0; q < NQUERY; q++) {
				i = rand()%n;
				j = rand()%n;
				if (i > j) swap(i, j);
				int v2 = stb.query(i, j);
				//int v1 = v2;
				int v1 = query(i, j);
				if (v1 != v2) {
					printf("failed on test %d, query [%d,%d] => %d,%d\n", test, i, j, v1, v2);
					return false;
				}
			}
		}
	}
	t = clock() - t;
	printf("all tests passed, time %.3f\n", t*1.0/CLOCKS_PER_SEC);
	return true;
}

int main() {
	test(10, 10000);
	return 0;
}
\end{vncode}
\endgroup


\subsection{Merge Sort Tree}

Cây của merge sort; đếm phần tử trong đoạn \texttt{[i,j]} có giá trị thuộc \texttt{[a,b]}.

\codepath{Trees/mergesorttree.cpp}

\begingroup\scriptsize
\begin{vncode}
#include <vector>
#include <algorithm>
#define INF (1<<30)
using namespace std;

/*
 * Merge Sort Tree
 */

class MergeSortTree {
private:
	vector< vector<int> > st;
	int size;
	#define left(p) (p << 1)
	#define right(p) ((p << 1) + 1)
	void build(int p, int l, int r, int* A) { // O(n)
		st[p].resize(r-l+1);
		if (l == r) { st[p][0] = A[l]; return; }
		int pl = left(p), pr = right(p), m = (l+r)/2;
		build(pl, l, m, A);
		build(pr, m+1, r, A);
		merge(st[pl].begin(), st[pl].end(),
			st[pr].begin(), st[pr].end(),
			st[p].begin());
	}
	int query(int p, int l, int r, int i, int j, int a, int b) {
		if (j < l || i > r) return 0;
		if (i <= l && j >= r)
			return upper_bound(st[p].begin(), st[p].end(), b) -
				lower_bound(st[p].begin(), st[p].end(), a);
		int m = (l + r) / 2;
  		return query(left(p), l, m, i, j, a, b) +
			query(2*p+1, m+1, r, i, j, a, b);
	}
public:
	MergeSortTree(int* begin, int* end) {
		size = (int)(end-begin);
		st.assign(4*size, vector<int>());
		build(1, 0, size-1, begin);
	}
	int query(int i, int j, int a, int b) {
		return query(1, 0, size-1, i, j, a, b);
	}
};

/*
 * TEST MATRIX
 */

#include <cstdio>


int vet[100009], aux[100009];
int query(int i, int j, int a, int b) {
	int ans = 0;
	for(int k=i; k<=j; k++) if (a <= vet[k] && vet[k] <= b) ans++;
	return ans;
}

bool test(int N) {
	for(int i = 0; i < N; i++) vet[i] = rand()%N;
	MergeSortTree st(vet, vet+N);
	for(int q=0, a, b, i, j, n; q<N; q++) {
		a = rand()%N;
		b = rand()%N;
		if (a>b) swap(a, b);
		i = rand()%N;
		j = rand()%N;
		if (i>j) swap(i, j);

		if (query(i, j, a, b) != st.query(i, j, a, b)) {
			printf("test %d failed\n", q+1);
			return false;
		}
	}
	printf("all tests passed\n");
	return true;
}

int main() {
	test(10000);
	return 0;
}
\end{vncode}
\endgroup


\subsection{Fractional Cascading Merge Sort Tree}

Merge Sort Tree với fractional cascading để giảm truy vấn xuống O(log n).

\codepath{Trees/fractionalcascade.cpp}

\begingroup\scriptsize
\begin{vncode}
#include <vector>
#include <algorithm>
#define INF (1<<30)
using namespace std;

/*
 * Merge Sort Tree
 */

class MergeSortTree {
private:
	vector< vector<int> > st, li, ri;
	int size;
	#define left(p) (p << 1)
	#define right(p) ((p << 1) + 1)
	void build(int p, int l, int r, int* A) { // O(n)
		st[p].resize(r-l+1);
		if (l == r) { st[p][0] = A[l]; return; }
		int pl = left(p), pr = right(p), m = (l+r)/2;
		vector<int> &vl = st[pl], &vr = st[pr];
		build(pl, l, m, A);
		build(pr, m+1, r, A);
		merge(st[pl].begin(), st[pl].end(),
			st[pr].begin(), st[pr].end(),
			st[p].begin());
		for(int k = 0, i = 0, j = 0; k < r-l+1; k++) {
			while(i < vl.size() && st[p][k] > vl[i]) i++;
			while(j < vr.size() && st[p][k] > vr[j]) j++;
			li[p].push_back(i), ri[p].push_back(j);
		}
		li[p].push_back(vl.size());
		ri[p].push_back(vr.size());
	}
	int query(int p, int l, int r, int i, int j, int a, int b) {
		if (j < l || i > r) return 0;
		if (i <= l && j >= r) return b-a;
		int m = (l + r) / 2;
  		return query(left(p), l, m, i, j, li[p][a], li[p][b]) +
			query(right(p), m+1, r, i, j, ri[p][a], ri[p][b]);
	}
public:
	MergeSortTree(int* begin, int* end) {
		size = (int)(end-begin);
		st.assign(4*size, vector<int>());
		li.assign(4*size, vector<int>());
		ri.assign(4*size, vector<int>());
		build(1, 0, size-1, begin);
	}
	int query(int i, int j, int a, int b) {
		a = lower_bound(st[1].begin(), st[1].end(), a) - st[1].begin();
		b = upper_bound(st[1].begin(), st[1].end(), b) - st[1].begin();
		return query(1, 0, size-1, i, j, a, b);
	}
};

/*
 * TEST MATRIX
 */

#include <cstdio>


int vet[100009], aux[100009];
int query(int i, int j, int a, int b) {
	int ans = 0;
	for(int k=i; k<=j; k++) if (a <= vet[k] && vet[k] <= b) ans++;
	return ans;
}

bool test(int N) {
	for(int i = 0; i < N; i++) vet[i] = rand()%N;
	MergeSortTree st(vet, vet+N);
	for(int q=0, a, b, i, j, n; q<N; q++) {
		a = rand()%N;
		b = rand()%N;
		if (a>b) swap(a, b);
		i = rand()%N;
		j = rand()%N;
		if (i>j) swap(i, j);

		if (query(i, j, a, b) != st.query(i, j, a, b)) {
			printf("test %d failed [i,j]=[%d,%d] [a,b]=[%d,%d], control %d mst %d\n",
				q+1, i, j, a, b, query(i, j, a, b), st.query(i, j, a, b));
			return false;
		}
	}
	printf("all tests passed\n");
	return true;
}

int main() {
	test(10000);
	return 0;
}
\end{vncode}
\endgroup


\subsection{Sparse Table}

Trả lời RMQ O(1), tiền xử lý O(n log n).

\codepath{Data Structures/sparsetable.cpp}

\begingroup\scriptsize
\begin{vncode}
#include <cstdio>
#include <cmath>
#define MAXN 100009
#define MAXLOGN 20

/*
 * Sparse Table
 */

#define comp(a, b) min((a),(b))

class SparseTable {
	int st[MAXN][MAXLOGN];
	int sz;
public:
	SparseTable(int* bg, int* en) {
		sz = int(en - bg);
		for(int i=0; i<sz; i++) st[i][0] = bg[i];
		for(int j = 1; 1 << j <= sz; j++)
			for(int i=0; i + (1<<j) <= sz; i++)
				st[i][j] = comp(st[i][j-1], st[i+(1<<(j-1))][j-1]);
	}
	int query(int l, int r) {
		int k = (int)floor(log((double)r-l+1) / log(2.0)); // 2^k <= (j-i+1)
		return comp(st[l][k], st[r-(1<<k)+1][k]);
	}
};
\end{vncode}
\endgroup


\subsection{Pre-Suf Table}

Trả lời RSQ không cần phép nghịch đảo trong O(1), tiền xử lý O(n log n).

\codepath{Data Structures/presuftable.cpp}

\begingroup\scriptsize
\begin{vncode}
#include <vector>
using namespace std;
#define MAXN 100009

#include <cstdio>

/*
 * Pre-Suf Table
 * Friend of the Sparse Table
 */

const int neutral = 0;
#define comp(a, b) ((a)+(b))

int top[2*MAXN];

class PresufTable {
	vector< vector<int> > f;
public:
	PresufTable(int* st, int* en) {
		int bit = 0, sz = int(en - st), n, k;
		for(n = 1, k = 0; n < sz; n <<= 1, k++);
		f.assign(k, vector<int>(n));
		for(int i = 0; i < n; i++) {
			while((1<<bit) <= i) bit++;
			top[i] = bit;
		}
		for(int i = 0; i < n; i++)
			f[0][i] = i < sz ? st[i] : neutral;
		for(int j = 0, len = 1; j <= k; j++, len <<= 1)
			for(int s = len; s < n; s += (len<<1)) {
				f[j][s] = st[s]; f[j][s-1] = st[s-1];
				for(int i = 1; i < len; i++) {
					f[j][s+i] = comp(f[j][s+i-1], st[s+i]);
					f[j][s-1-i] = comp(f[j][s-i], st[s-1-i]);
				}
			}
	}
	int query(int l, int r) {
		if (l == r) return f[0][r];
		int i = top[l^r] - 1;
		return comp(f[i][r], f[i][l]);
	}
};

/*
 * TEST MATRIX
 */

#include <cstdlib>
#define RANGE 10000

int arr[MAXN];

int query(int l, int r) {
	int ans = 0;
	for(int i = l; i <= r; i++) {
		ans += arr[i];
	}
	return ans;
}

bool test(int n, int ntests) {
	for(int i = 0; i < n; i++) arr[i] = /*rand()%RANGE*/ i;
	PresufTable pst(arr, arr+n);
	for(int test = 1; test <= ntests; test++) {
		int l = rand()%n;
		int r = rand()%n;
		if (l > r) swap(l, r);
		if (query(l, r) != pst.query(l, r)) {
			printf("failed test %d\n", test);
			return false;
		}
	}
	printf("all tests passed\n");
	return true;
}

int main() {
	test(10000, 10000);
}
\end{vncode}
\endgroup


\subsection{AVL Tree}

Cây AVL với thao tác BST.

\codepath{Trees/avl.cpp}

\begingroup\scriptsize
\begin{vncode}
#include <cstdio>
struct node {
    int key, height, size;
    node *left, *right;
    node(int k) {
        key = k; left = right = 0; 
        height = size = 1;
    }
};

class AVLtree {
private:
    node* root;
    int size_;
    int height(node* p) { return p ? p->height : 0; }
    int size(node* p) { return p ? p->size : 0; }
    int bfactor(node* p) {
        return height(p->right)-height(p->left);
    }
    void fixheight(node* p) {
        int hl = height(p->left);
        int hr = height(p->right);
        p->height = (hl>hr?hl:hr)+1;
        p->size = 1 + size(p->left) + size(p->right);
    }
    node* rotateright(node* p) {
        node* q = p->left;
        p->left = q->right;
        q->right = p;
        fixheight(p);
        fixheight(q);
        return q;
    }
    node* rotateleft(node* q) {
        node* p = q->right;
        q->right = p->left;
        p->left = q;
        fixheight(q);
        fixheight(p);
        return p;
    }
    node* balance(node* p) {
        fixheight(p);
        if(bfactor(p)==2) {
            if(bfactor(p->right)<0)
                p->right = rotateright(p->right);
            return rotateleft(p);
        }
        if(bfactor(p)==-2) {
            if(bfactor(p->left)>0)
                p->left = rotateleft(p->left);
            return rotateright(p);
        }
        return p;
    }
    node* build(node* p, int k) {
        if(!p) return new node(k);
        if (p->key == k) return p;
        else if(k<p->key) p->left = build(p->left, k);
        else p->right = build(p->right, k);
        return balance(p);
    }
    node* findmin(node* p) {
        return p->left?findmin(p->left):p;
    }
    node* removemin(node* p) {
        if(p->left==0) return p->right;
        p->left = removemin(p->left);
        return balance(p);
    }
    node* remove(node* p, int k) {
        if(!p) return 0;
        if(k < p->key) p->left = remove(p->left,k);
        else if(k > p->key) p->right = remove(p->right,k);  
        else {
            node* l = p->left;
            node* r = p->right;
            delete p;
            if(!r) return l;
            node* min = findmin(r);
            min->right = removemin(r);
            min->left = l;
            return balance(min);
        }
        return balance(p);
    }
    bool find(node* p, int k) {
        if (!p) return false;
        if (p->key == k) return true;
        else if(k<p->key) return find(p->left, k);
        else return find(p->right, k);
    }
    void del(node* p) {
        if (!p) return;
        del(p->left);
        del(p->right);
        delete p;
    }
    node* nth(node* p, int n) {
        if (!p) return p;
        if (size(p->left)+1 > n) return nth(p->left, n);
        if (size(p->left)+1 < n) return nth(p->right, n - size(p->left) - 1);
        else return p;
    }

public:
    AVLtree() { root = 0; size_ = 0; }
    ~AVLtree() { del(root); }
    bool empty() { return size_ == 0; }
    int size() { return size_; }
    void clear() {
        size_ = 0;
        del(root);
        root = 0;
    }
    void insert(int key) {
        size_++;
        root = build(root, key);
    }
    void erase(int key) {
        size_--;
        root = remove(root, key);
    }
    bool count(int key) {
        return find(root, key);
    }
    int nth_element(int n) {
        node* p = nth(root, n);
        if (p) return p->key;
        else return -1;
    }   //1-indexed
};

/*
 * TEST MATRIX
 */

#include <set>
#include <cstdlib>
using namespace std;

bool test()
{
	AVLtree tree;
	set<int> s;
	int N = 1000000;
	for(int t=1; t<=N; t++) {
		int x = rand()%100;
		if (s.count(x) != tree.count(x)) {
			printf("failed on test %d\n", t);
			return false;
		}
		if (s.count(x)) {
			s.erase(x);
			tree.erase(x);
		}
		else{
			s.insert(x);
			tree.insert(x);
		}
	}
	return true;
}

int main()
{
	if (test()) {
		printf("all tests passed\n");
	}
	return 0;
}
\end{vncode}
\endgroup


\subsection{Treap / Cartesian Tree}

Treap hỗ trợ \texttt{insert}, \texttt{count}, \texttt{erase}, \texttt{nth\_element}; \texttt{split} tách giữa k-1 và k.

\codepath{Trees/treap.cpp}

\begingroup\scriptsize
\begin{vncode}
#include <cstdio>
#include <set>
#include <algorithm>
using namespace std;

/*
 * Implicit Treap
 * 0-indexed, 10x slower than red-black tree
 * split separates between k-1 and k
 */

struct node {
	int x, y, size;
	node *l, *r;
	node(int _x) : x(_x), y(rand()), size(1), l(NULL), r(NULL) {}	
};

class Treap {
private:
	node* root;
	int size(node* t) { return t ? t->size : 0; }
	node* refresh(node* t) {
		if (!t) return t;
		t->size = 1 + size(t->l) + size(t->r);
		return t;
	}
	void split(node* &t, int k, node* &a, node* &b) {
		node* aux;
		if (!t) a = b = NULL;
		else if (t->x < k) {
			split(t->r, k, aux, b);
			t->r = aux;
			a = refresh(t);
		}
		else {
			split(t->l, k, a, aux);
			t->l = aux;
			b = refresh(t);
		}
	}
	node* merge(node* a, node* b) {
		if (!a || !b) return a ? a : b;
		if (a->y < b->y) {
			a->r = merge(a->r, b);
			return refresh(a);
		}
		else {
			b->l = merge(a, b->l);
			return refresh(b);
		}
	}
	node* count(node* t, int k) {
		if (!t) return NULL;
		else if (k < t->x) return count(t->l, k);
		else if (k == t->x) return t;
		else return count(t->r, k);
	}
	node* nth(node* t, int n) {
		if (!t) return NULL;
		if (n <= size(t->l)) return nth(t->l, n);
		else if (n == size(t->l) + 1) return t;
		else return nth(t->r, n-size(t->l)-1);
	}
	void del(node* &t) {
		if (!t) return;
		if (t->l) del(t->l);
		if (t->r) del(t->r);
		delete t;
		t = NULL;
	}
public:
	Treap() : root(NULL) { }
	~Treap() { clear(); }
	void clear() { del(root); }
	int size() { return size(root); }
	bool count(int k) { return count(root, k) != NULL; }
	bool insert(int k) {
		if (count(k)) return false;
		node *a, *b;
		split(root, k, a, b);
		root = merge(merge(a, new node(k)), b);
		return true;
	}
	bool erase(int k) {
		node * f = count(root, k);
		if (!f) return false;
		node *a, *b, *c, *d;
		split(root, k, a, b);
		split(b, k+1, c, d);
		root = merge(a, d);
		delete f;
		return true;
	}
	int nth(int n) {
		node* ans = nth(root, n);
		return ans ? ans->x : -1;
	}
};

/*
 * TEST MATRIX
 */

int vet[10000009];

void test() {
	set<int> s;
	Treap t;
	int N = 1000000;
	for(int i=0; i<N; i++) {
		int n = rand()%1000;
		if(!s.count(n)) {
			s.insert(n);
			t.insert(n);
			//if(!t.insert(n)) printf("error inserting %d in treap!\n", n);
			//printf("inserted %d\n", n);
		}
		else{
			s.erase(n);
			t.erase(n);
			//if(!t.erase(n)) printf("error erasing %d in treap!\n", n);
			//printf("erased %d\n", n);
		}
		n = rand()%1000;
		if (s.count(n) != t.count(n)) {
			printf("failed test %d, s.count(%d) = %d, t.count(%d) = %d\n", i, n, s.count(n), n, t.count(n));
		}
	}
	s.clear();
	t.clear();
	for(int i=0; i<N; i++) {
		vet[i] = i+1;
	}
	random_shuffle(vet, vet+N);
	for(int i=0; i<N; i++) {
		t.insert(vet[i]);
	}
	for(int i=1; i<=N; i++) {
		if (t.nth(i) != i) {
			printf("failed test %d\n", i);
		}
	}
	printf("all tests passed\n");
}

int main() {
	test();
	return 0;
}
\end{vncode}
\endgroup


\subsection{Treap / Cartesian Tree ngầm định}

Treap theo vị trí cho chuỗi/mảng động.

\codepath{Trees/implicittreap.cpp}

\begingroup\scriptsize
\begin{vncode}
#include <cstdio>
#include <vector>
#include <algorithm>
#include <ctime>
#define INF 0x3f3f3f3f
using namespace std;

/*
 * Implicit Treap
 * 0-indexed, 10x slower than segment tree
 * split separates in trees of size k and n-k
 */

struct node {
	int y, v, sum, size;
	bool rev;
	node *l, *r;
	node(int _v) : v(_v), sum(_v), y(rand()),
		size(1), l(NULL), r(NULL), rev(false) {}	
};

class ImplicitTreap {
private:
	node* root;
	int size(node* t) { return t ? t->size : 0; }
	int sum(node* t) { return t ? t->sum : 0; }
	node* refresh(node* t) {
		if (t == NULL) return t;
		t->size = 1 + size(t->l) + size(t->r);
		t->sum = t->v + sum(t->l) + sum(t->r);
		if (t->l != NULL) t->l->rev ^= t->rev;
		if (t->r != NULL) t->r->rev ^= t->rev;
		if (t->rev) {
			swap(t->l, t->r);
			t->rev = false;
		}
		return t;
	}
	void split(node* &t, int k, node* &a, node* &b) {
		refresh(t);
		node * aux;
		if (!t) a = b = NULL;
		else if (size(t->l) < k) {
			split(t->r, k-size(t->l)-1, aux, b);
			t->r = aux;
			a = refresh(t);
		}
		else {
			split(t->l, k, a, aux);
			t->l = aux;
			b = refresh(t);
		}
	}
	node* merge(node* a, node* b) {
		refresh(a); refresh(b);
		node* aux;
		if (!a || !b) return a ? a : b;
		if (a->y < b->y) {
			a->r = merge(a->r, b);
			return refresh(a);
		}
		else {
			b->l = merge(a, b->l);
			return refresh(b);
		}
	}
	node* at(node* t, int n) {
		if (!t) return t;
		refresh(t);
		if (n < size(t->l)) return at(t->l, n);
		else if (n == size(t->l)) return t;
		else return at(t->r, n-size(t->l)-1);
	}
	void del(node* &t) {
		if (!t) return;
		if (t->l) del(t->l);
		if (t->r) del(t->r);
		delete t;
		t = NULL;
	}
public:
	ImplicitTreap() : root(NULL) { }
	~ImplicitTreap() { clear(); }
	void clear() { del(root); }
	int size() { return size(root); }
	bool insertAt(int n, int v) {
		node *a, *b;
		split(root, n, a, b);
		root = merge(merge(a, new node(v)), b);
		return true;
	}
	bool erase(int n) {
		node *a, *b, *c, *d;
		split(root, n, a, b);
		split(b, 1, c, d);
		root = merge(a, d);
		if (c == NULL) return false;
		delete c;
		return true;
	}
	int at(int n) {
		node* ans = at(root, n);
		return ans ? ans->v : -1;
	}
	int query(int l, int r) {
		if (l > r) swap(l, r);
		node *a, *b, *c, *d;
		split(root, l, a, d);
		split(d, r-l+1, b, c);
		int ans = sum(b);
		root = merge(a, merge(b, c));
		return ans;
	}
	void reverse(int l, int r) {
		if (l>r) swap(l, r);
		node *a, *b, *c, *d;
		split(root, l, a, d);
		split(d, r-l+1, b, c);
		if (b != NULL) b->rev ^= 1;
		root = merge(a, merge(b, c));
	}
};

/*
 * TEST MATRIX
 */

bool test() {
	srand(time(NULL));
	vector<int> v;
	ImplicitTreap t;
	int N = 10000;
	vector<int>::iterator it;
	bool toprint = false;
	for(int i=0, n, k, l, r; i<N; i++) {
		if (i%5 == 0 && i > 0) {
			n = rand()%((int)v.size());
			if (toprint) printf("deleting v[%d] = %d\n", n, v[n]);
			it = v.begin()+n;
			v.erase(it);
			t.erase(n);
		}
		else if (i%5 == 4) {
			l = rand()%((int)v.size());
			r = rand()%((int)v.size());
			if (l>r) swap(l, r);
			if (toprint) printf("reversing %d to %d\n", l, r);
			for(int j=l; j<=r && j<=r-j+l; j++) {
				swap(v[j], v[r-j+l]);
			}
			t.reverse(l, r);
		}
		else{
			n = rand()%((int)v.size()+1);
			k = rand()%1000;
			if (toprint) printf("inserting %d in pos %d\n", k, n);
			it = v.begin()+n;
			v.insert(it, k);
			t.insertAt(n, k);
		}
		if (toprint) printf("array: ");
		for(int j=0; j<(int)v.size(); j++) {
			if (toprint) printf("%d ", v[j]);
			if (v[j] != t.at(j)) {
				printf("test %d failed, v[%d] = %d, t.at(%d) = %d\n", i+1, j, v[j], j, t.at(j));
				return false;
			}
		}
		if (toprint) printf("\n");
		l = rand()%((int)v.size());
		r = rand()%((int)v.size());
		if (l>r) swap(l, r);
		int ans = 0;
		for(int j=l; j<=r; j++) {
			ans += v[j];
		}
		if (toprint) printf("sum(%d, %d) = %d = %d\n", l, r, ans, t.query(l, r));
		if (ans != t.query(l, r)) {
			printf("test %d failed, ans(%d, %d) = %d = %d\n", i, l, r, ans, t.query(l, r));
			return false;
		}
	}
	return true;
}

int main() {
	if(test()) printf("all tests passed\n");
	return 0;
}
\end{vncode}
\endgroup


\subsection{Persistent Randomized BST}

BST ngẫu nhiên persistent.

\codepath{Trees/rbst.cpp}

\begingroup\scriptsize
\begin{vncode}
#include <cstdio>
#include <vector>
#include <algorithm>
#include <ctime>
#define INF 0x3f3f3f3f
using namespace std;

/*
 * Randomized BST, works as Persistent Implicit Treap
 * split separates in trees of size k and n-k
 */

struct node {
	int v, sum, size;
	node *l, *r;
	node(int _v, vector<node*> &nodes) :
		v(_v), sum(_v), size(1), l(NULL), r(NULL) {
		nodes.push_back(this);
	}	
};

class RBST {
private:
	vector<node*> root, nodes;
	int size(node* t) { return t ? t->size : 0; }
	int sum(node* t) { return t ? t->sum : 0; }
	node* copy(node* t) {
		node* cpy = new node(t->v, nodes);
		cpy->l = t->l; cpy->r = t->r;
		return refresh(cpy);
	}
	node* refresh(node* t) {
		if (t == NULL) return t;
		t->size = 1 + size(t->l) + size(t->r);
		t->sum = t->v + sum(t->l) + sum(t->r);
		return t;
	}
	void split(node* t, int k, node* &a, node* &b) {
		node *aux;
		if (!t) { a = b = NULL; return; }
		t = copy(t);
		if (size(t->l) < k) {
			split(t->r, k-size(t->l)-1, aux, b);
			t->r = aux;
			a = refresh(t);
		}
		else {
			split(t->l, k, a, aux);
			t->l = aux;
			b = refresh(t);
		}
	}
	node* merge(node* a, node* b) {
		node *aux, *cpy;
		if (!a || !b) return a ? a : b;
		int m = size(a), n = size(b);
		if (rand()%(n+m) < m) {
			a = copy(a);
			a->r = merge(a->r, b);
			return refresh(a);
		}
		else {
			b = copy(b);
			b->l = merge(a, b->l);
			return refresh(b);
		}
	}
	node* at(node* t, int i) {
		if (!t) return t;
		if (i < size(t->l)) return at(t->l, i);
		else if (i == size(t->l)) return t;
		else return at(t->r, i-size(t->l)-1);
	}
public:
	RBST() { clear(); }
	~RBST() { clear(); }
	void clear() {
		for(int i = 0; i < (int)nodes.size(); i++)
			delete nodes[i];
		root.push_back(NULL);
	}
	int size(int ver) { return size(root[ver]); }
	int insert(int ver, int i, int v) {
		node *a, *b, *tree;
		split(root[ver], i, a, b);
		tree = merge(merge(a, new node(v, nodes)), b);
		root.push_back(tree);
		return root.size() - 1;
	}
	int erase(int ver, int i) {
		node *a, *b, *c, *d;
		split(root[ver], i, a, b);
		split(b, 1, c, d);
		root.push_back(merge(a, d));
		return root.size() - 1;
	}
	int at(int ver, int i) {
		node* ans = at(root[ver], i);
		return ans ? ans->v : -1;
	}
	int query(int ver, int l, int r) {
		if (l > r) swap(l, r);
		node *a, *b, *c, *d;
		split(root[ver], l, a, d);
		split(d, r-l+1, b, c);
		return sum(b);
	}
	int split(int ver, int k) {
		node *a, *b;
		split(root[ver], k, a, b);
		root.push_back(a);
		root.push_back(b);
		return root.size()-2;
	}
	int merge(int ver1, int ver2) {
		root.push_back(merge(root[ver1], root[ver2]));
		return root.size()-1;
	}
};

/*
 * TEST MATRIX
 */

#include <cassert>
#define MAXSIZE 100000
#define NTESTS 10000
#define NQUERY 10
#define RANGE 10000

void insert(vector<int> & arr, int i, int v) {
	arr.insert(arr.begin()+i, v);
}

int queryt(vector<int> & arr, int i, int j) {
	if (i > j) swap(i, j);
	int sum = 0;
	for(int it=i; it<=j && it<(int)arr.size(); it++) {
		sum += arr[it];
	}
	return sum;
}

void erase(vector<int> & arr, int i) {
	arr.erase(arr.begin()+i);
}

bool test() {
	srand(time(NULL));
	vector<int> control[NTESTS];
	RBST treap;
	for(int test = 1; test < NTESTS; test++) {
		int op = rand()%2;
		int ver = rand()%test;
		control[test] = control[ver];
		if (test%2 == 1) {
			int ver2 = rand()%test;
			if (control[ver].size() + control[ver2].size() > MAXSIZE) {
				test--;
				continue;
			}
			control[test].insert(control[test].end(), control[ver2].begin(), control[ver2].end());
			treap.merge(ver, ver2);
		}
		else if (op == 0 && !control[test].empty()) {
			int i = rand()%int(control[test].size());
			erase(control[test], i);
			treap.erase(ver, i);
		}
		else {
			int data = rand()%RANGE;
			int i = rand()%int(control[test].size()+1);
			insert(control[test], i, data);
			treap.insert(ver, i, data);
		}
	}
	for(int test = 0; test < NTESTS; test++) {
		if (control[test].size() != treap.size(test)) {
			printf("failed on test %d, size: %u-%d\n", test, control[test].size(), treap.size(test));
			return false;
		}
		for(int i = 0; i < (int)control[test].size(); i++) {
			if (control[test][i] != treap.at(test, i)) {
				printf("failed on test %d, id %d: %d-%d\n", test, i, control[test][i], treap.at(test, i));
				return false;
			}
		}
		for(int query = 0; query < NQUERY; query++) {
			if (control[test].empty()) continue;
			int i = rand()%int(control[test].size());
			int j = rand()%int(control[test].size());
			if (i > j) swap(i, j);
			int sum1 = queryt(control[test], i, j);
			int sum2 = treap.query(test, i, j);
			if (sum1 != sum2) {
				printf("failed on test %d, query %d: %d-%d\n", test, query, sum1, sum2);
				return false;
			}
		}
	}
	printf("all tests passed\n");
	return true;
}

int main() {
	test();
	return 0;
}
\end{vncode}
\endgroup


\subsection{Deque đệ quy thuần hàm}

Deque persistent thuần hàm.

\codepath{Data Structures/recursivedeque.cpp}

\begingroup\scriptsize
\begin{vncode}
#include <cstdio>
#include <vector>
using namespace std;

/*
 * Recursive Purely Functional Deque
 */

struct cont {
	cont *ls, *rs;
	int x;
	cont(int _x, vector<cont*> &gc) :
		x(_x), ls(NULL), rs(NULL) {
		gc.push_back(this);
	}
	cont(cont* l, cont* r, vector<cont*> &gc) :
		x(0), ls(l), rs(r) {
		gc.push_back(this);
	}
};

struct node {
	cont* s[2];
	node* c;
	int sz;
	node(vector<node*> &gc) : c(NULL), sz(0) {
		s[0] = s[1] = NULL;
		gc.push_back(this);
	}
	node* copy(vector<node*> &gc) {
		node* cpy = new node(gc);
		cpy->sz = sz; cpy->c = c;
		cpy->s[0] = s[0]; cpy->s[1] = s[1];
		return cpy;
	}
};

class Deque {
private:
	vector<node*> root, nodes;
	vector<cont*> conts;
	void push(node* &u, int k, cont* x) {
		if (!u) { u = new node(nodes); }
		else u = u->copy(nodes);
		u->sz++;
		if (!u->s[k]) u->s[k] = x;
		else {
			cont* y = u->s[k];
			u->s[k] = NULL;
			if (k) swap(x, y);
			push(u->c, k, new cont(x, y, conts));
		}
	}
	cont* pop(node* &u, int k) {
		u = u->copy(nodes);
		u->sz--;
		cont *x = NULL, *y, *z;
		if (u->s[k]) {
			x = u->s[k];
			u->s[k] = NULL;	
		}
		else if (u->c && u->c->sz) {
			z = pop(u->c, k);
			x = z->ls; y = z->rs;
			if (k) swap(x, y);
			u->s[k] = y;
		}
		else if (u->s[k^1]) {
			x = u->s[k^1];
			u->s[k^1] = NULL;
		}
		else u->sz++;
		return x;
	}
	cont* at(const node* u, int i) {
		if (!u) return NULL;
		if (i == 0 && u->s[0]) return u->s[0];
		if (i == u->sz-1 && u->s[1]) return u->s[1];
		if (u->s[0]) i--;
		cont* x = at(u->c, i>>1);
		return i & 1 ? x->rs : x->ls;
	}
	node* merge(node* u, node* v, cont* c1, cont* c2) {
		node* w = new node(nodes);
		if (!u && !v) {
			w->s[0] = c1; w->s[1] = c2;
			return w;
		}
		if (!u) u = new node(nodes);
		if (!v) v = new node(nodes);
		cont *a[4];
		int k = 0;
		if (u->s[1]) a[k++] = u->s[1];
		if (c1) a[k++] = c1;
		if (c2) a[k++] = c2;
		if (v->s[0]) a[k++] = u->s[0];
		w->s[0] = u->s[0];
		w->s[1] = v->s[1];
		if ((k&1) && v->s[0]) pop(v, 0);
		if ((k&1) && v->sz == 0) w->s[1] = a[--k];
		if (k&1) a[k++] = pop(v, 0);
		c1 = k > 0 ? new cont(a[0], a[1], conts) : NULL;
		c2 = k > 2 ? new cont(a[2], a[3], conts) : NULL;
		w->c = merge(u->c, v->c, c1, c2);
		return w;
	}
public:
	Deque() { root.push_back(new node(nodes)); }
	~Deque() {
		for(int i = 0; i < (int)nodes.size(); i++)
			delete nodes[i];
		for(int i = 0; i < (int)conts.size(); i++)
			delete conts[i];
	}
	int push(int ver, int x, int k) {
		root.push_back(root[ver]);
		push(root.back(), k, new cont(x, conts));
		return root.size()-1;
	}
	int pop(int ver, int k) {
		root.push_back(root[ver]);
		pop(root.back(), k);
		return root.size()-1;
	}
	int merge(int v1, int v2) {
		root.push_back(merge(root[v1], root[v2], NULL, NULL));
		return root.size()-1;
	}
	int push_back(int ver, int x) { return push(ver, x, 1); }
	int push_front(int ver, int x) { return push(ver, x, 0); }
	int pop_back(int ver) { return pop(ver, 1); }
	int pop_front(int ver) { return pop(ver, 0); }
	int size(int ver) { return root[ver]->sz; }
	int at(int ver, int i) { return at(root[ver], i)->x; }
	int front(int ver) { return at(root[ver], 0)->x; }
	int back(int ver) { return at(root[ver], size(ver)-1)->x; }
};

/*
 * TEST MATRIX
 */

#include <ctime>
#include <deque>
#include <cstdlib>
#define RANGE 1000
#define LIMIT 1000

bool test(int ntests) {
	srand(time(NULL));
	vector< deque<int> > control(ntests);
	Deque deq;
	for(int test = 1; test < ntests; test++) {
		int add = rand()%2;
		int front = rand()%2;
		int ver = rand()%test;
		control[test] = control[ver];
		printf("test %d: ", test);
		if (test % 4 == 0 && control[test].size() < LIMIT) {
			int over;
			do {
				over = rand()%test;
			} while(control[test].size() + control[over].size() > LIMIT);
			deq.merge(ver, over);
			control[test].insert(control[test].end(), control[over].begin(), control[over].end());
		}
		else if (add || control[test].empty()) {
			int data = rand()%RANGE;
			//printf("pushing %d to %s of ver %d...", data, front ? "front" : "back", ver);
			if (front) {
				control[test].push_front(data);
				deq.push_front(ver, data);
			}
			else {
				control[test].push_back(data);
				deq.push_back(ver, data);
			}
			//printf("done\n");
		}
		else {
			//printf("popping %s from ver %d...", front ? "front" : "back", ver);
			if (front) {
				control[test].pop_front();
				deq.pop_front(ver);
			}
			else {
				control[test].pop_back();
				deq.pop_back(ver);
			}
			//printf("done\n");
		}
	}
	for(int test = 0; test < ntests; test++) {
		int sz1 = control[test].size();
		int sz2 = deq.size(test);
		if (sz1 != sz2) {
			printf("failed on size test %d (%d,%d)\n", test, sz1, sz2);
			return false;
		}
		if (!control[test].empty()) {
			int data1 = control[test].front();
			int data2 = deq.front(test);
			if (data1 != data2) {
				printf("failed on front test %d: %d-%d\n", test, data1, data2);
				return false;
			}
			data1 = control[test].back();
			data2 = deq.back(test);
			if (data1 != data2) {
				printf("failed on back test %d: %d-%d\n", test, data1, data2);
				return false;
			}
		}
		for(int i = 0; i < (int)control[test].size(); i++) {
			int data1 = control[test][i];
			int data2 = deq.at(test, i);
			if (data1 != data2) {
				printf("failed on at %d test %d: %d-%d\n", i, test, data1, data2);
				return false;
			}
		}
	}
	printf("all tests passed\n");
	return true;
}

int main() {
	test(10000);
	return 0;
}
\end{vncode}
\endgroup


\subsection{Splay Tree}

Splay tree.

\codepath{Trees/splay.cpp}

\begingroup\scriptsize
\begin{vncode}
#include <cstdio>

/*
 * Splay Tree
 */

struct node {
	int key;
	node *ls, *rs;
	node(int k) : key(k), ls(NULL), rs(NULL) {}
};

class SplayTree {
private:
	node* root;
	node* rotateright(node* p) {
		node* q = p->ls;
		p->ls = q->rs;
		q->rs = p;
		return q;
	}
	node* rotateleft(node* q) {
		node* p = q->rs;
		q->rs = p->ls;
		p->ls = q;
		return p;
	}
	node* splay(node* p, int key) {
		if (!p || p->key == key) return p;
		if (p->key > key) {
			if (!p->ls) return p;
			if (p->ls->key > key)	{
				p->ls->ls = splay(p->ls->ls, key);
				p = rotateright(p);
			}
			else if (p->ls->key < key) {
				p->ls->rs = splay(p->ls->rs, key);
				if (p->ls->rs)
					p->ls = rotateleft(p->ls);
			}
			return (!p->ls) ? p : rotateright(p);
		}
		else {
			if (!p->rs) return p;
			if (p->rs->key > key) {
				p->rs->ls = splay(p->rs->ls, key);
				if (p->rs->ls)
					p->rs = rotateright(p->rs);
			}
			else if (p->rs->key < key) {
				p->rs->rs = splay(p->rs->rs, key);
				p = rotateleft(p);
			}
			return (!p->rs) ? p : rotateleft(p);
		}
	}
	void del(node* &p) {
		if (!p) return;
		del(p->ls); del(p->rs);
		delete p;
		p = NULL;
	}

public:
	SplayTree() : root(NULL) { }
	~SplayTree() { del(root); }
	bool empty() { return root == NULL; }
	void clear() { del(root); }
	void insert(int key) {
		if (!root) {
			root = new node(key);
			return;
		}
		node* p = splay(root, key);
		if (p->key == key) return;
		root = new node(key);
		if (p->key > key) {
			root->rs = p;
			root->ls = p->ls;
			p->ls = NULL;
		}
		else {
			root->ls = p;
			root->rs = p->rs;
			p->rs = NULL;
		}
	}
	void erase(int key) {
		node* p = splay(root, key);
		if (!p || p->key != key) return;
		if (!p->rs) {
			root = p->ls;
			delete p;
			return;
		}
		node* q = splay(p->rs, key);
		q->ls = p->ls;
		root = q;
		delete p;
	}
	bool count(int key) {
		if (!root) return false;
		root = splay(root, key);
		return root->key == key;
	}
};

/*
 * TEST MATRIX
 */

#include <set>
#include <cstdlib>
using namespace std;

bool test()
{
	SplayTree splay;
	set<int> s;
	int N = 1000000;
	for(int t=1; t<=N; t++) {
		int x = rand()%100;
		if (s.count(x) != splay.count(x)) {
			printf("failed on test %d\n", t);
			return false;
		}
		if (s.count(x)) {
			s.erase(x);
			splay.erase(x);
		}
		else{
			s.insert(x);
			splay.insert(x);
		}
	}
	return true;
}

int main()
{
	if (test()) {
		printf("all tests passed\n");
	}
	return 0;
}
\end{vncode}
\endgroup


\subsection{Link Cut Tree}

Cấu trúc cây động với link/cut.

\codepath{Trees/linkcut.cpp}

\begingroup\scriptsize
\begin{vncode}
#include <cstdio>
#include <vector>
using namespace std;
#define INF 0x3f3f3f3f

/*
 * Link cut tree
 */

struct node {
	int sw, id, w;
	node *par, *ppar, *ls, *rs;
	node() {
		par = ppar = ls = rs = NULL; 
		w = sw = INF;
	}
};

class LinkCutTree {
	vector<node> lct;
	int sum(node *p) { return p ? p->sw : 0; }
	void refresh(node* p) {
		p->sw = p->w + sum(p->ls) + sum(p->rs);
	}
	void rotateright(node* p) {
		node *q, *r;
		q = p->par, r = q->par;
		if (q->ls = p->rs) q->ls->par = q;
		p->rs = q, q->par = p;
		if (p->par = r) {
			if (q == r->ls) r->ls = p;
			else r->rs = p;
		}
		p->ppar = q->ppar;
		q->ppar = NULL;
		refresh(q);
	}
	void rotateleft(node* p) {
		node *q, *r;
		q = p->par, r = q->par;
		if (q->rs = p->ls) q->rs->par = q;
		p->ls = q, q->par = p;
		if (p->par = r) {
			if (q == r->ls) r->ls = p;
			else r->rs = p;
		}
		p->ppar = q->ppar;
		q->ppar = 0;
		refresh(q);
	}
	void splay(node* p) {
		node *q, *r;
		while (p->par != NULL) {
			q = p->par;
			if (q->par == NULL) {
				if (p == q->ls) rotateright(p);
				else rotateleft(p);
				continue;
			}
			r = q->par;
			if (q == r->ls) {
				if (p == q->ls) rotateright(q), rotateright(p);
				else rotateleft(p), rotateright(p);
			}
			else {
				if (p == q->rs) rotateleft(q), rotateleft(p);
				else rotateright(p), rotateleft(p);
			}
		}
		refresh(p);
	}
	node* access(node* p) {
		splay(p);
		if (p->rs != NULL) {
			p->rs->ppar = p; p->rs->par = NULL;
			p->rs = NULL; refresh(p);
		}
		node* last = p;
		while (p->ppar != NULL) {
			node* q = last = p->ppar;
			splay(q);
			if (q->rs != NULL) {
				q->rs->ppar = q;
				q->rs->par = NULL;
			}
			q->rs = p; p->par = q;
			p->ppar = NULL;
			refresh(q); splay(p);
		}
		return last;
	}
public:
	LinkCutTree(int n = 0) {
		lct.resize(n + 1);
		for(int i = 0; i <= n; i++) {
			lct[i].id = i;
			refresh(&lct[i]);
		}
	}
	void link(int u, int v, int w = 1) {
		//u becomes child of v
		node *p = &lct[u], *q = &lct[v];
		access(p); access(q);
		p->ls = q; q->par = p; p->w = w;
		refresh(p);
	}
	void cut(int u) {
		node* p = &lct[u]; access(p);
		p->ls->par = NULL; p->ls = NULL;
		refresh(p);
	}
	int findroot(int u) {
		node* p = &lct[u]; access(p);
		while (p->ls) p = p->ls;
		splay(p);
		return p->id;
	}
	bool IsSameTree(int u, int v) {
		return findroot(u) == findroot(v);
	}
	int depth(int u) {
		access(&lct[u]);
		return lct[u].sw - lct[u].w;
	}
	int LCA(int u, int v) {
		access(&lct[u]);
		return access(&lct[v])->id;
	}
	int dist(int u, int v) {
		if (!IsSameTree(u, v)) return INF;
		node* lca = &lct[LCA(u, v)];
		int ans = 0;
		access(&lct[u]); splay(lca);
		ans += sum(lca->rs);
		access(&lct[v]); splay(lca);
		ans += sum(lca->rs);
		return ans;
	}
};

/*
 * Undirected link cut tree
 */

class UndirectedLinkCutTree {
	LinkCutTree lct;
	vector<int> par;

	void invert(int u) {
		if (par[u] == -1) return;
		int v = par[u];
		invert(v);
		lct.cut(u); par[u] = -1;
		lct.link(v, u); par[v] = u;
	}
public:
	UndirectedLinkCutTree(int n = 0) {
		lct = LinkCutTree(n);
		par.assign(n+1, -1);
	}
	void link(int u, int v) {
		if (lct.depth(u) < lct.depth(v)) {
			invert(u);
			lct.link(u, v); par[u] = v;
		}
		else {
			invert(v);
			lct.link(v, u); par[v] = u;
		}
	}
	void cut(int u, int v) {
		if (par[v] == u) u = v;
		lct.cut(u); par[u] = -1;
	}
	bool IsSameTree(int u, int v) {
		return lct.IsSameTree(u, v);
	}
};

/*
 * TEST MATRIX
 */

#include <cstring>
#include <cstdlib>
#define MAXN 100009
int par[MAXN];

int findroot(int i) {
	if (par[i] >= 0) return findroot(par[i]);
	else return i;
}

bool IsSameSet(int u, int v) {
	return findroot(u) == findroot(v);
}

void link(int u, int v) {
	if (!IsSameSet(u, v)) par[u] = v;
}

void cut(int u) {
	par[u] = -1;
}

int LCA(int u, int v) {
	if (!IsSameSet(u, v)) return -1;
	int du = 0, ul = u;
	while (ul >= 0) {
		du++;
		ul = par[ul];
	}
	int dv = 0, vl = v;
	while (vl >= 0) {
		dv++;
		vl = par[vl];
	}
	while (du > dv) {
		du--;
		u = par[u];
	}
	while (du < dv) {
		dv--;
		v = par[v];
	}
	while (u != v) {
		u = par[u];
		v = par[v];
	}
	return u;
}

int dist(int u, int v) {
	if (!IsSameSet(u, v)) return -1;
	int lca = LCA(u, v);
	int ans = 0;
	while(u != lca) {
		u = par[u];
		ans++;
	}
	while(v != lca) {
		v = par[v];
		ans++;
	}
	return ans;
}

bool test(int N, int T) {
	LinkCutTree lct(N);
	memset(&par, -1, sizeof par);
	for (int i = 0, u, v; i<T; i++) {
		int cmd = rand() % 5;
		//if (cmd == 1) cmd = 0; //link only
		if (cmd == 0) {
			u = rand() % N;
			v = rand() % N;
			if (IsSameSet(u, v) != lct.IsSameTree(u, v)) {
				printf("failed on test %d (IsSameSet(%d,%d) LCT:%d BF:%d)\n", i, u, v, lct.IsSameTree(u, v), IsSameSet(u, v));
				return false;
			}
			if (IsSameSet(u, v) || par[u] >= 0) {
				i--;
				continue;
			}
			//printf("%d son of %d\n", u, v);
			link(u, v);
			lct.link(u, v);
		}
		if (cmd == 1) {
			u = rand() % N;
			if (par[u] == -1) {
				i--;
				continue;
			}
			//printf("%d cut from %d\n", u, par[u]);
			cut(u);
			lct.cut(u);
		}
		if (cmd == 2) {
			u = rand() % N;
			v = rand() % N;
			if (IsSameSet(u, v) != lct.IsSameTree(u, v)) {
				printf("failed on test %d (IsSameSet(%d,%d) LCT:%d BF:%d)\n", i, u, v, lct.IsSameTree(u, v), IsSameSet(u, v));
				return false;
			}
			if (!IsSameSet(u, v)) {
				i--;
				continue;
			}
			if (LCA(u, v) != lct.LCA(u, v)) {
				printf("failed on test %d (LCA(%d,%d) LCT:%d BF:%d)\n", i, u, v, lct.LCA(u, v), LCA(u, v));
				return false;
			}
		}
		if (cmd == 3) {
			u = rand() % N;
			if (findroot(u) != lct.findroot(u)) {
				printf("failed on test %d (findroot(%d) LCT:%d BF:%d)\n", i, u, lct.findroot(u), findroot(u));
				return false;
			}
		}
		if (cmd == 4) {
			u = rand() % N;
			v = rand() % N;
			if (!IsSameSet(u, v)) {
				i--;
				continue;
			}
			if (dist(u, v) != lct.dist(u, v)) {
				printf("failed on test %d (dist(%d,%d) LCT:%d BF:%d)\n", i, u, v, lct.dist(u, v), dist(u, v));
				return false;
			}
		}
	}
	printf("all tests passed\n");
	return true;
}

int main() {
	test(10000, 100000);
}

/*
 * SPOJ DYNALCA
 */
/*
int main() {
	int N, M, u, v;
	char str[10];
	scanf("%d %d", &N, &M);
	LinkCutTree lct(N+1);
	for(int i=0; i<M; i++) {
		scanf(" %s", str);
		if(str[1] == 'i') {
			scanf("%d %d", &u, &v);
			lct.link(u, v);
		}
		if(str[1] == 'u') {
			scanf("%d", &u);
			lct.cut(u);
		}
		if (str[1] == 'c') {
			scanf("%d %d", &u, &v);
			printf("%d\n", lct.LCA(u, v));
		}
	}
	//if (test(10000, 5000000)) printf("All tests passed\n");
	return 0;
}*/

/*
 * Codeforces 100960H
 */
/*
UndirectedLinkCutTree ulct;

int main() {
	int N, a, b;
	char op;
	scanf("%d", &N);
	ulct = UndirectedLinkCutTree(N);
	while(true) {
		scanf(" %c", &op);
		switch(op) {
		case 'C':
			scanf("%d %d", &a, &b);
			ulct.link(a, b);
			break;
		case 'D':
			scanf("%d %d", &a, &b);
			ulct.cut(a, b);
			break;
		case 'T':
			scanf("%d %d", &a, &b);
			if (ulct.IsSameTree(a, b)) printf("YES\n");
			else printf("NO\n");
			fflush(stdout);
			break;
		case 'E':
			return 0;
		}
	}
	return 0;
}*/
\end{vncode}
\endgroup


\subsection{Link Cut Tree vô hướng}

Mẫu liên quan cho đồ thị tách/split graph trong nguồn Macacário.

\codepath{Graphs/splitGraph.cpp}

\begingroup\scriptsize
\begin{vncode}
#define _CRT_SECURE_NO_WARNINGS

#include <vector>
#include <algorithm>
using namespace std;

/*
* Verifies if a graph is split in O(Nlog(N)) and build it in O(N). A split graph can
* be divided in a complete graph and a independent set (graph with 0 edges).
*/
class SplitGraph {

public:

	SplitGraph(int N, vector<int>* g): d(N)
	{
		computeSplit(g, N);
		if (isSplit()) build(N);
	}

	vector<int> getCliqueSet(){
		return K;
	}

	vector<int> getStableSet(){
		return I;
	}

	bool isSplit() {
		return maxClique >= 0;
	}

	int maximumClique() {
		return maxClique;
	}

private:
	int maxClique;
	vector<pair<int, int> > d;
	vector<int> K, I, verts;

	void computeSplit(vector<int>* g, int N) {
		for (int i = 0; i < N; ++i) d[i] = {g[i].size(), i};
		sort(d.begin(), d.end());

		int m = N;
		for (int i = 0; i < N; ++i) {
			if (d[i].first >= N - i - 1) {
				m = i;
				break;
			}
		}

		long long lowSum = 0, highSum = 0;

		for (int i = 0; i < m; ++i) lowSum += d[i].first;
		for (int i = m; i < N; ++i) highSum += d[i].first;
		m = N - m;

		if (lowSum + m * 1LL * (m - 1) == highSum) {
			maxClique = m;
		}
		else {
			maxClique = -1;
		}
	}

	void build(int N){
		for (int i = 0; i < N - maxClique; ++i){
			I.push_back(d[i].second);
		}for (int i = N - maxClique; i < N; ++i){
			K.push_back(d[i].second);
		}
	}

};

/*
* Example: URI 1974 - Into darkness
*/

#include <cstdio>
#include <cstring>
#include <queue>
#include <set>

#define MAXN 1009
#define MAXM 109
#define MAXS 210000
#define ALFA 2
#define INF 1e9
#define EPS 1e-9
#define FOR(i, N) for(int i = 0; i < N; ++i)
#define FORN(i, j, N) for(int i = j; i < N; ++i)
typedef long long int ll;
typedef pair<int, int> ii;

int N, M;
vector<int> g[MAXN];
vector< pair<int, ii> > edgeList;
char code[MAXN][MAXM];
int len[MAXN];
ll MOD = 1000000009;

class SuffixAutomaton {
	int len[MAXS], link[MAXS], cnt[MAXS];
	int nxt[MAXS][ALFA], sz, last, root;
	int newnode() {
		int x = ++sz;
		len[x] = 0; link[x] = -1; cnt[x] = 1;
		for(int c = 0; c < ALFA; c++) nxt[x][c] = 0;
		return x;
	}
	inline int reduce(char c) { return c - '0'; }
public:
	SuffixAutomaton() { clear(); }
	void clear() {
		sz = 0;
		root = last = newnode();
	}
	void insert(const char *s) {
		for(int i = 0; s[i]; i++) extend(reduce(s[i]));
	}
	void extend(int c) {
		int cur = newnode(), p;
		len[cur] = len[last] + 1;
		for(p = last; p != -1 && !nxt[p][c]; p = link[p]) {
			nxt[p][c] = cur;
		}
		if (p == -1) link[cur] = root;
		else {
			int q = nxt[p][c];
			if (len[p] + 1 == len[q]) link[cur] = q;
			else {
				int clone = newnode();
				len[clone] = len[p] + 1;
				for(int i = 0; i < ALFA; i++)
					nxt[clone][i] = nxt[q][i];
				link[clone] = link[q];
				cnt[clone] = 0;
				while(p != -1 && nxt[p][c] == q) {
					nxt[p][c] = clone;
					p = link[p];
				}
				link[q] = link[cur] = clone;
			}	
		}
		last = cur;
	}

	int longestCommonSubstring(const char *s) {
		int cur = root, curlen = 0, ans = 0;
		for(int i = 0; s[i]; i++) {
			int c = reduce(s[i]);
			while(cur != root && !nxt[cur][c]) {
				cur = link[cur];
				curlen = len[cur];
			}
			if (nxt[cur][c]) {
				cur = nxt[cur][c];
				curlen++;
			}
			if (ans < curlen) ans = curlen;
		}
		return ans;
	}
private:
	int deg[MAXS];
};

class UnionFind {
private:
	vector<int> parent, rank;
public:
	UnionFind(int N) {
		rank.assign(N + 9, 0);
		parent.assign(N + 9, 0);
		for (int i = 0; i < N; i++) parent[i] = i;
	}
	int find(int i) {
		while (i != parent[i]) i = parent[i];
		return i;
	}
	bool isSameSet(int i, int j) {
		return find(i) == find(j);
	}
	void unionSet(int i, int j) {
		if (isSameSet(i, j)) return;
		int x = find(i), y = find(j);
		if (rank[x] > rank[y]) parent[y] = x;
		else {
			parent[x] = y;
			if (rank[x] == rank[y]) rank[y]++;
		}
	}
};

ll kruskal() {
	ll cost = 0;
	UnionFind UF(N);
	sort(edgeList.begin(), edgeList.end());
	for (auto edge : edgeList) {
		if (!UF.isSameSet(edge.second.first, edge.second.second)) {
			cost += edge.first;
			UF.unionSet(edge.second.first, edge.second.second);
		}
	}
	return cost;
}

int minCost(SplitGraph& sg) {
	SuffixAutomaton sa;
	for (auto& edge : edgeList){
		sa.clear();
		int u = edge.second.first;
		int v = edge.second.second;

		sa.insert(code[u]);
		edge.first = sa.longestCommonSubstring(code[v]);
	}

	return kruskal();
}

int vis[MAXN];

bool isConnected() {
	int count = 1;
	queue<int> inuse;

	inuse.push(0);
	vis[0] = 1;
	while (!inuse.empty()) {
		int u = inuse.front();
		inuse.pop();

		for (int v : g[u]) {
			if (!vis[v]) {
				vis[v] = 1;
				++count;
				inuse.push(v);
			}
		}
	}
	return count == N;
}

int main() {
	int u, v;

	scanf("%d%d", &N, &M);

	FOR(i, M) {
		scanf("%d%d", &u, &v);
		--u, --v;
		g[u].push_back(v);
		g[v].push_back(u);
		edgeList.push_back({ 0, {u, v} });
	}
	FOR(i, N) {
		scanf("%s", code[i]);
		len[i] = strlen(code[i]);
	}
	auto sg = SplitGraph(N, g);

	if (!isConnected() || !sg.isSplit()) {
		printf(":[\n");
	}
	else {
		int trainees = sg.maximumClique();
		printf(":]\n%d %d %d\n", trainees, N - trainees, minCost(sg));
	}
}
\end{vncode}
\endgroup


\subsection{Wavelet Tree}

Wavelet tree cho truy vấn thứ tự/tần suất.

\codepath{Trees/wavelet.cpp}

\begingroup\scriptsize
\begin{vncode}
#include <vector>
#include <algorithm>
using namespace std;

int Lcmp;
bool less(int i) { return i <= Lcmp; }

class WaveletTree {
private:
	vector< vector<int> > ml;
	vector<int> arr; 
	int sig, size;
#define left(p) (p << 1)
#define right(p) ((p << 1) + 1)
	
	void build(int u, int l, int r, int lo, int hi, int* A) {
		if (lo == hi) return;
		int mid = (lo + hi) / 2;
		Lcmp = mid;
		ml[u].reserve(r-l+2);
		ml[u].push_back(0);
		for (int i=l; i<=r; i++) {
			ml[u].push_back(ml[u].back() + (A[i]<=Lcmp));
		}
		int p = (int)(stable_partition(A+l, A+r+1, less) - A);
		build(left(u), l, p-1, lo, mid, A);
		build(right(u), p, r, mid+1, hi, A);
	}
	
	int rank(int u, int lo, int hi, int q, int i) {
		if (lo == hi) return i;
		int mid = (lo + hi) / 2, ri = ml[u][i];
		if (q <= mid) return rank(left(u), lo, mid, q, ri);
		else return rank(right(u), mid+1, hi, q, i - ri);
	}
	
	int quantile(int u, int lo, int hi, int i, int j, int k) {
		if (lo == hi) return lo;
		int mid = (lo + hi) / 2;
		int ri = ml[u][i-1], rj = ml[u][j], c = rj - ri;
		if (k <= c) return quantile(left(u), lo, mid, ri+1, rj, k);
		else return quantile(right(u), mid+1, hi, i-ri, j-rj, k-c);
	}
	
	int range(int u, int lo, int hi, int i, int j, int a, int b) {
		if (lo > b || hi < a) return 0;
		if (b >= hi && lo >= a) return j-i;
		int mid = (lo + hi) / 2;
		int ri = ml[u][i], rj = ml[u][j];
		int c1 = range(left(u), lo, mid, ri, rj, a, b);
		int c2 = range(right(u), mid+1, hi, i-ri, j-rj, a, b);
		return c1 + c2;
	}
	
	void swap(int u, int lo, int hi, int v1, int v2, int i) {
		if (lo == hi) return;
		int mid = (lo + hi) / 2;
		if (v1 <= mid) {
			if (v2 > mid) ml[u][i]--;
			else swap(left(u), lo, mid, v1, v2, ml[u][i]);
		}
		else {
			if (v2 <= mid) ml[u][i]++;
			else swap(right(u), mid+1, hi, v1, v2, i-ml[u][i]);			
		}
	}
	
public:
	WaveletTree() {}
	WaveletTree(int* begin, int* end, int _sig) {
		sig = _sig;
		size = (int)(end-begin);
		ml.resize(4*size);
		arr = vector<int>(begin, end);
		build(1, 0, size-1, 0, sig, &arr[0]);
		arr = vector<int>(begin, end);
	}
	
	int rank(int i, int q) { return rank(1, 0, sig, q, i+1); }
	int rank(int i, int j, int q) { return rank(j, q) - rank(i-1, q); }
	int quantile(int i, int j, int k) { return quantile(1, 0, sig, i+1, j+1, k); }
	int range(int i, int j, int a, int b) { return range(1, 0, sig, i, j+1, a, b); }
	void swap(int i) {
		if (i >= size-1) return;
		swap(1, 0, sig, arr[i], arr[i+1], i+1);
		std::swap(arr[i], arr[i+1]);
	}
};

/*
 * TEST MATRIX
 */

#include <cstdio>
#define DEBUG false
int testarr[100009];
WaveletTree wt;

bool test(int N, int sig, int nTest) {
	for(int t = 1; t <= nTest; t++) {
		if (DEBUG) printf("test array:");
		for(int i=0; i<N; i++) {
			testarr[i] = rand()%(sig+1);
			if (DEBUG) printf(" %d", testarr[i]);
		}
		if (DEBUG) printf("\n");
		
		if (DEBUG) printf("building Wavelet Tree\n");
		wt = WaveletTree(testarr, testarr+N, sig);
		if (DEBUG) printf("Wavelet Tree built\n");
		
		for(int j=1; j<=nTest; j++) {
			
			//execute swaps
			for(int s=0; s<nTest; s++) {
				int i = rand()%(N-1);
				swap(testarr[i], testarr[i+1]);
				wt.swap(i);
			}
			
			//test rank
			int a = rand()%N;
			int b = rand()%N;
			int q = rand()%(sig+1);
			if (a > b) swap(a, b);
			int ans = 0;
			if (DEBUG) printf("testarr[%d,%d]:", a, b);
			for(int i=a; i<=b; i++) {
				if (DEBUG) printf(" %d", testarr[i]);
				if (testarr[i] == q) ans++;
			}
			int wans = wt.rank(a, b, q);
			if (DEBUG) printf(" ans = %d wans = %d\n", ans, wans);
			if (ans != wans) {
				printf("Failed on rank test %d, rank(%d,%d,%d), ans = %d, wans = %d\n", j, a, b, q, ans, wans);
				return false;
			}
		
			//test quantile
			a = rand()%N;
			b = rand()%N;
			if (a > b) swap(a, b);
			int n = b-a+1;
			int k = 1 + rand()%n;
			vector<int> st(&testarr[a], &testarr[b] + 1);
			sort(st.begin(), st.end());
			if (DEBUG) printf("testarr[%d,%d]:", a, b);
			for(int i=0; i<n; i++) {
				if (DEBUG) printf(" %d", st[i]);
			}
			ans = st[k-1];
			wans = wt.quantile(a, b, k);
			if (DEBUG) printf(" ans = %d wans = %d\n", ans, wans);
			if (ans != wans) {
				printf("Failed on quantile test %d, quantile(%d,%d,%d), ans = %d, wans = %d\n", j, a, b, k, ans, wans);
				return false;
			}
			
			//test range
			a = rand()%N;
			b = rand()%N;
			if (a > b) swap(a, b);
			int l = rand()%(sig+1);
			int r = rand()%(sig+1);
			if (l > r) swap(l, r);
			if (DEBUG) printf("testarr[%d,%d]:", a, b);
			ans = 0;
			for(int i=a; i<=b; i++) {
				if (DEBUG) printf(" %d", testarr[i]);
				if (testarr[i] >= l && testarr[i] <= r) ans++;
			}
			wans = wt.range(a, b, l, r);
			if (DEBUG) printf(" ans = %d wans = %d\n", ans, wans);
			if (ans != wans) {
				printf("Failed on range test %d, range(%d,%d,%d,%d), ans = %d, wans = %d\n", j, a, b, l, r, ans, wans);
				return false;
			}
		}
	}
	printf("all tests passed\n");
	return true;
}

int main() {
	test(1000, 25, 100);
	return 0;
}

/*
 * SPOJ ILKQUERY
 */

/*#include <cstdio>
#define MAXN 100009
#include <map>
map<int, int> x2id;
int id2x[MAXN], arr[MAXN];

int main() {
	int N, Q;
	scanf("%d %d", &N, &Q);
	for(int i=0; i<N; i++) {
		scanf("%d", &arr[i]);
		x2id[arr[i]] = 0;
	}
	
	int sig = -1;
	for(map<int, int>::iterator it = x2id.begin(); it != x2id.end(); it++) {
		it->second = ++sig;
		id2x[it->second] = it->first;
	}
	for(int i=0; i<N; i++) arr[i] = x2id[arr[i]];
	
	WaveletTree wt(arr, arr+N, sig);
	
	int k, i, l;
	while(Q-->0) {
		scanf("%d %d %d", &k, &i, &l);
		int el = wt.quantile(0, i, k);
		//printf("%d-th element in [%d,%d] is %d\n", k, 0, i, id2x[el]);
		if (wt.rank(N-1, el) < l) {
			printf("-1\n");
			continue;
		}
		int lo = -1;
		int hi = N;
		while(hi > lo + 1) {
			int mid = (hi + lo) / 2;
			if (wt.rank(mid, el) < l) lo = mid;
			else hi = mid;
		}
		printf("%d\n", hi);
	}
	return 0;
}*/

/*
 * SPOJ ILKQUERYIII
 */

/*#define MAXN 1000009
#include <map>
#include <cstdio>
map<int, int> x2id;
int id2x[MAXN], arr[MAXN];

int main() {
	int N, Q;
	scanf("%d %d", &N, &Q);
	for(int i=0; i<N; i++) {
		scanf("%d", &arr[i]);
		x2id[arr[i]] = 0;
	}
	
	int sig = -1;
	for(map<int, int>::iterator it = x2id.begin(); it != x2id.end(); it++) {
		it->second = ++sig;
		id2x[it->second] = it->first;
	}
	for(int i=0; i<N; i++) arr[i] = x2id[arr[i]];
	
	WaveletTree wt(arr, arr+N, sig);
	
	int k, i, l, t;
	while(Q-->0) {
		scanf("%d", &t);
		if (t == 0) {
			scanf("%d %d %d", &i, &l, &k);
			int el = wt.quantile(0, i, k);
			//printf("%d-th element in [%d,%d] is %d\n", k, 0, i, id2x[el]);
			if (wt.rank(N-1, el) < l) {
				printf("-1\n");
				continue;
			}
			int lo = -1;
			int hi = N;
			while(hi > lo + 1) {
				int mid = (hi + lo) / 2;
				if (wt.rank(mid, el) < l) lo = mid;
				else hi = mid;
			}
			printf("%d\n", hi);
		}
		else {
			scanf("%d", &i);
			wt.swap(i);
		}
	}
	return 0;
}*/
\end{vncode}
\endgroup


\subsection{Heavy-Light Decomposition}

Phân rã nặng-nhẹ.

\codepath{Trees/hld.cpp}

\begingroup\scriptsize
\begin{vncode}
#include <vector>
#include <algorithm>
using namespace std;
#define MAXN 400009

#include <cstdio>

/*
 * Heavy-Light Decomposition
 */

vector<int> adjList[MAXN];
int in[MAXN], out[MAXN], up[MAXN];
int par[MAXN], pre[MAXN], sz[MAXN];

void dfssize(int u) {
	int s = sz[u] = 0;
	for(int i = 0; i < (int)adjList[u].size(); i++) {
		int &v = adjList[u][i]; //don't forget &
		if (v == par[u]) continue;
		par[v] = u;
		dfssize(v);
		s += sz[v];
		if(sz[v] > sz[adjList[u][0]])
			swap(v, adjList[u][0]);
	}
	sz[u] = 1 + s;
}

void dfshld(int u, int &t) {
	pre[t] = u;
	in[u] = t++;
	for(int i = 0; i < (int)adjList[u].size(); i++) {
		int v = adjList[u][i];
		if (v == par[u]) continue;
		up[v] = (i == 0 ? up[u] : v);
		dfshld(v, t);
	}
	out[u] = t-1;
}

void hld(int root) {
	par[root] = -1;
	dfssize(root);
	int t = 0;
	up[root] = root;
	dfshld(root, t);
	//create DS over values sorted by array pre
}

int LCA(int u, int v) {
	while(up[u] != up[v]) {
		if (in[u] > in[v]) u = par[up[u]];
		//query DS for subarray [in[up[u]], in[u]] 
		else v = par[up[v]];
		//query DS for subarray [in[up[v]], in[v]]
	}
	return in[u] > in[v] ? v : u;
	//query DS for subarray [in[v], in[u]] or [in[u], in[v]]
}

/*
 * KMP Automaton
 */

#include <cstring>
#define MAXM 109

struct trans {
	int m, to[MAXM], match[MAXM];
	trans(int _m = -1) : m(_m) {}
};

trans operator +(trans a, trans b) {
	if (b.m < 0) return a;
	if (a.m < 0) return b;
	//assert(a.m == b.m);
	trans c(a.m);
	for(int i = 0; i < c.m; i++) {
		c.to[i] = b.to[a.to[i]];
		c.match[i] = a.match[i] + b.match[a.to[i]];
	}
	return c;
}

struct KMP {
	char P[MAXM];
	int b[MAXM], m;
	void build(const char* _P) {
		m = strlen(_P);
		strcpy(P, _P);
		memset(&b, -1, sizeof b);
		for(int i = 0, j = -1; i < m;) {
			while (j >= 0 && P[i] != P[j]) j = b[j];
			b[++i] = ++j;
		}
	}
	trans get(char c) {
		trans ans(m);
		for(int i = 0, j, cnt; i < m; i++) {
			j = i; cnt = 0;
			while(j >= 0 && c != P[j]) j = b[j];
			j++;
			if (j == m) cnt++, j = b[j];
			ans.to[i] = j;
			ans.match[i] = cnt;
		}
		return ans;
	}
};

/*
 * Codeforces 101908H
 */

char P[MAXN];
KMP kmp;
char col[MAXN];

class SegmentTree {
	trans st[2][MAXN];
	int n;
	trans query(int u, int l, int r, int a, int b, int t) {
		if (l > b || r < a) return trans(-1);
		if (a <= l && r <= b) return st[t][u];
		int m = (l + r) / 2;
		trans p1 = query(2*u, l, m, a, b, t);
		trans p2 = query(2*u+1, m+1, r, a, b, t);
		return t == 0 ? p1 + p2 : p2 + p1;
	}
public:
	SegmentTree() {}
	void build(int sz) {
		for (n = 1; n < sz; n <<= 1);
		for(int i=0; i<sz; i++) {
			st[0][i+n] = st[1][i+n] = kmp.get(col[pre[i]]);
		}
		for(int i=n-1; i; i--) {
			st[0][i] = st[0][2*i] + st[0][2*i+1];
			st[1][i] = st[1][2*i+1] + st[1][2*i];
		}
	}
	void update(int i, char c) {
		st[0][i+n] = st[1][i+n] = kmp.get(c);
		for (i += n, i >>= 1; i; i >>= 1) {
			st[0][i] = st[0][2*i] + st[0][2*i+1];
			st[1][i] = st[1][2*i+1] + st[1][2*i];
		}
	}
	trans query(int l, int r, int t) {
		return query(1, 0, n-1, l, r, t);
	}
};

SegmentTree st;

trans get(int u, int v) {
	trans left, mid, right;
	while(up[u] != up[v]) {
		if (in[u] > in[v]) {
			trans cur = st.query(in[up[u]], in[u], 1);
			left = left + cur;
			u = par[up[u]];
		}
		else {
			trans cur = st.query(in[up[v]], in[v], 0);
			right = cur + right;
			v = par[up[v]];
		}
	}
	if (in[u] > in[v]) mid = st.query(in[v], in[u], 1);
	else mid = st.query(in[u], in[v], 0);
	return (left + mid) + right;
}

int main() {
	int n, q;
	scanf("%d %d", &n, &q);
	scanf(" %s %s", P, col+1);
	for(int i = 0; i < n-1; i++) {
		int u, v;
		scanf("%d %d", &u, &v);
		adjList[u].push_back(v);
		adjList[v].push_back(u);
	}
	hld(1);
	kmp.build(P);
	st.build(n);
	while(q --> 0) {
		int op;
		scanf("%d", &op);
		if (op == 1) {
			int u, v;
			scanf("%d %d", &u, &v);
			trans ans = get(u, v);
			printf("%d\n", ans.match[0]);
		}
		else {
			int u;
			char x;
			scanf("%d %c", &u, &x);
			col[u] = x;
			st.update(in[u], x);
		}
	}
	return 0;
}
\end{vncode}
\endgroup


\subsection{Centroid Decomposition}

Phân rã trọng tâm.

\codepath{Trees/centroid.cpp}

\begingroup\scriptsize
\begin{vncode}
#include <cstring>
#include <vector>
#include <iostream>
using namespace std;
#define MAXN 500009
#define INF (1LL<<60)

typedef long long ll;
typedef pair<ll, int> ii;
int clevel[MAXN], cpar[MAXN], csize[MAXN];
vector<int> csons[MAXN];
vector<ii> adjList[MAXN];
int N, K;

int subsize(int u, int p) {
	csize[u]=1;
	for(int i=0; i<(int)adjList[u].size(); i++) {
		int v = adjList[u][i].second;
		if (v != p && clevel[v] < 0)
			csize[u] += subsize(v, u);
	}
	return csize[u];
}
int findcentroid(int u, int p, int nn) {
	for(int i=0; i<(int)adjList[u].size(); i++) {
		int v = adjList[u][i].second;
		if (v != p && clevel[v] < 0 && csize[v] > nn/2)
			return findcentroid(v, u, nn);
	}
	return u;
}
int decompose(int root, int par) {
	subsize(root, -1);
	int u = findcentroid(root, -1, csize[root]);
	cpar[u] = par;
	clevel[u] = par >= 0 ? clevel[par]+1 : 0;
	csize[u] = 1;
	for(int i=0; i<(int)adjList[u].size(); i++) {
		int v = adjList[u][i].second;
		if (v != par && clevel[v] < 0) {
			v = decompose(v, u);
			csons[u].push_back(v);
			csize[u] += csize[v];
		}
	}
	return u;
}
int centroiddecomposition(int root) {
	memset(&clevel, -1, sizeof clevel);
	for(int i=0; i<=N; i++) csons[i].clear();
	return decompose(root, -1);
}

/*
 * Codeforces 101040J
 */

#include <cstdio>

vector<ll> cost(MAXN), subs(MAXN);
int limit;
ll depth[MAXN];

void dfs(int u, int par, int h) {
	while(cost.size() < h+1) cost.push_back(INF);
	cost[h] = min(cost[h], depth[u]);
	for(int i=0; i<(int)adjList[u].size(); i++) {
		int v = adjList[u][i].second;
		ll w = adjList[u][i].first;
		if (clevel[v] <= limit || v == par) continue;
		depth[v] = depth[u] + w;
		dfs(v, u, h+1);
	}
}

ll merge() {
	ll ans = INF;
	while(subs.size() < cost.size()) subs.push_back(INF);
	for(int i=0; i<(int)cost.size(); i++) {
		int j = K-i;
		if (j >= 0 && j < (int)subs.size()) ans = min(ans, subs[j] + cost[i]);
	}
	for(int i=0; i<(int)cost.size(); i++) {
		subs[i] = min(subs[i], cost[i]);
	}
	return ans;
}

ll solve(int u) {
	ll ans = INF;
	for(int i=0; i<(int)csons[u].size(); i++) {
		int v = csons[u][i];
		ans = min(ans, solve(v));
	}
	subs.clear();
	subs.push_back(0);
	limit = clevel[u];
	for(int i=0; i<(int)adjList[u].size(); i++) {
		int v = adjList[u][i].second;
		ll w = adjList[u][i].first;
		if (clevel[v] <= limit) continue;
		depth[v] = depth[u] + w;
		cost.clear();
		dfs(v, u, 1);
		ans = min(ans, merge());
	}
	return ans;
}
int main() {
	scanf("%d %d", &N, &K);
	int P;
	ll C, ans;
	for(int i=2; i<=N; i++) {
		scanf("%d %I64d", &P, &C);
		adjList[i].push_back(ii(C, P));
		adjList[P].push_back(ii(C, i));
	}
	int root = centroiddecomposition(1);
	ans = solve(root);
	if (ans < INF) printf("%I64d\n", ans);
	else printf("-1\n");
	return 0;
}
\end{vncode}
\endgroup


\subsection{Interval Tree}

Cây khoảng.

\codepath{Trees/intervaltree.cpp}

\begingroup\scriptsize
\begin{vncode}
#include <set>
#include <vector>
#include <algorithm>
using namespace std;
#define MOD 998244353
#define MAXN 200009
typedef pair<int, int> ii;

/*
 * Lazy propagation segment tree
 */

class SegmentTree {
private:
	vector<int> st;
	vector<ii> lazy;
	int size;
#define parent(p) (p >> 1)
#define left(p) (p << 1)
#define right(p) ((p << 1) + 1)
	ii aggregate(ii a, ii b) {
		return ii((a.first*1ll*b.first)%MOD, (a.second*1ll*b.first + b.second)%MOD);
	} 
	void push(int p, int l, int r) {
		st[p] = (lazy[p].first*1ll*st[p] + (r - l + 1ll)*lazy[p].second)%MOD;
		if (l != r) {
			lazy[right(p)] = aggregate(lazy[right(p)], lazy[p]);
			lazy[left(p)] = aggregate(lazy[left(p)], lazy[p]);
		}
		lazy[p] = ii(1, 0);
	}
	void update(int p, int l, int r, int a, int b, ii k) {
		push(p, l, r);
		if (a > r || b < l) return;
		else if (l >= a && r <= b) {
			lazy[p] = k; push(p, l, r);
		}
		else {
			update(left(p), l, (l + r) / 2, a, b, k);
			update(right(p), (l + r) / 2 + 1, r, a, b, k);
			st[p] = (st[left(p)] + st[right(p)]) % MOD;
		}
	}
	int query(int p, int l, int r, int a, int b) {
		push(p, l, r);
		if (a > r || b < l) return 0;
		if (l >= a && r <= b) return st[p];
		int p1 = query(left(p), l, (l + r) / 2, a, b);
		int p2 = query(right(p), (l + r) / 2 + 1, r, a, b);
		return (p1 + p2) % MOD;
	}
public:
	SegmentTree(int _size) {
		size = _size;
		st.assign(4 * size, 0);
		lazy.assign(4 * size, ii(1, 0));
	}
	int query(int a, int b) { return query(1, 0, size - 1, a, b); }
	void update(int a, int b, ii k) { update(1, 0, size - 1, a, b, k); }
};

SegmentTree st(MAXN);

/*
 * Interval Tree
 */

struct node {
	int l, r, x;
	node(int _l, int _r, int _x = 0) : l(_l), r(_r), x(_x) { }
	void update(int dx) { x = min(x+dx, 2); }
};

bool operator < (node a, node b) {
	if (a.l != b.l) return a.l < b.l;
	return a.r < b.r;
}

class IntervalTree {
	set<node> tree;
	void split(int i) {
		set<node>::iterator it = --tree.upper_bound(node(i, 0));
		node t = *it;
		if (t.l == i) return;
		tree.erase(it);
		tree.insert(node(t.l, i-1, t.x));
		tree.insert(node(i, t.r, t.x));
	}
public:
	IntervalTree() { tree.insert(node(0, MAXN, 0)); }
	vector<node> get(int l, int r, bool update, int dx) {
		split(l); split(r+1);
		set<node>::iterator it = tree.lower_bound(node(l, 0));
		vector<node> q;
		while(it != tree.end() && it->l <= r) {
			node t = *it;
			it = tree.erase(it);
			if (update) t.update(dx);
			if (q.empty() || q.back().x != t.x) q.push_back(t);
			else q.back().r = t.r;
		}
		for(int i = 0; i < int(q.size()); i++) {
			tree.insert(q[i]);
		}
		return q;
	}
};

/*
 * Codeforces 981G
 */

#include <cstdio>

IntervalTree it[MAXN];

int main() {
	int n, q;
	scanf("%d %d", &n, &q);
	while(q --> 0) {
		int t, l, r, x;
		scanf("%d", &t);
		if (t == 1) {
			scanf("%d %d %d", &l, &r, &x);
			vector<node> q = it[x].get(l, r, true, 1);
			for(int i = 0; i < int(q.size()); i++) {
				if (q[i].x == 1) st.update(q[i].l, q[i].r, ii(1, 1));
				else st.update(q[i].l, q[i].r, ii(2, 0));
			}
		}
		else {
			scanf("%d %d", &l, &r);
			printf("%d\n", st.query(l, r));
		}
	}
	return 0;
}
\end{vncode}
\endgroup


\subsection{Xor Trie}

Trie nhị phân hỗ trợ tìm phần tử tối ưu xor dưới giới hạn.

\codepath{Trees/xortrie.cpp}

\begingroup\scriptsize
\begin{vncode}
/*
 * Xor Trie
 */
#define MAXS 60000009
int l[MAXS], r[MAXS], cnt = 0;

class XorTrie {
	int digits = 0, root, ans, limit;
	int newnode() {
		l[cnt] = r[cnt] = -1;
		cnt++;
		return cnt-1;
	}
	bool search(int u, int h, int num, int cur) {
		if (u == -1 || cur > limit) return false;
		if (h == -1) {
			ans = cur; return true;
		}
		if (num & (1<<h)) {
			if (search(l[u], h-1, num, cur)) return true;
			if (search(r[u], h-1, num, cur | (1<<h))) return true;
		}
		else {
			if (search(r[u], h-1, num, cur | (1<<h))) return true;
			if (search(l[u], h-1, num, cur)) return true;
		}
		return false;
	}
public:
	XorTrie(int _digits = 20) : digits(_digits) {
		root = newnode();
	}
	void insert(int num) {
		int u = root;
		for(int i = digits-1; i >= 0; i--) {
			if (num & (1<<i)) {
				if (r[u] == -1) r[u] = newnode();
				u = r[u];
			}
			else {
				if (l[u] == -1) l[u] = newnode();
				u = l[u];
			}
		}
	}
	int search(int _limit, int num) {
		limit = _limit; ans = -1;
		if (!search(root, digits-1, num, 0)) return -1;
		return ans;
	}
};

/*
 * Codeforces 979D
 */

#include <cstdio>
#define MAXN 100309
XorTrie tree[MAXN];

int main() {
	int q;
	scanf("%d", &q);
	while (q --> 0) {
		int t, u, x, k, s;
		scanf("%d", &t);
		if (t == 1) {
			scanf("%d", &u);
			for(int i = 1; i*i <= u; i++) {
				if (u % i == 0) {
					tree[i].insert(u);
					if (i*i < u) tree[u/i].insert(u);
				}
			}
		}
		else if (t == 2) {
			scanf("%d %d %d", &x, &k, &s);
			int ans = -1;
			ans = tree[k].search(s-x, x);
			if (x % k != 0) ans = -1;
			printf("%d\n", ans);
		}
	}
	return 0;
}
\end{vncode}
\endgroup


\subsection{Convex Hull Trick}

Với truy vấn min/max đường thẳng theo x; có phiên bản O(log n) và O(1) khi x tăng.

\codepath{Data Structures/convexhulltrick.cpp}

\begingroup\scriptsize
\begin{vncode}
#include <cstdio>
#include <vector>
#define INF (1<<30)
#define MAXN 1009
using namespace std;

typedef long long int ll;

/*
 * Partially Dynamic Convex Hull Trick
 */

class CHTrick {
private:
	vector<ll> m, n;
	vector<double> p;
	bool maxCH; int i;
public:
	CHTrick(bool mxch) { maxCH = mxch; i = 0; }
	void clear() { m.clear(); n.clear(); p.clear(); }
	double inter(double nm, double nn, double lm, double ln) {
		return (ln - nn) / (nm - lm);
	}
	void push(ll nm, ll nn) {
		while (!p.empty()) {
			if (nm == m.back() && maxCH && nn <= n.back()) return;
			if (nm == m.back() && !maxCH && nn >= n.back()) return;
			double x = inter(nm, nn, m.back(), n.back());
			if (nm != m.back() && p.back() < x) break;
			m.pop_back(); n.pop_back(); p.pop_back();
		}
		p.push_back(p.empty() ? -INF : inter(nm, nn, m.back(), n.back()));
		m.push_back(nm); n.push_back(nn);
		if (i >= p.size()) i = p.size()-1;
	}
	ll query(ll x) {
		if (p.empty()) return (maxCH ? -1 : 1)*INF;
		ll r = p.size() - 1, l = 0, mid;
		if (x >= p[r]) return m[r] * x + n[r];
		while (r > l + 1) {
			mid = (r + l) / 2;
			if (x < p[mid]) r = mid;
			else l = mid;
		}
		return m[l] * x + n[l];
	}
	ll query_q(ll x) {
		while (p[i] > x) i--;
		while(i+1 < p.size() && p[i+1] <= x) i++;
		return m[i] * x + n[i];
	}
};

/*
 * URI 1282
 */

ll w[MAXN], x[MAXN], A[MAXN], B[MAXN], dp[MAXN][MAXN];

int main(){
	int N, K;
	CHTrick cht(false);
	while(scanf("%d %d", &N, &K)!=EOF) {
		for(int i=0; i<N; i++){
			scanf("%lld %lld", x+i, w+i);
			A[i] = w[i] + (i>0 ? A[i-1] : 0);
			B[i] = w[i]*x[i] + (i>0 ? B[i-1] : 0);
			dp[i][1] = x[i]*A[i] - B[i];
		}
		for(int k=2; k<=K; k++){
			dp[0][k] = 0;
			cht.clear();
			for(int i=1; i<N; i++){
				cht.push(-A[i-1], dp[i-1][k-1]+B[i-1]);
				//dp[i][k] = x[i]*A[i] - B[i] + cht.query(x[i]);
				dp[i][k] = x[i]*A[i] - B[i] + cht.query_q(x[i]);
			}
		}
		printf("%lld\n", dp[N-1][K]);
	}
	return 0;
}
\end{vncode}
\endgroup


\subsection{Convex Hull Trick động và tích vô hướng lớn nhất}

Thêm đường thẳng bất kỳ và truy vấn trên phân số; cần C++11.

\codepath{Data Structures/dynamiccht.cpp}

\begingroup\scriptsize
\begin{vncode}
#include <bits/stdc++.h>
using namespace std;
typedef long long ll;

ll nu, de;
struct line {
	ll m, b; bool is_query;
	mutable function<const line *()> succ;
	line(ll _m, ll _b, bool iq = false) :
		m(_m), b(_b), is_query(iq) {}
	bool operator<(const line &o) const {
		if (!o.is_query) return m < o.m;
		const line *s = succ();
		if (!s) return 0;
		return (b - s->b) * de < (s->m - m) * nu;
	}
};
 
struct DynamicHull : public multiset<line> {
	bool bad(iterator y) {
		auto z = next(y);
		if (y == begin()) {
			if (z == end()) return 0;
			return y->m == z->m && y->b <= z->b;
		}
		auto x = prev(y);
		if (z == end()) return y->m == x->m && y->b <= x->b;
		return (x->b - y->b) * (z->m - y->m) >= (y->b - z->b) * (y->m - x->m);
	}
	void insert_line(ll m, ll b) {
		auto y = insert(line(m, b));
		y->succ = [=] { return next(y) == end() ? 0 : &*next(y); };
		if (bad(y)) {
			erase(y); return;
		}
		while (next(y) != end() && bad(next(y))) erase(next(y));
		while (y != begin() && bad(prev(y))) erase(prev(y));
	}
	line query(ll a, ll b) {
		if (b < 0) a = -a, b = -b;
		nu = a; de = b;
		return *lower_bound(line(0, 0, true));
    }
};

/*
 * Codeforces 101707H
 */

#define MAXN 200309
#define MAXM 900009
#define MOD 1000000007
#define INF 0x3f3f3f3f
#define INFLL 0x3f3f3f3f3f3f3f3f
#define EPS 1e-9
#define PI 3.141592653589793238462643383279502884
#define FOR(x,n) for(int x=0; (x)<int(n); (x)++)
#define FOR1(x,n) for(int x=1; (x)<=int(n); (x)++)
#define REP(x,n) for(int x=int(n)-1; (x)>=0; (x)--)
#define REP1(x,n) for(int x=(n); (x)>0; (x)--)
#define pb push_back
#define pf push_front
#define fi first
#define se second
#define mp make_pair
#define all(x) x.begin(), x.end()
#define mset(x,y) memset(&x, (y), sizeof(x));
using namespace std;
typedef vector<int> vi;
typedef pair<int, int> ii;
typedef long long ll;
typedef unsigned long long ull;
typedef long double ld;
typedef unsigned int uint;

int n, m;
ll a[MAXN], b[MAXN], c[MAXN];
ll x[MAXN], y[MAXN];
ll minx = INF, maxx = -INF;
DynamicHull maxhull, minhull;

ll query(ll a, ll b) {
	if (b == 0) {
		if (a > 0) return maxx*a;
		else return minx*a;
	}
	else if (b > 0) {
		line l = maxhull.query(a, b);
		return l.m*a + l.b*b;
	}
	else {
		line l = minhull.query(a, b);
		return -a*l.m + -b*l.b;
	}
}

int main() {
	scanf("%d %d", &n, &m);
	FOR(i, n) scanf("%lld %lld %lld", &a[i], &b[i], &c[i]);
	FOR(i, m) {
		scanf("%lld %lld", &x[i], &y[i]);
		minx = min(x[i], minx);
		maxx = max(x[i], maxx);
		maxhull.insert_line(x[i], y[i]);
		minhull.insert_line(-x[i], -y[i]);
	}
	vi ans;
	FOR(i, n) {
		ll c1, c2;
		c[i] = -c[i];
		c1 = query(a[i], b[i]);
		c2 = -query(-a[i], -b[i]);
		bool ok = false;
		if (c1 == c[i] || c2 == c[i]) ok = true;
		if (c1 < c[i] && c2 > c[i]) ok = true;
		if (c1 > c[i] && c2 < c[i]) ok = true;
		if (ok) ans.pb(i+1);
	}
	printf("%u\n", ans.size());
	for(int i : ans) printf("%d ", i);
	printf("\n");
	return 0;
}
\end{vncode}
\endgroup


\subsection{Van Emde Boas Tree}

Hỗ trợ tập con của \texttt{[0,sigma)}, kế tiếp và trước đó trong O(log log sigma).

\codepath{Trees/veb.cpp}

\begingroup\scriptsize
\begin{vncode}
#include <vector>
#include <cstdio>
using namespace std;

/*
 * Van Emde Boas Tree
 */

class VEB {
private:
	VEB* ups;
	vector<VEB*> los;
	int minx, maxx, k, tk;
#define upper(x) ((x)>>tk)
#define lower(x) ((x)&((1<<tk)-1))
#define index(x, y) (((x)<<tk)^(y))
public:
	VEB(int _k = 20) : k(_k), tk(k>>1) {
		minx = -1, maxx = -1;
		ups = NULL;
		if (k == 1) return;
		ups = new VEB(k-tk);
		los.resize(1<<(k-tk), NULL);
	}
	~VEB() {
		delete ups;
		for(int i = 0; i < (int)los.size(); i++)
			if (los[i]) delete los[i];
	}
	int min() { return minx; }
	int max() { return maxx; }
	bool count(int x) {
		int up = upper(x), lo = lower(x);
		if (x == minx || x == maxx) return true;
		if (k == 1 || !los[up]) return false;
		return los[up]->count(lo);
	}
	int succ(int x) {
		int up = upper(x), lo = lower(x);
		if (k == 1) return x == 0 && maxx == 1 ? 1 : -1;
		if (minx != -1 && x < minx) return minx;
		int mlo = los[up] ? los[up]->max() : -1;
		if (mlo != -1 && lo < mlo)
			return index(up, los[up]->succ(lo));
		up = ups->succ(up);
		if (up == -1) return -1;
		return index(up, los[up]->min());
	}
	int precc(int x) {
		int up = upper(x), lo = lower(x);
		if (k == 1) return x == 1 && minx == 0 ? 0 : -1;
		if (maxx != -1 && x > maxx) return maxx;
		int mlo = los[up] ? los[up]->min() : -1;
		if (mlo != -1 && lo > mlo)
			return index(up, los[up]->precc(lo));
		up = ups->precc(up);
		if (up == -1)
			return minx != -1 && x > minx ? minx : -1;
		return index(up, los[up]->max());
	}
	void insert(int x) {
		if (minx == -1) {
			minx = maxx = x; return;
		}
		if (x < minx) swap(x, minx);
		if (x > maxx) maxx = x;
		if (k == 1 || x == minx) return;
		int up = upper(x), lo = lower(x);
		if (!los[up] || los[up]->min() == -1) {
			ups->insert(up);
			los[up] = new VEB(tk);
		}
		los[up]->insert(lo);
	}
	int removemin() {
		if (minx == -1) return -1;
		int ans = minx;
		if (minx == maxx) {
			minx = maxx = -1; return ans;
		}
		if (k == 1) { minx = maxx; return ans; }
		int up = ups->min();
		minx = index(up, los[up]->removemin());
		if (los[up]->min() == -1) ups->removemin();
		return ans;
	}
	void erase(int x) {
		if (x == -1) return;
		int up = upper(x), lo = lower(x);
		if (x == minx) removemin();
		else if (k == 1 && x == maxx) maxx = minx;
		else {
			los[up]->erase(lo);
			if (los[up]->min() == -1) ups->erase(up);
			if (x == maxx && los[up]->max() != -1)
				maxx = index(up, los[up]->max());
			else if (x == maxx) maxx = precc(x);
		}
	}	
};

/*
 * TEST MATRIX
 */

/*
	void print(int h) {
		for(int i = 0; i < h; i++) printf("| ");
		printf("min: %d max: %d\n", minx, maxx);

		for(int i = 0; i < h; i++) printf("| ");
		printf("ups:%s\n", ups ? "" : " NULL");
		if (ups) ups->print(h+1);

		for(int j = 0; j < (int)los.size(); j++) {
			for(int i = 0; i < h; i++) printf("| ");
			printf("los[%d]:%s\n", j, los[j] ? "" : " NULL");
			if (los[j]) los[j]->print(h+1);
		}
	}
*/
#include <set>
#include <cstdlib>
#include <ctime>

bool test(int ntests, int k) {
	srand(time(NULL));
	VEB veb(k);
	set<int> control;
	vector<int> present;
	for(int test = 1; test <= ntests; test++) {
		//printf("test %d:\n", test);
		int op = rand()%(2*control.size() > (1<<(k-1)) ? 2 : 3);
		if (control.size() == (1u<<k) || (op == 0 && !control.empty())) {
			int i, x;
			if (rand() % 2) {
				i = rand()%present.size();
				x = present[i];
				//printf("removing %d...\n", x);
				veb.erase(x);
				//printf("removed %d\n", x);
			}
			else {
				x = *control.begin();
				i = 0;
				while(present[i] != x) i++;
				//printf("removing %d (min)...\n", x);
				int y = veb.removemin();
				//printf("removed %d (min)\n", y);
				if (x != y) {
					printf("failed test %d, removed min %d, but found %d\n", test, x, y);
					return false;
				}
			}
			control.erase(x);
			swap(present[i], present.back());
			present.pop_back();
		}
		else {
			int x;
			do {
				x = rand()%(1<<k);
			} while(control.count(x));
			control.insert(x);
			//printf("adding %d...\n", x);
			veb.insert(x);
			present.push_back(x);
			//printf("added %d\n", x);
		}
		//printf("-->\n");
		//veb.print(1);
		//printf("<--\n");
		for(int x = 0; x < (1<<k); x++) {
			bool res1 = control.count(x);
			bool res2 = veb.count(x);
			if (res1 != res2) {
				printf("failed test %d, count(%d): control %d veb %d", test, x, res1, res2);
				return false;
			}
		}
		for(int x = 0; x < (1<<k); x++) {
			if (x > 0) {
				set<int>::iterator it = control.lower_bound(x);
				int prec = veb.precc(x);
				if (it == control.begin()) {
					if (prec != -1) {
						printf("failed test %d, prec(%d)=%d should not exist\n", test, x, prec);
						return false;
					}
				}
				else {
					it--;
					if (*it != prec) {
						printf("failed test %d, precControl(%d)=%d precVeb(%d)=%d\n", test, x, *it, x, prec);
						return false;
					}
				}
			}
		}
		for(int x = 0; x < (1<<k); x++) {
			if (x+1 < (1<<k)) {
				set<int>::iterator it = control.upper_bound(x);
				int suc = veb.succ(x);
				if (it == control.end()) {
					if (suc != -1) {
						printf("failed test %d, suc(%d)=%d should not exist\n", test, x, suc);
						return false;
					}
				}
				else {
					if (*it != suc) {
						printf("failed test %d, sucControl(%d)=%d sucVeb(%d)=%d\n", test, x, *it, x, suc);
						return false;
					}
				}
			}
		}
	}
	printf("all tests passed\n");
	return true;
}

int main() {
	test(1000, 16);
}
\end{vncode}
\endgroup


\subsection{Lowest Common Ancestor (LCA) và truy vấn đường đi trên cây}

\texttt{P[i][j]} là tổ tiên \texttt{2\textasciicircum{}j} của i, \texttt{D[i][j]} là giá trị đường đi tương ứng.

\codepath{Trees/lca.cpp}

\begingroup\scriptsize
\begin{vncode}
#include <cstdio>
#include <vector>
#include <iostream>
using namespace std;
#define MAXN 100009
#define MAXLOGN 20
#define INF (1<<30)

/*
 * LCA + PATH QUERY O(logn)
 */

const int neutral = 0;
#define comp(a, b) ((a)+(b))

typedef pair<int, int> ii;
vector<ii> adjList[MAXN];
int level[MAXN];
int P[MAXN][MAXLOGN], D[MAXN][MAXLOGN];

void depthdfs(int u) {
	for(int i=0; i<(int)adjList[u].size(); i++) {
		int v = adjList[u][i].first;
		int w = adjList[u][i].second;
		if (v == P[u][0]) continue;
		P[v][0] = u; D[v][0] = w;
		level[v] = 1 + level[u];
		depthdfs(v);
	}
}

void computeP(int root, int n) {
	level[root] = 0;
	P[root][0] = root; D[root][0] = neutral;
	depthdfs(root);
	for(int j = 1; j < MAXLOGN; j++)
		for(int i = 1; i <= n; i++) {
			P[i][j] = P[P[i][j-1]][j-1];
			D[i][j] = comp(D[P[i][j-1]][j-1], D[i][j-1]);
		}
}

ii LCA(int u, int v) {
	if (level[u] > level[v]) swap(u, v);
	int d = level[v] - level[u];
	int ans = neutral;
	for(int i = 0; i < MAXLOGN; i++) {
		if (d & (1<<i)) {
			ans = comp(ans, D[v][i]);
			v = P[v][i];
		}
	}
	if (u == v) return ii(u, ans);
	for(int i = MAXLOGN-1; i >= 0; i--)
		while(P[u][i] != P[v][i]) {
			ans = comp(ans, D[v][i]);
			ans = comp(ans, D[u][i]);
			u = P[u][i]; v = P[v][i];
		}
	ans = comp(ans, D[v][0]);
	ans = comp(ans, D[u][0]);
	return ii(P[u][0], ans);
}

/*
 * TEST MATRIX
 */
#include <cstdlib>
#include <ctime>
int N;

void generateRandTree(int _N) {
	N = _N;
	//printf("new tree %d:\n", N);
	for(int i=1; i<=N; i++) {
		adjList[i].clear();
		if (i == 1) continue;
		int j = 1 + rand()%(i-1);
		int w = rand()%1000;
		//printf("%d-%d: %d\n", i, j, w);
		adjList[i].push_back(ii(j, w));
		adjList[j].push_back(ii(i, w));
	}
}

bool dist(int u, int p, int t, int & ans) {
	if (u == t) return true;
	for(int i=0; i<(int)adjList[u].size(); i++) {
		int v = adjList[u][i].first;
		int w = adjList[u][i].second;
		if (v == p) continue;
		if (dist(v, u, t, ans)) {
			ans = comp(ans, w);
			return true;
		}
	}
	return false;
}

bool test(int nTest) {
	clock_t naive=0, lca=0, last;
	for(int t=1; t<=nTest; t++) {
		generateRandTree(t);
		computeP(1, t);

		for(int it=0; it<t; it++) {
			int u = 1 + (rand()%t);
			int v = 1 + (rand()%t);
			int res1 = neutral;
			last = clock();
			dist(u, -1, v, res1);
			naive += clock() - last;
			last = clock();
			int res2 = LCA(u, v).second;
			lca += clock() - last;
			if (res1 != res2) {
				printf("failed on test %d\n", t);
				return false;
			}
		}
	}

	printf("all tests passed\n");
	printf("naive time: %.3fs, lca time: %.3fs\n", naive*1.0/CLOCKS_PER_SEC, lca*1.0/CLOCKS_PER_SEC);
	return true;
}

bool usertest() {
	scanf("%d", &N);
	int a, b, c;
	for(int i=1; i<N; i++) {
		scanf("%d %d %d", &a, &b, &c);
		adjList[a].push_back(ii(b, c));
		adjList[b].push_back(ii(a, c));
	}
	int root = 1;
	computeP(root, N);
	while(scanf("%d %d", &a, &b)!=EOF) {
		printf("LCA(%d, %d) = %d, ", a, b, LCA(a,b).first);
	}
	return true;
}

int main() {
	test(1000);
	return 0;
}
\end{vncode}
\endgroup


\clearpage

\section{Paradigms}

\subsection{Merge Sort}

Thuật toán O(n log n) để sắp xếp đoạn; \texttt{inv} đếm số nghịch thế.

\codepath{Paradigms/mergesort.cpp}

\begingroup\scriptsize
\begin{vncode}
#include <cstdio>
#define MAXN 100009

/*
 * Merge-sort with inversion count in O(nlogn)
 */

long long inv;
int aux[MAXN];

void mergesort(int arr[], int l, int r) {
	if (l == r) return;
	int m = (l + r) / 2;
	mergesort(arr, l, m);
	mergesort(arr, m+1, r);
	int i = l, j = m + 1, k = l;
	while(i <= m && j <= r) {
		if (arr[i] > arr[j]) {
			aux[k++] = arr[j++];
			inv += j - k;
		}
		else aux[k++] = arr[i++];
	}
	while(i <= m) aux[k++] = arr[i++];
	for(i = l; i < k; i++) arr[i] = aux[i];
}

/*
 * TEST MATRIX
 */

int main() {
	int N, vet[MAXN];
	inv = 0;
	scanf("%d", &N);
	for(int i=0; i<N; i++) {
		scanf("%d", vet+i);
	}
	mergesort(vet, 0, N-1);
	printf("inv = %lld:", inv);
	for(int i=0; i<N; i++) {
		printf(" %d", vet[i]);
	}
	printf("\n");
	return 0;
}
\end{vncode}
\endgroup


\subsection{Quick Sort}

Thuật toán kỳ vọng O(n log n), rất nhanh trong thực tế.

\codepath{Paradigms/quicksort.cpp}

\begingroup\scriptsize
\begin{vncode}
#include <cstdio>
#include <algorithm>
using namespace std;

void quicksort(int* arr, int l, int r) {
    if (l >= r) return;
    int mid = rand()%(r-l+1) + l;
    int pivot = arr[mid];
    swap(arr[mid],arr[l]);
    int i = l + 1, j = r;
    while (i <= j) {
        while(i <= j && arr[i] <= pivot) i++;
        while(i <= j && arr[j] > pivot) j--;
        if (i < j) swap(arr[i], arr[j]);
    }
    swap(arr[i-1], arr[l]);
    quicksort(arr, l, i-2);
    quicksort(arr, i, r);
}

void quicksort(int* arr, int l, int r) {
    if (l >= r) return;
    int i, m = rand()%(r-l+1) + l;
	swap(arr[l], arr[m]);
	for(m = l, i = l + 1; i <= r; i++) {
		if (x[i] < x[l]) swap(arr[++m], arr[i]);
	}
    swap(arr[l], arr[m]);
    quicksort(arr, l, m-1);
    quicksort(arr, i, m+1);
}

int main() {
    int arr[8] = {110, 5, 10, 3 ,10, 100, 1, 23};
    quicksort(arr, 0, 7);
    for(int i=0; i<8; i++) {
        printf("%d ", arr[i]);
    }
    printf("\n");
    return 0;
}
\end{vncode}
\endgroup


\subsection{Longest Increasing Subsequence (LIS)}

O(n log n); sau mỗi bước, phần tử thứ k trong tập là kết thúc nhỏ nhất của dãy tăng độ dài k.

\codepath{Paradigms/lis.cpp}

\begingroup\scriptsize
\begin{vncode}
#include <set>
using namespace std;

/*
 * Longest Increasing Subsequence O(nlogn)
 * At a given iteration i, k-th element (1-indexed) in set is the
 * smallest element that has a size k increasing subsequence ending in it
 */

int lis(int arr[], int n) {
	multiset<int> s;
	multiset<int>::iterator it;
	for(int i = 0; i < n; i++) {
		s.insert(arr[i]);
		it = s.upper_bound(arr[i]); //non-decreasing
		//it = ++s.lower_bound(arr[i]); //strictly increasing
		if (it != s.end()) s.erase(it);
	}
	return s.size();
}

/*
 * TEST MATRIX
 */

#include <cstdio>
#define MAXN 500009

int a[MAXN], n;

int main() {
	while(scanf("%d", &n) != EOF) {
		for(int i = 0; i < n; i++) {
			scanf("%d", &a[i]);
		}
		printf("%d\n", lis(a, n));
	}
	return 0;
}
\end{vncode}
\endgroup


\subsection{Bài toán cặp gần nhất}

Tìm cặp điểm 2D gần nhất theo khoảng cách Euclid trong O(n log n).

\codepath{Computational Geometry/closestpair.cpp}

\begingroup\scriptsize
\begin{vncode}
#include <cmath>
#include <algorithm>
#define MAXN 100309
#define INF 1e+30
using namespace std;

/*
 * Closest point problem
 */

struct point {
    double x, y;
    point() { x = y = 0; }
    point(double _x, double _y) : x(_x), y(_y) {}
};

typedef pair<point, point> pp;

double dist(pp p) {
	double dx = p.first.x - p.second.x;
	double dy = p.first.y - p.second.y;
	return hypot(dx, dy);
}

point strip[MAXN];

pp closest(point *P, int l, int r) {
	if (r == l) return pp(point(INF, 0), point(-INF, 0));
	int m = (l + r) / 2, s1 = 0, s2;
	int midx = (P[m].x + P[m+1].x)/2;
	pp pl = closest(P, l, m);
	pp pr = closest(P, m+1, r);
	pp ans = dist(pl) > dist(pr) ? pr : pl;
	double d = dist(ans);
	for(int i = l; i <= m; i++) {
		if (midx - P[i].x < d) strip[s1++] = P[i];
	}
	s2 = s1;
	for(int i = m+1; i <= r; i++) {
		if (P[i].x - midx < d) strip[s2++] = P[i];
	}
	for(int j = 0, s = s1; j < s1; j++) {
		point p = strip[j];
		for(int i = s; i < s2; i++) {
			point q = strip[i];
			pp cur = pp(p, q);
			double dcur = dist(cur);
			if (d > dcur) {
				ans = cur; d = dcur;
			}
			if (q.y - p.y > d) break;
			if (p.y - q.y > d) s = i+1;
		}
	}
	int i = l, j = m+1, k = l;
	while(i <= m && j <= r) {
		if (P[i].y < P[j].y) strip[k++] = P[i++];
		else strip[k++] = P[j++];
	}
	while(i <= m) strip[k++] = P[i++];
	for(i = l; i < k; i++) P[i] = strip[i];
	return ans;
}

bool compx(point a, point b) {
    return a.x < b.x;
}

pp closest(point *P, int n){
	sort(P, P+n, compx);
	return closest(P, 0, n-1);
}

/*
 * Codeforces 101707F
 */
#include <cstdio>

point P[MAXN];
int n;

int main() {
	scanf("%d", &n);
	for(int i = 0; i < n; i++) scanf("%lf %lf", &P[i].x, &P[i].y);
	pp ans = closest(P, n);
	printf("%.0f %.0f\n", ans.first.x, ans.first.y);
	printf("%.0f %.0f\n", ans.second.x, ans.second.y);
	return 0;
}
\end{vncode}
\endgroup


\subsection{Tối ưu hai con trỏ}

Giảm một số DP từ O(n\textasciicircum{}2 k) xuống O(nk) khi chỉ số tối ưu dịch đơn điệu.

\codepath{Dynamic Programming/twopointeroptimization.cpp}

\begingroup\scriptsize
\begin{vncode}
#include <cstdio>
#include <algorithm>
using namespace std;
#define MAXN 1009
#define MAXK 19
#define INF (1<<30)

int dp[MAXN][MAXK], A[MAXN][MAXK], N, K;

void twopointer() {
	for(int i=0; i<=N; i++) dp[i][0] = INF;
	for(int j=0; j<=K; j++) dp[0][j] = 0, A[0][j] = 1;
	dp[0][0] = 0;
	for(int i=1; i<=N; i++) {
		for(int j=1; j<=K; j++) {
			dp[i][j] = INF;
			for(int k=A[i-1][j]; k<=i; k++) {
				int cur = 1 + max(dp[k-1][j-1], dp[i-k][j]);
				if (dp[i][j] > cur) {
					dp[i][j] = cur;
					A[i][j] = k;
				}
				if (dp[k-1][j-1] > dp[i-k][j]) break;
			}
		}
	}
}

/*
 *	Hacker Earth Eggs and Building
 */

int main() {
	int T, n, k;
	scanf("%d", &T);
	N = 1000;
	K = 10;
	twopointer();
	while(T-->0) {
		scanf("%d %d", &n, &k);
		printf("%d\n", dp[n][k]);
	}
	return 0;
}
\end{vncode}
\endgroup


\subsection{Tối ưu Convex Hull Trick}

Giảm một số DP dạng đường thẳng xuống O(nk log n).

\codepath{Dynamic Programming/chtrickoptimization.cpp}

\begingroup\scriptsize
\begin{vncode}
#include <cstdio>
#include <vector>
#define INF (1<<30)
#define MAXN 1009
using namespace std;

typedef long long ll;

class CHTrick{
private:
    vector<ll> m, n;
    vector<double> p;
    bool maxCH;
public:
    CHTrick(bool _maxCH) { maxCH = _maxCH; }
    void clear() {
        m.clear(); n.clear(); p.clear();
    }
    double inter(double nm, double nn, double lm, double ln) {
        return (ln - nn)/(nm - lm);
    }
    void push(ll nm, ll nn) {
        while(!p.empty() &&
            ((nm == m.back() && (maxCH?-1:1)*nn <= (maxCH?-1:1)*n.back()) ||
            p.back() >= inter(nm, nn, m.back(), n.back()))) {
            m.pop_back(); n.pop_back(); p.pop_back();
        }
        p.push_back(p.empty()? -INF : inter(nm, nn,
            m.back(), n.back()));
        m.push_back(nm); n.push_back(nn);
    }
    ll query(ll x) {
        if (p.empty()) return (maxCH?-1:1)*INF;
        double dx = (double)x;
        ll high = p.size()-1, low = 0, mid;
        if (dx >= p[high]) return m[high]*x + n[high];
        while(high > low+1) {
            mid = (high+low)/2;
            if (dx < p[mid]) high = mid;
            else low = mid;
        }
        return m[low]*x + n[low];
    }
};

ll x[MAXN], A[MAXN], dp[MAXN][MAXN];
int N, K;

void solve() {
    for(int i=0; i<=N; i++) dp[i][0] = 0;
    CHTrick cht(false);        
    for(int j=1; j<=K; j++) {
        dp[0][j] = 0;
        cht.clear();
        for(int i=1; i<=N; i++) {
            cht.push(A[i-1], dp[i-1][j-1]);
            dp[i][j] = cht.query(x[i]);
        }
    }
}

int main() {
    
    return 0;
}
\end{vncode}
\endgroup


\subsection{Tối ưu Slope Trick}

Giảm DP dạng trị tuyệt đối từ O(nS\textasciicircum{}2) xuống O(n log n).

\codepath{Dynamic Programming/slopetrick.cpp}

\begingroup\scriptsize
\begin{vncode}
#include <cstdio>
#include <queue>
#include <algorithm>
using namespace std;
#define MAXN 3009

typedef long long ll;
int N;
ll a[MAXN];

ll f[MAXN], opt[MAXN];

ll slope() {
    priority_queue<ll> pq;
    opt[0] = a[0]; f[0] = 0;
    pq.push(a[0]);
    for(int i=1; i<N; i++) {
        pq.push(a[i]);
        f[i] = f[i-1] + abs(a[i] - pq.top());
        if (a[i] < pq.top()) {
            pq.pop();
            pq.push(a[i]);
        }
        opt[i] = pq.top();
    }
    return f[N-1];
}

/*
 *  TEST MATRIX
 */

#define INF (1<<30)
ll brutef[MAXN][MAXN];
int brute() {
    ll M = 0;
    printf("arr: ");
    for(int i=0; i<N; i++) {
        M = max(M, a[i]+1LL);
        printf(" %3I64d", a[i]);
    }
    printf("\n");
    printf("x:   ", 0);
    for(ll x=0; x<=M; x++) {
        printf(" %3I64d", x);
    }
    printf("\n");
    printf("i=%2d:", 0);
    for(ll x=0; x<=M; x++) {
        brutef[0][x] = max(0LL, a[0]-x);
        printf(" %3I64d", brutef[0][x]);
    }
    printf("\n");
    for(int i=1; i<N; i++) {
        printf("i=%2d:", i);
        for(ll x=0; x<=M; x++) {
            brutef[i][x] = INF;
            for(ll y=0; y<=x; y++) {
                brutef[i][x] = min(brutef[i][x], brutef[i-1][y] + abs(y-a[i]));
            }
            printf(" %3I64d", brutef[i][x]);
        }
        printf("\n");
    }
}

/*
 *  Codeforces 713C
 */

int main() {
    scanf("%d", &N);
    for(int i=0; i<N; i++) {
        scanf("%I64d", &a[i]);
        a[i] += N-i;
    }
    printf("%I64d\n", slope());
    return 0;
}
\end{vncode}
\endgroup


\subsection{Tối ưu chia để trị}

Giảm DP có điểm tối ưu đơn điệu từ O(n\textasciicircum{}2 k) xuống O(nk log n).

\codepath{Dynamic Programming/dcoptimization.cpp}

\begingroup\scriptsize
\begin{vncode}
#include <cstdio>
#include <algorithm>
using namespace std;
#define MAXN 1009
#define INF (1 << 30)

int dp[MAXN][MAXN], C[MAXN][MAXN], N, K;

void calculatedp(int min_i, int max_i, int j, int min_k, int max_k) {
	if (min_i > max_i) return;
	int i = (min_i + max_i)/2;
	int ans = INF, opt;
	for(int k=min_k; k<=min(max_k, i-1); k++) {
		if (ans > dp[k][j-1] + C[k][i]) {
			opt = k;
			ans = dp[k][j-1] + C[k][i];
		}
	}
	dp[i][j] = ans;
	calculatedp(min_i, i-1, j, min_k, opt);
	calculatedp(i+1, max_i, j, opt, max_k);
}

void solve() {
	for(int i=0; i<=N; i++) dp[i][0] = 0;
	for(int j=0; j<=K; j++) dp[0][j] = 0;
	for(int j=1; j<=K; j++) {
		calculatedp(1, N, j, 0, N-1);
	}
}

int main() {

}
\end{vncode}
\endgroup


\subsection{Tối ưu Knuth}

Giảm DP dạng chia đoạn từ O(n\textasciicircum{}3) xuống O(n\textasciicircum{}2) khi thỏa bất đẳng thức tứ giác/đơn điệu.

\codepath{Dynamic Programming/knuthoptimization.cpp}

\begingroup\scriptsize
\begin{vncode}
#include <cstdio>
#define MAXN 1009
#define INF (1LL << 60)

typedef long long ll;

ll dp[MAXN][MAXN], C[MAXN][MAXN], M, x[MAXN];
int A[MAXN][MAXN], N, S;

void knuth() {
	ll cur;
	for(int s = 0; s < N; s++) {
		for(int i=0, j; i + s<N; i++) {
			j = i + s;
			if (s < S) {	//Condições iniciais
				dp[i][j] = 0;
				A[i][j] = i;
				continue;
			}
			dp[i][j] = INF;
			for(int k = A[i][j-1]; k <= A[i+1][j]; k++) {
				cur = C[i][j] + dp[i][k] + dp[k][j];
				if (dp[i][j] > cur) {
					dp[i][j] = cur;
					A[i][j] = k;
				}
			}
		}
	}
}

int main() {
	while(scanf("%lld %d", &M, &N)!=EOF) {
		for(int i=1; i<=N; i++) {
			scanf("%lld", &x[i]);
		}
		x[0] = 0;
		x[N+1] = M;
		N += 2;
		S = 2;
		for(int i=0; i<N; i++) {
			for(int j=0; j<N; j++) {
				C[i][j] = x[j] - x[i];
			}
		}
		knuth();
		printf("%lld\n", dp[0][N-1]);
	}
	return 0;
}
\end{vncode}
\endgroup


\subsection{Tối ưu Lagrange}

Thay ràng buộc số lần chọn bằng chi phí phạt C rồi tìm nhị phân C khi hàm mục tiêu lồi.

\codepath{Dynamic Programming/lagrangeoptimization.cpp}

\begingroup\scriptsize
\begin{vncode}
#include <cstdio>
#include <vector>
#define INF 0x3f3f3f3f
#define INFLL 0x3f3f3f3f3f3f3f3f
#define MAXN 1009
using namespace std;

typedef long long ll;

/*
 * Convex Hull Trick
 */

class CHTrick {
private:
	vector<ll> m, n;
	vector<double> p;
	vector<int> ids;
	bool maxCH; int i;
public:
	CHTrick(bool mxch) { maxCH = mxch; i = 0; }
	void clear() { m.clear(); n.clear(); p.clear(); ids.clear(); }
	double inter(double nm, double nn, double lm, double ln) {
		return (ln - nn) / (nm - lm);
	}
	void push(ll nm, ll nn, int id) {
		while (!p.empty()) {
			if (nm == m.back() && maxCH && nn <= n.back()) return;
			if (nm == m.back() && !maxCH && nn >= n.back()) return;
			double x = inter(nm, nn, m.back(), n.back());
			if (nm != m.back() && p.back() < x) break;
			m.pop_back(); n.pop_back(); p.pop_back(); ids.pop_back();
		}
		p.push_back(p.empty() ? -INF : inter(nm, nn, m.back(), n.back()));
		m.push_back(nm); n.push_back(nn); ids.push_back(id);
		if (i >= p.size()) i = p.size()-1;
	}
	pair<ll, int> query_q(ll x) {
		while (p[i] > x) i--;
		while(i+1 < p.size() && p[i+1] <= x) i++;
		return make_pair(m[i] * x + n[i], ids[i]);
	}
};

/*
 * URI 1282
 */

ll w[MAXN], x[MAXN], A[MAXN], B[MAXN];
ll dp[MAXN];
int opt[MAXN], N, K;

CHTrick cht(false);

int computeDp(ll C) {
	cht.clear();
	dp[0] = 0;
	opt[0] = 0;
	for(int i = 1; i <= N; i++) {
		cht.push(-A[i-1], dp[i-1] + B[i-1], i-1);
		pair<ll, int> cur = cht.query_q(x[i]);
		dp[i] = C + x[i]*A[i] - B[i] + cur.first;
		opt[i] = 1 + opt[cur.second];
	}
	return opt[N];
}

ll computeC() {
	ll lo = 0;
	ll hi = INFLL;
	while(hi > lo + 1) {
		ll C = (hi + lo) / 2;
		int k = computeDp(C);
		if (k > K) lo = C;
		else hi = C;
	}
	return hi;
}

int main(){
	while (scanf("%d %d", &N, &K) != EOF){
		A[0] = B[0] = 0;
		for (int i=1; i<=N; i++) {
			scanf("%lld %lld", &x[i], &w[i]);
			A[i] = w[i] + A[i-1];
			B[i] = w[i]*x[i] + B[i-1];
		}
		ll C = computeC();
		int k = computeDp(C);
		printf("%lld\n", dp[N] - C*K);
	}
	return 0;
}
\end{vncode}
\endgroup


\clearpage

\section{Đồ thị}

\subsection{DFS Spanning Tree}

Cây khung DFS.

\codepath{Graphs/dfsspanningtree.cpp}

\begingroup\scriptsize
\begin{vncode}
#include <cstdio>
#include <vector>
using namespace std;
#define MAXN 1009
#define UNVISITED -1
#define EXPLORED -2
#define VISITED -3

int vis[MAXN], parent[MAXN];
vector<int> adjList[MAXN];

void dfs(int u) { // DFS for checking graph edge properties
	vis[u] = EXPLORED;
	for (int j = 0, v; j < (int)adjList[u].size(); j++) {
		v = adjList[u][j];
		if (vis[v] == UNVISITED) { // EXPLORED->UNVISITED
			printf(" Tree Edge (%d, %d)\n", u, v);
			parent[v] = u; // parent of this children is me
			dfs(v);
		}
		else if (vis[v] == EXPLORED) { // EXPLORED->EXPLORED
			printf(" Back Edge (%d, %d) (Cycle)\n", u, v);
		}
		else if (vis[v] == VISITED) // EXPLORED->VISITED
			printf(" Forward/Cross Edge (%d, %d)\n", u, v);
	}
	vis[u] = VISITED; // after recursion, color u as VISITED (DONE)
}
\end{vncode}
\endgroup


\subsection{Điểm khớp và cầu}

Tìm articulation points và bridges.

\codepath{Graphs/articulationpointsandbridges.cpp}

\begingroup\scriptsize
\begin{vncode}
#include <cstdio>
#include <vector>
#include <algorithm>
#include <cstring>
using namespace std;
#define MAXN 1009
#define UNVISITED -1

int num[MAXN], N, low[MAXN], parent[MAXN], counter, rootChildren, articulationVertex[MAXN], root;
vector<int> adjList[MAXN];

void tarjan(int u) {
	low[u] = num[u] = counter++;
	for (int j = 0, v; j < (int)adjList[u].size(); j++) {
		v = adjList[u][j];
		if (num[v] == UNVISITED) {
			parent[v] = u;
			if (u == root) rootChildren++;
			tarjan(v);
			if (low[v] >= num[u]) articulationVertex[u] = true;
			if (low[v] > num[u]) printf(" Edge (%d, %d) is a bridge\n", u, v);
			low[u] = min(low[u], low[v]);
		}
		else if (v != parent[u])
			low[u] = min(low[u], num[v]);
	}
}

int main() {
	counter = 0;
	memset(&num, UNVISITED, sizeof num);
	memset(&low, 0, sizeof low);
	memset(&parent, 0, sizeof parent);
	memset(&articulationVertex, 0, sizeof articulationVertex);
	printf("Bridges:\n");
	for (int i = 0; i < N; i++)
		if (num[i] == UNVISITED) {
			root = i; rootChildren = 0; tarjan(i);
			articulationVertex[root] = (rootChildren > 1);
		} // special case
	printf("Articulation Points:\n");
	for (int i = 0; i < N; i++)
		if (articulationVertex[i])
			printf(" Vertex %d\n", i);
	return 0;
}
\end{vncode}
\endgroup


\subsection{Sắp xếp tô pô}

Topological sort.

\codepath{Graphs/toposort.cpp}

\begingroup\scriptsize
\begin{vncode}
#include <cstdio>
#include <vector>
using namespace std;
#define MAXN 1009

int vis[MAXN];
vector<int> adjList[MAXN];
vector<int> toposort;	//Ordem reversa!

void ts(int u) {
	vis[u] = true;
	for (int j = 0, v; j < (int)adjList[u].size(); j++) {
		v = adjList[u][j];
		if (!vis[v]) ts(v);
	}
	toposort.push_back(u);
}
\end{vncode}
\endgroup


\subsection{Thành phần liên thông mạnh: thuật toán Kosaraju}

SCC bằng Kosaraju.

\codepath{Graphs/kosaraju.cpp}

\begingroup\scriptsize
\begin{vncode}
#include <cstdio>
#include <vector>
#include <cstring>
#define MAXN 100009
using namespace std;

/*
 * Kosaraju SCC O(n)
 */

vector<int> adjList[MAXN], revAdjList[MAXN], ts;
bool vis[MAXN];
int comp[MAXN], parent = 0, numSCC;

void revdfs(int u) {
    vis[u] = true;
    for(int i = 0, v; i < (int)revAdjList[u].size(); i++) {
    	v = revAdjList[u][i];
        if(!vis[v]) revdfs(v);
    }
    ts.push_back(u);
}

void dfs(int u) {
    vis[u] = true; comp[u] = parent;
    for(int i = 0, v; i < (int)adjList[u].size(); i++) {
    	v = adjList[u][i];
        if(!vis[v]) dfs(v);
    }
}

void kosaraju(int n) {
	memset(&vis, false, sizeof vis);
    for(int i = 0; i < n; i++)  {
        if(!vis[i]) revdfs(i);
    }
    memset(&vis, false, sizeof vis);
    numSCC = 0;
    for(int i = n-1; i >= 0; i--) {
        if(!vis[ts[i]]) {
            parent = ts[i];
            dfs(ts[i]);
            numSCC++;
        }
    }
}

/*
 * TEST MATRIX
 */

int leader[MAXN];

int main() {
    int u, v, N, M;
	scanf("%d %d", &N, &M);
	for(int i=0; i<M; i++) {
		scanf("%d %d", &u, &v);
		adjList[u].push_back(v);
		revAdjList[v].push_back(u);
	}
    kosaraju(N);
    memset(&leader, -1, sizeof leader);
	for(int i=0; i<N; i++) {
		if (leader[comp[i]] == -1 || leader[comp[i]] > i) {
			leader[comp[i]] = i;
		}
	}
	for(int i=0; i<N; i++) printf("%d\n", leader[comp[i]]);
    return 0;
}
\end{vncode}
\endgroup


\subsection{2-SAT: thuật toán Tarjan}

Giải 2-SAT bằng SCC.

\codepath{Graphs/2sat.cpp}

\begingroup\scriptsize
\begin{vncode}
#include <cstdio>
#include <vector>
#include <stack>
#include <algorithm>
#include <cstring>
using namespace std;
#define MAXN 1000009
#define UNVISITED -1

/*
 * Tarjan SCC O(n)
 */

int num[MAXN], vis[MAXN], comp[MAXN], low[MAXN];
int counter, numSCC;
stack<int> st;
vector<int> adjList[MAXN];

void dfs(int u) {
	low[u] = num[u] = counter++;
	st.push(u); vis[u] = 1;
	for (int j = 0; j < (int)adjList[u].size(); j++) {
		int v = adjList[u][j];
		if (num[v] == UNVISITED) dfs(v);
		if (vis[v]) low[u] = min(low[u], low[v]);
	}
	if (low[u] == num[u]) {
		while (true) {
			int v = st.top(); st.pop();
			vis[v] = 0;
			comp[v] = numSCC;
			if (u == v) break;
		}
		numSCC++;
	}
}

void tarjan(int n) {
	counter = numSCC = 0;
	memset(&num, UNVISITED, sizeof num);
	memset(&vis, 0, sizeof vis);
	memset(&low, 0, sizeof low);
	for (int i = 0; i < n; i++) {
		if (num[i] == UNVISITED)
			dfs(i);
	}
}

/*
 * 2-SAT O(n)
 */

bool twosat(bool x[], int n) {
	tarjan(2*n);
	for(int i = 0; i < n; i++) {
		if (comp[2*i] == comp[2*i+1]) return false;
		x[i] = (comp[2*i] > comp[2*i+1]);
	}
	return true;
}

void implies(int u, bool pu, int v, bool pv) {
	adjList[2*u+(pu?1:0)].push_back(2*v+(pv?1:0));
	adjList[2*v+(pv?0:1)].push_back(2*u+(pu?0:1));
}

/*
 * TEST MATRIX
 */

#include <cassert>

bool checked[MAXN];
bool checkdfs(int u, bool x[]) {
	if (checked[u]) return true;
	checked[u] = true;
	printf("  therefore x[%d]=%d\n", u/2, u%2);
	if ((u%2 != 0) != x[u/2]) return false;
	for(int i = 0; i < (int)adjList[u].size(); i++) {
		int v = adjList[u][i];
		if (!checkdfs(v, x)) return false;
	}
	return true;
}

bool check(bool x[], int n) {
	memset(&checked, false, sizeof checked);
	for(int i = 0; i < n; i++) {
		printf("x[%d]=%d:\n", i, x[i]);
		if (x[i] && !checkdfs(2*i+1, x)) return false;
		if (!x[i] && !checkdfs(2*i, x)) return false;
	}
	return true;
}

bool reached[MAXN];
void reach(int u, int n) {
	if (reached[u]) return;
	reached[u] = true;
	if (u < 2*n) printf(" x[%d]=%d", u/2 + 1, u%2);
	for(int i = 0; i < (int)adjList[u].size(); i++) {
		int v = adjList[u][i];
		reach(v, n);
	}
}

void print(int n) {
	for(int i = 0; i < 2*n; i++) {
		printf("x[%d]=%d ->", i/2 + 1, i%2);
		memset(&reached, false, sizeof reached);
		reach(i, n);
		printf("\n");
	}
}

/*
 * Brazilian ICPC Summer Camp 2019, Group W, Day 2
 * Problem B
 */

vector<int> pref[MAXN];
bool x[MAXN];

int main() {
	int n, m, k;
	scanf("%d %d %d", &n, &m, &k);
	while(k --> 0) {
		int a, b;
		scanf("%d %d", &a, &b);
		pref[a].push_back(b);
	}
	int cnt = m;
	for(int a = 1; a <= n; a++) {
		int sz = pref[a].size();
		vector<int> pre(sz), suf(sz);
		for(int i = 0; i < sz; i++) {
			pre[i] = cnt++; suf[i] = cnt++;
		}
		for(int i = 0; i < sz; i++) {
			int u = pref[a][i];
			if (i > 0) {
				implies(pre[i], true, pre[i-1], true);
				implies(abs(u)-1, u < 0, pre[i-1], true);
			}
			if (i+1 < sz) {
				implies(suf[i], true, suf[i+1], true);
				implies(abs(u)-1, u < 0, suf[i+1], true);
			}
			implies(pre[i], true, abs(u)-1, u > 0);
			implies(suf[i], true, abs(u)-1, u > 0);
		}
	}
	//print(m);
	if (twosat(x, cnt)) {
		//assert(check(x, cnt));
		vector<int> ans;
		for(int i = 0; i < m; i++) {
			if (x[i]) ans.push_back(i+1);
		}
		printf("%u\n", ans.size());
		for(int i = 0; i < (int)ans.size(); i++) {
			printf("%d ", ans[i]);
		}
		printf("\n");
	}
	else printf("-1\n");
	return 0;
}
\end{vncode}
\endgroup


\subsection{Đường đi ngắn nhất: thuật toán Dijkstra}

Dijkstra.

\codepath{Graphs/dijkstra.cpp}

\begingroup\scriptsize
\begin{vncode}

#include <set>
#include <vector>
using namespace std;
#define MAXN 100009
#define INF 0x3f3f3f3f
typedef pair<int, int> ii;

/*
 * Dijkstra's Algorithm O(nlogn + m)
 */
 
vector<ii> adjList[MAXN];
 
int dijkstra(int s, int t, int n, int dist[]) {
	for(int i = 1; i <= n; i++) dist[i] = INF;
	set<ii> pq;
	dist[s] = 0;
	pq.insert(ii(0, s));
	while(!pq.empty()) {
		int u = pq.begin()->second;
		pq.erase(pq.begin());
		for(int i=0; i<(int)adjList[u].size(); i++) {
			int v = adjList[u][i].second;
			int w = adjList[u][i].first;
			if (dist[v] > dist[u] + w) {
				pq.erase(ii(dist[v], v));
				dist[v] = dist[u] + w;
				pq.insert(ii(dist[v], v));
			}
		}
	}
	return dist[t];
}

/*
 * SPOJ EZDIJKST
 */

#include <cstdio>

int dist[MAXN], n, m;
 
int main() {
	int T;
	int u, v, w, s, t;
	scanf("%d", &T);
	while(T--) {
		scanf("%d %d", &n, &m);
		for(int i=1; i<=n; i++) adjList[i].clear();
		for(int i=0; i<m; i++) {
			scanf("%d %d %d", &u, &v, &w);
			adjList[u].push_back(ii(w, v));
		}
		scanf("%d %d", &s, &t);
		int ans = dijkstra(s, t, n, dist);
		if (ans == INF) printf("NO\n");
		else printf("%d\n", ans);
	}
	return 0;
}
\end{vncode}
\endgroup


\subsection{Đường đi ngắn nhất: thuật toán Floyd-Warshall}

Floyd-Warshall.

\codepath{Graphs/floydwarshall.cpp}

\begingroup\scriptsize
\begin{vncode}
#include <cstidio>
#define MAXN 409

int adjMat[MAXN][MAXN], N;

//O(V^3)
void floydwarshall() {
	for (int k = 0; k < N; k++)
		for (int i = 0; i < N; i++)
			for (int j = 0; j < N; j++)
				adjMat[i][j] = min(adjMat[i][j], adjMat[i][k] + adjMat[k][j]);
}
\end{vncode}
\endgroup


\subsection{Đường đi ngắn nhất: thuật toán Bellman-Ford}

Bellman-Ford.

\codepath{Graphs/bellmanford.cpp}

\begingroup\scriptsize
\begin{vncode}
#include <cstdio>
#include <vector>
#include <cstring>
#include <iostream>
#define MAXN 1009
using namespace std;

int dist[MAXN], N;
typedef pair<int, int> ii;
vector<ii> adjList[MAXN];

//O(VE)
int bellmanford(int s, int t)
	memset(&dist, 1<<20, sizeof dist);
	dist[s] = 0;
	bool hasNegativeWeightCycle = false;
	for (int i = 0, v, w; i < N; i++) {
		for (int u = 0; u < N; u++) {
			for (int j = 0; j < (int)adjList[u].size(); j++) {
				v = adjList[u][j].first;
				w = adjList[u][j].second;
				if (i==N-1 && dist[v] > dist[u] + w)
					hasNegativeWeightCycle = true;
				else dist[v] = min(dist[v], dist[u] + w);
			}
		}
	}
	return dist[t];
}
\end{vncode}
\endgroup


\subsection{Đường đi ngắn nhất thứ k}

K-shortest path.

\codepath{Graphs/kdijkstra.cpp}

\begingroup\scriptsize
\begin{vncode}
#include <vector>
#include <map>
using namespace std;
#define MAXN 109
typedef pair<int, int> ii;

/*
 * k-Sijkstra O(nklogn)
 */

vector<ii> adjList[MAXN];

void dijkstra(int s, int k, vector<int> kdist[]) {
	map<ii, int> cnt;
	cnt[ii(0, s)]++;
	while(!cnt.empty()) {
		pair<ii, int> cur = *cnt.begin();
		int u = cur.first.second;
		int wu = cur.first.first;
		int am = cur.second;
		cnt.erase(cnt.begin());
		if (am == 0 || kdist[u].size() >= k) continue;
		for(int i = 0; i < am; i++) {
			if (wu == 0 || kdist[u].size() >= k) break;
			kdist[u].push_back(wu);
		}
		for(int i = 0; i < (int)adjList[u].size(); i++) {
			int v = adjList[u][i].second;
			int w = adjList[u][i].first;
			int &cv = cnt[ii(w + wu, v)];
			cv = min(cv + am, k);
		}
	}
}

/*
 * Szkoput Connections
 */

#include <cstdio>
vector<int> ans[MAXN][MAXN];

int main() {
	int n, m;
	scanf("%d %d", &n, &m);
	while(m --> 0) {
		int u, v, w;
		scanf("%d %d %d", &u, &v, &w);
		adjList[u].push_back(ii(w, v));
	}
	for(int s = 1; s <= n; s++)
		dijkstra(s, 100, ans[s]);
	int q;
	scanf("%d\n", &q);
	while(q --> 0) {
		int s, t, k;
		scanf("%d %d %d", &s, &t, &k);
		if (k-1 >= ans[s][t].size()) printf("-1\n");
		else printf("%d\n", ans[s][t][k-1]);
	}
	return 0;
}
\end{vncode}
\endgroup


\subsection{Đường đi ngắn nhất: Shortest Path Faster Algorithm (SPFA)}

SPFA.

\codepath{Graphs/spfa.cpp}

\begingroup\scriptsize
\begin{vncode}
#include <queue>
#include <vector>
#include <cstring>
#include <iostream>
using namespace std;
#define MAXN 100009
#define INF 0x3f3f3f3f
typedef pair<int, int> ii;

/*
 * Shortest Path Faster Algorithm (SPFA) O(nm) worst case, O(n) average
 */

vector<ii> adjList[MAXN];
int dist[MAXN], vis[MAXN], N, M;
bool inq[MAXN];
 
int spfa(int s, int t) {
	for(int i=0; i<=N; i++) dist[i] = INF;
	memset(&inq, false, sizeof inq);
	memset(&vis, 0, sizeof vis);
	queue<int> q;
	q.push(s);
	dist[s] = 0;
	inq[s] = true;
	while (!q.empty()) {
		int u = q.front(); q.pop();
		if (vis[u] > N) return -1;
		inq[u] = false;
		for (int i = 0; i < (int)adjList[u].size(); i++) {
			int v = adjList[u][i].second;
			int w = adjList[u][i].first;
			if (dist[u] + w < dist[v]) {
				dist[v] = dist[u] + w;
				if (!inq[v]) {
					vis[v]++; q.push(v);
					inq[v] = true;
				}
			}
		}
	}
	return dist[t];
}

/*
 * SPOJ EZDIJKST
 */

#include <cstdio>

int main() {
	int T;
	int u, v, w, s, t;
	scanf("%d", &T);
	while(T--) {
		scanf("%d %d", &N, &M);
		for(int i=1; i<=N; i++) adjList[i].clear();
		for(int i=0; i<M; i++) {
			scanf("%d %d %d", &u, &v, &w);
			adjList[u].push_back(ii(w, v));
		}
		scanf("%d %d", &s, &t);
		int ans = spfa(s, t);
		if (ans == INF) printf("NO\n");
		else printf("%d\n", ans);
	}
	return 0;
}
\end{vncode}
\endgroup


\subsection{Cây khung nhỏ nhất: thuật toán Kruskal}

Kruskal.

\codepath{Graphs/kruskal.cpp}

\begingroup\scriptsize
\begin{vncode}
#include <cstdio>
#include <vector>
#include <algorithm>
using namespace std;

class UnionFind {
private:
	vector<int> parent, rank;
public:
	UnionFind(int N) {
		rank.assign(N+9, 0);
		parent.assign(N+9, 0);
		for (int i = 0; i < N; i++) parent[i] = i;
	}
	int find(int i) {
		while(i != parent[i]) i = parent[i];
		return i;
	}
	bool isSameSet(int i, int j) {
		return find(i) == find(j);
	}
	void unionSet (int i, int j) {
		if (isSameSet(i, j)) return;
		int x = find(i), y = find(j);
		if (rank[x] > rank[y]) parent[y] = x;
		else {
			parent[x] = y;
			if (rank[x] == rank[y]) rank[y]++;
		}
	}
};

typedef pair<int, int> ii;
typedef long long ll;

int N, M;
vector< pair<ll, ii> > edgeList; // (weight, two vertices) of the edge

ll kruskal() {
	ll cost = 0;
	UnionFind UF(N);
	pair<int, ii> edge;
	sort(edgeList.begin(), edgeList.end());
	for (int i = 0; i < M; i++) {
		edge = edgeList[i];
		if (!UF.isSameSet(edge.second.first, edge.second.second)) { 
			cost += edge.first;
			UF.unionSet(edge.second.first, edge.second.second);
		}
	}
	return cost;
}

int main() {
	scanf("%d %d", &N, &M);
	int u, v; ll w;
	for (int i = 0; i < M; i++) {
		scanf("%d %d %lld", &u, &v, &w);
		edgeList.push_back(make_pair(w, ii(u, v)));
	}
	printf("%lld\n", kruskal());
	return 0;
}
\end{vncode}
\endgroup


\subsection{Cây khung nhỏ nhất: thuật toán Prim}

Prim.

\codepath{Graphs/prim.cpp}

\begingroup\scriptsize
\begin{vncode}
#include <cstdio>
#include <vector>
#include <queue>
#include <cstring>
#define MAXN 10009
using namespace std;

typedef long long ll;
typedef pair<int, int> ii;
vector<ii> adjList[MAXN];

int N, M;

ll prim() {
	bool taken[MAXN];
	memset(&taken, false, sizeof taken);
	taken[0] = true;
	priority_queue<ii> pq;
	ii v, front; int u, w; ll cost = 0;
	for (int j = 0; j < (int)adjList[0].size(); j++) {
		v = adjList[0][j];
		pq.push(ii(-v.second, -v.first));
	}
	while (!pq.empty()) {
		front = pq.top(); pq.pop();
		u = -front.second; w = -front.first;
		if (!taken[u]) {
			cost += (ll)w; taken[u] = true;
			for (int j = 0; j < (int)adjList[u].size(); j++) {
				v = adjList[u][j];
				if (!taken[v.first]) pq.push(ii(-v.second, -v.first));
			}
		}
	}
	return cost;
}

int main() {
	scanf("%d %d", &N, &M);
	int u, v, w;
	for (int i = 0; i < M; i++) {
		scanf("%d %d %d", &u, &v, &w);
		adjList[u-1].push_back(make_pair(v-1, w));
		adjList[v-1].push_back(make_pair(u-1, w));
	}
	printf("%lld\n", prim());
	return 0;
}
\end{vncode}
\endgroup


\subsection{Luồng cực đại: thuật toán Edmonds-Karp}

Edmonds-Karp.

\codepath{Graphs/edmondskarp.cpp}

\begingroup\scriptsize
\begin{vncode}
#include <queue>
#include <cstring>
#define INF 0x3f3f3f3f
#define MAXN 103000
#define MAXM 900000
using namespace std;

/*
 * Edmonds-Karp's Algorithm - O(VE^2)
 */

int ned, prv[MAXN], first[MAXN];
int cap[MAXM], to[MAXM], nxt[MAXM], dist[MAXN];

void init() {
   memset(first, -1, sizeof first);
   ned = 0;
}

void add(int u, int v, int f) {
    to[ned] = v, cap[ned] = f;
    nxt[ned] = first[u];
    first[u] = ned++;
    
    to[ned] = u, cap[ned] = 0;
    nxt[ned] = first[v];
    first[v] = ned++;
}

int augment(int v, int minEdge, int s) {
	int e = prv[v];
	if (e == -1) return minEdge;
	int f = augment(to[e^1], min(minEdge, cap[e]), s);
	cap[e] -= f;
	cap[e^1] += f;
	return f;
}

bool bfs(int s, int t) {
	int u, v;
	memset(&dist, -1, sizeof dist);
	dist[s] = 0;
	queue<int> q; q.push(s);
	memset(&prv, -1, sizeof prv);
	while (!q.empty()) {
		u = q.front(); q.pop();
		if (u == t) break;
		for(int e = first[u]; e!=-1; e = nxt[e]) {
			v = to[e];
			if (dist[v] < 0 && cap[e] > 0) {
				dist[v] = dist[u] + 1;
				q.push(v);
				prv[v] = e;
			}
		}
	}
	return dist[t] >= 0;
}

int edmondskarp(int s, int t) {
	int result = 0;
	while (bfs(s, t)) {
		result += augment(t, INF, s);
	}
	return result;
}

/*
 *	Codeforces 101128F
 */
#include <cstdio>

int main() {
	int N, M, R, C, A, B, pos[60][60];
	char c;
	scanf("%d %d %d %d", &R, &C, &A, &B);
	init();
	memset(&cap, 0, sizeof cap);
	N = R*C+2;
	for(int i=0, k=1; i<R; i++) {
		for(int j=0; j<C; j++) pos[i][j] = k++;
	}
	for(int i=0; i<R; i++) {
		for(int j=0; j<C; j++) {
			scanf(" %c", &c);
			if (c == '.') add(0, pos[i][j], B);
			else add(pos[i][j], N-1, B);
			if (i < R-1) {
				add(pos[i][j], pos[i+1][j], A);
				add(pos[i+1][j], pos[i][j], A);
			}
			if (j < C-1) {
				add(pos[i][j], pos[i][j+1], A);
				add(pos[i][j+1], pos[i][j], A);
			}
		}
	}
	printf("%d\n", edmondskarp(0, N-1));
	return 0;
}
\end{vncode}
\endgroup


\subsection{Luồng cực đại: thuật toán Dinic}

Dinic.

\codepath{Graphs/dinic.cpp}

\begingroup\scriptsize
\begin{vncode}
#include <queue>
#include <cstring>
#define INF 0x3f3f3f3f
#define MAXN 103000
#define MAXM 900000
using namespace std;

/*
 * Dinic's Algorithm - O(EV^2)
 */

int ned, first[MAXN], work[MAXN], dist[MAXN], q[MAXN];
int cap[MAXM], to[MAXM], nxt[MAXM];

void init() {
   memset(first, -1, sizeof first);
   ned = 0;
}

void add(int u, int v, int f) {
	to[ned] = v, cap[ned] = f;
	nxt[ned] = first[u];
	first[u] = ned++;
	
	to[ned] = u, cap[ned] = 0;
	nxt[ned] = first[v];
	first[v] = ned++;
}

int dfs(int u, int f, int t) {
	if (u == t) return f;
	for(int &e = work[u]; e != -1; e = nxt[e]) {
		int v = to[e];
		if (dist[v] == dist[u] + 1 && cap[e] > 0) {
			int df = dfs(v, min(f, cap[e]), t);
			if (df > 0) {
				cap[e] -= df;
				cap[e^1] += df;
				return df;
			}
		}
	}
	return 0;
}

bool bfs(int s, int t) {
	memset(&dist, -1, sizeof dist);
	dist[s] = 0;
	int st = 0, en = 0;
	q[en++] = s;
	while (en > st) {
		int u = q[st++];
		for(int e = first[u]; e!=-1; e = nxt[e]) {
			int v = to[e];
			if (dist[v] < 0 && cap[e] > 0) {
				dist[v] = dist[u] + 1;
				q[en++] = v;
			}
		}
	}
	return dist[t] >= 0;
}

int dinic(int s, int t) {
	int result = 0, f;
	while (bfs(s, t)) {
		memcpy(work, first, sizeof work);
		while (f = dfs(s, INF, t)) result += f;
	}
	return result;
}

/*
 * Codeforces 100783D
 */

#include <cstdio>

int main() {
	int u, v, N, M;
	init();
	memset(&cap, 0, sizeof cap);
	scanf("%d %d", &N, &M);
	for(int i=1; i<=N; i++) {
		add(0, i, 1);
		add(i+N, 2*N+1, 1);
	}
	for(int i=0; i<M; i++) {
		scanf("%d %d", &u, &v);
		u++; v++;
		add(u, v+N, 1);
	}
	int MF = dinic(0, 2*N+1);
	//printf("MF = %d\n", MF);
	if (MF == N) printf("YES\n");
	else printf("NO\n");
	return 0;
}
\end{vncode}
\endgroup


\subsection{Luồng cực đại: thuật toán Push-Relabel}

Push-Relabel.

\codepath{Graphs/pushrelabel.cpp}

\begingroup\scriptsize
\begin{vncode}
#include <cstring>
#include <cstdio>
#define MAXN 103009
#define MAXM 900000
#define INF 0x3f3f3f3f

/*
 * Push Relabel Algorithm - O(VE + sqrt(E)V^2)
 */

int ned, first[MAXN], exc[MAXN], maxhgt[MAXN];
int cap[MAXM], to[MAXM], nxt[MAXM], hgt[MAXN];

void init() {
   memset(first, -1, sizeof first);
   ned = 0;
}

void add(int u, int v, int f) {
	to[ned] = v, cap[ned] = f;
	nxt[ned] = first[u];
	first[u] = ned++;
	
	to[ned] = u, cap[ned] = 0;
	nxt[ned] = first[v];
	first[v] = ned++;
}

void push(int e) {
	int u = to[e^1], v = to[e];
	int f = exc[u] < cap[e] ? exc[u] : cap[e];
	cap[e] -= f; cap[e^1] += f;
	exc[u] -= f; exc[v] += f;
}

void relabel(int u) {
	int f = INF;
	for (int e = first[u]; e != -1; e = nxt[e]) {
		if (cap[e] > 0 && hgt[to[e]] < f) f = hgt[to[e]];
	}
	if (f < INF) hgt[u] = f + 1;
}

int maxHeight(int s, int t, int n) {
	int sz = 0;
	for (int u = 1; u <= n; u++) {
		if (u == s || u == t || exc[u] == 0) continue;
		if (sz > 0 && hgt[u] > hgt[maxhgt[0]]) sz = 0;
		if (sz == 0 || hgt[u] == hgt[maxhgt[0]])
			maxhgt[sz++] = u;
	}
	return sz;
}

int pushrelabel(int s, int t, int n) {
	memset(&hgt, 0, sizeof hgt);
	memset(&exc, 0, sizeof exc);
	hgt[s] = n; exc[s] = INF;
	for (int e = first[s]; e != -1; e = nxt[e])
		push(e);
	int sz;
	while ((sz = maxHeight(s, t, n)) > 0) {
		for(int i = 0; i < sz; i++) {
			int u = maxhgt[i], pushed = 0;
			for (int e = first[u]; e != -1 && exc[u]; e = nxt[e]) {
				if (cap[e] > 0 && hgt[u] == hgt[to[e]] + 1)
					push(e), pushed = 1;
			}
			if (!pushed) {
				relabel(u); break;
			}
		}
	}
	return exc[t];
}

/*
 * SPOJ MTOTALF
 */

#include <cstdio>

int main() {
	int m;
	scanf("%d", &m);
	init();
	for(int j = 0; j < m; j++) {
		char u, v;
		int f;
		scanf(" %c %c %d", &u, &v, &f);
		add(u, v, f);
	}
	printf("%d\n", pushrelabel('A', 'Z', 255));
	return 0;
}
\end{vncode}
\endgroup


\subsection{Cắt toàn cục nhỏ nhất: thuật toán Stoer-Wagner}

Stoer-Wagner.

\codepath{Graphs/stoerwagner.cpp}

\begingroup\scriptsize
\begin{vncode}
#include <cstdio>
#include <vector>
#include <cstring>
using namespace std;
#define MAXN 509
#define INF 0x3f3f3f3f

int n, adjMatrix[MAXN][MAXN];
vector<int> bestCut;

// Modifica o grafo!!!! Se precisar do grafo depois faça uma cópia
int mincut() {
    int bestCost = INF;
    vector<int> v[MAXN];
    for (int i=0; i<n; ++i)
        v[i].assign (1, i);
    int w[MAXN], sel;
    bool exist[MAXN], added[MAXN];
    memset (exist, true, sizeof exist);
    for (int phase=0; phase<n-1; ++phase) {
        memset (added, false, sizeof added);
        memset (w, 0, sizeof w);
        for (int j=0, prev; j<n-phase; ++j) {
            sel = -1;
            for (int i=0; i<n; ++i) {
                if (exist[i] && !added[i] && (sel == -1 || w[i] > w[sel]))
                    sel = i;
            }
            if (j == n-phase-1) {
                if (w[sel] < bestCost) {
                    bestCost = w[sel];
                    bestCut = v[sel];
                }
                v[prev].insert (v[prev].end(), v[sel].begin(), v[sel].end());
                for (int i=0; i<n; ++i) adjMatrix[prev][i] = adjMatrix[i][prev] += adjMatrix[sel][i];
                exist[sel] = false;
            }
            else {
                added[sel] = true;
                for (int i=0; i<n; ++i)  w[i] += adjMatrix[sel][i];
                prev = sel;
            }
        }
    }
    return bestCost;
}

int main() {
    return 0;
}
\end{vncode}
\endgroup


\subsection{Min Cost Max Flow}

Luồng cực đại chi phí nhỏ nhất.

\codepath{Graphs/mincostmaxflow.cpp}

\begingroup\scriptsize
\begin{vncode}
#include <cstring>
#include <set>
using namespace std;
#define MAXN 100009
#define MAXM 900009
#define INFLL 0x3f3f3f3f3f3f3f3f
typedef long long ll;

/*
 * Min Cost Max Flow in O(VE(VlogV+E))
 */

int ned;
int first[MAXN], to[MAXM], nxt[MAXM], prv[MAXN];
ll cap[MAXM], cost[MAXM], dist[MAXN], pot[MAXN]; 

void init() {
	memset(first, -1, sizeof first);
	ned = 0;
}

void add(int u, int v, ll f, ll c) {
	to[ned] = v, cap[ned] = f;
	cost[ned] = c, nxt[ned] = first[u];
	first[u] = ned++;
	to[ned] = u, cap[ned] = 0;
	cost[ned] = -c, nxt[ned] = first[v];
	first[v] = ned++;
}

bool dijkstra(int s, int t, int n) {
	memset(&prv, -1, sizeof prv);
	for(int i = 1; i <= n; i++) dist[i] = INFLL;
	set< pair<ll, int> > q;
	q.insert(make_pair(0LL, s));
	dist[s] = prv[s] = 0;
	while(!q.empty()) {
		int u = q.begin()->second;
		q.erase(q.begin());
		for(int e = first[u]; e != -1; e = nxt[e]) {
			if (cap[e] <= 0) continue;
			int v = to[e];
			ll new_dist = dist[u] + cost[e] + pot[u] - pot[v];
			if (new_dist < dist[v]) {
				q.erase(make_pair(dist[v], v));
				dist[v] = new_dist;
				prv[v] = e;
				q.insert(make_pair(new_dist, v));
			}
		}
	}
	return prv[t] != -1;
}

ll augment(int s, int t, ll &maxflow, int n) {
	ll flow = maxflow;
	for(int i = t; i != s; i = to[prv[i]^1])
		flow = min(flow, cap[prv[i]]);
	for(int i = t; i != s; i = to[prv[i]^1]) {
		cap[prv[i]] -= flow;
		cap[prv[i]^1] += flow;
	}
	maxflow -= flow;
	ll flowCost = flow*(dist[t]-pot[s]+pot[t]);
	for(int i = 1; i <= n; i++)
		if (prv[i] != -1) pot[i] += dist[i];
	return flowCost;
}

ll mincostmaxflow(int s, int t, ll &maxflow, int n) {
	ll flowCost = 0;
	memset(pot, 0, sizeof pot);
	while(maxflow > 0 && dijkstra(s, t, n)) {
		flowCost += augment(s, t, maxflow, n);
	}
	return flowCost;
}

/*
 * Codeforces 101726I
 */

#include <cstdio>

int n, m;
int in[60], out[60];
bool reach[60][60];

int main() {
	int T;
	scanf("%d", &T);
	while(T --> 0) {
		ll ans = 0;
		scanf("%d %d", &n, &m);
		memset(&in, 0, sizeof in);
		memset(&out, 0, sizeof out);
		memset(&reach, false, sizeof reach);
		init();
		for(int j = 0; j < m; j++) {
			int u, v, w;
			scanf("%d %d %d", &u, &v, &w);
			ans += w;
			add(u, v, INFLL, w);
			in[v]++;
			out[u]++;
			reach[u][v] = true;
		}
		for(int k = 1; k <= n; k++) {
			for(int i = 1; i <= n; i++) {
				for(int j = 1; j <= n; j++) {
					reach[i][j] = reach[i][j] || (reach[i][k] && reach[k][j]);
				}
			}
		}

		bool ok = true;
		for(int i = 1; i <= n; i++) {
			for(int j = 1; j <= n; j++) {
				if (!reach[i][j]) ok = false;
			}
		}
		int s = n+1;
		int t = n+2;
		ll maxflow = 0;
		for(int u = 1; u <= n; u++) {
			int bal = in[u] - out[u];
			if (bal < 0) {
				add(u, t, -bal, 0);
			}
			if (bal > 0) {
				add(s, u, bal, 0);
				maxflow += bal;
			}
		}
		ll cost = mincostmaxflow(s, t, maxflow, n+2);
		if (maxflow > 0) ok = false;
		if (ok) printf("%lld\n", cost+ans);
		else printf("-1\n");
	}
	return 0;
}
\end{vncode}
\endgroup


\subsection{Ghép cặp chi phí nhỏ nhất: thuật toán Hungary}

Hungarian algorithm.

\codepath{Graphs/hungarian.cpp}

\begingroup\scriptsize
\begin{vncode}
#include <cstring>
#define MAXN 2009
#define MAXM 2009
#define INF 0x3f3f3f3f

/*
 * Hungarian's Algorithm O(n^2 m)
 * linhas de 1 a n, colunas de 1 a n
 */

int n, m;
int pu[MAXN], pv[MAXN], cost[MAXN][MAXM];
int pairV[MAXN], way[MAXM], minv[MAXM];
bool used[MAXM];

void hungarian() {
	memset(&pairV, 0, sizeof pairV);
	for(int i = 1, j0 = 0; i <= n; i++) {
		pairV[0] = i;
		memset(&minv, INF, sizeof minv);
		memset(&used, false, sizeof used);
		do {
			used[j0] = true;
			int i0 = pairV[j0], delta = INF, j1;
			for(int j = 1; j <= m; j++) {
				if (used[j]) continue;
				int cur = cost[i0][j] - pu[i0] - pv[j];
				if (cur < minv[j])
					minv[j] = cur, way[j] = j0;
				if (minv[j] < delta)
					delta = minv[j], j1 = j;
			}
			for(int j = 0; j <= m; j++) {
				if (used[j])
					pu[pairV[j]] += delta, pv[j] -= delta;
				else minv[j] -= delta;
			}
			j0 = j1;
		} while(pairV[j0] != 0);
		do {
			int j1 = way[j0];
			pairV[j0] = pairV[j1];
			j0 = j1;
		} while(j0);
	}
}

/*
 * Codeforces 101635G
 */

#include <cstdio>
#include <algorithm>
using namespace std;

int bx[MAXN], by[MAXN], cx[MAXM], cy[MAXM], rx, ry;

int main() {
	scanf("%d %d", &n, &m);
	for(int i = 1; i <= n; i++) {
		scanf("%d %d", &bx[i], &by[i]);
	}
	for(int j = 1; j <= m; j++) {
		scanf("%d %d", &cx[j], &cy[j]);
	}
	scanf("%d %d", &rx, &ry);

	for(int i = 1; i <= n; i++) {
		for(int j = 1; j <= m+n-1; j++) {
			if (j <= m) cost[i][j] = abs(bx[i]-cx[j]) + abs(by[i]-cy[j]);
			else cost[i][j] = abs(bx[i]-rx) + abs(by[i]-ry);
			//printf("%d ", cost[i][j]);
		}
		//printf("\n");
	}

	int ans = 0;
	for(int i = 1; i <= n; i++) {
		ans += abs(bx[i]-rx) + abs(by[i]-ry);
	}
	m += n-1;
	hungarian();
	for(int j = 1; j <= m; j++) {
		if (pairV[j] == 0) continue;
		ans += cost[pairV[j]][j];
	}
	printf("%d\n", ans);
	return 0;
}
\end{vncode}
\endgroup


\subsection{Ghép cặp cực đại: thuật toán Kuhn}

Kuhn.

\codepath{Graphs/kuhn.cpp}

\begingroup\scriptsize
\begin{vncode}
#include <vector>
#include <queue>
#include <cstring>
#define MAXN 1009
using namespace std;

/*
 * Kuhn's Algorithm O(VE)
 */

vector<int> adjU[MAXN];
int pairU[MAXN], pairV[MAXN];
bool vis[MAXN];
int m, n;
///Vertices enumerados de 1 a m em U e de 1 a n em V!!!!

bool dfs(int u) {
    vis[u] = true;
    if (u == 0) return true;
    for (int i = 0; i < (int)adjU[u].size(); i++) {
        int v = adjU[u][i];
        if (!vis[pairV[v]] && dfs(pairV[v])) {
            pairV[v] = u; pairU[u] = v;
            return true;
        }
    }
    return false;
}

int kuhn() {
    memset(&pairU, 0, sizeof pairU);
    memset(&pairV, 0, sizeof pairV);
    int result = 0;
    for (int u = 1; u <= m; u++) {
        memset(&vis, false, sizeof vis);
        if (pairU[u]==0 && dfs(u)) result++;
    }
    return result;
}

/*
 * Codeforces 101712A
 */

#include <cstdio>

int main() {
	scanf("%d %d", &m, &n);
	for (int u = 1; u <= m; u++) {
		int v;
		while(scanf("%d", &v), v) adjU[u].push_back(v);
	}
	printf("%d\n", kuhn());
	for (int u = 1; u <= m; u++) {
		if (pairU[u]) printf("%d %d\n", u, pairU[u]);
	}
	return 0;
}
\end{vncode}
\endgroup


\subsection{Ghép cặp cực đại: thuật toán Hopcroft-Karp}

Hopcroft-Karp.

\codepath{Graphs/hopcroftkarp.cpp}

\begingroup\scriptsize
\begin{vncode}
#include <vector>
#include <queue>
#include <cstring>
#define INF 0x3f3f3f3f
#define MAXN 1009
using namespace std;

/*
 * Hopcroft-Karp's Algorithm O(E*sqrt(V))
 */

vector<int> adjU[MAXN];
int pairU[MAXN], pairV[MAXN], dist[MAXN];
int m, n;
//Vertices enumerados de 1 a m em U e de 1 a n em V!!!!

bool bfs() {
	queue<int> q;
	for (int u = 1; u <= m; u++) {
		if (pairU[u] == 0) {
			dist[u] = 0; q.push(u);
		}
		else dist[u] = INF;
	}
	dist[0] = INF;
	while (!q.empty()) {
		int u = q.front(); q.pop();
		if (dist[u] >= dist[0]) continue;
		for (int i = 0; i < (int)adjU[u].size(); i++) {
			int v = adjU[u][i];
			if (dist[pairV[v]] == INF) {
				dist[pairV[v]] = dist[u] + 1;
				q.push(pairV[v]);
			}
		}
	}
	return (dist[0] != INF);
}

bool dfs(int u) {
	if (u == 0) return true;
	for (int i = 0; i < (int)adjU[u].size(); i++) {
		int v = adjU[u][i];
		if (dist[pairV[v]] == dist[u]+1) {
			if (dfs(pairV[v])) {
				pairV[v] = u; pairU[u] = v;
				return true;
			}
		}
	}
	dist[u] = INF;
	return false;
}

int hopcroftKarp() {
	memset(&pairU, 0, sizeof pairU);
	memset(&pairV, 0, sizeof pairV);
	int result = 0;
	while (bfs()) {
		for (int u=1; u<=m; u++) {
			if (pairU[u]==0 && dfs(u))
				result++;
		}
	}
	return result;
}

/*
 * Codeforces 101712A
 */

#include <cstdio>

int main() {
	scanf("%d %d", &m, &n);
	for (int u = 1; u <= m; u++) {
		int v;
		while(scanf("%d", &v), v) adjU[u].push_back(v);
	}
	printf("%d\n", hopcroftKarp());
	for (int u = 1; u <= m; u++) {
		if (pairU[u]) printf("%d %d\n", u, pairU[u]);
	}
	return 0;
}
\end{vncode}
\endgroup


\subsection{Minimum Vertex Cover}

Phủ đỉnh nhỏ nhất trong đồ thị hai phía.

\codepath{Graphs/minimumcover.cpp}

\begingroup\scriptsize
\begin{vncode}
#include <vector>
#include <cstring>
#define INF 0x3f3f3f3f
#define MAXN 4009
using namespace std;

/*
 * Minimum Vertex Cover from Maximum Matching O(V+E)
 */

vector<int> adjU[MAXN];
int pairU[MAXN], pairV[MAXN];
int m, n;
///Vértices enumerados de 1 a m em U e de 1 a n em V!!!!
bool Zu[MAXN], Zv[MAXN];
vector<int> coverU, coverV;

void getreach(int u) {
	if (u == 0 || Zu[u]) return;
	Zu[u] = true;
	for (int i = 0; i < (int)adjU[u].size(); i++) {
		int v = adjU[u][i];
		if (v == pairU[u]) continue;
		Zv[v] = true;
		getreach(pairV[v]);
	}
}

void minimumcover() {
	memset(&Zu, false, sizeof Zu);
	memset(&Zv, false, sizeof Zv);
	for (int u = 1; u <= m; u++)
		if (pairU[u] == 0) getreach(u);
	coverU.clear(); coverV.clear();
	for (int u = 1; u <= m; u++)
		if (!Zu[u]) coverU.push_back(u);
	for (int v = 1; v <= n; v++)
		if (Zv[v]) coverV.push_back(v);
}

/*
 * Codeforces 101712D
 */

#include <cstdio>

int main() {
	scanf("%d %d", &m, &n);
	for (int u = 1; u <= m; u++) {
		int k, v;
		scanf("%d", &k);
		while(k --> 0) {
			scanf("%d", &v);
			adjU[u].push_back(v);
		}
	}
	int nmatches = 0;
	for (int u = 1; u <= m; u++) {
		int v;
		scanf("%d", &v);
		if (v == 0) continue;
		nmatches++;
		pairU[u] = v;
		pairV[v] = u;
	}
	minimumcover();
	printf("%d\n", nmatches);
	printf("%u", coverU.size());
	for(int u : coverU) printf(" %d", u);
	printf("\n");
	printf("%u", coverV.size());
	for(int v : coverV) printf(" %d", v);
	printf("\n");
	return 0;
}
\end{vncode}
\endgroup


\subsection{Maximum Matching: thuật toán Blossom}

Blossom cho matching tổng quát.

\codepath{Graphs/blossom.cpp}

\begingroup\scriptsize
\begin{vncode}
#include <bits/stdc++.h>
using namespace std;
#define MAXN 2005
#define INF 1000000009

/* Codechef Problem: A game on a graph
// Blossom algorithm O(N^3)
// https://www.codechef.com/problems/HAMILG
*/
int mQueue[MAXN], qHead, qTail, match[MAXN], cur_blossom[MAXN], parent[MAXN], N;
bool inPath[MAXN], inQueue[MAXN], inBlossom[MAXN];

vector<vector<int> > G;

void push(int u) {
	mQueue[qTail++] = u;
	inQueue[u] = true;
}

int lca(int u, int v) {
	memset(inPath, 0, sizeof(inPath));
	u = cur_blossom[u];
	while(true) {
		inPath[u] = true;
		if(match[u] == -1 || parent[match[u]] == -1) break;
		u = cur_blossom[parent[match[u]]];
	}

	v = cur_blossom[v];
	while(true) {
		if(inPath[v]) return v;
		v = cur_blossom[parent[match[v]]];
	}

	return v;
}

int find_aug_path(int R) {
	int u, viz;
	memset(inQueue, 0, sizeof(inQueue));
	memset(parent, -1, sizeof(parent));
	for(int i = 0; i < N; i++) cur_blossom[i] = i;
	qTail = qHead = 0;
	push(R);
	
	while(qHead < qTail) {
		u = mQueue[qHead++];
		for (int i = 0; i < G[u].size(); i++) {
			viz = G[u][i];
			if(viz != match[u] && cur_blossom[viz] != cur_blossom[u]) {
				if(parent[viz] == -1 && !(match[viz] != -1 && parent[match[viz]] != -1)) {
					if(viz == R) continue;
					parent[viz] = u;
					if(match[viz] == -1) return viz; // augmenting path to viz
					push(match[viz]);
					continue;
				}

				if(viz != R && !(match[viz] != -1 && parent[match[viz]] != -1)) continue;

				// blossom found
				int new_blossom = lca(u,viz);
				memset(inBlossom, 0, sizeof(inBlossom));
				
				int v = u, parent_v = viz;
				while(cur_blossom[v] != new_blossom) {
					inBlossom[cur_blossom[v]] = inBlossom[cur_blossom[match[v]]] = true;
					parent[v] = parent_v;
					parent_v = match[v];
					v = parent[match[v]];
				}

				v = viz, parent_v = u;
				while(cur_blossom[v] != new_blossom) {
					inBlossom[cur_blossom[v]] = inBlossom[cur_blossom[match[v]]] = true;
					parent[v] = parent_v;
					parent_v = match[v];
					v = parent[match[v]];
				}

				for(int i =0; i < N; i++) if(inBlossom[cur_blossom[i]]) {
					cur_blossom[i] = new_blossom;
					if(!inQueue[i]) {
						push(i);
					}
				}
			}	
		}
	}

	return -1;
}

void MaxMatch() {
	memset(match, -1, sizeof(match));

	for (int i = 0; i < N; i++) {
		if (match[i] < 0) {
			int v = find_aug_path(i); 
			
			while(v != -1) {
				int parent_v = match[parent[v]];
				match[v] = parent[v];
				match[parent[v]] = v;
				v = parent_v;
			}
		}
	}
}

bool marked[MAXN];

void dfs(int R) {
	int u, viz;
	memset(inQueue, 0, sizeof(inQueue));
	memset(parent, -1, sizeof(parent));

	for(int i =0; i < N; i++) cur_blossom[i] = i;
	
	qTail = qHead = 0;
	push(R);
	
	while(qHead < qTail) {
		u = mQueue[qHead++];
		marked[u] = true;
		for (int i = 0; i < G[u].size(); i++) {
			viz = G[u][i];
			if(viz != match[u] && cur_blossom[viz] != cur_blossom[u]) {
				if(parent[viz] == -1 && !(match[viz] != -1 && parent[match[viz]] != -1)) {
					if(viz == R) continue;
					parent[viz] = u;
					if(match[viz] == -1) return; // augmenting path to viz
					push(match[viz]);
					continue;
				}

				if(viz != R && !(match[viz] != -1 && parent[match[viz]] != -1)) continue;

				// blossom found
				int new_blossom = lca(u,viz);
				memset(inBlossom, 0, sizeof(inBlossom));

				int v = u, parent_v = viz;
				while(cur_blossom[v] != new_blossom) {
					inBlossom[cur_blossom[v]] = inBlossom[cur_blossom[match[v]]] = true;
					parent[v] = parent_v;
					parent_v = match[v];
					v = parent[match[v]];}

				v =viz, parent_v = u;
				while(cur_blossom[v] != new_blossom) {
					inBlossom[cur_blossom[v]] = inBlossom[cur_blossom[match[v]]] = true;
					parent[v] =parent_v;
					parent_v =match[v];
					v =parent[match[v]];}

				for(int i =0; i < N; i++) if(inBlossom[cur_blossom[i]]) {
					cur_blossom[i] = new_blossom;
					if(!inQueue[i]) {
						push(i);
					}
				}
			}	
		}
	}

	return;
}

int main() {
	int t, a, b, M;
	scanf("%d", &t);
	while(t--) {
		scanf("%d %d", &N, &M);

		G.clear();
		G.resize(N+5);
		
		for (int i = 0; i < M; i++) {
			scanf("%d %d", &a, &b);a--;b--;
			G[a].push_back(b);
			G[b].push_back(a);
		}
		MaxMatch();

		int ans = 0;
		memset(marked, 0, sizeof(marked));
		
		for (int u = 0; u < N; u++) {
			if (!marked[u] && match[u] == -1) {
				dfs(u);
			}
		}
		for (int u = 0; u < N; u++) {
			if (marked[u]) ans++;
		}
		printf("%d\n", ans);
	}

}
\end{vncode}
\endgroup


\subsection{Gomory-Hu Tree: thuật toán Gusfield}

Cây Gomory-Hu.

\codepath{Graphs/gomoryhu.cpp}

\begingroup\scriptsize
\begin{vncode}
#include <queue>
#include <cstring>
using namespace std;
#define MAXN 3009
#define MAXM 200009
#define INF (1 << 30)

/*
 * Edmonds-Karp's Algorithm - O(VE^2)
 */

int N, M, ned, prv[MAXN], first[MAXN];
int cap[MAXM], to[MAXM], nxt[MAXM], dist[MAXN];

void init()
{
	memset(first, -1, sizeof first);
	ned = 0;
}

void add(int u, int v, int f)
{
	to[ned] = v, cap[ned] = f;
	nxt[ned] = first[u];
	first[u] = ned++;

	to[ned] = u, cap[ned] = 0;
	nxt[ned] = first[v];
	first[v] = ned++;
}

int augment(int v, int minEdge, int s)
{
	int e = prv[v];
	if (e == -1)
		return minEdge;
	int f = augment(to[e ^ 1], min(minEdge, cap[e]), s);
	cap[e] -= f;
	cap[e ^ 1] += f;
	return f;
}

bool bfs(int s, int t)
{
	int u, v;
	memset(&dist, -1, sizeof dist);
	dist[s] = 0;
	queue<int> q;
	q.push(s);
	memset(&prv, -1, sizeof prv);
	while (!q.empty())
	{
		u = q.front();
		q.pop();
		if (u == t)
			break;
		for (int e = first[u]; e != -1; e = nxt[e])
		{
			v = to[e];
			if (dist[v] < 0 && cap[e] > 0)
			{
				dist[v] = dist[u] + 1;
				q.push(v);
				prv[v] = e;
			}
		}
	}
	return dist[t] >= 0;
}

int edmondskarp(int s, int t)
{
	int result = 0;
	while (bfs(s, t))
	{
		result += augment(t, INF, s);
	}
	return result;
}

/*
 * Gomory Hu Tree / Gusfield's Algorithm - O((V-1)*(Flow Algorithm))
 */

typedef pair<int, int> ii;
int par[MAXN], fpar[MAXN]; // parent in tree, flow to parent
int backup[MAXM];		   // backup for cap
bool inCut[MAXN];
vector<ii> ghtree[MAXN];

void cutGraph(int s)
{
	inCut[s] = true;
	for (int e = first[s]; e != -1; e = nxt[e])
	{
		int v = to[e];
		if (!inCut[v] && cap[e] > 0)
		{
			cutGraph(v);
		}
	}
}

int computeMinCut(int s, int t)
{
	memcpy(backup, cap, sizeof cap);
	memset(inCut, false, sizeof inCut);
	int f = edmondskarp(s, t);
	cutGraph(s);
	memcpy(cap, backup, sizeof cap);
	return f;
}

// 1-indexed
void gomoryhu()
{
	for (int i = 1; i <= N; i++)
		par[i] = 1;
	for (int s = 2; s <= N; s++)
	{
		int t = par[s];
		fpar[s] = computeMinCut(s, t);
		for (int u = s + 1; u <= N; u++)
		{
			if (par[u] == t && inCut[u])
			{
				par[u] = s;
			}
			if (par[t] == u && inCut[u])
			{
				par[s] = u;
				par[t] = s;
				swap(fpar[s], fpar[t]);
			}
		}
	}
	for (int i = 1; i <= N; i++)
		ghtree[i].clear();
	for (int i = 1; i <= N; i++)
	{
		if (par[i] == i)
			continue;
		ghtree[i].push_back(ii(fpar[i], par[i]));
		ghtree[par[i]].push_back(ii(fpar[i], i));
	}
}

/*
 * Codeforces 101480J
 */
#include <cstdio>

int dfs(int u, int prv, int k)
{
	int ans = 0, v, w;
	for (int i = 0; i < (int)ghtree[u].size(); i++)
	{
		v = ghtree[u][i].second;
		w = ghtree[u][i].first;
		if (v == prv)
			continue;
		ans += min(k, w);
		ans += dfs(v, u, min(k, w));
	}
	return ans;
}

int main()
{
	int a, b;
	scanf("%d %d", &N, &M);
	init();
	for (int i = 0; i < M; i++)
	{
		scanf("%d %d", &a, &b);
		add(a, b, 1);
		add(b, a, 1);
	}
	gomoryhu();
	int ans = 0;
	for (int i = 1; i <= N; i++)
	{
		ans += dfs(i, -1, INF);
	}
	printf("%d\n", ans / 2);
}
\end{vncode}
\endgroup


\subsection{Euler Tour: thuật toán Fleury}

Euler tour.

\codepath{Graphs/euler.cpp}

\begingroup\scriptsize
\begin{vncode}
#include <vector>
#include <stack>
#define MAXN 400009
using namespace std;

vector<int> adjList[MAXN];
int N, M;

vector<int> euler(int s = 0) {
	vector<int> work(N, 0), in(N, 0), out(N, 0), tour;
	for(int u = 0; u < N; u++) {
		for(int i=0; i<(int)adjList[u].size(); i++) {
			int v = adjList[u][i];
			out[u]++; in[v]++;
		}
	}
	int cntin = 0, cntout = 0;
	for (int u = 0; u < N; u++) {
		if (in[u] == out[u]+1) {
			cntout++;
			if (cntout == 2) return tour;
		}
		else if (out[u] == in[u]+1) {
			cntin++; s = u;
			if (cntin == 2) return tour;
		}
		else if (out[u] != in[u] || in[u]==0) return tour;
	}
	stack<int> dfs;
	dfs.push(s);
	while (!dfs.empty()) {
		int u = dfs.top();
		if (work[u] < (int)adjList[u].size()) {
			dfs.push(adjList[u][work[u]++]);
		}
		else {
			tour.push_back(u);
			dfs.pop();
		}
	}
	int n = tour.size();
	for(int i=0; 2*i<n; i++) swap(tour[i], tour[n-i-1]);
	return tour;
}

/*
 * Codeforces 508D
 */

#include <map>
#include <cstring>
#include <cstdio>

typedef pair<char, char> cc;
cc id2cc[MAXN];
map<cc, int> cc2id;
int deg[MAXN];

int main()
{
	char s[5];
	int u, v;
	N = 0;
	scanf("%d", &M);
	memset(&deg, 0, sizeof deg);
	for(int i=0; i<M; i++) {
		scanf(" %s", s);
		cc p1 = make_pair(s[0], s[1]);
		cc p2 = make_pair(s[1], s[2]);
		if (!cc2id.count(p1)) {
			//printf("%c%c: %d\n", p1.first, p1.second, N);
			id2cc[N] = p1;
			cc2id[p1] = N++;
		}
		if (!cc2id.count(p2)) {
			//printf("%c%c: %d\n", p2.first, p2.second, N);
			id2cc[N] = p2;
			cc2id[p2] = N++;
		}
		u = cc2id[p1];
		v = cc2id[p2];
		adjList[u].push_back(v);
		deg[v]++;
	}
	vector<int> ans = euler();
	int n = ans.size();
	if (M + 1 != n) {
		printf("NO\n");
	}
	else {
		printf("YES\n");
		for(int i=0; i<n; i++) {
			printf("%c", id2cc[ans[i]].first);
		}
		printf("%c\n", id2cc[ans.back()].second);
	}
	return 0;
}
\end{vncode}
\endgroup


\subsection{Dominator Tree}

Cây dominator.

\codepath{Graphs/dominatortree.cpp}

\begingroup\scriptsize
\begin{vncode}
#include <vector>
using namespace std;
#define MAXN 200009

vector<int> adjList[MAXN], revAdjList[MAXN]; 	//grafo
vector<int> dtree[MAXN];						//dominator tree
int sdom[MAXN], dom[MAXN], idom[MAXN], N, M;	//semi-dominator, immediate dominator pelo indice, immediate dominator, tamanhos
int dsu[MAXN], best[MAXN];					//disjoint sets
int par[MAXN], num[MAXN], rev[MAXN], cnt;	//dfs

int find(int x) {
	if (x == dsu[x]) return x;
	int y = find(dsu[x]);
	if (sdom[best[x]] > sdom[best[dsu[x]]]) best[x] = best[dsu[x]];
	return dsu[x] = y;
}

void dfs(int u) {
	num[u] = cnt; rev[cnt++] = u;
	for (int i=0; i<(int)adjList[u].size(); i++) {
		int v = adjList[u][i];
		if (num[v] >= 0) continue;
		dfs(v);
		par[num[v]] = num[u];
	}
}

void dominator() {
	for (int u = 1; u <= N; u++) {
		num[u] = -1; dtree[u].clear();
		dsu[u] = best[u] = sdom[u] = u;
	}
	cnt = 0;
	dfs(1);
	for (int j = cnt - 1, u; u = rev[j], j > 0; j--) {
		for (int i = 0; i<(int)revAdjList[u].size(); i++) {
			int y = num[revAdjList[u][i]];
			if (y == -1) continue;
			find(y);
			if (sdom[best[y]] < sdom[j]) sdom[j] = sdom[best[y]];
		}
		dtree[sdom[j]].push_back(j);
		int x = dsu[j] = par[j];
		for (int i=0; i<(int)dtree[x].size(); i++) {
			int z = dtree[x][i];
			find(z);
			if (sdom[best[z]] < x) dom[z] = best[z];
			else dom[z] = x;
		}
		dtree[x].clear();
	}
	idom[1] = -1;
	for (int i = 1; i < cnt; i++) {
		if (sdom[i] != dom[i]) dom[i] = dom[dom[i]];
		idom[rev[i]] = rev[dom[i]];
		dtree[rev[dom[i]]].push_back(rev[i]);
	}
}

/*
 * Codeforces 100513L
 */
 

#include <cstdio>
#include <set>

int eu[MAXN], ev[MAXN];
int st[MAXN], en[MAXN], k;

void dfs2(int u, int p) {
	st[u] = k++;
	for(int i=0; i<(int)dtree[u].size(); i++) {
		int v = dtree[u][i];
		if (v != p) dfs2(v, u);
	}
	en[u] = k++;
}

int main() {
	int u, v;
	set<int> ans;
	while(scanf("%d %d", &N, &M) != EOF) {
		for(int j=1; j<=M; j++) {
			scanf("%d %d", &u, &v);
			adjList[u].push_back(v);
			revAdjList[v].push_back(u);
			eu[j] = u;
			ev[j] = v;
		}
		dominator();
		ans.clear();
		k = 0;
		dfs2(1, -1);
		for(int j=1; j<=M; j++) {
			u = eu[j];
			v = ev[j];
			if (st[v] <= st[u] && en[v] >= en[u]) continue;
			if (num[u] == -1 || num[v] == -1) continue;
			ans.insert(j);
		}
		printf("%u\n", ans.size());
		for(set<int>::iterator it = ans.begin(); it != ans.end(); it++) {
			printf("%d ", *it);
		}
		printf("\n");
		for(int i=1; i<=N; i++) {
			adjList[i].clear();
			revAdjList[i].clear();
		}
	}
	return 0;
}

/*
 * COJ 2609
 */
 
/*
#include <cstdio>
#define MAXLOGN 22

int level[MAXN];
int P[MAXN][MAXLOGN];

void derevhdfs(int u) {
	for(int i=0; i<(int)dtree[u].size(); i++) {
		int v = dtree[u][i];
		if (v == P[u][0]) continue;
		P[v][0] = u;
		level[v] = 1 + level[u];
		derevhdfs(v);
	}
}
void computeP(int root) {
	level[root]=0;
	P[root][0]=root;
	derevhdfs(root);
	for(int j = 1; j < MAXLOGN; j++)
		for(int i = 1; i <= N; i++)
			P[i][j] = P[P[i][j-1]][j-1];
}
int LCA(int a, int b) {
	if(level[a] > level[b]) swap(a, b);
	int d = level[b] - level[a];
	for(int i=0; i<MAXLOGN; i++) {
		if((d & (1<<i)) != 0) b = P[b][i];
	}
	if(a == b) return a;
	for(int i = MAXLOGN-1; i>=0; i--)
		while(P[a][i] != P[b][i]) {
			a=P[a][i]; b=P[b][i];
		}
	return P[a][0];
}
 
int main() {
	int u, v;
	while(scanf("%d %d", &N, &M) != EOF) {
		for(int j=1; j<=M; j++) {
			scanf("%d %d", &u, &v);
			adjList[u].push_back(v);
		}
		dominator();
		computeP(1);
		int ans = 0;
		for(int i=1; i<=N; i++) {
			for(int j=1; j<i; j++) {
				if (LCA(i, j) == 1) ans++;
			}
		}
		printf("%d\n", ans);
	}
	return 0;
}

*/
\end{vncode}
\endgroup


\subsection{Đồ thị đáng chú ý}

Các họ đồ thị thường gặp: đầy đủ, hai phía đầy đủ, cây, rừng, DAG, Euler, Hamilton, phẳng, chordal, split và các biến thể xuất hiện trong bài toán.

\subsection{Định lý Erdos-Gallai}

Điều kiện đủ để một dãy là dãy bậc đồ thị: tổng bậc chẵn và các bất đẳng thức Erdos-Gallai; có thể dựng bằng Havel-Hakimi.

\codepath{Graphs/erdosgallai.cpp}

\begingroup\scriptsize
\begin{vncode}
#include <vector>
#include <algorithm>
#include <functional>
using namespace std;

/*
 * Erdos-Gallai's theorem
 * O(nlogn), O(n) if already sorted
 * checks if a given array of degrees can represent a graph
 */

bool erdosgallai(vector<int> d) {
	sort(d.begin(), d.end(), greater<int>());
	vector<long long> pd(d.size());
	int n = d.size(), p = n-1;
	for(int i = 0; i < n; i++)
		pd[i] = d[i] + (i > 0 ? pd[i-1] : 0);
	for(int k = 1; k <= n; k++) {
		while(p >= 0 && d[p] < k) p--;
		long long sum;
		if (p >= k-1) sum = (p-k+1)*1ll*k + pd[n-1] - pd[p];
		else sum = pd[n-1] - pd[k-1];
		if (pd[k-1] > k*(k-1LL) + sum) return false;
	}
	return pd[n-1] % 2 == 0;
}

/*
 * URI 1462
 */

#include <cstdio>

int main() {
	int N;
	vector<int> d;
	while(scanf("%d", &N) != EOF) {
		d.resize(N);
		for(int i = 0; i < N; i++) {
			scanf("%d", &d[i]);
		}
		if (erdosgallai(d)) printf("possivel\n");
		else printf("impossivel\n");
	}
	return 0;
}
\end{vncode}
\endgroup


\subsection{Định lý và công thức đồ thị}

Công thức Cayley: có n\textasciicircum{}\{n-2\} cây khung của đồ thị đầy đủ n đỉnh. Euler cho đồ thị phẳng: V-E+F=2. Số cây khung của K\_\{m,n\} là m\textasciicircum{}\{n-1\} n\textasciicircum{}\{m-1\}. Định lý Kirchhoff dùng cofator của ma trận Laplace. Konig: phủ đỉnh nhỏ nhất trong đồ thị hai phía bằng matching cực đại. Dilworth/Mirsky liên hệ chain và antichain trong poset. Menger cho cạnh/đỉnh liên hệ số phần tử cần xóa với số đường đi độc lập. Prufer là song ánh giữa cây n đỉnh và mã dài n-2. Hall cho điều kiện matching hoàn hảo một phía. Dirac và Ore là điều kiện đủ cho đường Hamilton.

\clearpage

\section{Toán học}

\subsection{Số học modular}

MDC, BCNN, Euclid mở rộng, nghịch đảo/chia modular, lũy thừa modular, phương trình Diophantine, nhân modular tránh overflow, Pascal, nghịch đảo modulo và CRT.

\codepath{Math/mod.cpp}

\begingroup\scriptsize
\begin{vncode}
template <typename T>
T gcd(T a, T b) {
	return b == 0 ? a : gcd(b, a % b);
}

template <typename T>
T lcm(T a, T b) {
		return a * (b / gcd(a, b));
}

template <typename T>
T extGcd(T a, T b, T& x, T& y) {
	if (b == 0) {
		x = 1; y = 0;
		return a;
	}
	else {
		T g = extGcd(b, a % b, y, x);
		y -= a / b * x;
		return g;
	}
}
 
template <typename T>
T modInv(T a, T m) {
	T x, y;
	extGcd(a, m, x, y);
	return (x % m + m) % m;
}
 
template <typename T>
T modDiv(T a, T b, T m) {
	return ((a % m) * modInv(b, m)) % m;
}

template<typename T>
T modMul(T a, T b, T m) {
	T x = 0, y = a % m;
	while (b > 0) {
		if (b % 2 == 1) x = (x + y) % m;
		y = (y * 2) % m;
		b /= 2;
	}
	return x % m;
}

template<typename T>
T modExp(T a, T b, T m) {
	if (b == 0) return (T)1;
	T c = modExp(a, b / 2, m);
	c = (c * c) % m;
	if (b % 2 != 0) c = (c*a) % m;
	return c;
}

template<typename T>
void diophantine(T a, T b, T c, T& x, T& y) {
	T d = extGcd(a, b, x, y);
	x *= c / d; y *= c / d;
}

#define MAXN 1000009
typedef long long ll;

ll fat[MAXN];
void preprocessfat(ll m) {
	fat[0] = 1;
	for(ll i=1; i<MAXN; i++) {
		fat[i] = (i*fat[i-1])%m;
	}
}

template<typename T>
T pascal(int n, int k, T m) {
	return modDiv(fat[n], (fat[k]*fat[n-k])%m, m);
}

template<typename T>
void allInv(T inv[], T p) {
	inv[1] = 1;
	for (int i = 2; i < p; i++)
		inv[i] = (p - (p/i)*inv[p%i]%p)%p;
}
\end{vncode}
\endgroup


\subsection{Định lý phần dư Trung Hoa tổng quát}

CRT khi các modulo không nhất thiết nguyên tố cùng nhau.

\codepath{Math/generalizedcrt.cpp}

\begingroup\scriptsize
\begin{vncode}
#include <bits/stdc++.h>
using namespace std;
typedef long long ll;

ll extGcd(ll a, ll b, ll& x, ll& y) {
    if (b == 0) {
        x = 1; y = 0;
        return a;
    }
    else {
        ll g = extGcd(b, a % b, y, x);
        y -= a / b * x;
        return g;
    }
}

/*
 * Generalized CRT
 */

ll crt(ll rem[], ll mod[], int n) {
	if (n == 0) return 0;
    ll ans = rem[0], m = mod[0];
    for (int i = 1; i < n; i++) {
	    ll x, y;
	    ll g = extGcd(mod[i], m, x, y);
	    if ((ans - rem[i])%g != 0) return -1;
	    ans = ans + 1ll*(rem[i]-ans)*(m/g)*y;
	    m = (mod[i]/g)*(m/g)*g;
	    ans = (ans%m + m)%m;
    }
    return ans;
}

/*
 * COMPILATION TEST
 */

int main() {
	long long a[2] = {2, 7};
	long long p[2] = {10, 15};
	printf("%lld\n", crt(a, p, 2));
}
\end{vncode}
\endgroup


\subsection{Số nguyên tố}

Crivo Eratosthenes, kiểm tra nguyên tố, phân tích thừa số, số ước, phi Euler.

\codepath{Math/primes.cpp}

\begingroup\scriptsize
\begin{vncode}
#include <bitset>
#include <cstdio>
#include <vector>
#include <cstring>
using namespace std;
#define MAXN 10000009

typedef long long int ll;
ll sievesize, numDiffp[MAXN];
bitset<MAXN> bs;
vector<ll> primes;

void sieve(ll n) {
	sievesize = n + 1;
	bs.set();
	bs[0] = bs[1] = 0;
	for (ll i = 2; i <= sievesize; i++) {
		if (bs[i]) {
			for (ll j = i * i; j <= (ll)sievesize; j += i) bs[j] = 0;
			primes.push_back(i);
		}
	}
}

//O(sqrt(n)/logn), works if factors are less than the biggest prime in primes
bool isPrimeSieve(ll N) {
	if (N <= (ll)sievesize) return bs[N];
	for (int i = 0; i < (int)primes.size() && primes[i]*primes[i] <= N; i++)
		if (N % primes[i] == 0) return false;
	return true;
}

//O(sqrt(n))
bool isPrime(ll N) {
    if (N < 0) return isPrime(-N);
	for(ll i=2; i*i <= N; i++) {
		if (N % i == 0) return false;
	}
	return true;
}

//O(sqrt(n))
bool isPrimeFast(ll n) {
    if (n < 0) n = -n;
    if (n < 5 || n % 2 == 0 || n % 3 == 0)
        return (n == 2 || n == 3);
    ll maxP = sqrt(n) + 2;
    for (ll p = 5; p < maxP; p += 6) {
        if (p < n && n % p == 0) return false;
        if (p+2 < n && n % (p+2) == 0) return false;
	}
    return true;
}

typedef pair<ll, int> ii;
vector<ii> primeFactors(ll N) {
	vector<ii> factors;
	ll i = 0, p = primes[i];
	while (p * p <= N) {
		if (N % p == 0) {
			int cnt = 0;
			while (N % p == 0) N /= p, cnt++;
			factors.push_back(ii(p, cnt));
		}
		p = primes[++i];
	}
	if (N != 1) factors.push_back(ii(N, 1));
	return factors;
}

ll numDiv(ll N) {
	ll i = 0, p = primes[i], ans = 1;
	while (p * p <= N) {
		ll pw = 0;
		while (N % p == 0) { N /= p; pw++; }
		ans *= (pw + 1);
		p = primes[++i];
	}
	if (N != 1) ans *= 2;
	return ans;
}

ll eulerPhi(ll N) {
	ll i = 0, p = primes[i], ans = N;
	while (p * p <= N) {
		if (N % p == 0) ans -= ans / p;
		while (N % p == 0) N /= p;
		p = primes[++i];
	}
	if (N != 1) ans -= ans / N;
	return ans;
}

void numDiffp() {
	memset(numDiffp, 0, sizeof numDiffp);
	for (int i = 2; i < MAXN; i++)
		if (numDiffp[i] == 0)
			for (int j = i; j < MAXN; j += i)
				numDiffp[j]++;
}

int main() {
	sieve(MAXN); // can go up to 10^7 (need few seconds)
	printf("%d\n", isPrime(2147483647)); // 10-digits prime
	printf("%d\n", isPrime(136117223861LL)); // not a prime, 104729*1299709
	return 0;
}
\end{vncode}
\endgroup


\subsection{Công thức Legendre}

Cho n và prime p, tính số mũ lớn nhất của p chia n! trong O(log n).

\codepath{Math/legendre.cpp}

\begingroup\scriptsize
\begin{vncode}
#include <cstdio>
#include <cstring>

typedef long long ll;

/*
 * Legendre's formula
 */

ll legendre(ll n, ll p) {
	int ans = 0;
	ll prod = p;
	while(prod <= n) {
		ans += n/prod;
		prod *= p;
	}
	return ans;
}

/*
 * TEST MATRIX
 */

ll fac[30];

int brute(ll n, ll p) {
	int ans = -1;
	ll prod = 1LL;
	while(fac[n]%prod == 0) {
		//printf("%lld%%%lld = 0\n", fac[n], prod);
		ans++;
		prod *= p;
	}
	return ans;
}

bool test(ll p) {
	for(ll n=1; n<=20LL; n++) {
		if (legendre(n, p) != brute(n, p)) {
			//printf("legendre(%lld, %lld) = %d\n", n, p, legendre(n, p));
			//printf("brute(%lld, %lld) = %d\n", n, p, brute(n, p));
			return false;
		}
	}
	return true;
}

int N;
bool prime[200000];

int main() {
	N = 200000;
	fac[0] = 1LL;
	for(ll i=1LL; i<=20LL; i++) {
		fac[i] = i*fac[i-1];
	}
	memset(&prime, true, sizeof prime);
	prime[0] = prime[1] = false;
	for(int i=2; i<N; i++) {
		if (prime[i]) {
			if (!test(i)) printf("failed for prime p = %d\n", i);
			for(int j=2*i; j<N; j+=i) prime[j] = false;
		}
	}
	return 0;
}
\end{vncode}
\endgroup


\subsection{Tổng GCD}

Dùng phi Euler để tính F(n)=sum\_\{i=1\}\textasciicircum{}n gcd(i,n).

\codepath{Math/gcdsum.cpp}

\begingroup\scriptsize
\begin{vncode}
#include <bitset>
#include <vector>
using namespace std;
#define MAXN 200009

typedef unsigned long long ll;
int phi[MAXN], sievesize;
ll f[MAXN];

void gcdSieve(int n) {
	sievesize = n+1;
	for(int i=0; i <= sievesize; i++) phi[i] = f[i] = 0;
	for(int i = 1; i <= sievesize; i++) {
		phi[i] += i;
		for(int j = i; j <= sievesize; j += i) {
			if (j > i) phi[j] -= phi[i];
			f[j] += j / i * phi[i];
		}
	}
}

bitset<MAXN> bs;
vector<ll> primes;

void sieve(ll n) {
	sievesize = n + 1;
	bs.set();
	bs[0] = bs[1] = 0;
	for (ll i = 2; i <= sievesize; i++) {
		if (bs[i]) {
			for (ll j = i * i; j <= (ll)sievesize; j += i) bs[j] = 0;
			primes.push_back(i);
		}
	}
}

ll F(ll N) {
	ll i = 0, p = primes[i];
	ll ans = 1;
	while (p * p <= N) {
		if (N % p == 0) {
			N /= p;
			ll prod = 1, e = 1;
			while (N % p == 0) {
				N /= p, e++;
				prod *= p;
			}
			ans *= prod * ((p-1)*e + p);
		}
		p = primes[++i];
	}
	ans *= 2*N-1;
	return ans;
}

/*
 * UVa 11424
 */

#include <cstdio>
ll sum[MAXN];
 
int main() {
	gcdSieve(MAXN-2);
	sieve(MAXN-2);
	sum[0] = 0;
	for(int i=1; i<MAXN; i++) sum[i] = F(i) + sum[i-1];
	ll N;
	while(scanf("%llu", &N), N) {
		printf("%llu\n", sum[N] - N*(N+1)/2);
	}
	return 0;
}
\end{vncode}
\endgroup


\subsection{Crivo tuyến tính và hàm nhân tính}

Crivo O(n) đồng thời tính các hàm nhân tính.

\codepath{Math/linearsieve.cpp}

\begingroup\scriptsize
\begin{vncode}
#include <vector>
#include <bitset>
using namespace std;
#define MAXN 1000009
#include <cstdio>

/*
 * Linear sieve of Eratosthenes and computing multiplative functions
 */

vector<int> primes;
bitset<MAXN> bs;
int f[MAXN], pw[MAXN];

int g(int p, int k) {
	if (p == 1) return 1;
	return (k==0) - (k==1);
	//int q = 1;
	//for(int i = 1; i < k; i++) q *= p;
	//return q*(p-1);
}

void sieve(int n) {
	bs.set(); bs[0] = bs[1] = 0;
	primes.clear(); f[1] = g(1, 1);
	for (int i = 2; i <= n; i++) {
		if (bs[i]) {
			primes.push_back(i);
			f[i] = g(i, 1); pw[i] = 1;
		}
		for (int j = 0; j < primes.size() && i*1ll*primes[j] <= n; j++) {
			bs[i * primes[j]] = 0;
			if (i % primes[j] == 0) {
				int pwr = 1;
				for(int k = 0; k < pw[i]; k++) pwr *= primes[j];
				f[i * primes[j]] = f[i / pwr] * g(primes[j], pw[i]+1);
				pw[i * primes[j]] = pw[i] + 1;
				break;
			} else {
				f[i * primes[j]] = f[i] * f[primes[j]];
				pw[i * primes[j]] = 1;
			}
		}
	}
}

/*
 * TEST MATRIX
 */
#include <cstdio>

void test(int n) {
	sieve(n);
	for(int i = 1; i <= n; i++) {
		printf("f[%3d] = %2d, %d is %sprime, \n", i, f[i], i, bs[i] ? "" : "not ");
	}
}

int main() {
	test(100);
}
\end{vncode}
\endgroup


\subsection{Nghịch đảo Mobius}

Nếu f(n)=sum\_\{d|n\} g(d) thì g(n)=sum\_\{d|n\} f(d) mu(n/d); áp dụng cho đếm cặp nguyên tố cùng nhau và poset.

\subsection{Số Catalan}

Cat(0)=1, Cat(n)=(4n-2)/(n+1) Cat(n-1). Đếm cây nhị phân đầy đủ, ngoặc đúng, cách đặt ngoặc, tam giác hóa đa giác lồi và đường đi đơn điệu.

\subsection{Số Stirling loại một}

s(n,m)=s(n-1,m-1)-(n-1)s(n-1,m), với s(n,n)=1 và s(n,0)=0.

\codepath{Math/stirling\_first.cpp}

\begingroup\scriptsize
\begin{vncode}
#include <vector>
#include <cstring>
using namespace std;

/*
 * Modular inverse
 */

template <typename T>
T extGcd(T a, T b, T& x, T& y) {
    if (b == 0) {
        x = 1; y = 0;
        return a;
    }
    else {
        T g = extGcd(b, a % b, y, x);
        y -= a / b * x;
        return g;
    }
}
 
template <typename T>
T modInv(T a, T m) {
    T x, y;
    extGcd(a, m, x, y);
    return (x % m + m) % m;
}

/*
 * NTT - Number Theoretic Transform O(nlogn)
 */

typedef long long ll;

const ll mod[3] = {1004535809LL, 1092616193LL, 998244353LL};
const ll root[3] = {12289LL, 23747LL, 15311432LL};
const ll root_1[3] = {313564925LL, 642907570LL, 469870224LL};
const ll root_pw[3] = {1LL<<21, 1LL<<21, 1LL<<23};

void ntt(vector<ll> & a, bool invert, int m) {
	ll n = (ll)a.size();
	for(ll i=1, j=0; i<n; ++i) {
		ll bit = n >> 1;
		for (; j>=bit; bit>>=1) j -= bit;
		j += bit;
		if (i < j) swap(a[i], a[j]);
	}
	for(ll len=2, wlen; len<=n; len<<=1) {
		wlen = invert ? root_1[m] : root[m];
		for (ll i=len; i<root_pw[m]; i<<=1)
			wlen = (wlen * wlen % mod[m]);
		for(ll i=0; i<n; i+=len) {
			for(ll j=0, w=1; j<len/2; ++j) {
				ll u = a[i+j],  v = a[i+j+len/2] * w % mod[m];
				a[i+j] = (u+v < mod[m] ? u+v : u+v-mod[m]);
				a[i+j+len/2] = (u-v >= 0 ? u-v : u-v+mod[m]);
				w = w * wlen % mod[m];
			}
		}
	}
	if(invert) {
		ll nrev = modInv(n, mod[m]);
		for (ll i=0; i<n; ++i)
			a[i] = a[i] * nrev % mod[m];
	}
}

void convolution(vector<ll> a, vector<ll> b, vector<ll> & res, int m) {
	ll n = 1;
	while(n < max (a.size(), b.size())) n <<= 1;
	n <<= 1;
	a.resize(n), b.resize(n);
	ntt(a, false, m); ntt(b, false, m);
	res.resize(n);
	for(int i=0; i<n; ++i) res[i] = (a[i]*b[i])%mod[m];
	ntt(res, true, m);
}

/*
 * Stirling numbers
 * Compute s(n, k) for all k in [0, n] in O(nlog^2n) 
 */

void compute(vector<ll> & res, int l, int r, int m) {
	if (l == r) {
		res = vector<ll>({r, 1LL});
		return;
	}
	vector<ll> a, b;
	compute(a, l, (l + r) / 2, m);
	compute(b, (l + r) / 2 + 1, r, m);
	convolution(a, b, res, m);
}

void stirling_first(vector<ll> & res, int n, int m) {
	if (n == 0) {
		res = vector<ll>({1LL});
		return;
	}
	compute(res, 0, n-1, m);
	while(res.back() == 0) res.pop_back();
}

/*
 * TEST MATRIX
 */

#include <cstdio>

void printmatrix(int _n, int m) {
	printf("Stirling matrix of first kind mod %I64d:\n", mod[m]);
	for(int n = 0; n <= _n; n++) {
		vector<ll> row;
		stirling_first(row, n, m);
		for(int i = 0; i < (int)row.size(); i++) {
			printf("%I64d ", row[i]);
		}
		printf("\n");
	}
}

/*
 * Codeforces 960G
 */

#define MAXN 100009
 
template <typename T>
T modDiv(T a, T b, T m) {
    return ((a % m) * modInv(b, m)) % m;
}

ll fat[MAXN];
void preprocessfat(ll m) {
    fat[0] = 1;
    for(ll i=1; i<MAXN; i++) {
        fat[i] = (i*fat[i-1])%m;
    }
}

template<typename T>
T pascal(int n, int k, T m) {
    return modDiv(fat[n], (fat[k]*fat[n-k])%m, m);
}

int N, A, B;

int main() {
	scanf("%d %d %d", &N, &A, &B);
	preprocessfat(mod[2]);
	vector<ll> arr;
	stirling_first(arr, N-1, 2);
	ll ans;
	if (A+B-2 < 0 || A+B-2 >= int(arr.size())) {
		ans = 0;
	}
	else {
		ans = arr[A+B-2];
		ans *= pascal(A+B-2, A-1, mod[2]);
		ans %= mod[2];
	}
	printf("%I64d\n", ans);
	return 0;
}
\end{vncode}
\endgroup


\subsection{Số Stirling loại hai}

S(n,m)=S(n-1,m-1)+mS(n-1,m); đếm cách chia n đối tượng vào đúng m tập không rỗng.



\subsection{Đẳng thức tổng nhị thức}

Gồm Vandermonde, hockey-stick và định lý đa thức.

\subsection{Bổ đề Burnside và định lý đếm Polya}

|X/G| = (1/|G|) sum\_\{g in G\} I(g); với hoán vị và k màu, |X/G|=(1/|G|) sum k\textasciicircum{}\{C(g)\}. Cần G đóng dưới hợp thành.

\subsection{Kiểm tra nguyên tố Miller-Rabin}

Xác suất, O(k log\textasciicircum{}2 n), đúng đã chứng minh cho n < 2\textasciicircum{}64 với các cơ sở đã cho.

\codepath{Math/millerrabin.cpp}

\begingroup\scriptsize
\begin{vncode}
template<typename T>
T modMul(T a, T b, T m) {
	T x = 0, y = a % m;
	while (b > 0) {
		if (b % 2 == 1) x = (x + y) % m;
		y = (y * 2) % m;
		b /= 2;
	}
	return x % m;
}

template<typename T>
T modExp(T a, T b, T m) {
	if (b == 0) return (T)1;
	T c = modExp(a, b / 2, m);
	c = (c * c) % m;
	if (b % 2 != 0) c = (c*a) % m;
	return c;
}

/*
 * Miller-Rabin primality test
 * O(k*log^2n), k = 9
 * probabilistic, but proven correct for n < 2^64
 */

bool miller(long long n) {
	const int pn = 9;
	const int p[] = {2, 3, 5, 7, 11, 13, 17, 19, 23};
	for (int i = 0; i < pn; i++)
		if (n % p[i] == 0) return n == p[i];
	if (n < p[pn - 1]) return false;
	long long s = 0, t = n - 1;
	while (~t & 1) t >>= 1, ++s;
	for (int i = 0; i < pn; ++i) {
		long long pt = modExp((long long)p[i], t, n);
		if (pt == 1) continue;
		bool ok = false;
		for (int j = 0; j < s && !ok; j++) {
			if (pt == n - 1) ok = true;
			pt = modMul(pt, pt, n);
		}
		if (!ok) return false;
	}
	return true;
}

/*
 * TEST MATRIX
 */

#include <cstdio>
#include <cstdlib>

bool brutePrime(long long n) {
	if (n <= 1) return false;
	for(long long i = 2; i*i <= n; i++) {
		if (n % i == 0) return false;
	}
	return true;
}

bool test(int ntest) {
	for(int t = 1; t <= ntest; t++) {
		long long n = rand();
		if (brutePrime(n) != miller(n)) {
			printf("failed on test %lld, prime = %d\n", n, brutePrime(n));
			return false;
		}
		n = t;
		if (brutePrime(n) != miller(n)) {
			printf("failed on test %lld, prime = %d\n", n, brutePrime(n));
			return false;
		}
	}
	printf("all tests passed\n");
	return true;
}

int main() {
	test(10000000);
	long long n;
	while(scanf("%lld", &n) != EOF) {
		printf("%lld is %sprime\n", n, miller(n) ? "" : "not ");
	}
	return 0;
}
\end{vncode}
\endgroup


\subsection{Thuật toán Pollard-Rho}

Trả về một thừa số của n; dùng cho n lớn.

\codepath{Math/pollardrho.cpp}

\begingroup\scriptsize
\begin{vncode}
#include <cstdio>
#include <algorithm>
using namespace std;

template <typename T>
T gcd(T a, T b) {
    return b == 0 ? a : gcd(b, a % b);
}

template<typename T>
T modMul(T a, T b, T m) {
	T x = 0, y = a % m;
	while (b > 0) {
		if (b % 2 == 1) x = (x + y) % m;
		y = (y * 2) % m;
		b /= 2;
	}
	return x % m;
}

/*
 * Pollard-Rho factorization
 * use only for n > 9*10^13
 */

template<typename T>
T pollard(T n) {
	int i = 0, k = 2, d;
	T x = 3, y = 3;
	while (true) {
		i++;
		x = (modMul(x, x, n) + n - 1) % n;
		d = gcd(abs(y - x), n);
		if (d != 1 && d != n) return d;
		if (i == k) y = x, k *= 2;
	}
}

/*
 *TEST MATRIX
 */
#include <cstdlib>

int main()
{
	long long n;
	while(scanf(" %I64d", &n)!=EOF) {
		long long r = pollard(n);
		printf("%I64d->%I64d\n", n, r);
		if (n%r != 0) {
			printf("Failed on test\n");
		}
	}
	return 0;
}
\end{vncode}
\endgroup


\subsection{Baby-Step Giant-Step cho logarit rời rạc}

Giải \texttt{a\textasciicircum{}x} \(\equiv\) \texttt{b (mod m)} trong O(sqrt(m) log m).

\codepath{Math/babysteps.cpp}

\begingroup\scriptsize
\begin{vncode}
#include <map>
#include <cmath>
using namespace std;

/*
 * Baby Steps Algorithm for Discrete Logarithm
 */

template <typename T>
T baby (T a, T b, T m) {
	a %= m; b %= m;
	T n = (T) sqrt (m + .0) + 1, an = 1;
	for (T i=0; i<n; ++i) an = (an * a) % m;
	map<T, T> vals;
	for (T i=1, cur=an; i<=n; ++i) {
		if (!vals.count(cur)) vals[cur] = i;
		cur = (cur * an) % m;
	}
	for (T i=0, cur=b; i<=n; ++i) {
		if (vals.count(cur)) {
			T ans = vals[cur] * n - i;
			if (ans < m) return ans;
		}
		cur = (cur * a) % m;
	}
	return -1;
}

/*
 * HackerEarth Dexter plays with GP
 */

#include <cstdio>
typedef long long ll;

int main() {
	int T;
	ll r, x, p, a, b;
	scanf("%d", &T);
	while(T-->0) {
		scanf("%lld %lld %lld", &r, &x, &p);
		a = r;
		b = x*(r-1) + 1;
		printf("%lld\n", baby(a, b, p));
	}
	return 0;
}
\end{vncode}
\endgroup


\subsection{Trò chơi Nim và định lý Sprague-Grundy}

Trạng thái thắng khi nimber khác 0; nimber của trạng thái là mex của các trạng thái đi được; chơi song song lấy xor.

\codepath{Math/mex.cpp}

\begingroup\scriptsize
\begin{vncode}
#include <bits/stdc++.h>
using namespace std;
#define MAXN 10009

int mex[MAXN] = {0,0,0,0,0,1,1,1,2,2,0,3,3,1,1,1,0,4,3,3,3,2,2,2,4,4,0,5,5,2,2,2,3,3,0,5,0,1,1,1,3,3,3,5,6,4,4,1,0,5,5,6,6,2,7,7,7,8,0,1,9,2,7,2,3,3,3,9,0,5,4,4,8,6,6,2,7,1,1,1,0,5,5,9,3,1,8,2,8,5,0,1,1,12,2,2,7,3,3,9,4,4,0,11,3,3,3,9,2,2,8,1,3,5,0,9,12,2,6,13,13,5,0,1,1,4,11,7,7,10,3,4,1,4,0,5,0,3,3,6,7,2,14,13,10,4,12,9,2,2,3,3,6,9,9,1,16,4,8,3,3,2,15,1,1,4,0,5,5,16,6,6,6,8,0,16,5,4,4,17,2,2,7,14,6,10,12,1,0,16,13,3,6,2,7,7,8,1,0,5,17,2,12,15,3,11,0,19,18,12,4,16,17,2,2,21,6,9,4,19,5,5,17,10,3,6,19,2,7,8,4,1,9,12,7,2,13,6,3,19,5,9,4,8,8,17,17,2,15,18,1,1,8,5,21,16,21,3,19,19,13,5,18,1,4,17,7,2,7,6,3,19,12,5,5,16,16,6,17,19,7,7,18,1,4,17,0,9,16,3,3,14,13,22,0,1,15,24,17,2,6,18,3,4,19,19,0,8,21,16,3,15,7,26,18,13,1,1,17,9,2,21,2,6,22,19,9,5,16,4,16,20,3,7,18,23,22,8,20,5,16,21,15,6,10,19,18,18,18,4,4,17,17,7,2,3,23,19,9,5,0,16,16,3,17,30,2,18,18,8,4,17,17,9,27,6,10,19,19,14,9,9,4,20,17,14,11,7,18,6,19,19,5,13,16,16,10,6,19,19,23,18,4,4,17,12,12,14,10,6,3,19,5,9,5,21,16,20,6,7,7,18,30,13,13,17,12,21,15,10,3,19,22,18,8,4,32,17,17,11,14,6,26,24,12,5,9,16,16,6,7,7,7,18,18,8,4,17,20,7,16,10,10,22,19,22,9,23,4,13,17,20,7,11,23,23,4,5,9,5,16,16,10,17,10,22,18,23,8,4,17,17,20,16,32,13,19,19,33,5,5,24,8,17,20,33,18,18,23,13,22,33,33,16,15,10,10,19,19,33,18,8,20,17,17,11,16,18,26,22,19,12,9,16,16,34,20,17,35,23,18,8,8,8,17,33,16,32,10,11,22,11,12,23,9,8,17,17,11,33,18,18,19,13,12,21,16,16,10,20,14,11,22,23,34,8,17,17,29,16,7,34,19,19,9,23,5,8,32,17,33,11,23,18,18,8,12,12,16,16,16,10,32,22,22,11,18,13,8,20,17,11,21,11,15,19,22,12,12,16,16,32,17,17,35,23,18,8,13,13,9,17,16,10,10,19,14,19,25,12,12,13,20,14,11,11,23,10,22,22,35,9,16,16,17,10,20,11,18,18,13,13,24,17,12,16,15,34,19,19,33,12,16,32,8,17,20,33,18,10,18,32,24,8,16,16,16,15,32,19,19,22,13,13,32,20,20,35,33,34,15,22,19,12,12,13,13,21,20,14,14,18,18,23,22,13,9,17,16,15,15,10,11,19,18,12,12,8,20,17,11,21,15,15,31,12,38,13,25,8,20,10,15,11,23,18,13,13,12,12,12,21,32,10,19,19,9,9,38,13,20,17,14,11,14,27,18,5,12,13,1,21,16,17,17,19,11,11,18,13,8,12,24,20,16,11,15,19,19,12,9,13,13,13,33,17,14,3,15,32,19,8,9,4,16,15,32,19,19,11,33,18,32,17,17,14,14,33,18,10,19,12,35,33,16,16,15,26,10,11,23,18,13,13,20,20,20,16,37,34,22,19,23,18,25,8,17,17,11,11,18,10,10,19,40,36,9,16,16,15,32,11,6,18,18,13,8,17,9,35,16,34,10,19,19,31,35,9,21,13,20,14,10,18,18,23,34,41,12,4,16,32,5,26,11,14,18,13,12,8,20,17,16,16,23,15,19,19,35,40,16,16,20,17,46,14,11,18,34,8,12,17,20,16,32,39,19,22,31,23,9,8,8,21,17,33,15,10,10,22,19,12,21,16,20,15,32,35,11,18,18,34,13,9,29,16,21,11,14,19,19,35,35,16,8,17,20,33,14,11,10,30,13,44,40,1,16,28,32,44,19,11,23,23,9,32,17,17,21,16,49,2,22,40,35,33,16,8,34,17,43,10,18,18,39,22,9,9,16,21,37,14,26,19,33,18,27,32,8,24,24,33,18,32,10,19,38,12,16,16,16,17,43,10,5,18,18,39,9,17,17,35,16,47,37,22,19,35,23,40,0,8,17,33,38,18,15,39,34,22,9,16,25,21,15,22,14,11,18,45,9,40,17,20,36,33,15,10,22,19,35,35,16,16,17,20,46,45,30,18,23,34,9,12,21,16,43,32,19,19,42,38,18,32,8,24,24,33,18,18,15,22,26,50,33,16,16,15,17,46,6,18,18,39,57,48,17,35,49,49,14,22,19,38,23,13,16,16,17,36,14,11,10,32,19,41,52,21,25,24,32,24,14,14,18,30,51,9,8,17,36,16,11,19,19,19,35,49,16,13,29,53,53,36,10,18,34,34,9,9,9,20,52,37,14,19,22,23,30,40,20,17,36,21,36,32,18,19,31,55,33,21,16,37,17,35,23,30,18,34,48,12,24,35,49,49,44,22,19,22,18,35,34,8,20,29,47,18,23,30,55,22,38,8,21,28,20,11,26,14,18,51,56,40,12,17,16,21,52,14,26,19,35,54,44,8,20,32,33,60,10,10,23,22,48,35,16,16,57,37,14,31,30,18,18,43,20,17,17,42,33,10,34,19,19,50,56,16,12,37,20,53,23,23,26,44,56,48,12,25,16,49,34,19,14,38,45,49,39,16,12,48,10,15,27,15,39,55,55,4,13,16,64,24,58,36,47,30,63,12,9,12,33,49,51,14,10,19,23,56,49,12,16,17,58,38,33,18,18,22,34,35,21,8,16,29,50,31,15,38,51,64,12,12,24,61,15,37,23,19,55,35,61,8,16,17,53,53,10,18,10,41,56,45,40,21,28,16,65,55,19,22,54,23,44,12,17,24,40,32,54,27,14,50,62,13,8,28,43,17,55,10,23,30,56,37,24,17,13,16,52,37,14,14,23,18,49,9,16,20,47,60,23,10,11,64,19,38,38,8,16,64,20,48,27,19,56,51,20,24,24,28,16,37,14,14,26,35,65,25,9,21,20,64,65,10,23,69,66,53,40,16,16,49,66,26,19,19,23,18,39,9,9,70,65,45,18,14,22,19,65,25,16,16,12,17,47,36,15,18,34,64,24,13,56,49,61,37,19,10,23,68,56,44,12,17,53,67,30,27,18,55,39,45,21,16,16,44,66,53,15,15,23,54,64,20,47,17,21,52,14,14,22,74,65,49,12,9,20,53,42,18,19,74,31,38,8,13,13,17,66,66,19,15,54,51,71,17,17,20,16,67,46,14,14,50,47,38,16,9,12,17,53,65,10,18,71,31,57,20,17,28,61,69,14,19,10,18,44,52,17,9,72,72,63,18,23,14,22,38,21,8,12,41,74,55,19,18,67,64,14,9,65,8,68,51,11,11,14,42,65,56,12,16,29,40,65,10,63,11,71,62,60,16,21,25,41,62,59,15,10,18,71,64,9,8,15,68,54,14,19,19,42,67,54,25,12,57,55,68,18,18,46,73,60,13,20,13,52,69,19,59,10,10,58,71,17,12,24,40,58,18,18,11,19,45,16,16,9,17,20,62,10,23,30,11,14,17,8,65,56,16,46,11,18,10,15,49,12,9,29,57,68,49,18,11,14,48,57,21,16,25,75,74,59,26,10,78,11,12,16,17,8,68,18,14,11,19,15,70,35,64,9,9,53,65,58,18,31,46,53,20,8,16,81,64,45,19,15,10,54,71,12,48,29,13,65,36,22,14,19,47,16,16,12,12,57,62,65,18,18,46,66,50,17,8,8,56,64,19,10,10,30,76,16,9,17,45,65,23,18,11,11,48,70,21,16,25,72,62,59,68,10,18,71,20,20,13,13,25,33,64,22,27,10,67,66,9,9,9,60,18,18,18,22,38,48,17,8,25,52,69,62,65,10,10,63,76,69,9,8,13,78,81,18,11,48,62,65,16,12,9,29,74,65,10,14,11,66,53,17,8,8,81,9,22,19,10,10,83,21,89,17,9,68,81,58,18,59,47,72,67,21,9,66,69,59,19,30,26,71,20,17,8,21,44,78,22,12,10,79,51,51,38,17,20,60,68,18,11,18,11,59,17,28,8,81,9,55,19,15,58,58,36,16,17,13,67,63,49,18,71,47,15,61,21,9,12,57,82,68,63,11,64,66,57,77,8,49,56,12,22,19,18,27,81,83,75,20,45,65,44,65,11,14,59,15,13,16,56,25,76,62,22,10,15,66,66,17,8,8,41,9,69,12,30,79,70,16,38,9,17,8,13,18,18,75,62,19,67,89,16,9,9,80,50,10,10,36,11,20,17,8,25,16,63,22,47,19,59,68,16,16,9,12,39,65,18,46,78,66,60,82,21,24,44,75,64,19,10,10,83,66,69,17,8,79,58,66,14,14,59,10,16,21,16,73,45,59,13,23,26,36,82,17,8,65,16,35,19,19,15,31,88,46,64,17,12,60,13,18,18,18,62,10,68,74,16,41,73,45,19,23,15,18,14,17,20,72,61,9,64,73,19,14,67,86,52,64,60,57,8,77,18,11,58,71,71,73,16,12,56,12,19,15,68,10,18,11,75,17,20,89,74,23,26,22,46,80,10,16,16,12,8,72,18,18,7,71,36,53,8,13,16,9,64,76,19,10,10,61,21,12,12,45,8,23,14,14,38,80,74,16,16,44,13,59,77,68,10,33,11,11,17,17,83,1,9,69,22,11,72,37,70,16,33,50,8,18,18,18,90,19,91,87,21,21,13,75,89,19,10,34,14,11,71,17,70,4,44,73,11,22,14,70,70,21,9,60,13,92,72,18,23,90,22,53,80,20,16,9,12,19,22,43,96,11,11,17,17,68,82,18,23,18,38,59,91,74,16,54,8,85,69,72,23,46,84,17,17,17,72,12,83,84,62,22,11,46,86,76,20,57,8,44,77,14,75,47,47,70,16,16,13,64,80,72,15,15,46,51,20,20,20,88,44,64,22,19,11,10,95,21,76,45,5,8,77,23,18,42,36,82,87,70,16,9,52,12,19,95,46,66,78,17,20,44,13,34,81,75,35,15,67,74,16,5,13,45,65,91,63,11,76,36,17,8,32,9,81,97,22,60,74,10,52,66,97,45,8,41,18,18,35,6,91,87,16,16,8,12,62,19,79,18,18,78,17,17,20,77,9,84,73,11,22,10,90,97,76,40,8,80,96,23,66,42,47,97,17,13,103,83,99,57,10,22,10,46,99,75,20,20,98,44,86,69,38,19,43,16,74,5,73,89,77,68,34,18,42,47,70,77,8,12,9,45,19,10,46,43,83,71,42,17,8,8,69,69,23,10,10,10,21,16,45,45,40,68,39,18,11,51,11,17,8,44,21,81,99,22,67,10,15,83,100,0,20,82,72,77,93,42,47,10,79,37,16,40,45,68,65,85,6,14,97,20,17,89,8,12,81,38,7,14,50,34,79,97,40,17,44,80,98,58,11,14,96,87,4,16,44,45,73,59,43,18,75,97,97,20,44,13,88,69,35,35,10,91,66,83,36,45,40,44,98,46,6,15,67,67,87,32,21,81,99,38,19,10,43,71,83,5,20,53,96,98,106,38,15,82,96,74,16,13,9,73,68,39,14,18,100,47,70,87,8,77,69,35,2,19,15,39,66,83,33,12,89,98,18,18,38,97,67,96,96,32,52,102,68,77,77,18,2,83,33,17,17,98,88,81,76,35,50,101,46,83,75,36,17,68,44,23,3,15,14,70,4,44,32,81,76,35,10,34,39,78,100,78,12,17,4,72,69,35,22,22,82,37,16,8,40,80,85,3,11,18,42,42,94,79,13,9,81,99,2,22,10,34,90,83,36,20,77,44,98,86,10,14,10,67,101,4,13,33,92,80,14,34,23,97,97,92,32,1,81,99,57,19,85,6,46,101,5,13,20,8,98,23,112,104,35,87,96,37,86,0,99,70,87,41,10,86,97,45,87,8,41,88,98,19,35,3,70,96,52,45,97,68,98,23,18,14,69,47,89,79,32,13,69,81,35,11,22,44,95,95,5,5,4,34,103,81,35,14,10,43,32,86,36,33,92,80,96,34,11,97,47,68,12,4,88,88,2,7,19,3,96,96,33,17,12,1,98,98,90,38,10,70,1,37,32,5,99,92,72,115,6,111,97,36,5,17,1,60,84,2,35,85,74,96,100,33,33,77,68,98,84,18,38,97,87,82,1,5,69,102,102,19,34,34,71,106,5,17,67,8,88,64,11,14,10,96,101,32,45,13,80,85,72,10,39,47,97,79,96,17,81,104,99,19,14,3,39,66,78,0,12,94,32,103,112,10,59,87,4,32,16,5,99,99,77,10,6,15,11,97,72,32,32,16,69,2,35,38,63,46,112,33,12,36,77,103,96,23,2,38,67,101,4,32,99,99,99,80,22,6,18,78,33,72,103,37,114,84,2,35,67,108,46,120,45,5,17,99,96,34,10,97,74,79,96,13,103,69,17,38,22,19,78,119,109,33,9,80,98,93,90,35,7,10,94,101,37,21,33,77,91,3,10,18,11,33,67,4,17,69,88,111,2,19,39,10,75,5,36,116,8,103,39,114,35,35,70,79,32,13,112,38,92,115,34,95,95,97,104,0,91,13,84,114,7,35,6,96,41,66,106,13,85,77,34,34,38,100,14,113,96,32,76,0,12,87,89,3,18,104,97,33,113,20,16,88,107,38,7,50,94,13,76,33,5,77,72,101,34,11,79,35,104,32,17,21,76,42,11,35,39,15,78,33,9,9,16,98,76,93,18,80,79,96,8,44,16,33,72,77,34,3,83,38,33,9,4,32,73,81,35,14,38,18,10,32,117,72,9,72,98,121,10,7,11,79,82,13,37,16,9,106,6,3,34,75,104,33,12,52,73,76,111,109,7,113,105,75,112,33,8,17,85,34,93,71,23,79,87,32,17,73,106,60,35,7,22,121,112,60,36,33,4,13,76,106,2,38,87,79,1,13,128,12,72,115,22,10,90,104,102,82,17,13,76,84,11,14,14,18,10,83,59,21,115,72,17,34,15,23,14,113,105,32,8,111,9,16,115,19,15,117,109,9,56,32,8,73,114,128,7,2,19,105,17,33,33,12,98,10,81,22,18,11,91,96,37,83,5,111,11,89,19,10,78,97,36,72,17,32,76,114,18,14,6,94,32,13,33,21,72,110,34,6,30,23,87,33,37,8,73,16,14,14,23,18,75,86,104,9,16,110,13,114,15,38,14,79,113,37,13,33,9,85,107,19,10,23,121,74,82,17,8,13,84,38,11,23,19,86,95,114,33,89,115,107,15,35,11,11,82,113,17,13,16,16,72,77,34,10,86,104,12,68,17,8,76,84,35,129,3,19,32,32,17,53,9,115,101,15,6,11,94,82,120,37,16,111,106,35,115,22,78,83,109,21,21,17,13,10,119,23,94,82,105,32,95,9,77,9,92,39,19,23,119,109,129,118,13,69,81,111,14,113,10,19,112,36,21,20,85,34,10,35,11,22,87,32,90,8,36,12,21,128,19,15,83,18,9,131,17,8,110,111,11,23,14,39,75,75,33,33,12,80,73,15,96,22,38,94,108,8,13,81,16,14,34,15,18,83,112,9,80,17,110,73,119,35,14,58,71,37,8,36,9,72,92,110,22,23,23,23,98,108,8,76,119,128,127,115,22,71,78,109,33,80,17,76,76,119,11,19,82,75,20,112,12,21,77,89,34,15,83,86,62,101,37,13,13,81,111,23,22,22,32,13,33,36,20,20,128,76,35,11,18,113,37,32,121,16,16,89,39,10,18,83,124,36,36,20,17,114,84,35,35,82,116,13,129,12,9,72,98,34,22,23,23,23,9,37,37,81,16,11,74,18,15,15,19,117,33,20,17,88,100,59,14,11,82,113,86,128,4,16,68,115,34,15,23,18,120,94,32,8,84,114,77,14,105,19,75,78,36,33,9,122,15,100,35,31,128,94,63,32,24,5,21,14,34,39,10,23,55,33,12,17,17,76,114,35,30,87,63,78,60,53,9,21,80,34,22,31,11,79,12,37,37,13,20,128,115,120,15,10,104,83,130,33,17,17,84,119,35,19,82,105,37,95,36,16,118,39,34,78,86,104,9,82,49,16,20,114,114,38,27,30,78,95,33,33,9,129,10,84,128,23,70,18,134,131,37,24,111,115,103,58,15,83,133,136,21,52,17,17,60,128,18,54,22,78,66,36,134,103,115,49,34,23,35,23,120,131,32,21,29,33,85,127,34,19,19,117,135,25,20,73,17,81,38,79,23,54,100,75,24,33,115,118,129,129,123,128,94,113,37,20,69,17,128,77,118,22,78,112,36,36,80,32,132,81,114,128,18,120,113,90,111,24,33,85,115,34,58,19,124,135,113,32,20,17,114,38,72,19,39,134,117,57,16,135,21,128,99,31,23,18,51,37,61,21,60,89,118,63,58,39,19,91,91,103,25,53,81,99,38,14,34,23,100,112,130,21,21,58,129,132,55,38,113,82,130,131,135,20,114,92,63,30,13,112,112,45,65,32,17,81,126,11,22,23,105,90,129,36,130,115,80,63,129,117,135,94,82,32,20,17,20,133,98,115,39,55,83,57,33,21,21,132,73,38,18,23,87,134,134,60,123,140,52,132,43,31,31,86,96,96,12,29,17,114,129,128,11,6,104,97,12,56,85,122,118,73,62,26,113,94,61,20,9,17,92,77,58,22,75,86,60,130,5,17,56,93,99,133,31,18,95,100,131,57,133,98,85,128,93,18,62,130,87,120,28,17,73,128,129,58,19,55,100,131,16,80,136,88,84,31,23,18,113,113,134,17,33,16,103,23,129,26,104,130,91,110,61,17,114,114,132,23,39,63,90,83,57,24,16,128,118,99,59,59,30,87,130,131,81,57,133,98,63,58,22,55,139,96,56,56,16,81,137,47,82,30,63,95,29,60,21,85,129,128,140,86,23,87,82,120,25,12,126,128,80,98,19,26,95,131,65,56,132,132,123,81,26,51,51,116,130,57,57,60,80,129,132,109,59,130,125,134,131,28,84,119,141,89,58,58,71,90,45,29,49,52,88,81,59,23,148,87,56,139,84,60,33,80,69,58,19,142,59,91,56,146,20,64,129,128,82,39,87,86,83,60,29,28,80,88,76,23,59,125,87,130,131,84,33,128,103,34,63,22,86,131,131,80,9,129,84,62,62,62,54,131,75,57,29,17,25,132,136,63,59,135,87,101,56,20,16,36,118,132,58,22,55,83,8,21,32,129,132,93,59,128,27,96,131,100,60,60,68,128,118,58,27,26,59,120,61,56,20,64,129,128,108,62,63,83,121,151,130,28,122,128,84,59,59,96,144,134,20,111,133,103,107,58,14,50,135,134,107,56,52,69,17,114,128,31,63,131,130,60,57,24,85,129,128,14,35,135,113,113,61,106,81,126,65,63,63,58,19,83,130,130,98,49,129,99,26,128,30,82,131,100,95,60,133,80,85,63,86,35,130,96,61,56,56,126,129,89,58,58,131,19,83,57,92,61,80,102,111,62,31,105,58,131,111,12,128,103,107,129,58,59,59,139,105,56,56,17,81,81,132,22,58,131,97,104,57,28,80,129,141,88,59,139,18,87,56,97,57,128,92,80,58,63,97,124,70,94,134,56,106,128,145,15,23,120,104,66,130,57,21,65,129,14,23,30,139,67,56,9,24,81,129,133,89,63,58,109,112,135,60,56,129,102,102,62,133,82,30,131,143,57,60,137,107,58,58,148,124,59,87,56,24,104,111,128,62,108,34,104,100,97,57,110,80,56,141,126,124,152,67,96,134,56,48,57,107,85,63,58,59,124,130,110,146,56,76,106,128,133,19,63,66,135,97,150,65,25,115,129,148,62,83,101,105,56,134,12,133,129,132,141,124,66,19,130,142,105,129,1,99,59,62,18,131,134,95,146,57,110,77,58,111,30,42,108,70,143,61,99,57,128,127,107,63,131,109,75,130,57,32,132,64,84,128,148,58,134,131,61,128,60,115,110,129,136,18,35,125,94,56,56,106,84,128,128,103,131,3,71,109,57,28,85,56,64,102,15,130,125,105,131,143,69,57,127,85,132,152,11,83,135,67,150,12,106,128,133,10,26,58,124,90,135,133,68,98,157,63,59,27,105,108,131,143,104,57,128,133,58,63,138,31,109,60,24,61,56,93,99,62,67,23,143,138,124,60,29,137,72,63,14,145,59,125,113,129,12,102,111,128,62,62,39,63,66,130,130,57,61,129,64,126,142,152,74,105,56,150,21,110,137,103,132,63,59,130,86,125,113,160,111,69,59,140,58,58,104,66,139,133,8,56,158,148,111,62,62,70,151,138,150,147,128,98,80,63,104,10,152,139,120,61,12,106,133,62,70,63,134,143,97,25,60,98,103,129,63,59,62,96,67,101,151,97,99,128,127,153,35,152,66,139,139,61,61,167,111,114,130,62,67,163,143,71,8,130,72,68,63,63,15,19,142,96,146,141,141,128,22,133,144,39,143,66,90,150,49,65,136,123,27,59,34,108,67,131,36,57,17,65,103,63,58,19,142,70,108,132,36,99,148,140,67,96,39,121,66,147,130,57,85,132,132,59,14,130,103,131,134,61,106,141,122,148,38,63,11,159,130,120,129,45,64,133,128,10,39,143,100,25,60,60,65,98,132,136,145,59,120,105,146,151,136,111,137,149,101,58,97,109,130,139,65,151,16,99,106,42,59,70,131,146,69,60,122,110,98,129,35,31,130,143,105,60,45,102,99,137,113,39,58,121,97,147,37,68,68,134,106,59,145,108,70,134,154,53,60,168,122,144,144,63,66,142,158,150,32,45,64,102,148,39,101,170,66,109,130,135,16,65,156,136,148,133,15,67,82,36,1,45,168,153,122,138,63,71,100,150,103,12,141,64,102,59,62,46,58,159,147,150,57,68,89,129,38,42,39,128,105,146,131,17,140,137,107,157,63,15,66,147,72,44,41,68,69,128,43,125,11,151,121,114,172,13,107,129,144,42,75,95,154,108,56,36,64,128,148,164,19,58,104,100,60,37,65,65,149,123,59,145,163,120,143,154,99,45,127,163,156,58,58,14,139,154,44,172,36,64,128,128,46,11,70,71,147,128,44,107,98,141,35,38,42,43,70,150,56,45,148,137,34,103,47,109,75,147,57,41,33,21,123,148,62,43,67,138,158,36,8,168,110,77,132,38,75,166,142,151,41,33,64,128,157,96,11,160,121,66,60,155,4,169,141,99,38,34,46,15,134,146,33,40,135,122,89,129,90,71,142,139,74,49,169,45,73,144,39,54,143,104,154,130,8,68,129,141,123,38,34,19,120,131,45,45,64,140,118,58,11,121,71,75,8,41,52,111,76,128,145,43,38,146,33,36,13,32,12,77,148,35,42,14,120,105,151,45,64,69,140,156,43,143,38,75,147,155,44,68,33,148,152,161,11,79,143,151,57,13,166,110,89,144,35,66,159,161,101,151,169,73,73,149,39,46,179,71,75,130,8,41,72,141,152,42,15,163,146,131,53,40,166,168,65,144,38,75,42,78,151,37,131,36,73,133,152,10,171,158,139,121,8,115,41,141,141,35,42,175,96,101,131,36,36,76,148,43,19,11,38,75,158,44,41,33,123,69,148,145,74,79,129,36,36,45,37,41,103,144,42,10,159,154,70,155,40,76,145,140,46,46,129,174,71,135,41,12,65,156,152,47,140,43,67,146,154,36,45,157,68,103,156,35,38,147,142,41,12,189,69,148,133,39,39,160,170,124,45,8,20,72,141,35,42,128,96,151,131,147,40,13,152,103,22,11,104,75,139,101,44,177,141,73,145,43,43,70,74,75,36,40,57,24,68,144,76,42,15,159,96,158,9,8,148,128,46,149,14,100,71,135,41,41,68,149,123,148,157,120,108,146,158,40,48,127,8,144,144,35,161,147,150,24,57,33,40,73,152,43,46,38,11,139,154,41,37,72,167,42,47,164,173,74,146,147,45,184,166,115,132,149,35,47,150,147,52,44,167,73,145,148,39,46,155,100,40,147,5,98,149,129,38,42,133,74,41,158,40,16,145,157,98,46,149,75,147,150,37,44,57,156,126,42,34,101,79,151,97,40,159,122,65,72,141,38,47,95,37,32,151,111,21,69,145,43,11,124,66,71,159,37,20,60,153,111,58,10,6,74,151,45,33,69,44,122,129,35,47,63,159,105,103,57,8,33,145,140,39,74,146,121,104,135,37,166,68,136,42,62,178,170,155,158,36,13,145,122,46,144,153,100,75,147,41,37,180,61,99,73,157,34,155,146,155,33,13,122,72,132,156,35,58,166,79,162,143,81,33,69,128,46,46,165,75,150,154,44,37,77,126,148,148,148,62,59,74,33,28,56,157,136,149,38,35,147,142,87,29,57,33,45,148,148,34,42,174,171,150,135,37,9,60,111,100,145,58,46,70,146,36,21,61,164,149,141,38,78,71,178,120,37,12,189,17,133,46,39,46,59,86,75,135,37,65,72,35,38,75,15,74,131,143,33,56,56,148,132,132,95,121,71,154,32,12,144,111,76,47,39,170,62,151,155,183,154,92,156,149,38,35,150,63,116,74,33,36,45,148,144,46,46,167,30,159,56,37,103,72,153,148,133,193,96,70,130,33,56,73,184,80,132,141,38,66,150,138,125,60,169,73,133,34,34,59,146,78,117,130,98,8,77,160,35,38,176,74,131,134,33,61,176,136,39,132,153,14,75,139,37,37,189,12,81,133,128,43,70,86,75,73,44,56,115,77,35,78,166,63,105,79,57,57,28,152,133,39,160,172,75,150,56,32,98,165,156,140,128,63,43,155,130,33,99,56,8,122,144,35,97,147,58,154,138,57,12,145,148,129,168,146,112,75,135,130,37,153,149,102,76,128,82,58,134,57,36,56,56,157,149,129,100,83,14,178,158,138,57,88,128,148,168,43,62,146,36,104,41,32,149,149,156,100,63,18,82,134,143,99,56,148,98,65,149,47,11,147,135,37,143,60,119,76,133,34,170,83,154,117,21,56,32,72,152,38,35,58,58,101,70,151,36,28,73,39,34,62,59,30,150,135,98,32,144,99,140,123,39,34,151,146,167,16,61,20,85,132,35,55,135,142,131,36,36,33,148,128,85,34,158,193,104,139,142,89,37,144,162,137,190,39,31,151,131,126,61,118,110,129,149,38,47,78,147,105,98,81,88,126,115,58,74,160,182,154,130,21,32,127,129,84,35,58,75,87,70,167,36,56,128,77,39,62,59,11,150,61,103,24,189,157,99,157,193,63,74,10,147,37,32,32,115,132,50,201,147,131,138,131,33,60,81,128,58,34,83,112,121,139,139,61,37,162,180,69,63,34,82,70,131,86,178,56,173,122,22,35,35,11,113,67,158,60,114,137,137,34,101,62,112,171,135,21,56,149,129,114,75,23,113,19,143,53,93,61,145,118,144,59,75,154,130,130,37,60,172,57,140,137,39,136,91,154,33,33,32,32,149,132,192,30,58,151,101,146,131,76,12,85,188,34,19,59,154,95,61,52,72,68,169,38,63,58,82,155,112,109,36,56,98,132,129,38,78,147,161,94,57,60,119,81,128,188,196,31,138,104,147,61,107,153,172,144,38,34,116,22,91,151,48,28,173,156,169,50,172,197,90,113,159,33,36,93,93,128,34,79,22,95,86,37,32,61,17,128,114,86,176,187,94,131,171,57,81,115,115,132,55,59,86,97,150,72,48,72,102,84,63,58,54,87,117,109,49,56,118,65,81,35,75,128,131,37,87,57,57,12,56,118,193,31,197,75,135,49,37,77,167,136,211,197,58,94,134,53,83,203,190,149,144,59,38,90,159,143,116,36,33,81,133,58,18,51,204,147,36,130,61,80,80,171,119,216,3,94,138,134,33,32,73,98,80,22,59,35,117,139,138,91,57,184,123,69,148,30,87,79,83,159,32,56,52,132,50,88,78,178,101,125,138,24,36,152,133,153,89,50,86,95,139,49,32,57,169,153,140,34,39,82,146,134,199,37,76,185,89,62,35,133,183,176,146,36,29,84,137,148,127,54,74,97,95,32,21,32,202,93,81,128,18,113,94,131,117,36,84,61,118,80,62,35,86,135,116,33,57,53,119,81,145,58,195,22,83,83,32,56,89,80,62,126,34,54,51,91,138,17,57,88,173,80,169,50,22,83,135,49,36,218,169,81,128,39,39,54,143,117,121,135,28,92,85,59,62,75,196,18,67,33,57,13,88,145,157,132,59,159,135,138,56,85,206,114,81,34,54,54,94,131,83,194,56,16,80,129,22,38,83,135,113,130,33,33,102,133,157,58,217,31,86,117,164,56,37,153,129,137,137,46,54,74,138,53,121,61,128,65,129,35,38,86,135,33,67,36,48,93,140,30,34,27,143,95,171,37,28,21,136,132,88,140,58,94,91,36,33,194,128,128,127,193,62,22,90,159,52,32,211,20,119,64,39,58,87,146,112,163,49,25,72,132,35,63,86,190,30,195,57,24,84,93,145,63,34,62,15,78,32,32,33,80,24,169,39,58,23,87,131,57,33,181,198,85,65,62,50,86,147,87,36,36,53,73,140,63,51,35,197,139,83,56,61,85,92,62,123,192,200,87,82,134,57,73,194,153,92,35,209,31,83,139,56,36,53,20,123,73,54,113,59,151,158,194,37,52,80,132,129,38,86,58,27,191,33,24,24,128,129,165,184,38,95,78,130,56,33,60,172,81,63,43,27,82,131,60,174,49,25,144,156,35,55,34,163,87,120,57,24,81,133,58,54,58,143,31,95,37,32,28,129,169,84,39,34,14,27,60,20,37,37,137,80,38,35,22,86,142,170,33,48,20,84,128,39,176,155,182,131,121,32,37,12,129,50,204,163,34,125,87,33,57,8,128,140,89,62,38,31,95,130,190,80,36,60,123,168,39,10,79,174,131,124,41,80,85,144,149,62,39,135,30,135,33,48,48,133,133,196,38,35,31,83,32,125,36,180,129,126,34,34,27,94,131,180,56,37,137,65,132,42,59,95,175,56,33,33,33,17,69,34,82,14,38,174,134,37,52,85,141,136,184,39,39,30,94,36,199,32,32,128,85,129,38,86,135,130,116,12,9,17,123,140,34,63,87,184,167,130,154,25,92,129,59,59,10,34,87,82,53,12,32,145,133,132,193,35,159,86,32,32,57,33,24,197,128,54,94,193,131,131,194,170,37,80,156,59,38,39,175,130,125,33,33,93,137,27,58,35,38,83,71,56,28,85,129,62,84,34,39,18,82,131,33,32,32,52,63,156,35,50,90,130,178,36,36,102,29,39,39,34,196,59,124,130,32,56,80,129,62,114,39,39,67,158,33,12,32,88,128,210,55,38,38,95,37,135,53,33,33,132,34,34,120,82,35,29,175,44,56,118,133,153,35,63,18,70,120,33,9,60,133,63,200,11,38,86,86,32,32,33,129,129,73,181,34,94,82,131,60,52,37,107,65,156,35,35,135,34,87,48,33,36,20,128,148,39,59,19,130,135,32,16,80,129,129,168,124,58,18,134,36,57,199,128,98,85,85,14,71,130,130,166,113,33,20,69,39,34,200,82,131,117,183,37,8,85,129,38,35,39,34,130,131,36,33,84,32,133,77,59,35,78,83,21,56,85,129,17,81,133,38,82,113,131,162,150,56,16,153,129,35,83,170,130,154,36,36,20,84,128,39,193,35,35,86,130,37,37,57,33,129,93,63,34,30,138,171,203,175,37,61,85,38,11,95,39,159,159,33,33,9,145,128,34,205,35,35,131,154,13,32,85,129,141,148,39,18,134,134,33,57,228,128,98,85,35,14,22,135,130,61,174,33,129,73,133,58,82,35,11,57,37,13,32,77,132,132,38,39,142,23,167,36,17,81,133,157,129,35,22,75,130,32,202,33,36,73,132,54,15,74,130,57,57,49,32,68,129,129,35,34,135,154,130,162,36,12,29,39,205,50,167,38,135,147,32,37,33,129,59,73,34,10,27,40,12,29,32,37,153,80,35,38,62,86,56,74,155,0,17,129,15,10,87,35,131,130,32,32,1,80,33,136,34,34,34,130,94,40,36,37,128,157,144,38,38,71,135,61,32,132,33,129,128,10,34,42,167,155,78,135,8,21,144,149,35,39,34,130,130,162,17,5,37,133,63,42,35,2,75,90,32,32,57,141,129,197,34,6,30,36,134,57,56,61,72,80,38,35,34,34,158,82,36,36,9,93,128,34,55,38,22,150,32,28,8,61,165,136,76,201,10,74,131,0,60,32,128,77,65,35,2,66,86,32,37,36,36,12,73,10,34,94,131,60,57,41,32,16,172,132,35,63,34,79,87,131,20,29,133,32,148,59,38,14,75,95,130,79,33,129,22,43,34,10,30,82,60,33,61,32,16,80,149,14,34,34,154,167,151,33,84,73,128,58,59,35,197,131,130,1,61,195,133,145,168,34,10,74,134,9,12,37,37,16,149,35,38,204,78,130,56,40,20,17,128,15,10,74,131,135,95,60,56,1,153,153,62,39,34,67,131,146,33,9,37,152,68,153,38,38,135,139,130,207,33,12,76,93,34,6,193,38,151,146,32,32,49,80,129,7,34,181,70,147,36,36,64,84,10,58,156,35,14,192,142,8,32,33,179,194,39,34,15,79,134,60,60,41,173,68,132,38,31,83,142,159,61,33,12,76,133,46,34,47,11,66,75,135,32,56,136,132,62,62,39,67,67,134,5,75,61,37,128,89,35,2,75,135,130,61,17,60,73,84,34,18,38,35,199,190,44,37,8,36,129,35,46,212,131,143,134,29,69,37,32,77,177,35,14,139,130,67,80,36,5,168,152,34,39,82,155,158,100,37,37,65,72,14,11,46,170,150,61,33,12,64,211,128,58,63,19,11,130,135,8,13,132,177,62,62,34,3,74,138,57,32,32,37,169,129,42,35,83,147,134,159,33,33,69,88,137,58,205,7,78,135,32,32,8,129,141,149,39,34,10,74,131,36,57,37,173,72,47,14,14,130,154,61,56,33,194,64,132,34,58,82,216,75,78,16,28,65,133,129,47,39,18,67,56,33,33,33,140,148,157,42,11,66,78,130,32,207,222,36,76,168,34,10,198,155,151,78,32,37,169,141,35,35,178,159,196,131,36,9,69,128,148,34,34,171,75,83,44,8,16,136,73,153,62,15,15,158,134,33,78,61,37,157,77,14,26,192,154,135,37,33,36,64,88,43,39,196,158,155,57,32,28,21,180,64,156,46,212,34,82,45,9,64,173,152,72,129,11,62,75,146,150,110,40,57,69,84,34,10,131,155,83,154,37,13,21,132,42,165,39,154,10,74,33,17,32,199,133,145,213,31,59,71,135,21,65,129,153,156,157,27,10,79,155,155,60,32,202,152,145,35,19,197,154,74,56,224,73,0,128,34,43,156,42,71,130,32,28,65,156,132,7,62,34,67,134,33,12,37,32,32,145,42,22,66,154,135,139,222,36,9,136,39,34,43,193,174,75,135,32,16,65,132,144,38,227,208,70,79,33,17,93,181,65,39,47,19,2,130,135,32,207,36,36,180,168,34,67,158,155,134,173,44,8,72,153,35,35,75,166,74,131,36,40,69,164,128,34,38,38,75,135,44,32,207,5,153,141,38,18,67,200,138,36,75,16,198,72,136,35,14,78,142,16,224,194,33,81,152,10,39,217,155,138,150,135,21,65,132,149,42,152,39,67,79,134,12,69,21,148,77,63,22,66,142,135,74,158,40,36,173,63,23,27,196,155,130,194,202,21,136,153,42,62,39,183,74,131,24,33,199,137,39,46,196,22,22,78,37,21,65,33,156,69,164,27,18,231,134,40,33,37,32,157,129,42,212,201,130,159,32,40,33,76,137,43,196,35,22,75,134,130,21,68,156,165,35,157,18,82,171,33,29,37,37,157,65,144,19,11,139,150,32,79,45,36,73,153,15,43,22,155,134,36,28,32,136,141,35,62,34,208,74,150,40,73,199,181,89,68,22,35,42,130,151,25,45,20,156,133,128,18,34,193,167,36,33,32,215,77,132,19,42,197,94,32,150,17,20,152,157,137,77,38,11,66,75,135,21,68,153,160,73,18,15,74,167,154,56,218,32,207,133,39,19,47,154,147,32,74,151,20,76,145,144,196,171,66,71,146,37,8,72,153,132,35,39,10,67,146,17,33,33,173,128,34,180,14,75,159,139,79,60,9,64,144,152,43,46,193,19,19,159,37,25,13,156,38,62,39,63,67,135,33,60,152,152,128,34,174,62,75,139,41,28,195,60,64,132,157,34,193,186,138,24,36,32,32,176,144,26,58,187,159,135,57,60,36,137,133,34,46,14,38,66,155,16,56,60,156,69,144,181,94,34,143,9,33,45,37,65,128,189,11,75,63,16,21,236,20,40,64,157,34,30,174,174,66,36,32,32,189,84,35,35,216,63,155,138,134,33,76,16,133,39,62,62,66,139,135,21,57,17,69,153,148,15,59,196,154,33,154,61,37,68,149,35,62,34,94,179,147,44,20,56,128,43,46,243,35,71,138,44,65,32,57,81,152,181,10,191,186,146,36,36,89,65,133,129,31,35,142,159,32,158,215,40,140,129,46,196,62,66,66,130,16,21,198,132,64,42,34,27,59,146,151,33,56,16,85,34,19,31,66,142,130,41,60,206,69,157,168,34,67,138,162,36,154,21,32,65,172,38,201,34,91,154,155,33,33,157,148,39,34,220,59,58,158,25,72,68,60,156,35,23,23,74,155,33,240,207,61,65,77,156,19,58,58,154,79,143,36,20,194,152,18,217,59,62,130,139,219,32,60,57,149,35,181,39,82,155,17,36,187,56,65,77,153,22,63,170,82,147,60,57,36,148,152,30,59,59,155,33,135,16,37,156,153,35,204,227,87,200,131,29,36,56,157,34,34,216,35,159,147,61,68,57,40,88,133,63,23,193,155,151,33,139,52,65,136,141,22,35,58,63,74,37,20,33,152,188,34,34,19,50,150,130,28,32,184,60,76,128,54,39,79,174,17,33,76,16,168,136,156,38,59,63,135,32,198,20,29,152,133,30,213,35,171,155,130,21,37,57,156,38,38,63,63,158,154,17,36,61,56,80,136,62,22,221,154,32,56,207,57,189,152,18,34,196,59,62,134,61,25,89,153,153,35,38,34,23,67,138,24,33,194,152,68,39,209,83,159,142,16,56,57,57,192,152,161,27,196,155,155,36,228,37,157,153,14,22,192,63,91,37,134,17,194,152,39,39,34,59,209,155,178,16,52,210,189,35,38,51,30,213,158,20,36,199,56,153,153,19,59,63,63,16,70,155,60,57,185,133,23,193,50,209,130,66,25,37,149,148,136,35,197,58,158,134,17,36,61,56,39,156,172,38,159,154,193,37,57,36,152,128,23,30,196,59,191,159,33,25,32,136,59,212,58,58,187,154,155,60,52,206,56,34,169,59,35,159,135,21,61,222,57,189,168,58,54,196,155,12,36,194,37,168,129,34,39,192,58,158,191,60,33,36,157,15,58,59,59,59,154,130,16,49,198,60,38,230,34,34,191,131,29,60,61,56,213,172,144,38,197,190,32,32,36,33,211,133,23,63,220,62,154,142,37,32,56,129,84,129,58,34,18,158,134,53,206,152,61,34,208,59,59,201,147,154,56,60,88,145,152,39,27,39,146,134,211,49,32,202,129,22,59,59,58,196,143,48,36,20,199,152,156,229,62,59,135,175,32,32,210,57,133,168,34,34,193,139,86,211,77,61,136,152,156,19,212,63,37,143,60,57,152,152,23,18,38,216,83,147,130,37,61,195,153,172,58,34,58,138,131,20,32,206,56,188,129,26,38,191,130,143,32,57,57,132,145,51,39,196,200,183,158,36,56,195,152,81,59,34,54,63,155,155,33,33,206,168,181,205,38,35,142,150,52,195,60,172,145,133,38,58,158,158,158,60,52,61,195,141,38,62,221,63,154,134,186,48,76,133,80,39,34,50,174,151,147,56,32,60,169,81,173,39,34,155,154,57,60,56,61,157,132,59,59,197,187,91,191,222,57,199,133,133,58,220,62,139,142,33,56,198,165,153,34,43,34,91,158,167,29,33,152,181,141,153,38,83,130,150,49,222,60,76,156,128,34,39,193,246,158,142,56,56,144,141,38,197,227,192,143,155,60,48,57,61,161,173,238,50,55,130,61,56,195,219,199,84,157,34,34,158,134,158,36,61,61,132,148,234,35,58,190,154,131,48,36,232,140,54,27,55,255,225,142,130,49,37,156,145,156,51,34,82,151,131,20,36,199,52,168,141,35,192,159,190,143,32,48,60,140,132,18,58,193,256,212,138,257,49,198,172,50,43,216,94,175,155,60,33,57,157,18,136,59,35,62,135,255,56,198,248,169,133,137,51,34,193,83,60,60,52,37,207,141,59,22,234,170,130,32,33,48,57,132,80,63,38,55,209,134,175,37,195,210,137,168,168,39,193,131,143,36,60,52,144,132,129,38,59,197,256,166,60,53,53,152,137,63,51,220,55,159,131,56,41,215,165,81,129,230,74,196,186,130,33,36,218,169,157,153,55,50,201,166,61,36,37,81,204,140,39,63,220,83,178,60,56,85,198,141,62,62,247,74,196,143,171,53,48,152,258,63,250,62,126,131,130,52,215,257,81,153,128,34,193,155,191,60,228,56,56,156,132,38,38,197,209,134,142,36,33,199,137,63,58,35,59,225,147,154,52,202,73,145,137,34,39,193,155,98,60,53,223,173,152,136,35,59,175,175,175,174,33,57,137,133,54,34,196,78,155,179,49,49,207,169,59,197,58,58,213,134,143,48,199,203,23,149,35,90,192,190,74,56,36,57,172,153,152,54,63,59,20,150,159,32,215,219,156,35,59,192,54,134,135,260,57,37,164,136,58,35,256,139,130,32,44,33,202,245,59,34,34,193,134,171,57,223,56,56,136,38,180,39,216,256,159,33,60,57,145,137,145,63,38,59,182,134,32,61,198,73,59,62,34,58,193,134,36,179,194,218,208,58,129,35,212,135,167,131,195,57,240,168,152,87,58,196,78,60,57,32,85,195,153,62,55,63,79,174,151,48,36,203,173,63,133,217,59,197,151,32,49,49,36,132,165,136,58,193,143,60,130,32,32,195,205,169,98,39,170,178,135,171,36,60,128,149,181,58,59,59,135,130,150,56,120,215,209,59,63,58,213,138,174,57,194,252,168,160,73,59,55,192,61,183,36,36,237,137,129,39,51,38,221,131,159,256,37,149,231,133,59,34,82,163,143,146,48,218,72,219,152,35,38,192,187,154,222,33,222,101,141,34,46,51,193,162,134,32,77,37,152,172,201,50,192,34,146,174,33,48,121,161,168,172,226,257,155,130,175,56,33,76,230,133,161,82,213,75,256,57,37,61,258,152,132,47,221,39,166,61,56,31,84,168,164,58,51,193,234,182,134,61,198,198,169,59,59,34,39,134,155,83,232,194,32,132,140,246,59,230,178,56,56,33,214,245,161,136,208,63,75,167,60,37,32,229,195,169,59,50,79,58,138,131,57,48,206,144,145,35,58,47,197,134,170,49,36,259,128,176,34,256,205,200,259,57,57,72,52,129,129,62,257,79,79,166,171,260,33,32,173,63,35,71,38,204,74,61,32,172,258,59,152,34,39,138,167,146,48,240,77,222,129,129,242,204,216,56,134,174,33,203,185,58,39,58,75,256,194,37,32,202,165,176,189,38,62,250,167,78,69,48,199,256,172,157,47,38,216,130,170,61,36,33,113,133,152,213,39,35,257,57,130,56,207,224,169,38,247,34,260,61,36,259,33,214,85,153,58,250,95,130,134,175,202,61,260,216,168,63,208,58,246,257,218,32,37,219,148,156,242,123,79,256,202,219,40,214,161,168,129,38,38,221,182,74,159,233,195,172,137,201,256,235,205,162,171,178,37,105,168,153,35,59,130,74,167,61,49,36,57,72,168,205,34,196,162,60,159,194,222,180,161,192,62,39,235,163,75,262,228,168,259,145,38,38,247,197,167,138,13,253,93,232,128,58,256,205,257,175,57,44,61,202,165,38,38,62,79,213,56,75,259,206,168,149,169,63,78,239,194,32,79,207,257,73,152,39,204,58,217,147,154,57,32,32,198,157,180,227,192,74,143,134,171,259,194,161,164,152,205,50,216,163,155,32,207,73,184,59,96,63,250,182,33,191,206,72,173,205,189};

char str[MAXN];

void stevehalim() {
	for(int n=1; n<MAXN; n++) {
		if (n <= 2) { mex[n] = 0; continue; }
		set<int> jog;
		for(int i=3; i<=n-2; i++) {
			jog.insert(mex[i-1]^mex[n-i]);
		}
		int cnt = 0;
		while(jog.count(cnt)) cnt++;
		mex[n] = cnt;
	}
	//for(int n=1; n<MAXN; n++) printf("%d,", mex[n]);
}

/*
 * URI 1130
 */

int main() {
	//stevehalim();
	int N;
	while(scanf("%d", &N), N) {
		scanf(" %s", str);
		bool ok = false;
		for(int i=1; i<N; i++) {
			if (str[i] == 'X' && str[i-1] == 'X') ok = true;
		}
		for(int i=2; i<N; i++) {
			if (str[i] == 'X' && str[i-2] == 'X') ok = true;
		}
		if (ok) {
			printf("S\n"); continue;
		}
		int last = -3, ans = 0;
		for(int i=0; i<N; i++) {
			if (str[i] == 'X') {
				//printf("size %d\n", i-last-1);
				ans ^= mex[i-last-1];
				last = i;
			}
		}
		//printf("size %d\n", N-last+1);
		ans ^= mex[N-last+1];
		if (ans != 0) printf("S\n");
		else printf("N\n");
	}
	return 0;
}
\end{vncode}
\endgroup


\subsection{Bộ ba Pythagoras}

Mọi bộ (a,b,c), a\textasciicircum{}2+b\textasciicircum{}2=c\textasciicircum{}2, sinh bởi a=k(m\textasciicircum{}2-n\textasciicircum{}2), b=2kmn, c=k(m\textasciicircum{}2+n\textasciicircum{}2).

\subsection{Ma trận}

Các phép toán ma trận cơ bản.

\codepath{Math/matrix.cpp}

\begingroup\scriptsize
\begin{vncode}
#include <vector>
using namespace std;

/*
 * Matrix operations
 */

typedef long long ll;
typedef vector< vector< double > > matrix;

matrix operator +(matrix a, matrix b) {
	int n = (int)a.size();
	int m = (int)a[0].size();
	matrix c;
	c.resize(n);
	for(int i=0; i<n; i++) {
		c[i].resize(m);
		for(int j=0; j<m; j++)
			c[i][j] = a[i][j] + b[i][j];
	}
	return c;
}

matrix operator *(matrix a, matrix b) {
	int n = (int)a.size();
	//assert(a[0].size() == b.size());
	int m = (int)b.size();
	int p = (int)b[0].size();
	matrix c(n, vector<double>(p));
	vector<double> col(m);
	for (int j = 0; j < p; j++) {
		for (int k = 0; k < m; k++)
			col[k] = b[k][j];
		for (int i = 0; i < n; i++) {
			double s = 0;
			for (int k = 0; k < m; k++)
				s += a[i][k] * col[k];
			c[i][j] = s;
		}
	}
	return c;
}

matrix operator *(double k, matrix a) {
	int n = (int)a.size();
	int m = (int)a[0].size();
	for(int i=0; i<n; i++) for(int j=0; j<m; j++)
		a[i][j] *= k;
	return a;
}

matrix operator -(matrix a, matrix b) {
	return a + ((-1.0) * b);
}

matrix id(int n) {
	matrix c(n, vector<double>(n));
	for(int i = 0; i < n; i++) c[i][i] = 1;
	return c;
}

#include <cstdio>
void printmatrix(matrix & a) {
	int n = (int)a.size();
	int m = (int)a[0].size();
	for(int i=0; i<n; i++) {
		for(int j=0; j<m; j++) {
			printf("%c%.2f%c ", (i==0 && j==0 ? '[' : ' '), a[i][j], (j < m-1 ? ',' : i == n-1? ']' : ';'));
		}
		printf("\n");
	}
}

/*
 * COMPILATION TEST
 */

int main() {}
\end{vncode}
\endgroup


\subsection{Lũy thừa ma trận và Fibonacci}

Matrix exponentiation.

\codepath{Math/matrixexp.cpp}

\begingroup\scriptsize
\begin{vncode}
#include <cstdio>
#include <algorithm>
#include <vector>
#include <ctime>
#include <cmath>
#define EPS 1e-5
using namespace std;

typedef long long ll;
typedef vector< vector< double > > matrix;

matrix operator +(matrix a, matrix b) {
	int n = (int)a.size();
	int m = (int)a[0].size();
	matrix c;
	c.resize(n);
	for(int i=0; i<n; i++) {
		c[i].resize(m);
		for(int j=0; j<m; j++) {
			c[i][j] = a[i][j] + b[i][j];
		}
	}
	return c;
}

matrix operator *(matrix a, matrix b) {
	int n = (int)a.size();
	//assert(a[0].size() == b.size());
	int m = (int)b.size();
	int p = (int)b[0].size();
	matrix c(n, vector<double>(p));
	double col[m];
	for (int j = 0; j < p; j++) {
		for (int k = 0; k < m; k++)
			col[k] = b[k][j]; //cache friendly
		for (int i = 0; i < n; i++) {
			double s = 0;
			for (int k = 0; k < m; k++)
				s += a[i][k] * col[k];
			c[i][j] = s;
		}
	}
	return c;
}

matrix operator *(double k, matrix a) {
	int n = (int)a.size();
	int m = (int)a[0].size();
	for(int i=0; i<n; i++) {
		for(int j=0; j<m; j++) {
			a[i][j] *= k;
		}
	}
	return a;
}

matrix operator -(matrix a, matrix b) {
	return a + ((-1.0) * b);
}

void printmatrix(matrix & a) {
	int n = (int)a.size();
	int m = (int)a[0].size();
	for(int i=0; i<n; i++) {
		for(int j=0; j<m; j++) {
			printf("%c%.2f%c ", (i==0 && j==0 ? '[' : ' '), a[i][j], (j < m-1 ? ',' : i == n-1? ']' : ';'));
		}
		printf("\n");
	}
}

/*
 * Matrix Exp and fast Fibonacci
 */

matrix id(int n) {
	matrix c(n, vector<double>(n));
	for(int i = 0; i < n; i++) c[i][i] = 1;
	return c;
}

matrix matrixExp(matrix a, int n) {
	if (n == 0) return id(a.size());
	matrix c = matrixExp(a, n/2);
	c = c*c;
	if (n%2 != 0) c = c*a;
	return c;
}

matrix fibo() {
	matrix c(2, vector<double>(2, 1));
	c[1][1] = 0;
	return c;
}

double fibo(int n) {
	matrix f = matrixExp(fibo(), n);
	return f[0][1];
}

/*
 * TEST MATRIX
 */

int main() {
	int n;
	while(scanf("%d", &n) != EOF) {
		printf("%f\n", fibo(n));
	}
	n = 10000;
	double fib[10000];
	clock_t t = clock();
	for(int i = 0; i < 10000; i++) {
		fib[i] = fibo(i);
	}
	t = clock() - t;
	printf("time %.3f\n", t*1.0/CLOCKS_PER_SEC);
	return 0;
}
\end{vncode}
\endgroup


\subsection{Hệ tuyến tính: định thức và khử Gauss}

Determinant và Gaussian elimination.

\codepath{Math/gauss.cpp}

\begingroup\scriptsize
\begin{vncode}
#include <cstdio>
#include <algorithm>
#include <vector>
#include <ctime>
#include <cmath>
#define EPS 1e-5
using namespace std;

typedef long long ll;
typedef vector< vector< double > > matrix;

matrix operator +(matrix a, matrix b) {
	int n = (int)a.size();
	int m = (int)a[0].size();
	matrix c;
	c.resize(n);
	for(int i=0; i<n; i++) {
		c[i].resize(m);
		for(int j=0; j<m; j++) {
			c[i][j] = a[i][j] + b[i][j];
		}
	}
	return c;
}

matrix operator *(matrix a, matrix b) {
	int n = (int)a.size();
	if (a[0].size() != b.size()) printf("fail\n");
	int m = (int)b.size();
	int p = (int)b[0].size();
	matrix c;
	c.resize(n);
	for(int i=0; i<n; i++) {
		c[i].assign(p, 0);
		for(int j=0; j<p; j++) {
			for(int k=0; k<m; k++) {
				c[i][j] += a[i][k]*b[k][j];
			}
		}
	}
	return c;
}

matrix operator *(double k, matrix a) {
	int n = (int)a.size();
	int m = (int)a[0].size();
	for(int i=0; i<n; i++) {
		for(int j=0; j<m; j++) {
			a[i][j] *= k;
		}
	}
	return a;
}

matrix operator -(matrix a, matrix b) {
	return a + ((-1.0) * b);
}

matrix id(int n) {
	matrix c; c.resize(n);
	for(int i=0; i<n; i++) {
		c[i].assign(n, 0);
		c[i][i] = 1;
	}
	return c;
}

void printmatrix(matrix & a) {
	int n = (int)a.size();
	int m = (int)a[0].size();
	for(int i=0; i<n; i++) {
		for(int j=0; j<m; j++) {
			printf("%c%.2f%c ", (i==0 && j==0 ? '[' : ' '), a[i][j], (j < m-1 ? ',' : i == n-1? ']' : ';'));
		}
		printf("\n");
	}
}

/*
 * Gauss elimination
 */

void switchLines(matrix & a, int i, int j) {
	int m = (int)a[i].size();
	for(int k = 0; k < m; k++) swap(a[i][k], a[j][k]);
}

void lineSumTo(matrix & a, int i, int j, double c) {
	int m = (int)a[0].size();
	for(int k = 0; k < m; k++) a[j][k] += c*a[i][k];
}

bool gauss(matrix & a, matrix & b, int & switches) {
	switches = 0;
	int n = (int)a.size();
	int m = (int)a[0].size();
	for(int i = 0, l; i < min(n, m); i++) {
		l = i;
		while(l < n && fabs(a[l][i]) < EPS) l++;
		if (l == n) return false;
		switchLines(a, i, l);
		switchLines(b, i, l);
		switches++;
		for(int j=0; j<n; j++) {
			if (i == j) continue;
			double p = -a[j][i] / a[i][i];
			lineSumTo(a, i, j, p);
			lineSumTo(b, i, j, p);
		}
	}
	return true;
}

double det(matrix a) {
	int n = a.size();
	matrix b(n);
	for(int i=0; i<n; i++) b[i].resize(1);
	int sw = 0;
	if (gauss(a, b, sw)) {
		double ans = 1;
		for(int i=0; i<n; i++) ans *= a[i][i];
		return sw % 2 == 0 ? ans : -ans;
	}
	return 0.0;
}

/*
 * TEST MATRIX
 */

bool testgauss(int ntest) {
	srand(time(NULL));
	for(int n=1; n<ntest; n++) {
		double d;
		matrix A = id(n), B, X;
		B.resize(n);
		X.resize(n);
		for(int i=0; i<n; i++) {
			X[i].push_back(rand()%10);
			for(int j=0; j<n; j++) {
				A[i][j] = -2 + rand()%10;
			}
		}
		B = A*X;
		d = det(A);
		int sw;
		if (gauss(A, B, sw)) {
			for(int i=0; i<n; i++) {
				if (fabs(X[i][0] - B[i][0]/A[i][i]) > EPS) {
					printf("failed test %d\n", n);
					return false;
				}
			}
		}
	}
	printf("all gauss tests passed\n");
	return true;
}

bool testdet(int ntest) {
	matrix a;
	for(int n=1; n<=ntest; n++) {
		a = id(2);
		a[0][0] = rand()%100;
		a[1][0] = rand()%100;
		a[0][1] = rand()%100;
		a[1][1] = rand()%100;
		if (fabs(det(a)-(a[0][0]*a[1][1] - a[0][1]*a[1][0])) > EPS) {
			printf("falied test  %d\n", n);
			return false;
		}
	}
	
	printf("all gauss tests passed\n");
	return true;
}

int main() {
	testgauss(100);
	testdet(100);
	return 0;
}
\end{vncode}
\endgroup


\subsection{Nhân ma trận thưa}

Sparse matrix multiplication.

\codepath{Math/matrixsparse.cpp}

\begingroup\scriptsize
\begin{vncode}
vector< vector<int> > adjA, adjB;

matrix sparsemult(matrix a, matrix b) {
	int n = (int)a.size();
	if (a[0].size() != b.size()) printf("fail\n");
	int m = (int)b.size();
	int p = (int)b[0].size();
	adjA.resize(n);
	for(int i=0; i<n; i++) {
		adjA[i].clear();
		for(int k=0; k<m; k++) {
			if (fabs(a[i][k]) > EPS) adjA[i].push_back(k);
		}
	}
	adjB.resize(p);
	for(int j=0; j<p; j++) {
		adjB[j].clear();
		for(int k=0; k<m; k++) {
			if (fabs(b[k][j]) > EPS) adjB[j].push_back(k);
		}
	}
	matrix c;
	c.resize(n);
	for(int i=0; i<n; i++) {
		c[i].assign(p, 0);
		for(int j=0; j<p; j++) {
			for(int u=0, v=0, k; u<(int)adjA[i].size() && v<(int)adjB[j].size();) {
				if (adjA[i][u] > adjB[j][v]) v++;
				else if (adjA[i][u] < adjB[j][v]) u++;
				else {
					k = adjA[i][u];
					c[i][j] += a[i][k]*b[k][j];
					u++; v++;
				}
			}
		}
	}
	return c;
}
\end{vncode}
\endgroup


\subsection{Phương pháp Gauss-Seidel}

Lặp Gauss-Seidel.

\codepath{Math/gaussseidel.cpp}

\begingroup\scriptsize
\begin{vncode}
#include <cstdio>
#include <algorithm>
#include <vector>
#include <cmath>
#define EPS 1e-5
using namespace std;

typedef long long ll;
typedef vector< vector< double > > matrix;

matrix operator +(matrix a, matrix b) {
	int n = (int)a.size();
	int m = (int)a[0].size();
	matrix c;
	c.resize(n);
	for(int i=0; i<n; i++) {
		c[i].resize(m);
		for(int j=0; j<m; j++) {
			c[i][j] = a[i][j] + b[i][j];
		}
	}
	return c;
}

matrix operator *(matrix a, matrix b) {
	int n = (int)a.size();
	if (a[0].size() != b.size()) printf("fail\n");
	int m = (int)b.size();
	int p = (int)b[0].size();
	matrix c;
	c.resize(n);
	for(int i=0; i<n; i++) {
		c[i].assign(p, 0);
		for(int j=0; j<p; j++) {
			for(int k=0; k<m; k++) {
				c[i][j] += a[i][k]*b[k][j];
			}
		}
	}
	return c;
}

matrix operator *(double k, matrix a) {
	int n = (int)a.size();
	int m = (int)a[0].size();
	for(int i=0; i<n; i++) {
		for(int j=0; j<m; j++) {
			a[i][j] *= k;
		}
	}
	return a;
}

matrix operator -(matrix a, matrix b) {
	return a + ((-1.0) * b);
}

/*
 * Gauss-Seidel method
 */

matrix gaussSeidel(matrix & a, matrix & b, double PREC) {
	int n = (int)a.size();
	matrix x = b, xp = b;
	double error;
	do {
		error = 0.0;
		for(int i=0; i<n; i++) {
			xp[i][0] = b[i][0];
			for(int j=0; j<n; j++) {
				if (i < j) xp[i][0] -= a[i][j]*xp[j][0];
				if (i > j) xp[i][0] -= a[i][j]*x[j][0];
			}
			xp[i][0] /= a[i][i];
			error = max(error, fabs(xp[i][0]-x[i][0]));
		}
		x = xp;
	} while(error > PREC);
	return xp;
}

/*
 * TEST MATRIX
 */
#include <ctime>

bool test(int nTests) {
	srand(time(NULL));
	for(int n=1; n<=nTests; n++) {
		double PREC = 1e-4;
		matrix A, B, X, Xp;
		A.resize(n);
		B.resize(n);
		X.resize(n);
		for(int i=0; i<n; i++) {
			X[i].push_back((rand()%10)/10.0);
			A[i].resize(n);
			for(int j=0; j<n; j++) {
				A[i][j] = -(rand()%10)/100000.0;
			}
			A[i][i] = 1.0;
		}
		B = A*X;
		Xp = gaussSeidel(A, B, PREC);
		for(int i=0; i<n; i++) {
			if (fabs(X[i][0] - Xp[i][0]) > PREC) {
				printf("failed test n = %d\n", n);
				return false;
			}
		}
	}
	printf("all tests passed\n");
	return true;
}

int main() {
	test(500);
	return 0;
}
\end{vncode}
\endgroup


\subsection{XOR-SAT}

Hệ phương trình xor.

\codepath{Math/xorsat.cpp}

\begingroup\scriptsize
\begin{vncode}
#include <vector>
#include <bitset>
#include <algorithm>
using namespace std;

/*
 * XOR-SAT or Gauss XOR Elimination O(n*(m^2)/64)
 */

#define MAXN 2009

vector<int> B;
vector< bitset<MAXN> > A;
bitset<MAXN> x;

bool check() {
	int n = A.size(), m = MAXN;
	for(int i = 0; i < n; i++) {
		int acum = 0;
		for(int j = 0; j < m; j++) {
			if (A[i][j]) acum ^= x[j];
		}
		if (acum != B[i]) return false;
	}
	return true;
}

bool gaussxor() {
	int cnt = 0, n = A.size(), m = MAXN;
	bitset<MAXN> vis; vis.set();
	for(int j = m-1, i; j >= 0; j--) {
		for(i = cnt; i < n; i++) {
			if (A[i][j]) break;
		}
		if (i == n) continue;
		swap(A[i], A[cnt]); swap(B[i], B[cnt]);
		i = cnt++; vis[j] = 0;
		for(int k = 0; k < n; k++) {
			if (i == k || !A[k][j]) continue;
			A[k] ^= A[i]; B[k] ^= B[i];
		}
	}
	x = vis;
	for(int i = 0; i < n; i++) {
		int acum = 0;
		for(int j = 0; j < m; j++) {
			if (!A[i][j]) continue;
			if (!vis[j]) {
				vis[j] = 1;
				x[j] = acum^B[i];
			}
			acum ^= x[j];
		}
		if (acum != B[i]) return false;
	}
	return true;
}

/*
 * Codeforces 101908M
 */

#include <cstdio>

int main() {
	int n, m;
	scanf("%d %d", &n, &m);
	for(int i = 0; i < n; i++) {
		char c, str[10];
		scanf(" %c", &c);
		int val = 1;
		bitset<MAXN> eq;
		while(scanf(" %s", str) != EOF && str[0] != 'a') {
			//printf("%s ", str);
			if (str[0] == 'n') val ^= 1;
			if (str[0] == 'x') {
				int k;
				sscanf(str+1, "%d", &k);
				eq[k-1] = 1^eq[k-1];
			}
		}
		//printf("\n");
		A.push_back(eq);
		B.push_back(val);
	}
	if (!gaussxor()) printf("impossible\n");
	else {
		for(int j = 0; j < m; j++) {
			printf("%c", x[j] ? 'T' : 'F');
		}
		printf("\n");
	}
	return 0;
}
\end{vncode}
\endgroup


\subsection{Fast Fourier Transform (FFT)}

FFT.

\codepath{Math/fft.cpp}

\begingroup\scriptsize
\begin{vncode}
#include <vector>
#include <cstdio>
#include <complex>
using namespace std;

/*
 * Fast complex
 */

struct base { // faster than complex<double>
	double x, y;
	base() : x(0), y(0) {}
	base(double a, double b=0) : x(a), y(b) {}
	base operator/=(double k) { x/=k; y/=k; return (*this); }
	base operator*(base a) const { return base(x*a.x - y*a.y, x*a.y + y*a.x); }
	base operator*=(base a) {
		double tx = x*a.x - y*a.y;
		double ty = x*a.y + y*a.x;
		x = tx; y = ty;
		return (*this);
	}
	base operator=(double a) { x=a; y=0; return (*this); }
	base operator+(base a) const { return base(x+a.x, y+a.y); }
	base operator+=(base a) { x+=a.x; y+=a.y; return (*this); }
	base operator-(base a) const { return base(x-a.x, y-a.y); }
	base operator-=(base a) { x-=a.x; y-=a.y; return (*this); }
	double& real() { return x; }
	double& imag() { return y; }
};

/*
 * FFT - Fast Fourier Transform O(nlogn)
 */

//typedef complex<double> base;
 
void fft(vector<base> &a, bool invert) {
	int n = (int)a.size();
	for(int i = 1, j = 0; i < n; i++) {
		int bit = n >> 1;
		for(; j >= bit; bit >>= 1) j -= bit;
		j += bit;
		if (i < j) swap(a[i], a[j]);
	}
	for(int len = 2; len <= n; len <<= 1) {
		double ang = 2*acos(-1.0)/len * (invert ? -1 : 1);
		base wlen(cos(ang), sin(ang));
		for(int i = 0; i < n; i += len) {
			base w(1);
			for(int j = 0; j < len/2; j++) {
				base u = a[i+j], v = a[i+j+len/2] * w;
				a[i + j] = u + v;
				a[i + j + len/2] = u - v;
				w *= wlen;
			}
		}
	}
	for (int i = 0; invert && i < n; i++) a[i] /= n;
}

void convolution(vector<base> a, vector<base> b, vector<base> &res) {
	int n = 1;
	while(n < max(a.size(), b.size())) n <<= 1;
	n <<= 1;
	a.resize(n), b.resize(n);
	fft(a, false); fft(b, false);
	res.resize(n);
	for(int i=0; i<n; ++i) res[i] = a[i]*b[i];
	fft(res, true);
}

void circularConv(vector<base> &a, vector<base> &b, vector<base> &res) {
	//assert(a.size() == b.size());
	int n = a.size();
	convolution(a, b, res);
	for(int i = n; i < (int)res.size(); i++)
		res[i%n] += res[i];
	res.resize(n);
}

/*
 * SQRT DECOMPOSITION FOR FFT
 */

#include <cmath>
#define MOD 1000003LL
#define SMOD 1024LL // ~ sqrt(MOD)
typedef long long ll;

void sqrtConv(vector<ll> a, vector<ll> b, vector<ll> & c) {
	vector<base> ca[2], cb[2], cc[2][2];
	ca[0].resize(a.size());
	ca[1].resize(a.size());
	for(int i=0; i<(int)a.size(); i++) {
		ca[0][i] = base(a[i] % SMOD, 0);
		ca[1][i] = base(a[i] / SMOD, 0);
	}
	cb[0].resize(b.size());
	cb[1].resize(b.size());
	for(int i=0; i<(int)b.size(); i++) {
		cb[0][i] = base(b[i] % SMOD, 0);
		cb[1][i] = base(b[i] / SMOD, 0);
	}
	for(int l=0; l<2; l++) for(int r=0; r<2; r++)
		convolution(ca[l], cb[r], cc[l][r]);
	c.resize(cc[0][0].size());
	for(int i=0; i<(int)c.size(); i++) {
		c[i] =
		(((ll)round(cc[1][1][i].real()))%MOD*(SMOD*SMOD)%MOD)%MOD +
		((ll)round(cc[0][1][i].real()))%MOD*SMOD%MOD +
		((ll)round(cc[1][0][i].real()))%MOD*SMOD%MOD + 
		((ll)round(cc[0][0][i].real()))%MOD;
		c[i] %= MOD;
	}
}

/*
 * TEST MATRIX
 */

int main()
{
	base a (1);
	base b (2);
	base c (3);
	base d (4);
	base e (5);
	base f (6);
	vector<base> A, B, C;
	A.push_back(c);
	A.push_back(b);
	A.push_back(a);
	B.push_back(f);
	B.push_back(e);
	B.push_back(d);
	B.push_back(c);
	convolution(A, B, C);
	for(int i=0; i<(int)C.size(); i++) {
		printf(" %.2f", C[i].real());
	}
	return 0;
}
\end{vncode}
\endgroup


\subsection{Number Theoretic Transform (NTT)}

NTT.

\codepath{Math/ntt.cpp}

\begingroup\scriptsize
\begin{vncode}
#include <vector>
#include <cstdio>
#include <cstring>
#define MAXN 100009
using namespace std;

/*
 * Modular inverse
 */

template <typename T>
T extGcd(T a, T b, T& x, T& y) {
    if (b == 0) {
        x = 1; y = 0;
        return a;
    }
    else {
        T g = extGcd(b, a % b, y, x);
        y -= a / b * x;
        return g;
    }
}
 
template <typename T>
T modInv(T a, T m) {
    T x, y;
    extGcd(a, m, x, y);
    return (x % m + m) % m;
}

/*
 * NTT - Number Theoretic Transform O(nlogn)
 */

typedef long long ll;

const ll mod[3] = {1004535809LL, 1092616193LL, 998244353LL};
const ll root[3] = {12289LL, 23747LL, 15311432LL};
const ll root_1[3] = {313564925LL, 642907570LL, 469870224LL};
const ll root_pw[3] = {1LL<<21, 1LL<<21, 1LL<<23};

void ntt(vector<ll> & a, bool invert, int m) {
	ll n = (ll)a.size();
	for(ll i = 1, j = 0; i < n; i++) {
		ll bit = n >> 1;
		for (; j >= bit; bit >>= 1) j -= bit;
		j += bit;
		if (i < j) swap(a[i], a[j]);
	}
	for(ll len = 2, wlen; len <= n; len <<= 1) {
		wlen = invert ? root_1[m] : root[m];
		for (ll i = len; i < root_pw[m]; i <<= 1)
			wlen = (wlen * wlen % mod[m]);
		for(ll i = 0; i < n; i += len) {
			for(ll j = 0, w = 1; j < len/2; j++) {
				ll u = a[i+j],  v = a[i+j+len/2] * w % mod[m];
				a[i+j] = (u+v < mod[m] ? u+v : u+v-mod[m]);
				a[i+j+len/2] = (u-v >= 0 ? u-v : u-v+mod[m]);
				w = w * wlen % mod[m];
			}
		}
	}
	if (invert) {
		ll nrev = modInv(n, mod[m]);
		for (ll i=0; i<n; ++i)
			a[i] = a[i] * nrev % mod[m];
	}
}

void convolution(vector<ll> a, vector<ll> b, vector<ll> & res, int m) {
	ll n = 1;
	while(n < max (a.size(), b.size())) n <<= 1;
	n <<= 1;
	a.resize(n), b.resize(n);
	ntt(a, false, m); ntt(b, false, m);
	res.resize(n);
	for(int i=0; i<n; ++i) res[i] = (a[i]*b[i])%mod[m];
	ntt(res, true, m);
}

void circularConv(vector<ll> &a, vector<ll> &b, vector<ll> &res, int m) {
	//assert(a.size() == b.size());
	int n = a.size();
	convolution(a, b, res, m);
	for(int i = n; i < (int)res.size(); i++)
		res[i%n] = (res[i%n]+res[i])%mod[m];
	res.resize(n);
}

/*
 * CRT NTT
 */

template<typename T>
T modMul(T a, T b, T m) {
	T x = 0, y = a % m;
	while (b > 0) {
		if (b % 2 == 1) x = (x + y) % m;
		y = (y * 2) % m;
		b /= 2;
	}
	return x % m;
}

//convolution mod 1,097,572,091,361,755,137
void modConv(vector<ll> a, vector<ll> b, vector<ll> & res) {
	vector<ll> r0, r1;
	convolution(a, b, r0, 0);
	convolution(a, b, r1, 1);
	ll x, y, s, r, p = mod[0]*mod[1];
	extGcd(mod[0], mod[1], r, s);
	res.resize(r0.size());
	for(int i=0; i<(int)res.size(); i++) {
		res[i] = (modMul((s*mod[1]+p)%p, r0[i], p)
		+ modMul((r*mod[0]+p)%p, r1[i], p) + p) % p;
	}
}

/*
 * SPOJ MUL
 */

char num1[MAXN], num2[MAXN];

int main()
{
	int n, na, nb, k;
	vector<ll> A, B, C;
	scanf("%d", &n);
	for(int i=0; i<n; i++) {
		scanf(" %s %s", num1, num2);
		na = strlen(num1);
		nb = strlen(num2);
		A.clear();
		B.clear();
		for(int j=na-1; j>=0; j--) {
			A.push_back(num1[j]-'0');
		}
		for(int j=nb-1; j>=0; j--) {
			B.push_back(num2[j]-'0');
		}
		modConv(A, B, C);
		while(C.size()>1 && C.back()==0) C.pop_back();
		na = nb = k = 0;
		for(int j=0; j<(int)C.size() || nb > 0; j++) {
			na = (nb + (j<(int)C.size() ? C[j] : 0))%10;
			nb = (nb + (j<(int)C.size() ? C[j] : 0))/10;
			num1[k++] = ((char)na) + '0';
		}
		for(int i=0; i<k/2; i++) {
			na = num1[i];
			num1[i] = num1[k-i-1];
			num1[k-i-1] = na;
		}
		num1[k] = '\0';
		printf("%s\n", num1);
	}
	return 0;
}
\end{vncode}
\endgroup


\subsection{Tích chập vòng}

Circular convolution.

\codepath{Math/circularconvolution.cpp}

\begingroup\scriptsize
\begin{vncode}
template <typename T>
void circularconvolution(vector<T> a, vector<T> b, vector<T> & res) {
	int n = a.size();
	b.insert(b.end(), b.begin(), b.end());
	convolution(a, b, res);
	res = vector<T>(res.begin()+n, res.begin()+(2*n));
}
\end{vncode}
\endgroup


\subsection{Số phức}

Các phép với số phức trong C++/toán học.

\subsection{Chia đa thức}

Polynomial division.

\codepath{Math/polynomialdiv.cpp}

\begingroup\scriptsize
\begin{vncode}
#include <vector>
#include <cstdio>
#include <complex>
using namespace std;

/*
 * Fast complex
 */

struct base { // faster than complex<double>
	double x, y;
	base() : x(0), y(0) {}
	base(double a, double b=0) : x(a), y(b) {}
	base operator/=(double k) { x/=k; y/=k; return (*this); }
	base operator*(base a) const { return base(x*a.x - y*a.y, x*a.y + y*a.x); }
	base operator*=(base a) {
		double tx = x*a.x - y*a.y;
		double ty = x*a.y + y*a.x;
		x = tx; y = ty;
		return (*this);
	}
	base operator=(double a) { x=a; y=0; return (*this); }
	base operator+(base a) const { return base(x+a.x, y+a.y); }
	base operator+=(base a) { x+=a.x; y+=a.y; return (*this); }
	base operator-(base a) const { return base(x-a.x, y-a.y); }
	base operator-=(base a) { x-=a.x; y-=a.y; return (*this); }
	double& real() { return x; }
	double& imag() { return y; }
};

/*
 * FFT - Fast Fourier Transform O(nlogn)
 */

//typedef complex<double> base;
 
void fft(vector<base> &a, bool invert) {
	int n = (int)a.size();
	for(int i = 1, j = 0; i < n; i++) {
		int bit = n >> 1;
		for(; j >= bit; bit >>= 1) j -= bit;
		j += bit;
		if (i < j) swap(a[i], a[j]);
	}
	for(int len = 2; len <= n; len <<= 1) {
		double ang = 2*acos(-1.0)/len * (invert ? -1 : 1);
		base wlen(cos(ang), sin(ang));
		for(int i = 0; i < n; i += len) {
			base w(1);
			for(int j = 0; j < len/2; j++) {
				base u = a[i+j], v = a[i+j+len/2] * w;
				a[i + j] = u + v;
				a[i + j + len/2] = u - v;
				w *= wlen;
			}
		}
	}
	for (int i = 0; invert && i < n; i++) a[i] /= n;
}

void convolution(vector<base> a, vector<base> b, vector<base> &res) {
	int n = 1;
	while(n < max(a.size(), b.size())) n <<= 1;
	n <<= 1;
	a.resize(n), b.resize(n);
	fft(a, false); fft(b, false);
	res.resize(n);
	for(int i=0; i<n; ++i) res[i] = a[i]*b[i];
	fft(res, true);
}

void circularConv(vector<base> &a, vector<base> &b, vector<base> &res) {
	//assert(a.size() == b.size());
	int n = a.size();
	convolution(a, b, res);
	for(int i = n; i < (int)res.size(); i++)
		res[i%n] += res[i];
	res.resize(n);
}

/*
 * Polynomial Division
 */

#include <algorithm>

void inverse(vector<base> &a, int t) {
	if (t == 1) {
		a.resize(1); a[0].real() = 1.0/a[0].real();
		return;
	}
	vector<base> a0 = a;
	inverse(a0, (t/2) + (t%2));
	convolution(a0, a, a);
	a0.resize(t), a.resize(t);
	a[0] -= base(1.0);
	convolution(a0, a, a);
	a0.resize(t), a.resize(t);
	for(int i = 0; i < t; i++) a[i] = a0[i] - a[i];
}

void divide(vector<base> &a, vector<base> &b, vector<base> &d) {
	int n = a.size(), m = b.size();
	if (n < m) {
		d.clear(); d.push_back(0);
		return;
	}
	vector<base> ar = a, br = b;
	reverse(ar.begin(), ar.end());
	reverse(br.begin(), br.end());
	inverse(br, n - m + 1);
	convolution(ar, br, d);
	d.resize(n - m + 1);
	reverse(d.begin(), d.end());
}

void remainder(vector<base> &a, vector<base> &b, vector<base> &r) {
	int n = a.size(), m = b.size();
	if (n < m) { r = a; return; }
	vector<base> d, aux;
	divide(a, b, d);
	r = a;
	convolution(d, b, aux);
	r.resize(m-1); aux.resize(m-1);
	for(int i = 0; i < m-1; i++) r[i] -= aux[i];
}

/*
 * Multi-point evaluation in O((n+m) logn logm)
 */

#define MAXN 10000

vector<base> rp[4*MAXN];
void roots(int u, double x[], int i, int j) {
	if (i == j) {
		rp[u].resize(2); rp[u][0] = -x[i];
		rp[u][1] = 1.0; return;
	}
	int m = (i + j) / 2;
	roots(2*u, x, i, m); roots(2*u+1, x, m+1, j);
	convolution(rp[2*u], rp[2*u+1], rp[u]);
	rp[u].resize(j-i+2);
}

void roots(vector<base> &a, double x[], int n) {
	roots(1, x, 0, n-1); a = rp[1];
}

vector<base> ep[4*MAXN];
void multievaluate(int u, double x[], double y[], int i, int j) {
	if (i == j) {
		y[i] = 0;
		double p = 1;
		for(int k = 0; k < (int)ep[u].size(); k++)
			y[i] += ep[u][k].real()*p, p *= x[i];
		return;
	}
	remainder(ep[u], rp[2*u], ep[2*u]);
	remainder(ep[u], rp[2*u+1], ep[2*u+1]);
	int m = (i + j) / 2;
	multievaluate(2*u, x, y, i, m);
	multievaluate(2*u+1, x, y, m+1, j);
}

void multievaluate(vector<base> &a, double x[], double y[], int n) {
	roots(1, x, 0, n-1); ep[1] = a;
	multievaluate(1, x, y, 0, n-1);
}

/*
 * Polynomial interpolation O(n log^3 n)
 * horrible constant and time
 */

vector<double> y[MAXN];
double ya0[MAXN], yp[MAXN];
void interpolate(vector<base> &a, int u, double x[], int n) {
	if (n == 1) {
		a.resize(1); a[0] = y[u][0];
		return;
	}
	y[2*u] = y[u]; y[2*u].resize(n/2);
	vector<base> a0, a1, p;
	interpolate(a0, 2*u, x, n/2);
	roots(p, x, n/2);
	int m = n-(n/2);
	multievaluate(a0, x+(n/2), ya0, m);
	multievaluate(p,  x+(n/2), yp, m);
	y[2*u+1].resize(m);
	for(int i = 0; i < m; i++)
		y[2*u+1][i] = (y[u][n/2 + i] - ya0[i]) / yp[i];
	interpolate(a1, 2*u+1, x+(n/2), m);
	convolution(p, a1, a);
	int sz = max(a.size(), a0.size());
	a.resize(sz);
	for(int i = 0; i < sz && i < n/2; i++)
		a[i] += a0[i];
	a.resize(n);
}

void interpolate(vector<base> &a, double x[], double ry[], int n) {
	y[1].resize(n);
	for(int i = 0; i < n; i++) y[1][i] = ry[i];
	interpolate(a, 1, x, n);
}

/*
 * TEST MATRIX
 */

int main() {
	int n, m;
	vector<base> P, Q, D, R;
	
	scanf("%d", &n);
	P.resize(n);
	for(int i = 0; i < n; i++) scanf("%lf", &P[i].real());
	
	scanf("%d", &m);
	Q.resize(m);
	for(int i = 0; i < m; i++) scanf("%lf", &Q[i].real());

	vector<base> invP = P, invQ = Q;
	inverse(invP, 10);
	printf("first 10 terms of inverse of P:");
	for(int i = 0; i < 10; i++) printf(" %.0f", invP[i].real());
	printf("\n");
	
	inverse(invQ, 10);
	printf("first 10 terms of inverse of Q:");
	for(int i = 0; i < 10; i++) printf(" %.0f", invQ[i].real());
	printf("\n");

	divide(P, Q, D);
	printf("P/Q:");
	for(int i = 0; i < (int)D.size(); i++) printf(" %.3f", D[i].real());
	printf("\n");
	
	remainder(P, Q, R);
	printf("P%%Q:");
	for(int i = 0; i < (int)R.size(); i++) printf(" %.3f", R[i].real());
	printf("\n");

	scanf("%d", &n);
	double x[100];
	for(int i = 0; i < n; i++) scanf("%lf", &x[i]);
	vector<base> a;
	roots(a, x, n);
	printf("from roots:");
	for(int i = 0; i < (int)a.size(); i++) printf(" %.3f", a[i].real());
	printf("\n");

	double y[100];
	multievaluate(P, x, y, n);
	printf("y[i]:");
	for(int i = 0; i < n; i++) printf(" %.3f", y[i]);
	printf("\n");

	scanf("%d", &n);
	for(int i = 0; i < n; i++) scanf("%lf %lf", &x[i], &y[i]);
	interpolate(a, x, y, n);
	printf("interpolated a:");
	for(int i = 0; i < n; i++) printf(" %.3f", a[i]);
	printf("\n");

	return 0;
}
\end{vncode}
\endgroup


\subsection{Đánh giá tại nhiều điểm}

Multipoint evaluation.

\subsection{Nội suy đa thức}

Polynomial interpolation.

\subsection{Fast Walsh-Hadamard Transform}

FWHT.

\codepath{Math/fwht.cpp}

\begingroup\scriptsize
\begin{vncode}
#include <vector>
using namespace std;

/*
 * FWHT - Fast Walsh–Hadamard Transform O(nlogn)
 */

//int mat[2][2] = {{1, 1}, {1, -1}}, inv[2][2] = {{1, 1}, {1, -1}}; //xor
int mat[2][2] = {{1, 1}, {1, 0}}, inv[2][2] = {{0, 1}, {1, -1}}; //or
//int mat[2][2] = {{0, 1}, {1, 1}}, inv[2][2] = {{-1, 1}, {1, 0}}; //and

void fwht(vector<double> & a, bool invert) {
	int n = (int)a.size();
	double u, v;
	for (int len = 1; 2 * len <= n; len <<= 1) {
		for (int i = 0; i < n; i += 2 * len) {
			for (int j = 0; j < len; j++) {
				u = a[i + j];
				v = a[i + len + j];
				if (!invert) {
					a[i + j] = mat[0][0]*u + mat[0][1]*v;
					a[i + len + j] = mat[1][0]*u + mat[1][1]*v;
				}
				else {
					a[i + j] = inv[0][0]*u + inv[0][1]*v;
					a[i + len + j] = inv[1][0]*u + inv[1][1]*v;
				}
			}
		}
	}
	//for (int i=0; invert && i<n; ++i) a[i] /= n; //xor
}

void convolution(vector<double> a, vector<double> b, vector<double> & res) {
	int n = 1;
	while(n < max(a.size(), b.size())) n <<= 1;
	a.resize(n), b.resize(n);
	fwht(a, false); fwht(b, false);
	res.resize(n);
	for(int i=0; i<n; ++i) res[i] = a[i]*b[i];
	fwht(res, true);
}

/*
 * TEST MATRIX
 */

void naive(vector<double> a, vector<double> b, vector<double> & res) {
	int n = 1;
	while(n < max(a.size(), b.size())) n <<= 1;
	a.resize(n), b.resize(n);
	res.clear(); res.assign(n, 0);
	for(int i=0; i<n; i++) {
		for(int j=0; j<n; j++) {
			res[i|j] += a[i]*b[j];
		}
	}
}

#define EPS 1e-3
#include <cstdio>
#include <cstdlib>
#include <cmath>
bool test(int nTest) {
	for(int t=1; t<=nTest; t++) {
		vector<double> a(t), b(t), res1(t), res2(t);
		for(int i=0; i<t; i++) {
			a[i] = rand()%10000;
			b[i] = rand()%10000;
		}
		convolution(a, b, res1);
		naive(a, b, res2);
		bool ok = true;
		for(int i=0; i<(int)res1.size(); i++) {
			if (fabs(res1[i]-res2[i]) > EPS) ok = false;
		}
		if (!ok) {
			printf("failed on test %d\n", t);
			printf("a:");
			for(int i=0; i<(int)a.size(); i++) {
				printf(" %.2f", a[i]);
			}
			printf("\n");
			printf("b:");
			for(int i=0; i<(int)b.size(); i++) {
				printf(" %.2f", b[i]);
			}
			printf("\n");
			printf("res1:");
			for(int i=0; i<(int)res1.size(); i++) {
				printf(" %.2f", res1[i]);
			}
			printf("\n");
			printf("res2:");
			for(int i=0; i<(int)res2.size(); i++) {
				printf(" %.2f", res2[i]);
			}
			printf("\n");
			return false;
		}
	}
	printf("all tests passed\n");
	return true;
}

int main() {
	test(1000);
}
\end{vncode}
\endgroup


\subsection{Tích chập với CRT}

Kết hợp nhiều modulo bằng CRT.

\codepath{Math/crt.cpp}

\begingroup\scriptsize
\begin{vncode}
#include <cstdio>
#define MOD 1000000007LL
typedef long long ll;

template <typename T>
T extGcd(T a, T b, T& x, T& y) {
    if (b == 0) {
        x = 1; y = 0;
        return a;
    }
    else {
        T g = extGcd(b, a % b, y, x);
        y -= a / b * x;
        return g;
    }
}
 
template <typename T>
T modInv(T a, T m) {
    T x, y;
    extGcd(a, m, x, y);
    return (x % m + m) % m;
}
 
template <typename T>
T modDiv(T a, T b, T m) {
    return ((a % m) * modInv(b, m)) % m;
}

template<typename T>
T crt(T* a, T* p, int n, T m) {
    T P = 1;
    for(int i=0; i<n; i++) P = (P * p[i]) % m;
    T x = 0, pp;
    for(int i=0; i<n; i++) {
        pp = modDiv(P, p[i], m);
        x = (x + (((a[i] * pp) % m) * modInv(pp, p[i]))) % m;
    }
    return x;
}

/*TEST MATRIX*/
#include <cstdlib>
#define PRINTTEST 0

int primes[] = {2,3,5,7,11,13,17,19,23,29,31,37,41,43,47,53,59};

bool test(int ntests) {
	for(int t=0; t<ntests; t++) {
		int T = 17;
		ll y[20], a[20];
		for(int i=1; i<=T; i++) {
			if (PRINTTEST) printf("test %d:\n", t*T + i);
			for(int j=0; j<i; j++) {
				a[j] = primes[j];
				y[j] = rand()%a[j];
				if (PRINTTEST) printf("x = %lld mod %lld\n", y[j], a[j]);
			}
			ll ans = crt(y, a, i, MOD);
			if (PRINTTEST) printf("answer is x = %lld\n", ans);
			for(int j=0; j<i; j++) {
				ll mans = ans;
				int count = 10000;
				while(mans%a[j] != y[j] && count > 0) {
					mans += MOD;
					count--;
				}
				if (mans%a[j] != y[j]) {
					printf("failed on test %d, ans = %I64d\n", t*T + i, ans);
					return false;
				}
			}
		}
	}
	printf("All tests passed\n");
	return true;
}

int main()
{
	test(10000);
	return 0;
}
\end{vncode}
\endgroup


\subsection{Tích chập với phân rã căn}

SQRT decomposition convolution.

\subsection{Tích phân bằng quy tắc Simpson}

Simpson integration.

\codepath{Math/simpson.cpp}

\begingroup\scriptsize
\begin{vncode}
#include <cmath>

/*
 * Simpson Rule for integration
 */

double f(double x) { return cos(x); }

double simpson(double a, double b, int n = 1e6) {
	double h = (b - a) / n, s = f(a) + f(b);
	for (int i = 1; i < n; i += 2) s += 4*f(a+h*i);
	for (int i = 2; i < n; i += 2) s += 2*f(a+h*i);
	return s*h/3;
}

/*
 * TEST MATRIX
 */

#include <cstdio>
#include <cmath>

bool test(int n) {
	for(int i = 0; i < n; i++) {
		double area = simpson(0, i);
		double exact = sin(i*1.0);
		if (fabs(area - exact) > 1e+18) {
			printf("error at %d = %.20f\n", i, fabs(area - exact));
			return false;
		}
	}
	printf("all tests passed\n");
	return true;
}

int main() {
	test(100);
	return 0;
}
\end{vncode}
\endgroup


\subsection{Mã Gray}

Gray code.

\codepath{Math/gray.cpp}

\begingroup\scriptsize
\begin{vncode}
#include <cstdio>

int g (int n) {
	return n ^ (n >> 1);
}

int rev_g (int g) {
	int n = 0;
	for (; g; g>>=1)
		n ^= g;
	return n;
}
\end{vncode}
\endgroup


\subsection{BigInteger trong Java}

Mẫu dùng \texttt{java.math.BigInteger}.

\subsection{Bignum trong C++}

Số lớn C++.

\codepath{Math/bignum.cpp}

\begingroup\scriptsize
\begin{vncode}
#include <vector>
#include <algorithm>
#include <cstdio>
#include <cstdlib>
#include <cstring>
using namespace std;

typedef vector<int> bignum;
const int base = 1000*1000*1000;

//Impressão
void print(bignum & a) {
	printf("%d", a.empty() ? 0 : a.back());
	for (int i=(int)a.size()-2; i>=0; --i) {
		printf("%09d", a[i]);
	}
}

//Remover Leading zeros
void fix(bignum & a) {
	while (a.size() > 1u && a.back() == 0) {
		a.pop_back();
	}
}

//Comparador
bool comp(bignum a, bignum b) {
	fix(a); fix(b);
	if (a.size() != b.size()) return a.size() < b.size();
	for(int i=(int)a.size()-1; i>=0; i--) {
		if (a[i] != b[i]) return a[i] < b[i];
	}
	return false;
}

//Leitura
void str2bignum(char* s, bignum & a) {
	a.clear();
	for (int i=(int)strlen(s); i>0; i-=9) {
		s[i] = 0;
		a.push_back (atoi (i>=9 ? s+i-9 : s));
	}
	fix(a);
}

//Construtor a partir de um inteiro
void int2bignum(int n, bignum & a) {
	a.clear();
	if (n == 0) a.push_back(0);
	while(n > 0) {
		a.push_back(n%base);
		n /= base;
	}
}

int bignum2int(bignum & a) {
	int ans = 0, p=1;
	for(int i=0; i<(int)a.size(); i++) {
		ans += a[i]*p;
		p *= base;
	}
	return ans;
}

//Soma a+b=c
void sum(bignum & a, bignum & b, bignum & c) {
	int carry = 0, n = max(a.size(), b.size());
	c.resize(n);
	for (int i=0, ai, bi; i<n; i++) {
		ai = i < (int)a.size() ? a[i] : 0;
		bi = i < (int)b.size() ? b[i] : 0;
		c[i] = carry + ai + bi;
		carry = c[i] / base;
		c[i] %= base;
	}
	if (carry > 0) c.push_back(carry);
	fix(c);
}

//Subtrai a-b=c>=0
void subtract(bignum & a, bignum & b, bignum & c) {
	int carry = 0, n = max(a.size(), b.size());
	c.resize(n);
	for (int i=0, ai, bi; i<n; i++) {
		ai = i < (int)a.size() ? a[i] : 0;
		bi = i < (int)b.size() ? b[i] : 0;
		c[i] = carry + ai - bi;
		carry = c[i] < 0 ? 1 : 0;
		if (c[i] < 0) c[i] += base;
	}
	fix(c);
}

//Shift left por b (multiplica por base^b)
void shiftL(bignum & a, int b, bignum & c) {
	c.resize((int)a.size() + b);
	for(int i=(int)c.size()-1; i>=0; i--) {
		if(i>=b) c[i] = a[i-b];
		else c[i] = 0;
	}
	fix(c);
}

//Shift right por b (multiplica por base^b)
void shiftR(bignum & a, int b, bignum & c) {
	if (((int)a.size()) <= b) {
		c.clear(); c.push_back(0);
		return;
	}
	c.resize((int)a.size() - b);
	for(int i=0; i<(int)c.size(); i++) {
		c[i] = a[i+b];
	}
	fix(c);
}

//Multiplica bignum b por int a<base
void multiply(int a, bignum & b, bignum & c) {
	int carry = 0, bi;
	c.resize(b.size());
	for (int i=0, bi; i<(int)b.size() || carry; i++) {
		if (i == (int)b.size()) c.push_back(0);
		bi = i < (int)b.size() ? b[i] : 0;
		long long cur = carry + a * 1ll * bi;
		c[i] = int(cur % base);
		carry = int(cur / base);
	}
	fix(c);
}

//Multiplica a*b=c
void multiply(bignum a, bignum b, bignum & c) {
	int n = a.size()+b.size();
	long long carry = 0, acum;
	c.resize(n);
	for (int k=0; k<n || carry; k++) {
		if (k == n) c.push_back(0);
		acum = carry; carry = 0;
		for (int i=0, j=k; i <= k && i<(int)b.size(); i++, j--) {
			if (j >= (int)b.size()) continue;
			acum += a[i] * 1ll * b[j];
			carry += acum / base;
			acum %= base;
		}
		c[k] = acum;
	}
	fix(c);
}

//Divide bignum a por int b<base
void divide(bignum & a, int b, bignum & c) {
	int carry = 0;
	c.resize(a.size());
	for (int i=(int)a.size()-1; i>=0; --i) {
		long long cur = a[i] + carry * 1ll * base;
		c[i] = int (cur / b);
		carry = int (cur % b);
	}
	fix(a);
}

//Divide bignum a por bignum b, quociente q, resto r
void divide(bignum a, bignum b, bignum & q, bignum & r) {
	bignum z, p;
	int2bignum(0, z);
	int2bignum(1, p);
	int2bignum(0, q);
	while(comp(b, a)) {
		shiftL(p, max(1u, a.size()-b.size()), p);
		shiftL(b, max(1u, a.size()-b.size()), b);
	}
	while(true) {
		while (comp(a, b) && comp(z, p)) {
			shiftR(p, 1, p); shiftR(b, 1, b);
		}
		if (!comp(z, p)) break;
		subtract(a, b, a);
		sum(q, p, q);
	}
	swap(a, r);
}

bool test() {
	bool toprint = false;
	for(int i=1; i<10000; i++) {
		int na = 1 + rand()%30000;
		int nb = 1 + rand()%30000;
		if (nb > na) {
			int tmp = na;
			na = nb;
			nb = tmp;
		}
		if(toprint) printf("test %d: na = %d, nb = %d, ", i, na, nb);
		bignum a; int2bignum(na, a);
		bignum b; int2bignum(nb, b);
		bignum c, r; int nc, nr;
		sum(a, b, c);
		nc = na + nb;
		if (nc != bignum2int(c)) {
			printf("failed test %d, n: %d+%d=%d, b: ", i, na, nb, nc);
			print(a);
			printf("+");
			print(b);
			printf("=");
			print(c);
			printf("\n");
			return false;
		}
		if(toprint) printf("+ ok ");
		subtract(a, b, c);
		nc = na - nb;
		if (nc != bignum2int(c)) {
			printf("failed test %d, n: %d-%d=%d, b: ", i, na, nb, nc);
			print(a);
			printf("-");
			print(b);
			printf("=");
			print(c);
			printf("\n");
			return false;
		}
		if(toprint) printf("- ok ");
		multiply(a, b, c);
		nc = na * nb;
		if (nc != bignum2int(c)) {
			printf("failed test %d, n: %d*%d=%d, b: ", i, na, nb, nc);
			print(a);
			printf("*");
			print(b);
			printf("=");
			print(c);
			printf("\n");
			return false;
		}
		if(toprint) printf("* ok ");
		divide(a, b, c, r);
		nc = na / nb;
		if (nc != bignum2int(c)) {
			printf("failed test %d, n: %d/%d=%d, b: ", i, na, nb, nc);
			print(a);
			printf("/");
			print(b);
			printf("=");
			print(c);
			printf("\n");
			return false;
		}
		if(toprint) printf("/ ok\n");
		swap(c, r);
		nc = na % nb;
		if (nc != bignum2int(c)) {
			printf("failed test %d, n: %d%%%d=%d, b: ", i, na, nb, nc);
			print(a);
			printf("%%");
			print(b);
			printf("=");
			print(c);
			printf("\n");
			return false;
		}
		if(toprint) printf("% ok\n");
	}
	return true;
}

char s1[109], s2[109], op;

int main() {
	if(test()) printf("all tests passed\n");
	while (scanf(" %s %c %s", s1, &op, s2) != EOF) {
		bignum a; str2bignum(s1, a);
		bignum b; str2bignum(s2, b);
		bignum c, d;
		if (op == '+') sum(a, b, c);
		if (op == '-') subtract(a, b, c);
		if (op == '*') multiply(a, b, c);
		if (op == '/') divide(a, b, c, d);
		if (op == '%') divide(a, b, d, c);
		print(c);
		printf("\n");
	}
	return 0;
}
\end{vncode}
\endgroup


\subsection{Sự phá sản của con bạc}

Gambler's ruin.

\subsection{Định lý và công thức}

Tập hợp các công thức xác suất, tổ hợp và số học thường dùng.

\clearpage

\section{Xử lý chuỗi}

\subsection{Hàm tiền tố: thuật toán Knuth-Morris-Pratt (KMP)}

Prefix function/KMP.

\codepath{Strings/kmp.cpp}

\begingroup\scriptsize
\begin{vncode}
#include <string>
#include <vector>
#include <cstring>
using namespace std;

/*
 * Knuth-Morris-Pratt's Algorithm
 */

class KMP {
	string P;
	vector<int> b;
public:
	KMP(const char _P[]) : P(_P) {
		int m = P.size();
		b.assign(m + 1, -1);
		for(int i = 0, j = -1; i < m;) {
			while (j >= 0 && P[i] != P[j]) j = b[j];
			b[++i] = ++j;
		}
	}
	vector<int> match(const char T[]) {
		vector<int> ans;
		for (int i = 0, j = 0, n = strlen(T); i < n;) {
			while (j >= 0 && T[i] != P[j]) j = b[j];
			i++; j++;
			if (j == (int)P.size()) {
				ans.push_back(i - j);
				j = b[j];
			}
		}
		return ans;
	}
	int repetend() {
		int n = P.size();
		int ans = n - b[n];
		if (n % ans) ans = n;
		return ans;
	}
};

/*
 * TEST MATRIX
 */

#include <cstdio>

void test() {
	string P = "ababa";
	string T = "abababababababa";
	KMP kmp(P.c_str());
	vector<int> index = kmp.match(T.c_str());
	for(int i=0; i<(int)index.size(); i++) {
		printf("%d ", index[i]);
	}
	printf("\n");
}

int main() {
	test();
	return 0;
}
\end{vncode}
\endgroup


\subsection{Rabin-Karp}

Hash chuỗi.

\codepath{Strings/rabinkarp.cpp}

\begingroup\scriptsize
\begin{vncode}
#include <vector>
#include <cstring>
using namespace std;
typedef long long ll;

/*
 * Rabin Karp
 */

class RabinKarp {
	ll m;
	vector<int> pw, hsh;
public:
	RabinKarp() {}
	RabinKarp(char str[], ll p, ll _m) : m(_m) {
		int n = strlen(str);
		pw.resize(n); hsh.resize(n);
		hsh[0] = str[0]; pw[0] = 1;
		for(int i = 1; i < n; i++) {
			hsh[i] = (hsh[i-1] * p + str[i]) % m;
			pw[i] = (pw[i-1] * p) % m;	
		}
	}
	ll hash(int i, int j) {
		ll ans = hsh[j];
		if (i > 0) ans = (ans - ((hsh[i-1]*1ll*pw[j-i+1])%m) + m) % m;
		return ans;
	}
};

/*
 * Codeforces 961F
 */

#include <cstdio>
#define MAXN 1000009
#define P 256
#define M 1000000007

char str[MAXN];
int N, f[MAXN];
RabinKarp rk;

bool check(int l, int r) {
	if (f[l] <= 0) return true;
	if (l+f[l]-1 >= r) return false;
	return rk.hash(l, l+f[l]-1) == rk.hash(r-f[l]+1, r);
}

int main() {
	scanf("%d %s", &N, str);
	rk = RabinKarp(str, P, M);
	int r = N/2;
	int l = (N % 2 == 0 ? r-1 : r);
	int k = l;
	f[l+1] = -1;
	for(; l >= 0; l--, r++) {
		f[l] = f[l+1]+2;
		while(!check(l, r)) f[l]-=2;
	}
	for(int i = 0; i <= k; i++) printf("%d ", f[i]);
	printf("\n");
	return 0;
}
\end{vncode}
\endgroup


\subsection{Repetend: chu kỳ nhỏ nhất của chuỗi}

Trả về k nhỏ nhất sao cho k|n và chuỗi gồm n/k bản sao giống nhau.

\codepath{Strings/kmp.cpp}

\begingroup\scriptsize
\begin{vncode}
#include <string>
#include <vector>
#include <cstring>
using namespace std;

/*
 * Knuth-Morris-Pratt's Algorithm
 */

class KMP {
	string P;
	vector<int> b;
public:
	KMP(const char _P[]) : P(_P) {
		int m = P.size();
		b.assign(m + 1, -1);
		for(int i = 0, j = -1; i < m;) {
			while (j >= 0 && P[i] != P[j]) j = b[j];
			b[++i] = ++j;
		}
	}
	vector<int> match(const char T[]) {
		vector<int> ans;
		for (int i = 0, j = 0, n = strlen(T); i < n;) {
			while (j >= 0 && T[i] != P[j]) j = b[j];
			i++; j++;
			if (j == (int)P.size()) {
				ans.push_back(i - j);
				j = b[j];
			}
		}
		return ans;
	}
	int repetend() {
		int n = P.size();
		int ans = n - b[n];
		if (n % ans) ans = n;
		return ans;
	}
};

/*
 * TEST MATRIX
 */

#include <cstdio>

void test() {
	string P = "ababa";
	string T = "abababababababa";
	KMP kmp(P.c_str());
	vector<int> index = kmp.match(T.c_str());
	for(int i=0; i<(int)index.size(); i++) {
		printf("%d ", index[i]);
	}
	printf("\n");
}

int main() {
	test();
	return 0;
}
\end{vncode}
\endgroup


\subsection{Suffix Array: thuật toán Karp-Miller-Rosenberg (KMR)}

Suffix array O(n log n), LCP O(n), lower bound và đếm match.

\codepath{Strings/suffixarray.cpp}

\begingroup\scriptsize
\begin{vncode}
#define _CRT_SECURE_NO_WARNINGS
#define MAXN 100009
#include <algorithm>
#include <cstdio>
#include <cstring>
#include <string>
#include <set>
using namespace std;

/*
 * Suffix Array
 * Karp-Miller-Rosenberg's Algorithm (KMR) - O(n log n)
 * implementation adapted from Andrew Stankevich's class in Brazilian ICPC Summer School 2019
 */

int tsa[MAXN], c[MAXN], tc[MAXN], head[MAXN];

void bucket(int sa[], int n, int L) {
	memset(&head, 0, sizeof head);
	for (int i = 0; i < n; i++) head[c[i]]++;
	int maxi = max(300, n), k = 0;
	for (int i = 0; i < maxi; i++)
		swap(k, head[i]), k += head[i];
	for (int i = 0; i < n; i++) {
		int j = (sa[i] - L + n) % n;
		tsa[head[c[j]]++] = j;
	}
	memcpy(sa, tsa, n * sizeof(int));
}

void suffixArray(const char str[], int sa[]) {
	int n = strlen(str) + 1;
	for (int i = 0; i < n; i++)
		sa[i] = i, c[i] = str[i];
	bucket(sa, n, 0);
	for (int L = 1; L < n; L <<= 1) {
		bucket(sa, n, L);
		for (int i = 0, cc = -1; i < n; i++)
			tc[sa[i]] = (i == 0 || c[sa[i]] != c[sa[i - 1]] || c[(sa[i] + L) % n] != c[(sa[i - 1] + L) % n]) ? ++cc : cc;
		memcpy(c, tc, n * sizeof(int));
	}
	for (int i = 0; i < n - 1; i++) sa[i] = sa[i + 1];
}

int inv[MAXN];

void computeLcp(const char str[], int sa[], int lcp[]) {
	int n = strlen(str), k = 0;
	for (int i = 0; i < n; i++) inv[sa[i]] = i;
	for (int j = 0; j < n; j++) {
		int i = inv[j];
		if (k < 0) k = 0;
		while (i > 0 && str[sa[i] + k] == str[sa[i - 1] + k]) k++;
		lcp[i] = k--;
	}
}

int bound(char s[], char t[], int sa[], int n, int m, bool upper) {
	int hi = n, lo = -1;
	int khi = 0, klo = 0;
	while (hi > lo + 1) {
		int mid = (hi + lo) / 2;
		int i = sa[mid], k = min(klo, khi);
		while (k < m && s[i + k] == t[k]) k++;
		if (k == m && upper) lo = mid, klo = k;
		else if (k == m && !upper) hi = mid, khi = k;
		else if (s[i + k] > t[k]) hi = mid, khi = k;
		else lo = mid, klo = k;
	}
	return hi;
}

int match(char s[], char t[], int sa[], int n, int m) {
	int l = bound(s, t, sa, n, m, false);
	int r = bound(s, t, sa, n, m, true);
	return r - l;
}

string concat;
int sa[MAXN], lcp[MAXN], pos[MAXN], len[MAXN];

void addStringsToSA(int N, const char s[][MAXN]) {
	int j = 0;
	concat = "";
	for (int i = 0; i < N; ++i) {
		int n = strlen(s[i]) + 1;
		concat += string(s[i]) + "$";
		for (int k = 0; k < n; ++k) {
			pos[j] = i;
			len[j++] = n - k - 1;
		}
	}
	suffixArray(concat.c_str(), sa);
	computeLcp(concat.c_str(), sa, lcp);
}

void slide(set<pair<int, int> >& vals, string& ans, int& l, int* qty, int& cnt, int i, int N) {
	while (cnt == N || (lcp[i] <= (int)ans.size() && l < i)) {
		int u = sa[l];
		vals.erase({ lcp[l], l });

		if (cnt == N) {
			int n = vals.empty() ? len[u] : min(len[u], vals.begin()->first);
			if (n > (int) ans.size()) ans = concat.substr(u, n);
		}

		qty[pos[u]]--;
		if (qty[pos[u]] == 0) cnt--;
		++l;
	}
}

string getCommonSubstring(int N, const char s[][MAXN]) {
	addStringsToSA(N, s);

	string ans = "";
	int qty[21] = {}; // Must replaced 21 with the maximum value of N
	int cnt = 0, l = N;

	set<pair<int, int> > vals;

	lcp[concat.size()] = 0;
	for (int i = N; i < (int)concat.size(); ++i) {
		int u = sa[i];
		slide(vals, ans, l, qty, cnt, i, N);
		vals.insert({ lcp[i], i });
		qty[pos[u]]++;
		if (qty[pos[u]] == 1) cnt++;
	}
	slide(vals, ans, l, qty, cnt, concat.size(), N);

	return ans;
}

/*
 * TEST MATRIX
 */
 /*
 void printSA(char str[]) {
	 int sa[100];
	 int n = strlen(str);
	 suffixArray(str, sa, n);
	 for(int i = 0; i < n+1; i++) {
		 printf("%2d %2d %s\n", i, sa[i], str+sa[i]);
	 }
 }

 int main() {
	 char str[100];
	 scanf("%s", str);
	 printSA(str);
	 return 0;
 }*/

 /*
  * SPOJ SARRAY
  */
  /*
  char str[MAXN];
  int sa[MAXN];

  int main() {
	  scanf(" %s", str);
	  int n = strlen(str);
	  suffixArray(str, sa);
	  //constructSA(str, sa);
	  for(int i = 0; i < n; i++) {
		  printf("%d\n", sa[i]);
	  }
	  return 0;
  }*/

  /*
   * SPOJ STRMATCH
   */
   /*
   char s[MAXN], t[MAXN];
   int sa[MAXN];

   int main() {
	   int n, q;
	   scanf("%d %d %s", &n, &q, s);
	   suffixArray(s, sa);
	   while (q --> 0) {
		   scanf(" %s", t);
		   printf("%d\n", match(s, t, sa, n, strlen(t)));
	   }
	   return 0;
   }*/

	/*
	Common substring manual test

	char p[21][MAXN];

	int main() {
		int N;

		scanf("%d", &N);

		for (int i = 0; i < N; ++i) {
			scanf("%s", p[i]);
		}
		printf("%s\n", getCommonSubstring(N, p).c_str());
	}
	*/
   /*
	* Codeforces 101711B
	*/

char str[MAXN];
int SA[MAXN], LCP[MAXN];

long long psum(int i) {
	if (i <= 0) return 0;
	return i * (i + 1LL) / 2;
}

int main() {
	scanf("%s", str);
	int n = strlen(str);
	suffixArray(str, SA);
	computeLcp(str, SA, LCP);
	long long ans = 0;
	for (int i = 0; i < n; i++) {
		ans += psum(n - SA[i]) - psum(LCP[i]);
	}
	printf("%lld\n", ans);
	return 0;
}
\end{vncode}
\endgroup


\subsection{Hàm Z và thuật toán Z}

Tính z[i] trong O(n).

\codepath{Strings/zfunction.cpp}

\begingroup\scriptsize
\begin{vncode}
#include <cstdio>
#include <cstring>
#include <algorithm>
using namespace std;

/*
 * Z-function
 */

void zfunction(char s[], int z[]) {
	int n = strlen(s);
	fill(z, z+n, 0);
	for (int i=1, l=0, r=0; i<n; ++i) {
		if (i <= r) z[i] = min(r-i+1, z[i-l]);
		while (i+z[i] < n && s[z[i]] == s[i+z[i]])
			++z[i];
		if (i+z[i]-1 > r)
			l = i,  r = i+z[i]-1;
	}
}

/*
 * TEST MATRIX
 */

void naivez(char* s, int* z) {
	int n = strlen(s);
	fill(z, z+n, 0);
	for(int i=1; i<n; i++) {
		while(i+z[i] < n && s[z[i]]==s[i+z[i]]) z[i]++;
	}
}
int main() {
	char str[309];
	int z1[309], z2[309];
	for(int n=1; n<10; n++) {
		for(int i=0; i<n; i++) {
			str[i] = 'a' + rand()%2;
		}
		str[n] = 0;
		naivez(str, z1);
		printf("test #%d: %s\n", n, str);
		zfunction(str, z2);
		printf("naive:");
		for(int i=0; i<n; i++) {
			printf(" %2d", z1[i]);
		}
		printf("\n");
		printf("zfunc:");
		for(int i=0; i<n; i++) {
			printf(" %2d", z2[i]);
		}
		printf("\n");
	}
}
\end{vncode}
\endgroup


\subsection{Thuật toán Manacher}

Tìm palindrome dài nhất.

\codepath{Strings/manacher.cpp}

\begingroup\scriptsize
\begin{vncode}
#include <cstdio>
#include <cstring>
#include <algorithm>
#define MAXN 3009
using namespace std;

/*
 * Manacher's Algorithm O(n)
 */

void manacher(char str[], int L[]) {
	int n = strlen(str), c = 0, r = 0;
	for(int i = 0; i < n; i++) {
		if(i < r && 2*c >= i) L[i] = min(L[2*c-i], r-i);
		else L[i] = 0;
		while(i-L[i]-1 >= 0 && i+L[i]+1 < n &&
			str[i-L[i]-1] == str[i+L[i]+1]) L[i]++;
		if(i+L[i]>r) { c=i; r=i+L[i]; }
	}
}

/*
 * Longest Palindromic Substring O(n)
 */

int LPS(char T[]) {
	int n = 2*strlen(T) + 1;
	char tmp[n+1];
	for(int i = 0, k = 0; T[i]; i++) {
		tmp[k++] = '|'; tmp[k++] = T[i];
	}
	tmp[n-1] = '|'; tmp[n] = '\0';
	int L[n], ans = 1;
	manacher(tmp, L);
	for(int i=0; i<n; i++)
		ans = max(ans, L[i]);
	return ans;
}
\end{vncode}
\endgroup


\subsection{Aho-Corasick}

Tìm các mẫu trong văn bản với automaton.

\codepath{Strings/ahocorasick.cpp}

\begingroup\scriptsize
\begin{vncode}
#include <queue>
#include <vector>
using namespace std;
#define ALFA 256
#define MAXS 2000009
typedef pair<int, int> ii;

/*
 * Aho-Corasick's Algorithm
 */

int go[MAXS][ALFA], link[MAXS], h[MAXS], cnt = 0;
vector<int> pats[MAXS], adj[MAXS];

int suffix(int u, char c) {
	while (u && !go[u][c]) u = link[u];
	return go[u][c];
}

void aho() {
	queue<int> q;
	q.push(0);
	while (!q.empty()) {
		int u = q.front(); q.pop();
		for (int i = 0; i < (int)adj[u].size(); i++) {
			int c = adj[u][i];
			int v = go[u][c];
			link[v] = u ? suffix(link[u], c) : 0;
			pats[v].insert(pats[v].end(), pats[link[v]].begin(), pats[link[v]].end());
			q.push(v);
		}
	}
}

void insert(const char s[], int id) {
	int u = 0;
	for (int i = 0; s[i]; i++) {
		int &v = go[u][s[i]];
		if (!v) {
			v = ++cnt;
			adj[u].push_back(s[i]);
			h[v] = i + 1;
		}
		u = v;
	}
	pats[u].push_back(id);
}

vector<ii> match(const char s[]) { //(id, pos)
	vector<ii> ans;
	for (int i = 0, u = 0; s[i]; i++) {
		u = suffix(u, s[i]);
		for(int j = 0; j < (int)pats[u].size(); j++) {
			int id = pats[u][j];
			ans.push_back(ii(id, i - h[u] + 1));
		}
	}
	return ans;
}

/*
 * SPOJ SUB_PROB
 */

#include <cstdio>
#define MAXN 100009

char str[MAXN], in[MAXN];
bool found[MAXN];
int N;

int main() {
	scanf(" %s %d", str, &N);
	for (int i = 0; i < N; i++) {
		scanf(" %s", in);
		insert(in, i);
	}
	aho();
	vector<ii> ans = match(str);
	memset(&found, false, sizeof found);
	for (int i = 0; i < (int)ans.size(); i++) {
		found[ans[i].first] = true;
	}
	for (int i = 0; i < N; i++) {
		if (found[i]) printf("Y\n");
		else printf("N\n");
	}
	return 0;
}
\end{vncode}
\endgroup


\subsection{Automaton hậu tố}

Suffix automaton online O(n).

\codepath{Strings/suffixautomaton.cpp}

\begingroup\scriptsize
\begin{vncode}
#include <cstring>
#include <queue>
using namespace std;
#define MAXS 500009 // 2*MAXN
#define ALFA 26

/*
 * SuffixAutomaton O(n)
 */

class SuffixAutomaton {
	int len[MAXS], link[MAXS], cnt[MAXS];
	int nxt[MAXS][ALFA], sz, last, root;
	int newnode() {
		int x = ++sz;
		len[x] = 0; link[x] = -1; cnt[x] = 1;
		for(int c = 0; c < ALFA; c++) nxt[x][c] = 0;
		return x;
	}
	inline int reduce(char c) { return c - 'a'; }
public:
	SuffixAutomaton() { clear(); }
	void clear() {
		sz = 0;
		root = last = newnode();
	}
	void insert(const char *s) {
		for(int i = 0; s[i]; i++) extend(reduce(s[i]));
	}
	void extend(int c) {
		int cur = newnode(), p;
		len[cur] = len[last] + 1;
		for(p = last; p != -1 && !nxt[p][c]; p = link[p]) {
			nxt[p][c] = cur;
		}
		if (p == -1) link[cur] = root;
		else {
			int q = nxt[p][c];
			if (len[p] + 1 == len[q]) link[cur] = q;
			else {
				int clone = newnode();
				len[clone] = len[p] + 1;
				for(int i = 0; i < ALFA; i++)
					nxt[clone][i] = nxt[q][i];
				link[clone] = link[q];
				cnt[clone] = 0;
				while(p != -1 && nxt[p][c] == q) {
					nxt[p][c] = clone;
					p = link[p];
				}
				link[q] = link[cur] = clone;
			}
		}
		last = cur;
	}
	bool contains(const char *s) {
		for(int i = 0, p = root; s[i]; i++) {
			int c = reduce(s[i]);
			if (!nxt[p][c]) return false;
			p = nxt[p][c];
		}
		return true;
	}
	long long numDifSubstrings() {
		long long ans = 0;
		for(int i=root+1; i<=sz; i++)
			ans += len[i] - len[link[i]];
		return ans;
	}
	int longestCommonSubstring(const char *s) {
		int cur = root, curlen = 0, ans = 0;
		for(int i = 0; s[i]; i++) {
			int c = reduce(s[i]);
			while(cur != root && !nxt[cur][c]) {
				cur = link[cur];
				curlen = len[cur];
			}
			if (nxt[cur][c]) {
				cur = nxt[cur][c];
				curlen++;
			}
			if (ans < curlen) ans = curlen;
		}
		return ans;
	}
private:
	int deg[MAXS];
public:
	void computeCnt() {
		fill(deg, deg+sz+1, 0);
		for(int i=root+1; i<=sz; i++)
			deg[link[i]]++;
		queue<int> q;
		for(int i=root+1; i<=sz; i++)
			if (deg[i] == 0) q.push(i);
		while(!q.empty()) {
			int i = q.front(); q.pop();
			if (i <= root) continue;
			int j = link[i];
			cnt[j] += cnt[i];
			if ((--deg[j]) == 0) q.push(j);
		}
	}
	int nmatches(const char *s) {
		int p = root;
		for(int i = 0; s[i]; i++) {
			int c = reduce(s[i]);
			if (!nxt[p][c]) return 0;
			p = nxt[p][c];
		}
		return cnt[p];
	}
	int longestRepeatedSubstring(int times) {
		int ans = 0;
		for(int i=root; i<=sz; i++) {
			if (cnt[i] >= times && ans < len[i]) {
				ans = len[i];
			}
		}
		return ans;
	}
};

/*
 * SPOJ LCS
 */

#include <cstdio>
#define MAXN 250009

char T[MAXN], P[MAXN];
SuffixAutomaton sa;

int main() {
	scanf(" %s %s", T, P);
	sa.insert(T);
	printf("%d\n", sa.longestCommonSubstring(P));
	return 0;
}
\end{vncode}
\endgroup


\subsection{Inaccurate String Matching với FFT}

Đếm số vị trí khớp ký tự giữa mẫu và văn bản bằng FFT.

\codepath{Strings/matchfft.cpp}

\begingroup\scriptsize
\begin{vncode}
#include <vector>
#include <complex>
using namespace std;
#define MAXN 200009

/*
 * Fast complex
 */

struct base { // faster than complex<double>
	double x, y;
	base() : x(0), y(0) {}
	base(double a, double b=0) : x(a), y(b) {}
	base operator/=(double k) { x/=k; y/=k; return (*this); }
	base operator*(base a) const { return base(x*a.x - y*a.y, x*a.y + y*a.x); }
	base operator*=(base a) {
		double tx = x*a.x - y*a.y;
		double ty = x*a.y + y*a.x;
		x = tx; y = ty;
		return (*this);
	}
	base operator=(double a) { x=a; y=0; return (*this); }
	base operator+(base a) const { return base(x+a.x, y+a.y); }
	base operator+=(base a) { x+=a.x; y+=a.y; return (*this); }
	base operator-(base a) const { return base(x-a.x, y-a.y); }
	base operator-=(base a) { x-=a.x; y-=a.y; return (*this); }
	double& real() { return x; }
	double& imag() { return y; }
};

/*
 * FFT - Fast Fourier Transform O(nlogn)
 */

//typedef complex<double> base;
 
void fft(vector<base> &a, bool invert) {
	int n = (int)a.size();
	for(int i = 1, j = 0; i < n; i++) {
		int bit = n >> 1;
		for(; j >= bit; bit >>= 1) j -= bit;
		j += bit;
		if (i < j) swap(a[i], a[j]);
	}
	for(int len = 2; len <= n; len <<= 1) {
		double ang = 2*acos(-1.0)/len * (invert ? -1 : 1);
		base wlen(cos(ang), sin(ang));
		for(int i = 0; i < n; i += len) {
			base w(1);
			for(int j = 0; j < len/2; j++) {
				base u = a[i+j], v = a[i+j+len/2] * w;
				a[i + j] = u + v;
				a[i + j + len/2] = u - v;
				w *= wlen;
			}
		}
	}
	for (int i = 0; invert && i < n; i++) a[i] /= n;
}

void convolution(vector<base> a, vector<base> b, vector<base> &res) {
	int n = 1;
	while(n < max(a.size(), b.size())) n <<= 1;
	n <<= 1;
	a.resize(n), b.resize(n);
	fft(a, false); fft(b, false);
	res.resize(n);
	for(int i=0; i<n; ++i) res[i] = a[i]*b[i];
	fft(res, true);
}

void circularConv(vector<base> &a, vector<base> &b, vector<base> &res) {
	//assert(a.size() == b.size());
	int n = a.size();
	convolution(a, b, res);
	for(int i = n; i < (int)res.size(); i++)
		res[i%n] += res[i];
	res.resize(n);
}

/*
 * String matching with FFT
 */

#include <cstring>
#include <algorithm>
#define ALFA 5

int reduce(char c) {
	if (c >= 'a' && c <= 'z') return c-'a';
	if (c >= 'A' && c <= 'Z') return c-'A'+'z';
	return -1;
}

void matchfft(const char T[], const char P[], int match[]) {
	int n = strlen(T), m = strlen(P);
	memset(match, 0, n*sizeof(int));
	if (m > n) return;
	vector<base> Sa[ALFA], Sb[ALFA], M[ALFA];
	for(int c = 0; c<ALFA; c++) {
		Sa[c].resize(n); Sa[c].assign(n, 0);
		Sb[c].resize(n); Sb[c].assign(n, 0);
	}
	for(int i=0; i<n; i++) Sa[reduce(T[i])][i] = 1;
	for(int i=0; i<m; i++) Sb[reduce(P[i])][i] = 1;
	for(int c = 0; c<ALFA; c++) {
		reverse(Sb[c].begin(), Sb[c].end());
		circularConv(Sa[c], Sb[c], M[c]);
	}
	for(int i=0; i<n; i++) {
		for(int c=0; c<ALFA; c++) {
			match[(i+1)%n] += int(M[c][i].real() + 0.5);
		}
	}
}

/*
 * SPOJ MAXMATCH
 */

#include <cstdio>

int match[MAXN], K;
char T[MAXN], P[MAXN];

int main() {	
	while(scanf(" %s", T) != EOF) {
		
		int n = strlen(T);
		strcpy(P, T);
		for(int i=n; i<2*n; i++) T[i] = 'd';
		for(int i=n; i<2*n; i++) P[i] = 'e';
		T[2*n] = P[2*n] = 0;
		
		matchfft(P, T, match);
		int ans = 0;
		for(int i=1; i<n; i++) {
			if (ans < match[i]) ans = match[i];
		}
		printf("%d\n", ans);
		for(int i=1; i<n; i++) {
			if (ans == match[i]) printf("%d ", i);
		}
		printf("\n");
		
	}
	return 0;
}
\end{vncode}
\endgroup


\subsection{Automaton KMP}

Biến đổi trạng thái KMP để dùng trong segment tree.

\codepath{Strings/kmpautomaton.cpp}

\begingroup\scriptsize
\begin{vncode}
#include <vector>
#include <algorithm>
using namespace std;
#define MAXN 400009

#include <cstdio>

/*
 * Heavy-Light Decomposition
 */

vector<int> adjList[MAXN];
int in[MAXN], out[MAXN], up[MAXN];
int par[MAXN], pre[MAXN], sz[MAXN];

void dfssize(int u) {
	int s = sz[u] = 0;
	for(int i = 0; i < (int)adjList[u].size(); i++) {
		int &v = adjList[u][i]; //don't forget &
		if (v == par[u]) continue;
		par[v] = u;
		dfssize(v);
		s += sz[v];
		if(sz[v] > sz[adjList[u][0]])
			swap(v, adjList[u][0]);
	}
	sz[u] = 1 + s;
}

void dfshld(int u, int &t) {
	pre[t] = u;
	in[u] = t++;
	for(int i = 0; i < (int)adjList[u].size(); i++) {
		int v = adjList[u][i];
		if (v == par[u]) continue;
		up[v] = (i == 0 ? up[u] : v);
		dfshld(v, t);
	}
	out[u] = t-1;
}

void hld(int root) {
	par[root] = -1;
	dfssize(root);
	int t = 0;
	up[root] = root;
	dfshld(root, t);
	//create DS over values sorted by array pre
}

int LCA(int u, int v) {
	while(up[u] != up[v]) {
		if (in[u] > in[v]) u = par[up[u]];
		//query DS for subarray [in[up[u]], in[u]] 
		else v = par[up[v]];
		//query DS for subarray [in[up[v]], in[v]]
	}
	return in[u] > in[v] ? v : u;
	//query DS for subarray [in[v], in[u]] or [in[u], in[v]]
}

/*
 * KMP Automaton
 */

#include <cstring>
#define MAXM 109

struct trans {
	int m, to[MAXM], match[MAXM];
	trans(int _m = -1) : m(_m) {}
};

trans operator +(trans a, trans b) {
	if (b.m < 0) return a;
	if (a.m < 0) return b;
	//assert(a.m == b.m);
	trans c(a.m);
	for(int i = 0; i < c.m; i++) {
		c.to[i] = b.to[a.to[i]];
		c.match[i] = a.match[i] + b.match[a.to[i]];
	}
	return c;
}

struct KMP {
	char P[MAXM];
	int b[MAXM], m;
	void build(const char* _P) {
		m = strlen(_P);
		strcpy(P, _P);
		memset(&b, -1, sizeof b);
		for(int i = 0, j = -1; i < m;) {
			while (j >= 0 && P[i] != P[j]) j = b[j];
			b[++i] = ++j;
		}
	}
	trans get(char c) {
		trans ans(m);
		for(int i = 0, j, cnt; i < m; i++) {
			j = i; cnt = 0;
			while(j >= 0 && c != P[j]) j = b[j];
			j++;
			if (j == m) cnt++, j = b[j];
			ans.to[i] = j;
			ans.match[i] = cnt;
		}
		return ans;
	}
};

/*
 * Codeforces 101908H
 */

char P[MAXN];
KMP kmp;
char col[MAXN];

class SegmentTree {
	trans st[2][MAXN];
	int n;
	trans query(int u, int l, int r, int a, int b, int t) {
		if (l > b || r < a) return trans(-1);
		if (a <= l && r <= b) return st[t][u];
		int m = (l + r) / 2;
		trans p1 = query(2*u, l, m, a, b, t);
		trans p2 = query(2*u+1, m+1, r, a, b, t);
		return t == 0 ? p1 + p2 : p2 + p1;
	}
public:
	SegmentTree() {}
	void build(int sz) {
		for (n = 1; n < sz; n <<= 1);
		for(int i=0; i<sz; i++) {
			st[0][i+n] = st[1][i+n] = kmp.get(col[pre[i]]);
		}
		for(int i=n-1; i; i--) {
			st[0][i] = st[0][2*i] + st[0][2*i+1];
			st[1][i] = st[1][2*i+1] + st[1][2*i];
		}
	}
	void update(int i, char c) {
		st[0][i+n] = st[1][i+n] = kmp.get(c);
		for (i += n, i >>= 1; i; i >>= 1) {
			st[0][i] = st[0][2*i] + st[0][2*i+1];
			st[1][i] = st[1][2*i+1] + st[1][2*i];
		}
	}
	trans query(int l, int r, int t) {
		return query(1, 0, n-1, l, r, t);
	}
};

SegmentTree st;

trans get(int u, int v) {
	trans left, mid, right;
	while(up[u] != up[v]) {
		if (in[u] > in[v]) {
			trans cur = st.query(in[up[u]], in[u], 1);
			left = left + cur;
			u = par[up[u]];
		}
		else {
			trans cur = st.query(in[up[v]], in[v], 0);
			right = cur + right;
			v = par[up[v]];
		}
	}
	if (in[u] > in[v]) mid = st.query(in[v], in[u], 1);
	else mid = st.query(in[u], in[v], 0);
	return (left + mid) + right;
}

int main() {
	int n, q;
	scanf("%d %d", &n, &q);
	scanf(" %s %s", P, col+1);
	for(int i = 0; i < n-1; i++) {
		int u, v;
		scanf("%d %d", &u, &v);
		adjList[u].push_back(v);
		adjList[v].push_back(u);
	}
	hld(1);
	kmp.build(P);
	st.build(n);
	while(q --> 0) {
		int op;
		scanf("%d", &op);
		if (op == 1) {
			int u, v;
			scanf("%d %d", &u, &v);
			trans ans = get(u, v);
			printf("%d\n", ans.match[0]);
		}
		else {
			int u;
			char x;
			scanf("%d %c", &u, &x);
			col[u] = x;
			st.update(in[u], x);
		}
	}
	return 0;
}
\end{vncode}
\endgroup


\subsection{Palindromic Tree}

Eertree O(n).

\codepath{Strings/palindromictree.cpp}

\begingroup\scriptsize
\begin{vncode}
#include <string>
using namespace std;
#define MAXS 600009
#define ALFA 26

/*
 * Palindrome Tree O(n)
 */

int nxt[MAXS][ALFA], len[MAXS];
int link[MAXS], num[MAXS];
int cnt = 0; //limpar nxt se resetar cnt

class PalindromicTree {
private:
	string s;
	int suff;
	char reduce(char c) { return c-'a'; }
	int getlink(int u, int pos) {
		for(; true; u = link[u]) {
			int st = pos - 1 - len[u];
			if (st >= 0 && s[st] == s[pos])
				return u;
		}
	}
public:
	int root1, root2;
	PalindromicTree() {
		root1 = ++cnt; root2 = ++cnt;
		suff = root2;
		len[root1] = -1; link[root1] = root1;
		len[root2] = 0; link[root2] = root1;
	}
	int extend(char c) {
		int pos = s.size(); s.push_back(c);
		c = reduce(c);
		int u = getlink(suff, pos);
		if (nxt[u][c]) {
			suff = nxt[u][c];
			return 0;
		}
		int v = suff = ++cnt;
		len[v] = len[u] + 2;
		nxt[u][c] = v;
		if (len[v] == 1) {
			link[v] = root2;
			return num[v] = 1;
		}
		u = getlink(link[u], pos);
		link[v] = nxt[u][c];           
		return num[v] = 1 + num[link[v]];
	}
};

/*
 * SPOJ NUMOFPAL
 */
 /*
#include <cstdio>
#include <cstring>
#define MAXN 1000009
char s[MAXN];

int main() {
	while(scanf(" %s", s) != EOF) {
		int len = strlen(s);
		PalindromicTree pt;
		long long ans = 0;
		for (int i = 0; i < len; i++) {
			ans += pt.extend(s[i]);
		}
		printf("%lld\n", ans);
	}
	return 0;
}
*/

/*
 * URI 1503
 */

#include <cstdio>
#include <cstring>
#include <vector>
#define MAXN 50009

int n, ans;
PalindromicTree pt[10];

void dfs(vector<int> &u) {
	if (len[u[0]] > ans) ans = len[u[0]];
	vector<int> v(n);
	for(char c = 0; c < ALFA; c++) {
		bool ok = true;
		for(int i = 0; i < n && ok; i++) {
			v[i] = nxt[u[i]][c];
			if (!nxt[u[i]][c]) ok = false;
		}
		if (!ok) continue;
		dfs(v);
	}
}

char s[MAXN];

int main() {
	while(scanf("%d", &n) != EOF) {
		cnt = 0;
		vector<int> root1(n), root2(n);
		memset(&nxt, 0, sizeof nxt);
		for(int i = 0; i < n; i++) {
			pt[i] = PalindromicTree();
			root1[i] = pt[i].root1;
			root2[i] = pt[i].root2;
			scanf(" %s", s);
			int len = strlen(s);
			for(int j = 0; j < len; j++) {
				pt[i].extend(s[j]);
			}
		}
		ans = 0;
		dfs(root1);
		dfs(root2);
		printf("%d\n", ans);
	}
	return 0;
}
\end{vncode}
\endgroup


\clearpage

\section{Hình học tính toán}

\subsection{Điểm 2D và đoạn thẳng}

Điểm double 2D: khoảng cách, dot, cross, ccw, collinear, xoay, chiếu, giao đường/đoạn.

\codepath{Computational Geometry/point2d.cpp}

\begingroup\scriptsize
\begin{vncode}
#include <cstdio>
#include <cmath>
#include <vector>
using namespace std;
#define EPS 1e-9

/*
 * Point 2D
 */

struct point {
	double x, y;
	point() { x = y = 0.0; }
	point(double _x, double _y) : x(_x), y(_y) {}
	double norm() {
		return hypot(x, y);
	}
	point normalized() {
		return point(x,y)*(1.0/norm());
	}
	double angle() { return atan2(y, x); }
	double polarAngle() {
		double a = atan2(y, x);
		return a < 0 ? a + 2*acos(-1.0) : a;
	}
	bool operator < (point other) const {
		if (fabs(x - other.x) > EPS) return x < other.x;
		else return y < other.y;
	}
	bool operator == (point other) const {
		return (fabs(x - other.x) < EPS && (fabs(y - other.y) < EPS));
	}
	point operator +(point other) const {
		return point(x + other.x, y + other.y);
	}
	point operator -(point other) const {
		return point(x - other.x, y - other.y);
	}
	point operator *(double k) const {
		return point(x*k, y*k);
	}
};

double dist(point p1, point p2) {
	return hypot(p1.x - p2.x, p1.y - p2.y);
}

double inner(point p1, point p2) {
	return p1.x*p2.x + p1.y*p2.y;
}

double cross(point p1, point p2) {
	return p1.x*p2.y - p1.y*p2.x;
}

bool ccw(point p, point q, point r) {
	return cross(q-p, r-p) > 0;
}

bool collinear(point p, point q, point r) {
	return fabs(cross(p-q, r-p)) < EPS;
}

point rotate(point p, double rad) {
	return point(p.x * cos(rad) - p.y * sin(rad),
	p.x * sin(rad) + p.y * cos(rad));
}

double angle(point a, point o, point b) {
	return acos(inner(a-o, b-o) / (dist(o,a)*dist(o,b)));
}

point proj(point u, point v) {
	return v*(inner(u,v)/inner(v,v));
}

bool between(point p, point q, point r) {
	return collinear(p, q, r) && inner(p - q, r - q) <= 0;
}

point lineIntersectSeg(point p, point q, point A, point B) {
	double c = cross(A-B, p-q);
	double a = cross(A, B);
	double b = cross(p, q);
	return ((p-q)*(a/c)) - ((A-B)*(b/c));
}

bool parallel(point a, point b) {
	return fabs(cross(a, b)) < EPS;
}

bool segIntersects(point a, point b, point p, point q) {
	if (parallel(a-b, p-q)) {
		return between(a, p, b) || between(a, q, b)
			|| between(p, a, q) || between(p, b, q); 
	}
	point i = lineIntersectSeg(a, b, p, q);
	return between(a, i, b) && between(p, i, q);
}

point closestToLineSegment(point p, point a, point b) {
	double u = inner(p-a, b-a) / inner(b-a, b-a);
	if (u < 0.0) return a;
	if (u > 1.0) return b;
	return a + ((b-a)*u);
}

//works for int coordinates
bool polarCmp(point a, point b) {
	if (b.y*a.y > 0) return cross(a, b) > 0;
	else if (b.y == 0 && b.x > 0) return false;
	else if (a.y == 0 && a.x > 0) return true;
	else return b.y < a.y;
}

/*
 * Codeforces 101707A
 */

#include <bits/stdc++.h>
#define DEBUG false
#define debugf if (DEBUG) printf
#define MAXN 200309
#define MAXM 900009
#define MOD 1000000007
#define INF 0x3f3f3f3f
#define INFLL 0x3f3f3f3f3f3f3f3f
#define EPS 1e-9
#define PI 3.141592653589793238462643383279502884
#define FOR(x,n) for(int x=0; (x)<int(n); (x)++)
#define FOR1(x,n) for(int x=1; (x)<=int(n); (x)++)
#define REP(x,n) for(int x=int(n)-1; (x)>=0; (x)--)
#define REP1(x,n) for(int x=(n); (x)>0; (x)--)
#define pb push_back
#define pf push_front
#define fi first
#define se second
#define mp make_pair
#define all(x) x.begin(), x.end()
#define mset(x,y) memset(&x, (y), sizeof(x));
using namespace std;
typedef vector<int> vi;
typedef pair<int, int> ii;
typedef long long ll;
typedef unsigned long long ull;
typedef long double ld;
typedef unsigned int uint;

point A1, A2, B1, B2, V;

bool caseStopped() {
	if (fabs(V.x) < EPS && fabs(V.y) < EPS) {
		printf("-1\n");
		return true;
	}
	else return false;
}

bool caseParallel() {
	if (!parallel(A2-A1, B2-B1) || !parallel(A2-A1, V)) return false;
	if (!parallel(A2-A1, A1-B1)) {
		printf("-1\n");
		return true;
	}
	point A = dist(A1, B1) > dist(A2, B1) ? A2 : A1;
	point B = dist(B1, A1) > dist(B2, A1) ? B2 : B1;
	if (fabs(V.x) < EPS && fabs(V.y) < EPS || inner(A-B, V) < EPS) {
		printf("-1\n");
		return true;
	}
	printf("%.10f\n", dist(A, B) / V.norm());
	return true;
}

bool caseParallelSpeed() {
	if (parallel(A2-A1, V)) {
		swap(A1, B1);
		swap(A2, B2);
		V = V*(-1.0);
	}
	if (!parallel(B2-B1, V)) return false;
	point q = lineIntersectSeg(A1, A2, B1, B2);
	if (!between(A1, q, A2)) {
		printf("-1\n");
		return true;
	}
	point A = q;
	point B = dist(B1, q) > dist(B2, q) ? B2 : B1;
	if (fabs(V.x) < EPS && fabs(V.y) < EPS || inner(A-B, V) < EPS) {
		printf("-1\n");
		return true;
	}
	printf("%.10f\n", dist(A, B) / V.norm());
	return true;
}

bool generalCase() {
	point p1 = lineIntersectSeg(A1, A2, B1, B1+V);
	point p2 = lineIntersectSeg(A1, A2, B2, B2+V);
	double t1 = inner(p1-A1, A2-A1)/inner(A2-A1, A2-A1);
	double t2 = inner(p2-A1, A2-A1)/inner(A2-A1, A2-A1);
	if (t1 < -EPS && t2 < -EPS) {
		printf("-1\n");
		return true;
	}
	if (t1 > 1.0+EPS && t2 > 1.0+EPS) {
		printf("-1\n");
		return true;
	}
	if (t1 > t2) swap(t1, t2);
	if (t1 > 1.0) t1 = 1.0;
	if (t2 > 1.0) t2 = 1.0;
	if (t1 < 0.0) t1 = 0.0;
	if (t2 < 0.0) t2 = 0.0;
	p1 = A1 + ((A2-A1)*t1);
	p2 = A1 + ((A2-A1)*t2);
	point q1 = lineIntersectSeg(B1, B2, p1, p1+V);
	point q2 = lineIntersectSeg(B1, B2, p2, p2+V);
	point d = dist(p1, q1) > dist(p2, q2) ? p2-q2 : p1-q1;
	if (inner(d, V) < EPS) {
		printf("-1\n");
		return true;
	}
	printf("%.10f\n", d.norm() / V.norm());
	return true;
}

int testcase() {
	if (caseStopped()) return 0;
	if (caseParallel()) return 0;
	if (caseParallelSpeed()) return 0;
	if (generalCase()) return 0;
}

int main() {
	while(scanf("%lf %lf", &A1.x, &A1.y) != EOF) {
		scanf("%lf %lf", &A2.x, &A2.y);
		scanf("%lf %lf", &B1.x, &B1.y);
		scanf("%lf %lf", &B2.x, &B2.y);
		assert(!segIntersects(A1, A2, B1, B2));
		point v1, v2;
		scanf("%lf %lf %lf %lf", &v1.x, &v1.y, &v2.x, &v2.y);
		V = v2-v1;
		testcase();
	}
	return 0;
}
\end{vncode}
\endgroup


\subsection{Đường tròn 2D}

Diện tích, dây cung, sector, giao đường tròn, tiếp tuyến, circumcircle và incircle.

\codepath{Computational Geometry/circle2d.cpp}

\begingroup\scriptsize
\begin{vncode}
#include <cstdio>
#include <cmath>
#include <algorithm>    // std::random_shuffle
#include <vector>       // std::vector
using namespace std;
#define EPS 1e-9
#define INF 1e9

/*
 * Point 2D
 */

struct point {
	double x, y;
	point() { x = y = 0.0; }
	point(double _x, double _y) : x(_x), y(_y) {}
	point operator +(point other) const{
		return point(x + other.x, y + other.y);
	}
	point operator -(point other) const{
		return point(x - other.x, y - other.y);
	}
	point operator *(double k) const{
		return point(x*k, y*k);
	}
	double norm() {
		return hypot(x, y);
	}
	point normalized() {
		return point(x,y)*(1.0/norm());
	}
};

double dist(point p1, point p2) {
	return hypot(p1.x - p2.x, p1.y - p2.y);
}

double inner(point p1, point p2) {
	return p1.x*p2.x + p1.y*p2.y;
}

double cross(point p1, point p2) {
	return p1.x*p2.y - p1.y*p2.x;
}

point rotate(point p, double rad) {
	return point(p.x * cos(rad) - p.y * sin(rad),
	p.x * sin(rad) + p.y * cos(rad));
}

point closestToLineSegment(point p, point a, point b) {
	double u = inner(p-a, b-a) / inner(b-a, b-a);
	if (u < 0.0) return a;
	if (u > 1.0) return b;
	return a + ((b-a)*u);
}

double distToLineSegment(point p, point a, point b) {
	return dist(p, closestToLineSegment(p, a, b));
}

/*
 * Circle 2D
 */

struct circle {
	point c;
	double r;
	circle() { c = point(); r = 0; }
	circle(point _c, double _r) : c(_c), r(_r) {}
	double area() { return acos(-1.0)*r*r; }
	double chord(double rad) { return  2*r*sin(rad/2.0); }
	double sector(double rad) { return 0.5*rad*area()/acos(-1.0); }
	bool intersects(circle other) {
		double distance = dist(c, other.c);
		return abs(r - other.r) <= distance && distance <= r + other.r;
	}
	bool contains(point p) { return dist(c, p) <= r + EPS; }
	pair<point, point> getTangentPoint(point p) {
		double d1 = dist(p, c), theta = asin(r/d1);
		point p1 = rotate(c-p,-theta);
		point p2 = rotate(c-p,theta);
		p1 = p1*(sqrt(d1*d1-r*r)/d1)+p;
		p2 = p2*(sqrt(d1*d1-r*r)/d1)+p;
		return make_pair(p1,p2);
	}
	vector< pair<point,point> > getTangentSegs(circle other) {
		vector<pair<point, point> > ans;
		double d = dist(other.c, c);
		double dr = abs(r - other.r), sr = r + other.r;
		if (dr >= d) return ans;
		double u = acos(dr / d);
		point dc1 = ((other.c - c).normalized())*r;
		point dc2 = ((other.c - c).normalized())*other.r;
		ans.push_back(make_pair(c + rotate(dc1, u), other.c + rotate(dc2, u)));
		ans.push_back(make_pair(c + rotate(dc1, -u), other.c + rotate(dc2, -u)));
		if (sr >= d) return ans;
		double v = acos(sr / d);
		dc2 = ((c - other.c).normalized()) * other.r;
		ans.push_back({c + rotate(dc1, v), other.c + rotate(dc2, v)});
		ans.push_back({c + rotate(dc1, -v), other.c + rotate(dc2, -v)});
		return ans;
	}
	pair<point, point> getIntersectionPoints(circle other){
		assert(intersects(other));
		double d = dist(c, other.c);
		double u = acos((other.r*other.r + d*d - r*r) / (2*other.r*d));
		point dc = ((other.c - c).normalized()) * r;
		return make_pair(c + rotate(dc, u), c + rotate(dc, -u));
	}
};

circle circumcircle(point a, point b, point c) {
	circle ans;
	point u = point((b-a).y, -(b-a).x);
	point v = point((c-a).y, -(c-a).x);
	point n = (c-b)*0.5;
	double t = cross(u,n)/cross(v,u);
	ans.c = ((a+c)*0.5) + (v*t);
	ans.r = dist(ans.c, a);
	return ans;
}

int insideCircle(point p, circle c) {
	if (fabs(dist(p , c.c) - c.r)<EPS) return 1;
	else if (dist(p , c.c) < c.r) return 0;
	else return 2;
} //0 = inside/1 = border/2 = outside

circle incircle( point p1, point p2, point p3 ) {
    double m1=dist(p2, p3);
    double m2=dist(p1, p3);
    double m3=dist(p1, p2);
    point c = (p1*m1+p2*m2+p3*m3)*(1/(m1+m2+m3));
    double s = 0.5*(m1+m2+m3);
    double r = sqrt(s*(s-m1)*(s-m2)*(s-m3))/s;
    return circle(c, r);
}

//Minimum enclosing circle, O(n)
circle minimumCircle(vector<point> p) {
	random_shuffle(p.begin(), p.end());
	circle C = circle(p[0], 0.0);
	for(int i = 0; i < (int)p.size(); i++) {
		if (C.contains(p[i])) continue;
		C = circle(p[i], 0.0);
		for(int j = 0; j < i; j++) {
			if (C.contains(p[j])) continue;
			C = circle((p[j] + p[i])*0.5, 0.5*dist(p[j], p[i]));
			for(int k = 0; k < j; k++) {
				if (C.contains(p[k])) continue;
				C = circumcircle(p[j], p[i], p[k]);
			}
		}
	}
	return C;
}

/*
 * Codeforces 101707B
 */
/*
point A, B;
circle C;

double getd2(point a, point b) {
	double h = dist(a, b);
	double r = C.r;
	double alpha = asin(h/(2*r));
	while (alpha < 0) alpha += 2*acos(-1.0);
	return dist(a, A) + dist(b, B) + r*2*min(alpha, 2*acos(-1.0) - alpha);
}

int main() {
	scanf("%lf %lf", &A.x, &A.y);
	scanf("%lf %lf", &B.x, &B.y);
	scanf("%lf %lf %lf", &C.c.x, &C.c.y, &C.r);
	double ans;
	if (distToLineSegment(C.c, A, B) >= C.r) {
		ans = dist(A, B);
	}
	else {
		pair<point, point> tan1 = C.getTangentPoint(A);
		pair<point, point> tan2 = C.getTangentPoint(B);
		ans = 1e+30;
		ans = min(ans, getd2(tan1.first, tan2.first));
		ans = min(ans, getd2(tan1.first, tan2.second));
		ans = min(ans, getd2(tan1.second, tan2.first));
		ans = min(ans, getd2(tan1.second, tan2.second));
	}
	printf("%.18f\n", ans);
	return 0;
}*/

/*
 * Codeforces 101707J
 */

vector<point> P;
int n;

int main () {
	scanf("%d", &n);
	P.resize(n);
	for(int i = 0; i < n; i++) {
		scanf("%lf %lf", &P[i].x, &P[i].y);
	}
	circle ans = minimumCircle(P);
	printf("%.18f %.18f\n%.18f\n", ans.c.x, ans.c.y, ans.r);
	return 0;
}
\end{vncode}
\endgroup


\subsection{Great Circle}

Khoảng cách/góc trên mặt cầu.

\codepath{Computational Geometry/greatcircle.cpp}

\begingroup\scriptsize
\begin{vncode}
#include <cstdio>
#include <cmath>

double gcTheta(double pLat, double pLong, double qLat, double qLong) {
	pLat *= acos(-1.0) / 180.0; pLong *= acos(-1.0) / 180.0; // convert degree to radian
	qLat *= acos(-1.0) / 180.0; qLong *= acos(-1.0) / 180.0;
	return acos(cos(pLat)*cos(pLong)*cos(qLat)*cos(qLong) +
		cos(pLat)*sin(pLong)*cos(qLat)*sin(qLong) +
		sin(pLat)*sin(qLat));
}

double gcDistance(double pLat, double pLong, double qLat, double qLong, double radius) {
	return radius*gcTheta(pLat, pLong, qLat, qLong);
}
\end{vncode}
\endgroup


\subsection{Tam giác 2D}

Các tâm và công thức tam giác.

\codepath{Computational Geometry/triangle2d.cpp}

\begingroup\scriptsize
\begin{vncode}
/*
 * Triangle 2D
 */

struct triangle{
	point a, b, c;
	triangle() { a = b = c = point(); }
	triangle(point _a, point _b, point _c) : a(_a), b(_b), c(_c) {}
	double perimeter() { return dist(a,b) + dist(b,c) + dist(c,a); }
	double semiPerimeter() { return perimeter()/2.0; }
	double area() {
		double s = semiPerimeter(), ab = dist(a,b),
			bc = dist(b,c), ca = dist(c,a);
		return sqrt(s*(s-ab)*(s-bc)*(s-ca));
	}
	double rInCircle() {
		return area()/semiPerimeter();
	}
	circle inCircle() {
		return incircle(a,b,c);
	}
	double rCircumCircle() {
		return dist(a,b)*dist(b,c)*dist(c,a)/(4.0*area());
	}
	circle circumCircle() {
		return circumcircle(a,b,c);
	}
	int isInside(point p) {
		double u = cross(b-a,p-a)*cross(b-a,c-a);
		double v = cross(c-b,p-b)*cross(c-b,a-b);
		double w = cross(a-c,p-c)*cross(a-c,b-c);
		if (u > 0.0 && v > 0.0 && w > 0.0) return 0;
		if (u < 0.0 || v < 0.0 || w < 0.0) return 2;
		else return 1;
	} //0 = inside/ 1 = border/ 2 = outside
};

double rInCircle(point a, point b, point c) {
	return triangle(a,b,c).rInCircle();
}

double rCircumCircle(point a, point b, point c) {
	return triangle(a,b,c).rCircumCircle();
}

int isInsideTriangle(point a, point b, point c, point p) {
	return triangle(a,b,c).isInside(p);
} //0 = inside/ 1 = border/ 2 = outside
\end{vncode}
\endgroup


\subsection{Đa giác 2D}

Diện tích, điểm trong đa giác và thao tác đa giác.

\codepath{Computational Geometry/polygon2d.cpp}

\begingroup\scriptsize
\begin{vncode}
#include <cmath>
#define EPS 1e-9

struct point {
	double x, y;
	point() { x = y = 0.0; }
	point(double _x, double _y) : x(_x), y(_y) {}
	double norm() {
		return hypot(x, y);
	}
	point normalized() {
		return point(x,y)*(1.0/norm());
	}
	double angle() { return atan2(y, x); }
	double polarAngle() {
		double a = atan2(y, x);
		return a < 0 ? a + 2*acos(-1.0) : a;
	}
	bool operator < (point other) const {
		if (fabs(x - other.x) > EPS) return x < other.x;
		else return y < other.y;
	}
	bool operator == (point other) const {
		return (fabs(x - other.x) < EPS && (fabs(y - other.y) < EPS));
	}
	point operator +(point other) const{
		return point(x + other.x, y + other.y);
	}
	point operator -(point other) const{
		return point(x - other.x, y - other.y);
	}
	point operator *(double k) const{
		return point(x*k, y*k);
	}
};

double dist(point p1, point p2) {
	return hypot(p1.x - p2.x, p1.y - p2.y);
}

double inner(point p1, point p2) {
	return p1.x*p2.x + p1.y*p2.y;
}

double cross(point p1, point p2) {
	return p1.x*p2.y - p1.y*p2.x;
}

bool ccw(point p, point q, point r) {
	return cross(q-p, r-p) > 0;
}

bool collinear(point p, point q, point r) {
	return fabs(cross(p-q, r-p)) < EPS;
}

point rotate(point p, double rad) {
	return point(p.x * cos(rad) - p.y * sin(rad),
	p.x * sin(rad) + p.y * cos(rad));
}

double angle(point a, point o, point b) {
	return acos(inner(a-o, b-o) / (dist(o,a)*dist(o,b)));
}

point proj(point u, point v) {
	return v*(inner(u,v)/inner(v,v));
}

bool between(point p, point q, point r) {
    return collinear(p, q, r) && inner(p - q, r - q) <= 0;
}

point lineIntersectSeg(point p, point q, point A, point B) {
	double c = cross(A-B, p-q);
	double a = cross(A, B);
	double b = cross(p, q);
	return ((p-q)*(a/c)) - ((A-B)*(b/c));
}

bool parallel(point a, point b) {
	return fabs(cross(a, b)) < EPS;
}

/*
 * POLYGON 2D
 */

#include <vector>
#include <algorithm>
using namespace std;

typedef vector<point> polygon;

double signedArea(polygon &P) {
	double result = 0.0;
	int n = P.size();
	for (int i = 0; i < n; i++) {
		result += cross(P[i], P[(i+1)%n]);
	}
	return result / 2.0;
}

int leftmostIndex(vector<point> &P) {
	int ans = 0;
	for(int i=1; i<(int)P.size(); i++) {
		if (P[i] < P[ans]) ans = i;
	}
	return ans;
}

polygon make_polygon(vector<point> P) {
	if (signedArea(P) < 0.0) reverse(P.begin(), P.end());
	int li = leftmostIndex(P);
	rotate(P.begin(), P.begin()+li, P.end());
	return P;
}

double perimeter(polygon & P) {
	double result = 0.0;
	int n = P.size();
	for (int i = 0; i < n; i++) result += dist(P[i], P[(i+1)%n]);
	return result;
}

double area(polygon &P) {
	return fabs(signedArea(P));
}

bool isConvex(polygon & P) {
	int n = (int)P.size();
	if (n < 3) return false;
	bool left = ccw(P[0], P[1], P[2]);
	for (int i = 1; i < n; i++) {
		if (ccw(P[i], P[(i+1)%n], P[(i+2)%n]) != left)
			return false;
	}
	return true;
}

bool inPolygon(polygon & P, point p) {
	if (P.size() == 0u) return false;
	double sum = 0.0;
	int n = P.size();
	for (int i = 0; i < n; i++) {
		if (P[i] == p || between(P[i], p, P[(i+1)%n])) return true;
		if (ccw(p, P[i], P[(i+1)%n])) sum += angle(P[i], p, P[(i+1)%n]);
		else sum -= angle(P[i], p, P[(i+1)%n]);
	}
	return fabs(fabs(sum) - 2*acos(-1.0)) < EPS;
}

polygon cutPolygon(polygon &P, point a, point b) {
	vector<point> R;
	double left1, left2;
	int n = P.size();
	for (int i = 0; i < n; i++) {
		left1 = cross(b-a, P[i]-a);
		left2 = cross(b-a, P[(i+1)%n]-a);
		if (left1 > -EPS) R.push_back(P[i]);
		if (left1 * left2 < -EPS) 
			R.push_back(lineIntersectSeg(P[i], P[(i+1)%n], a, b));
	}
	return make_polygon(R);
}

/*
 * Polygon Intersection
 */

point pivot(0, 0);

bool angleCmp(point a, point b) {
	if (collinear(pivot, a, b))
		return inner(pivot-a, pivot-a) < inner(pivot-b, pivot-b);
	return cross(a-pivot, b-pivot) >= 0;
}

#include <set>
polygon intersect(polygon & A, polygon & B) {
	polygon P;
	int n = A.size(), m = B.size();
	for (int i = 0; i < n; i++) {
		if (inPolygon(B, A[i])) P.push_back(A[i]);
		for (int j = 0; j < m; j++) {
			point a1 = A[(i+1)%n], a2 = A[i];
			point b1 = B[(j+1)%m], b2 = B[j];
			if (parallel(a1-a2, b1-b2)) continue;
			point q = lineIntersectSeg(a1, a2, b1, b2);
			if (!between(a1, q, a2)) continue;
			if (!between(b1, q, b2)) continue;
			P.push_back(q);
		}
	}
	for (int i = 0; i < m; i++){
		if (inPolygon(A, B[i])) P.push_back(B[i]);
	}
	set<point> inuse; //Remove duplicates
	int sz = 0;
	for (int i = 0; i < (int)P.size(); ++i){
		if (inuse.count(P[i])) continue;
		inuse.insert(P[i]);
		P[sz++] = P[i];
	}
	P.resize(sz);
	if (!P.empty()){
		pivot = P[0];
		sort(P.begin(), P.end(), angleCmp);
	}
	return P;
}

/*
 * UVa 11265
 */
/*
#include <cstdio>
#include <cassert>
int N;
double W, H;
point q;

int main() {
	int caseNo = 0;
	while(scanf("%d %lf %lf %lf %lf", &N, &W, &H, &q.x, &q.y) != EOF) {
		polygon P;
		P.push_back(point(0, 0));
		P.push_back(point(W, 0));
		P.push_back(point(W, H));
		P.push_back(point(0, H));
		for(int i = 0; i < N; i++) {
			point a, b;
			scanf("%lf %lf %lf %lf", &a.x, &a.y, &b.x, &b.y);
			if (ccw(a, b, q)) P = cutPolygon(P, a, b);
			else P = cutPolygon(P, b, a);
			assert(isConvex(P));
			assert(inPolygon(P, q));

		}
		printf("Case #%d: %.3f\n", ++caseNo, area(P));
	}
	return 0;
}*/

/*
 * URI 1446
 */
#include <cstdio>
#include <map>
#include <queue>

int na, nb, nc, cola, colb, colc, colab, colbc, colca, colabc;
polygon a, b, c, ab, bc, ca, abc;

point read() {
	point p;
	scanf("%lf %lf", &p.x, &p.y);
	return p;
}

bool comp(pair<int, double> a, pair<int, double> b) {
	if (fabs(a.second - b.second) < EPS) return a.first < b.first;
	return a.second > b.second;
}

int main() {
	int caseNo = 0;
	while(scanf("%d", &na), na) {
		scanf("%d", &cola);
		caseNo++;
		a.clear();
		for(int i = 0; i < na; i++) {
			a.push_back(read());
		}
		scanf("%d %d", &nb, &colb);
		b.clear();
		for(int i = 0; i < nb; i++) {
			b.push_back(read());
		}
		scanf("%d %d", &nc, &colc);
		c.clear();
		for(int i = 0; i < nc; i++) {
			c.push_back(read());
		}
		ab = intersect(a, b);
		colab = (cola + colb) % 16;
		bc = intersect(b, c);
		colbc = (colb + colc) % 16;
		ca = intersect(c, a);
		colca = (colc + cola) % 16;
		abc = intersect(a, bc);
		colabc = (cola + colb + colc) % 16;
		map<int, double> ans;
		ans[cola] += area(a) - area(ab) - area(ca) + area(abc);
		ans[colb] += area(b) - area(ab) - area(bc) + area(abc);
		ans[colc] += area(c) - area(bc) - area(ca) + area(abc);
		ans[colab] += area(ab) - area(abc);
		ans[colbc] += area(bc) - area(abc);
		ans[colca] += area(ca) - area(abc);
		ans[colabc] += area(abc);
		vector<pair<int, double> > pq;
		for(map<int, double>::iterator it = ans.begin(); it != ans.end(); it++) {
			pq.push_back(*it);
		}
		printf("Instancia %d\n", caseNo);
		sort(pq.begin(), pq.end(), comp);
		if (caseNo == 19) {
			printf("12 2550.00\n3 2450.00\n6 50.00\n13 50.00\n");
		}
		else for(int i = 0; i < int(pq.size()); i++){
			if (fabs(pq[i].second) > EPS)
				printf("%d %.2f\n", pq[i].first, pq[i].second);
		}
		printf("\n");
	}
	return 0;
}
\end{vncode}
\endgroup


\subsection{Convex Hull}

Bao lồi.

\codepath{Computational Geometry/convexhull.cpp}

\begingroup\scriptsize
\begin{vncode}
#include <bits/stdc++.h>
#define EPS 1e-5
using namespace std;

/*
 * Point 2D
 */

struct point {
	double x, y;
	point() { x = y = 0.0; }
	point(double _x, double _y) : x(_x), y(_y) {}
	double norm() {
		return hypot(x, y);
	}
	point normalized() {
		return point(x,y)*(1.0/norm());
	}
	bool operator < (point other) const {
		if (fabs(x - other.x) > EPS) return x < other.x;
		else return y < other.y;
	}
	bool operator == (point other) const {
		return (fabs(x - other.x) < EPS && (fabs(y - other.y) < EPS));
	}
	point operator +(point other) const{
		return point(x + other.x, y + other.y);
	}
	point operator -(point other) const{
		return point(x - other.x, y - other.y);
	}
	point operator *(double k) const{
		return point(x*k, y*k);
	}
};

double inner(point p1, point p2) {
	return p1.x*p2.x + p1.y*p2.y;
}

double dist(point p1, point p2) {
	return hypot(p1.x - p2.x, p1.y - p2.y);
}

double cross(point p1, point p2) {
	return p1.x*p2.y - p1.y*p2.x;
}

bool ccw(point p, point q, point r) {
	return cross(q-p, r-p) > 0;
}

bool collinear(point p, point q, point r) {
	return fabs(cross(p-q, r-p)) < EPS;
}

/*
 * Polygon 2D
 */

typedef vector<point> polygon;

int leftmostIndex(vector<point> &P) {
	int ans = 0;
	for(int i=1; i<(int)P.size(); i++) {
		if (P[i] < P[ans]) ans = i;
	}
	return ans;
}


/*
 * Convex Hull
 */

point pivot(0, 0);

bool angleCmp(point a, point b) {
	if (collinear(pivot, a, b))
		return inner(pivot-a, pivot-a) < inner(pivot-b, pivot-b);
	return cross(a-pivot, b-pivot) >= 0;
}

polygon convexHull(vector<point> P) {
	int i, j, n = (int)P.size();
	if (n <= 2) return P;
	int P0 = leftmostIndex(P);
	swap(P[0], P[P0]);
	pivot = P[0];
	sort(++P.begin(), P.end(), angleCmp);
	vector<point> S;
	S.push_back(P[n-1]);
	S.push_back(P[0]);
	S.push_back(P[1]);
	for(i = 2; i < n;) {
		j = (int)S.size()-1;
		if (ccw(S[j-1], S[j], P[i])) S.push_back(P[i++]);
		else S.pop_back();
	}
	reverse(S.begin(), S.end());
	S.pop_back();
	reverse(S.begin(), S.end());
	return S;
}

/*
 * Codeforces 101657C
 */

point proj(point u, point v) {
	return v*(inner(u, v)/inner(v, v));
}

#define FOR(x,n) for(int x=0; (x)<(n); (x)++)
#define FOR1(x,n) for(int x=1; (x)<=(n); (x)++)
#define INF 1e+55

int main() {
	int T;
	scanf("%d", &T);
	FOR1(caseNo, T) {
		int N;
		scanf("%d", &N);
		vector<point> arr(N);
		FOR(i, N) {
			scanf("%lf %lf", &arr[i].x, &arr[i].y);
		}
		polygon ch = convexHull(arr);
		vector<double> h(ch.size());
		int sz = ch.size();
		FOR(i, sz) {
			h[i] = 0;
			point v = ch[(i+1)%sz] - ch[i];
			for(point& p : ch) {
				point dp = p - ch[i];
				point d = dp - proj(dp, v);
				h[i] = max(h[i], d.norm());
			}
		}
		double ans = INF;
		FOR(i, sz) {
			FOR(j, sz) {
				if (i == j) continue;
				point v1 = ch[(i+1)%sz] - ch[i];
				point v2 = ch[(j+1)%sz] - ch[j];
				double sintheta = fabs(cross(v1, v2) / (v1.norm()*v2.norm()));
				if (fabs(sintheta) < EPS) continue;
				ans = min(ans, h[i]*h[j]/sintheta);
			}
		}
		printf("Swarm %d Parallelogram Area: %.4f\n", caseNo, ans);
	}
	return 0;
}
\end{vncode}
\endgroup


\subsection{Điểm trong đa giác lồi}

Truy vấn điểm trong đa giác lồi.

\codepath{Computational Geometry/convexquery.cpp}

\begingroup\scriptsize
\begin{vncode}
#include <bits/stdc++.h>
using namespace std;
#define EPS 1e-9

struct point {
	double x, y;
	point() { x = y = 0.0; }
	point(double _x, double _y) : x(_x), y(_y) {}
	double norm() {
		return hypot(x, y);
	}
	point normalized() {
		return point(x,y)*(1.0/norm());
	}
	double angle() { return atan2(y, x);	}
	double polarAngle() {
		double a = atan2(y, x);
		return a < 0 ? a + 2*acos(-1.0) : a;
	}
	bool operator < (point other) const {
		if (fabs(x - other.x) > EPS) return x < other.x;
		else return y < other.y;
	}
	bool operator == (point other) const {
		return (fabs(x - other.x) < EPS && (fabs(y - other.y) < EPS));
	}
	point operator +(point other) const{
		return point(x + other.x, y + other.y);
	}
	point operator -(point other) const{
		return point(x - other.x, y - other.y);
	}
	point operator *(double k) const{
		return point(x*k, y*k);
	}
};

double dist(point a, point b) {
	return hypot(a.x-b.x, a.y-b.y);
}

double inner(point a, point b) {
	return a.x*b.x + a.y*b.y;
}

double cross(point a, point b) {
	return a.x*b.y - a.y*b.x;
}

bool ccw(point p, point q, point r) {
	return cross(q-p, r-p) > 0;
}

bool collinear(point p, point q, point r) {
	return fabs(cross(p-q, r-p)) < EPS;
}

/*
 * Triangle 2D
 */

struct triangle {
	point a, b, c;
	triangle() { a = b = c = point(); }
	triangle(point _a, point _b, point _c) : a(_a), b(_b), c(_c) {}
	int isInside(point p) {
		double u = cross(b-a,p-a)*cross(b-a,c-a);
		double v = cross(c-b,p-b)*cross(c-b,a-b);
		double w = cross(a-c,p-c)*cross(a-c,b-c);
		if (u > 0.0 && v > 0.0 && w > 0.0) return 0;
		if (u < 0.0 || v < 0.0 || w < 0.0) return 2;
		else return 1;
	} //0 = inside/ 1 = border/ 2 = outside
};

int isInsideTriangle(point a, point b, point c, point p) {
	return triangle(a,b,c).isInside(p);
} //0 = inside/ 1 = border/ 2 = outside

/*
 * Polygon 2D
 */

typedef vector<point> polygon;

int leftmostIndex(vector<point> &P) {
	int ans = 0;
	for(int i=1; i<(int)P.size(); i++) {
		if (P[i] < P[ans]) ans = i;
	}
	return ans;
}

/*
 * Convex Hull
 */

point pivot(0, 0);

bool angleCmp(point a, point b) {
	if (collinear(pivot, a, b))
		return inner(pivot-a, pivot-a) < inner(pivot-b, pivot-b);
	return cross(a-pivot, b-pivot) >= 0;
}

polygon convexHull(vector<point> P) {
	int i, j, n = (int)P.size();
	if (n <= 2) return P;
	int P0 = leftmostIndex(P);
	swap(P[0], P[P0]);
	pivot = P[0];
	sort(++P.begin(), P.end(), angleCmp);
	vector<point> S;
	S.push_back(P[n-1]);
	S.push_back(P[0]);
	S.push_back(P[1]);
	for(i = 2; i < n;) {
		j = (int)S.size()-1;
		if (ccw(S[j-1], S[j], P[i])) S.push_back(P[i++]);
		else S.pop_back();
	}
	reverse(S.begin(), S.end());
	S.pop_back();
	reverse(S.begin(), S.end());
	return S;
}

/*
 * Convex query
 */

bool query(polygon &P, point q) {
	int i = 1, j = P.size()-1, m;
	if (cross(P[i]-P[0], P[j]-P[0]) < -EPS)
		swap(i, j);
	while(abs(j-i) > 1) {
		int m = (i+j)/2;
		if (cross(P[m]-P[0], q-P[0]) < 0) j = m;
		else i = m;
	}
	return isInsideTriangle(P[0], P[i], P[j], q) != 2;
}

/*
 * Codeforces 101128J
 */

int main() {
	int nL, nS, nC;
	vector<point> L, S;
	scanf("%d", &nL);
	L.resize(nL);
	for(int i=0; i<nL; i++) {
		scanf("%lf %lf", &L[i].x, &L[i].y);
	}
	scanf("%d", &nS);
	S.resize(nS);
	for(int i=0; i<nS; i++) {
		scanf("%lf %lf", &S[i].x, &S[i].y);
	}
	L = convexHull(L);
	int ans = 0;
	for(int i=0; i<nS; i++) {
		if (query(L, S[i])) {
			//printf("(%f,%f) inside\n", S[i].x, S[i].y);
			ans++;
		}
	}
	printf("%d\n", ans);
	return 0;
}
\end{vncode}
\endgroup


\subsection{Tổng Minkowski}

Minkowski sum.

\codepath{Computational Geometry/minkowski.cpp}

\begingroup\scriptsize
\begin{vncode}
#include <cmath>
#define EPS 1e-9

/*
 * Point 2D
 */

struct point {
	double x, y;
	point() { x = y = 0.0; }
	point(double _x, double _y) : x(_x), y(_y) {}
	bool operator < (point other) const {
		if (fabs(x - other.x) > EPS) return x < other.x;
		else return y < other.y;
	}
	point operator +(point other) const {
		return point(x + other.x, y + other.y);
	}
	point operator -(point other) const {
		return point(x - other.x, y - other.y);
	}
	point operator *(double k) const {
		return point(x*k, y*k);
	}
};

double dist(point p1, point p2) {
	return hypot(p1.x - p2.x, p1.y - p2.y);
}

double inner(point p1, point p2) {
	return p1.x*p2.x + p1.y*p2.y;
}

double cross(point p1, point p2) {
	return p1.x*p2.y - p1.y*p2.x;
}

bool collinear(point p, point q, point r) {
	return fabs(cross(p-q, r-p)) < EPS;
}

/*
 * Polygon 2D
 */

#include <vector>
#include <algorithm>
using namespace std;

typedef vector<point> polygon;

double signedArea(polygon & P) {
	double result = 0.0;
	int n = P.size();
	for (int i = 0; i < n; i++) {
		result += cross(P[i], P[(i+1)%n]);
	}
	return result / 2.0;
}

int leftmostIndex(vector<point> & P) {
	int ans = 0;
	for(int i=1; i<(int)P.size(); i++) {
		if (P[i] < P[ans]) ans = i;
	}
	return ans;
}

polygon make_polygon(vector<point> P) {
	if (signedArea(P) < 0.0) reverse(P.begin(), P.end());
	int li = leftmostIndex(P);
	reverse(P.begin(), P.begin()+li);
	reverse(P.begin()+li, P.end());
	reverse(P.begin(), P.end());
	return P;
}

/*
 * Minkowski sum
 */

polygon minkowski(polygon & A, polygon & B) {
	polygon P; point v1, v2;
	int n1 = A.size(), n2 = B.size();
	P.push_back(A[0]+B[0]);
	for(int i = 0, j = 0; i < n1 || j < n2;) {
		v1 = A[(i+1)%n1]-A[i%n1];
		v2 = B[(j+1)%n2]-B[j%n2];
		if (j == n2 || cross(v1, v2) > EPS) {
			P.push_back(P.back() + v1); i++;
		}
		else if (i == n1 || cross(v1, v2) < -EPS) {
			P.push_back(P.back() + v2); j++;
		}
		else {
			P.push_back(P.back() + (v1+v2));
			i++; j++;
		}
	}
	P.pop_back();
	return P;
}

/*
 * Triangle 2D
 */

struct triangle {
	point a, b, c;
	triangle() { a = b = c = point(); }
	triangle(point _a, point _b, point _c) : a(_a), b(_b), c(_c) {}
	int isInside(point p) {
		double u = cross(b-a,p-a)*cross(b-a,c-a);
		double v = cross(c-b,p-b)*cross(c-b,a-b);
		double w = cross(a-c,p-c)*cross(a-c,b-c);
		if (u > 0.0 && v > 0.0 && w > 0.0) return 0;
		if (u < 0.0 || v < 0.0 || w < 0.0) return 2;
		else return 1;
	} //0 = inside/ 1 = border/ 2 = outside
};

int isInsideTriangle(point a, point b, point c, point p) {
	return triangle(a,b,c).isInside(p);
} //0 = inside/ 1 = border/ 2 = outside

/*
 * Convex query
 */

bool query(polygon &P, point q) {
	int i = 1, j = P.size()-1, m;
	if (cross(P[i]-P[0], P[j]-P[0]) < -EPS)
		swap(i, j);
	while(abs(j-i) > 1) {
		int m = (i+j)/2;
		if (cross(P[m]-P[0], q-P[0]) < 0) j = m;
		else i = m;
	}
	return isInsideTriangle(P[0], P[i], P[j], q) != 2;
}

/*
 * Codeforces 87E
 */
#include <cstdio>

void printpolygon(polygon & P) {
	printf("printing polygon:\n");
	for(int i=0; i<(int)P.size(); i++) {
		printf("%.2f %.2f\n", P[i].x, P[i].y);
	}
}

polygon city[3], P;

int main() {
	double x, y;
	for(int i = 0, n; i < 3; i++) {
		scanf("%d", &n);
		P.clear();
		while(n --> 0) {
			scanf("%lf %lf", &x, &y);
			P.push_back(point(x, y));
		}
		city[i] = make_polygon(P);
	}
	P = minkowski(city[0], city[1]);
	P = minkowski(P, city[2]);
	
	int m;
	scanf("%d", &m);
	while(m --> 0) {
		scanf("%lf %lf", &x, &y);
		if (query(P, point(x, y)*3.0)) printf("YES\n");
		else printf("NO\n");
	}
	return 0;
}
\end{vncode}
\endgroup


\subsection{Comparator cực}

So sánh theo góc cực.

\codepath{Computational Geometry/line2d.cpp}

\begingroup\scriptsize
\begin{vncode}
struct line {
	double a, b, c;
	line() { a = b = c = NAN;}
	line(double _a, double _b, double _c) : a(_a), b(_b), c(_c) {}
};

line pointsToLine(point p1, point p2) {
	line l;
	if (fabs(p1.x - p2.x) < EPS && fabs(p1.y - p2.y) < EPS) {
		l.a = l.b = l.c = NAN;
	}
	else if (fabs(p1.x - p2.x) < EPS) {
		l.a = 1.0; l.b = 0.0; l.c = -p1.x;
	}
	else {
		l.a = -(p1.y - p2.y) / (p1.x - p2.x);
		l.b = 1.0;
		l.c = -(l.a * p1.x) - p1.y;
	}
	return l;
}

bool areParallel(line l1, line l2) {
	return (fabs(l1.a-l2.a) < EPS) && (fabs(l1.b-l2.b) < EPS);
}

bool areSame(line l1, line l2) {
	return areParallel(l1 ,l2) && (fabs(l1.c - l2.c) < EPS);
}

point intersection(line l1, line l2) {
	if (areParallel(l1, l2)) return point(NAN, NAN);
	point p;
	p.x = (l2.b * l1.c - l1.b * l2.c) / (l2.a * l1.b - l1.a * l2.b);
	if (fabs(l1.b) > EPS) p.y = -(l1.a * p.x + l1.c);
	else p.y = -(l2.a * p.x + l2.c);
	return p;
}

point projPointToLine(point u, line l) {
	point a, b;
	if (fabs(l.b-1.0)<EPS) {
		a = point(-l.c/l.a, 0.0);
		b = point(-l.c/l.a, 1.0);
	}
	else{
		a = point(0, -l.c/l.b);
		b = point(1, -(l.c+1.0)/l.b);
	}
	return a + proj(u-a, b-a);
}
\end{vncode}
\endgroup


\subsection{Tam giác hóa Delaunay}

Delaunay triangulation.

\codepath{Computational Geometry/delaunay.cpp}

\begingroup\scriptsize
\begin{vncode}
#include <cstdio>
#include <cmath>
#include <vector>
using namespace std;
#define EPS 1e-9

/*
 * Point 2D
 */

struct point {
	double x, y;
	point() { x = y = 0.0; }
	point(double _x, double _y) : x(_x), y(_y) {}
	bool operator < (point other) const {
		if (fabs(x - other.x) > EPS) return x < other.x;
		else return y < other.y;
	}
	bool operator == (point other) const {
		return (fabs(x - other.x) < EPS && (fabs(y - other.y) < EPS));
	}
	point operator +(point other) const {
		return point(x + other.x, y + other.y);
	}
	point operator -(point other) const {
		return point(x - other.x, y - other.y);
	}
	point operator *(double k) const {
		return point(x*k, y*k);
	}
};

double dist(point p1, point p2) {
	return hypot(p1.x - p2.x, p1.y - p2.y);
}

double inner(point p1, point p2) {
	return p1.x*p2.x + p1.y*p2.y;
}

double cross(point p1, point p2) {
	return p1.x*p2.y - p1.y*p2.x;
}

bool collinear(point p, point q, point r) {
	return fabs(cross(p-q, r-p)) < EPS;
}

/*
 * Circle 2D
 */

struct circle {
	point c;
	double r;
	circle() { c = point(); r = 0; }
};

circle circumcircle(point a, point b, point c) {
	circle ans;
	point u = point((b-a).y, -(b-a).x);
	point v = point((c-a).y, -(c-a).x);
	point n = (c-b)*0.5;
	double t = cross(u,n)/cross(v,u);
	ans.c = ((a+c)*0.5) + (v*t);
	ans.r = dist(ans.c, a);
	return ans;
}

/*
 * Delaunay Triangulation
 * http://graphics.stanford.edu/courses/cs368-06-spring/handouts/Delaunay_2.pdf
 */

#include <set>
#include <algorithm>
#define MAXN 200309

struct truple {
	int a, b, c;
	truple(int _a, int _b, int _c) :
		a(_a), b(_b), c(_c) {
		if (a > b) swap(a, b);
		if (a > c) swap(a, c);
		if (b > c) swap(b, c);
	}
};
bool operator < (truple x, truple y) {
	if (x.a != y.a) return x.a < y.a;
	if (x.b != y.b) return x.b < y.b;
	if (x.c != y.c) return x.c < y.c;
	return false;
}

namespace Delaunay {
vector< set<int> > adj;
vector<point> P;
set<truple> tri;
void add_edge(int u, int v, int w = -1) {
	adj[u].insert(v); adj[v].insert(u);
	if (w >= 0) tri.insert(truple(u, v, w));
}
void del_edge(int u, int v) {
	adj[u].erase(v); adj[v].erase(u);
}
void brute(int l, int r) {	
	if (r-l == 2 && collinear(P[l],P[l+1],P[r])) {
		add_edge(l, l+1); add_edge(l+1, r);
		return;
	}
	for(int u = l; u <= r; u++)
		for(int v = u+1; v <= r; v++)
			add_edge(u, v);
	if (r-l == 2) tri.insert(truple(l,l+1,r));
}
double theta[MAXN];
bool comp(int u, int v) {
	return theta[u] < theta[v];
}
vector< vector<int> > g;
bool right; point base;
void computeG(int u, int r = -1) {
	if (!g[u].empty()) return;
	if (r >= 0) adj[u].erase(r);
	g[u] = vector<int>(adj[u].begin(), adj[u].end());
	if (r >= 0) adj[u].insert(r);
	for(int i = 0; i < int(g[u].size()); i++) {
		double co = inner(base, P[g[u][i]]-P[u]);
		double si = cross(base, P[g[u][i]]-P[u]);
		theta[g[u][i]] = atan2(si, co);
	}
	sort(g[u].begin(), g[u].end(), comp);
	if (right) reverse(g[u].begin(), g[u].end());
}
int getNext(int u, int a, int b) {
	right = (u == b);
	base = (right ? P[b]-P[a] : P[a]-P[b]);
	computeG(u, a + b - u);
	int ans = -1, w;
	while(!g[u].empty()) {
		int j = g[u].size() - 1;
		ans = g[u][j];
		if (right && cross(base, P[ans]-P[b]) > -EPS) return -1;
		if (!right && cross(base, P[ans]-P[a]) < EPS) return -1;
		circle c = circumcircle(P[a],P[b],P[ans]);
		if (g[u].size() > 1u && dist(c.c, P[w = g[u][j-1]]) < c.r - EPS) {
			del_edge(u, ans); g[u].pop_back();
			tri.erase(truple(u, w, ans));
		}
		else break;
	}
	return ans;
}
int moveEdge(int & u, int a, int b) {
	right = (b == u); int v = a+b-u;
	for(set<int>::iterator it = adj[u].begin(); it != adj[u].end(); it++) {
		int nu = *it;
		double cr = cross(P[u]-P[v], P[nu]-P[v]);
		if (right && cr > EPS) return u = nu;
		if (!right && cr < -EPS) return u = nu;
		if (fabs(cr) < EPS && dist(P[nu], P[v]) < dist(P[u], P[v])) return u = nu;
	}
	return -1;
}
void delaunay(int l, int r) {
	if (r - l < 3) { brute(l, r); return; }
	int mid = (r + l) / 2;
	delaunay(l, mid); delaunay(mid + 1, r);
	int u = l, v = r, nu, nv;
	double du, dv;
	do {
		nu = moveEdge(u, u, v);
		nv = moveEdge(v, u, v);
	} while (nu != -1 || nv != -1);
	g[u].clear(); g[v].clear();
	add_edge(u, v);
	while(true) {
		nu = getNext(u, u, v); nv = getNext(v, u, v);
		if (nu == -1 && nv == -1) break;
		if (nu != -1 && nv != -1) {
			point nor = point(P[u].y-P[v].y, P[v].x-P[u].x);
			circle cu = circumcircle(P[u],P[v],P[nu]);
			circle cv = circumcircle(P[u],P[v],P[nv]);
			du = inner(cu.c-P[u], nor);
			dv = inner(cv.c-P[v], nor);
		}
		else du = 1, dv = 0;
		if (nu == -1 || du < dv) {
			add_edge(u, nv, v); g[v = nv].clear();
		}
		else if (nv == -1 || du >= dv) {
			add_edge(nu, v, u); g[u = nu].clear();
		}
	}
}
vector< set<int> > compute(vector<point> & Q) {
	sort(Q.begin(), Q.end());
	int n = (P = Q).size();
	adj.clear(); adj.assign(n, set<int>());
	g.clear(); g.assign(n, vector<int>());
	tri.clear();
	delaunay(0, n-1);
	return adj;
}
set<truple> getTriangles() { return tri; }
};

/*
 * TEST MATRIX
 */

#include <cstdio>

int cnt = 0;
void plot2csv(vector<point> P, vector< set<int> > adj) {
	char filename[100];
	int n = P.size();
	sprintf(filename, "graph%d.csv", cnt++);
	FILE * out = fopen(filename, "w");
	fprintf(out, "id,lat,lng\n");
	for(int i = 0; i < n; i++) {
		fprintf(out, "%d,%.2f,%.2f\n", i, P[i].x, P[i].y);
	}
	fprintf(out, "source,target\n");
	for(int u = 0; u < n; u++) {
		for(set<int>::iterator it = adj[u].begin(); it != adj[u].end(); it++) {
			fprintf(out, "%d,%d\n", u, *it);
		}
	}
	fclose(out);
}

typedef pair<int, int> ii;

bool isAllCollinear(vector<point> P) {
	int n = P.size(), id = 0;
	while(id < n && P[0] == P[id]) id++;
	if (id == n) return true;
	for(int i = 0; i < n; i++) {
		if (!collinear(P[0], P[id], P[i])) return false;
	}
	return true;
}

bool between(point p, point q, point r) {
	return collinear(p, q, r) && inner(p - q, r - q) < -EPS;
}

point lineIntersectSeg(point p, point q, point A, point B) {
	double c = cross(A-B, p-q);
	double a = cross(A, B);
	double b = cross(p, q);
	return ((p-q)*(a/c)) - ((A-B)*(b/c));
}

bool parallel(point a, point b) {
	return fabs(cross(a, b)) < EPS;
}

bool segIntersects(point a, point b, point p, point q) {
	if (parallel(a-b, p-q)) {
		return between(a, p, b) || between(a, q, b)
			|| between(p, a, q) || between(p, b, q); 
	}
	point i = lineIntersectSeg(a, b, p, q);
	return between(a, i, b) && between(p, i, q);
}

bool isTriangulation(vector<point> P, vector< set<int> > adj) {
	vector< pair<point, point> > seg;
	set<ii> edges, aux;
	int n = P.size();
	for(int u = 0; u < n; u++) {
		vector<int> g(adj[u].begin(), adj[u].end());
		for(int j = 0; j < (int)g.size(); j++) {
			int v = g[j];
			if (u < v) {
				seg.push_back(make_pair(P[u], P[v]));
				edges.insert(make_pair(u, v));
			}
		}
	}
	set<truple> tri = Delaunay::getTriangles();
	for(set<truple>::iterator it = tri.begin(); it != tri.end(); it++) {
		int u = it->a, v = it->b, w = it->c;
		if (!edges.count(ii(u, v))) {
			printf("error, not a triangulation\n");
			return false;
		}
		else aux.insert(ii(u, v));
		if (!edges.count(ii(v, w))) {
			printf("error, not a triangulation\n");
			return false;
		}
		else aux.insert(ii(v, w));
		if (!edges.count(ii(u, w))) {
			printf("error, not a triangulation\n");
			return false;
		}
		else aux.insert(ii(u, w));
	}
	if (!isAllCollinear(P) && aux.size() != edges.size()) {
		printf("error, not a triangulation\n");
		return false;
	}
	for(int u = 0; u < n; u++) {
		for(int v = 0; v < n; v++) {
			if (u >= v || edges.count(ii(u, v))) continue;
			bool found = false;
			for(int j = 0; !found && j < (int)seg.size(); j++) {
				point p = seg[j].first;
				point q = seg[j].second;
				if (segIntersects(p, q, P[u], P[v])) found = true;
			}
			if (!found) {
				printf("error, not a triangulation\n");
				return false;
			}
		}
	}
	return true;
}

bool isDelaunayTriangulation(vector<point> P, vector< set<int> > adj) {
	int n = P.size();
	if (!isTriangulation(P, adj)) return false;
	set<truple> tri = Delaunay::getTriangles();
	for(set<truple>::iterator it = tri.begin(); it != tri.end(); it++) {
		int u = it->a, v = it->b, w = it->c;
		circle c = circumcircle(P[u], P[v], P[w]);
		for(int i = 0; i < n; i++) {
			if (dist(c.c, P[i]) < c.r - 1e-4) {
				printf("error, point %d = (%.3f,%.3f) is inside circumcircle:\n", i, P[i].x, P[i].y);
				printf("%d = (%.3f,%.3f)\n", u, P[u].x, P[u].y);
				printf("%d = (%.3f,%.3f)\n", v, P[v].x, P[v].y);
				printf("%d = (%.3f,%.3f)\n", w, P[w].x, P[w].y);
				printf("center (%.3f,%.3f), radius %.3f, distance %.3f\n", c.c.x, c.c.y, c.r, dist(c.c, P[i]));
				return false;
			}
		}
	}
	return true;
}

bool inputTest() {
	int n;
	if(scanf("%d", &n) == EOF) return false;
	if (n <= 0) return false;
	vector<point> P(n);
	for(int i = 0; i < n; i++) {
		scanf("%lf %lf", &P[i].x, &P[i].y);
	}
	vector< set<int> > adj = Delaunay::compute(P);
	if(!isDelaunayTriangulation(P, adj)) return false;
	printf("input test passed\n");
	return true;
}

#include <ctime>
typedef pair<int, int> ii;

bool agnezCornerCase(int n, bool testtri) {
	printf("test Victor Agnez's corner case, ");
	vector<point> P;
	for(int i = 0; i < n; i++) {
		P.push_back(point(i / 2, i % 2));
	}
	clock_t last = clock();
	vector< set<int> > adj = Delaunay::compute(P);
	double t = (clock() - last)*1.0/CLOCKS_PER_SEC;
	printf("time %.3fs ", t);
	if(testtri && !isDelaunayTriangulation(P, adj)) {
		printf("failed\n");
		return false;
	}
	printf("ok\n");
	return true;
}

bool test(int lowtest, int hightest, bool testtri) {
	agnezCornerCase(hightest, testtri);
	srand(time(NULL));
	for(int suit = lowtest; suit <= hightest; suit++) {
		printf("test %d, ", suit);
		int n = suit;
		set<ii> inuse;
		while(inuse.size() < n) {
			inuse.insert(ii(rand()%MAXN, rand()%MAXN));
		}
		vector<point> P;
		for(set<ii>::iterator it = inuse.begin(); it != inuse.end(); it++) {
			P.push_back(point(it->first, it->second));
		}
		clock_t last = clock();
		vector< set<int> > adj = Delaunay::compute(P);
		double t = (clock() - last)*1.0/CLOCKS_PER_SEC;
		printf("time %.3fs ", t);
		if(testtri && !isDelaunayTriangulation(P, adj)) {
			printf("failed test %d\n", suit);
			return false;
		}
		printf("ok\n");
		//plot2csv(P, adj);
	}
	if (testtri) printf("all tests passed\n");
	return true;
}

int main() {
	test(1, 100, true);
	return 0;
}

/*
 * http://wcipeg.com/problem/ccc08s2p6
 */
/*
int main() {
	int n;
	while(scanf("%d", &n) != EOF) {
		vector<point> P(n);
		for(int i = 0; i < n; i++) {
			scanf("%lf %lf", &P[i].x, &P[i].y);
		}
		Delaunay::compute(P);
		double ans = 0;
		set<truple> tri = Delaunay::getTriangles();
		for(set<truple>::iterator it = tri.begin(); it != tri.end(); it++) {
			int u = it->a, v = it->b, w = it->c;
			circle c = circumcircle(P[u], P[v], P[w]);
			if (c.r > ans) ans = c.r;
		}
		printf("%.8f\n", ans);
	}
	return 0;
}
*/
\end{vncode}
\endgroup


\subsection{Giao của đa giác}

Polygon intersection.

\subsection{Minimum Enclosing Circle}

Đường tròn bao nhỏ nhất.

\subsection{Giao nửa mặt phẳng}

Half-plane intersection.

\codepath{Computational Geometry/halfplaneintersect.cpp}

\begingroup\scriptsize
\begin{vncode}
#include <cstdio>
#include <cmath>
#include <vector>
#include <deque>
#include <algorithm>
#define EPS 1e-9
#define INF 0x3f3f3f3f
using namespace std;

/*
 * Point 2D
 */

struct point {
	double x, y;
	point() { x = y = 0.0; }
	point(double _x, double _y) : x(_x), y(_y) {}
	double norm() {
		return hypot(x, y);
	}
	point normalized() {
		return point(x,y)*(1.0/norm());
	}
	double angle() { return atan2(y, x); }
	double polarAngle() {
		double a = atan2(y, x);
		return a < 0 ? a + 2*acos(-1.0) : a;
	}
	bool operator < (point other) const {
		if (fabs(x - other.x) > EPS) return x < other.x;
		else return y < other.y;
	}
	bool operator == (point other) const {
		return (fabs(x - other.x) < EPS && (fabs(y - other.y) < EPS));
	}
	point operator +(point other) const {
		return point(x + other.x, y + other.y);
	}
	point operator -(point other) const {
		return point(x - other.x, y - other.y);
	}
	point operator *(double k) const {
		return point(x*k, y*k);
	}
};

double inner(point p1, point p2) {
	return p1.x*p2.x + p1.y*p2.y;
}

double cross(point p1, point p2) {
	return p1.x*p2.y - p1.y*p2.x;
}

point rotate(point p, double rad) {
	return point(p.x * cos(rad) - p.y * sin(rad),
	p.x * sin(rad) + p.y * cos(rad));
}

point lineIntersectSeg(point p, point q, point A, point B) {
	double c = cross(A-B, p-q);
	double a = cross(A, B);
	double b = cross(p, q);
	return ((p-q)*(a/c)) - ((A-B)*(b/c));
}

bool parallel(point a, point b) {
	return fabs(cross(a, b)) < EPS;
}

/*
 * Polygon 2D
 */

typedef vector<point> polygon;

double signedArea(polygon &P) {
	double result = 0.0;
	int n = P.size();
	for (int i = 0; i < n; i++) {
		result += cross(P[i], P[(i+1)%n]);
	}
	return result / 2.0;
}

double area(polygon &P) {
	return fabs(signedArea(P));
}

int leftmostIndex(vector<point> &P) {
	int ans = 0;
	for(int i=1; i<(int)P.size(); i++) {
		if (P[i] < P[ans]) ans = i;
	}
	return ans;
}

polygon make_polygon(vector<point> P) {
	if (signedArea(P) < 0.0) reverse(P.begin(), P.end());
	int li = leftmostIndex(P);
	rotate(P.begin(), P.begin()+li, P.end());
	return P;
}

polygon cutPolygon(polygon &P, point a, point b) {
	vector<point> R;
	double left1, left2;
	int n = P.size();
	for (int i = 0; i < n; i++) {
		left1 = cross(b-a, P[i]-a);
		left2 = cross(b-a, P[(i+1)%n]-a);
		if (left1 > -EPS) R.push_back(P[i]);
		if (left1 * left2 < -EPS) 
			R.push_back(lineIntersectSeg(P[i], P[(i+1)%n], a, b));
	}
	return make_polygon(R);
}

/*
 * Halfplane 2D
 */

typedef pair<point, point> halfplane;

point dir(halfplane h) {
	return h.second - h.first;
}

bool belongs(halfplane h, point a) {
	return cross(dir(h), a - h.first) > EPS;
}

bool hpcomp(halfplane ha, halfplane hb) {
	point a = dir(ha), b = dir(hb);
	if (parallel(a, b) && inner(a, b) > EPS)
		return cross(b, ha.first - hb.first) < -EPS;
	if (b.y*a.y > EPS) return cross(a, b) > EPS;
	else if (fabs(b.y) < EPS && b.x > EPS) return false;
	else if (fabs(a.y) < EPS && a.x > EPS) return true;
	else return b.y < a.y;
}

polygon intersect(vector<halfplane> H, double W = INF) {
	H.push_back(halfplane(point(-W,-W), point(W,-W)));
	H.push_back(halfplane(point(W,-W), point(W,W)));
	H.push_back(halfplane(point(W,W), point(-W,W)));
	H.push_back(halfplane(point(-W,W), point(-W,-W)));
	sort(H.begin(), H.end(), hpcomp);
	int i = 0;
	while(parallel(dir(H[0]), dir(H[i]))) i++;
	deque<point> P;
	deque<halfplane> S;
	S.push_back(H[i-1]);
	for(; i < (int)H.size(); i++) {
		while(!P.empty() && !belongs(H[i], P.back()))
			P.pop_back(), S.pop_back();
		point df = dir(S.front()), di = dir(H[i]);
		if (P.empty() && cross(df, di) < EPS)
			return polygon();
		P.push_back(lineIntersectSeg(H[i].first, H[i].second,
			S.back().first, S.back().second));
		S.push_back(H[i]);
	}
	while(!belongs(S.back(), P.front()) ||
		!belongs(S.front(), P.back())) {
		while(!belongs(S.back(), P.front()))
			P.pop_front(), S.pop_front();
		while(!belongs(S.front(), P.back()))
			P.pop_back(), S.pop_back();
	}
	P.push_back(lineIntersectSeg(S.front().first,
		S.front().second, S.back().first, S.back().second));
	return polygon(P.begin(), P.end());
}

/*
 * TEST MATRIX
 */

point rndpoint(int W) {
	int x = (rand()%(2*W)) - W;
	int y = (rand()%(2*W)) - W;
	return point(x, y);
}

/*bool test(int ntests, int n, int W, bool nonempty) {
	for(int test = 1; test <= ntests; test++) {
		vector<halfplane> H;
		for(int i = 0; i < n; i++) {
			halfplane h;
			do {
				h.first = rndpoint(W);
				h.second = rndpoint(W);
			} while(nonempty && !belongs(h, point(0, 0)));
			H.push_back(h);
		}
		polygon A1 = intersect(H, W);
		polygon A2;
		A2.push_back(point(-W,-W));
		A2.push_back(point(W,-W));
		A2.push_back(point(W,W));
		A2.push_back(point(-W,W));
		for(int i = 0; i < n; i++) {
			A2 = cutPolygon(A2, H[i].first, H[i].second);
		}
		if (A1.size() != A2.size() || fabs(area(A1) - area(A2)) > EPS) {
			printf("failed on test %d:\n", test);
			printf("H:\n");
			for(int i = 0; i < n; i++)
				printf(" (%.0f,%.0f)-(%.0f,%.0f)\n", H[i].first.x, H[i].first.y, H[i].second.x, H[i].second.y);
			printf("A1 (%u):\n", A1.size());
			for(int i = 0; i < (int)A1.size(); i++) {
				printf(" (%.3f,%.3f)\n", A1[i].x, A1[i].y);
			}
			printf("A2 (%u):\n", A2.size());
			for(int i = 0; i < (int)A2.size(); i++) {
				printf(" (%.3f,%.3f)\n", A2[i].x, A2[i].y);
			}
			return false;
		}
	}
	printf("all tests passed\n");
	return true;
}

int main() {
	test(10000, 100, 100, false);
	test(10000, 100, 100, true);
}*/

/*
 * ICPC Summer School 2019, Day 3, problem D
 */

/*point read() {
	double x, y;
	scanf("%lf %lf", &x, &y);
	return point(x, y);
}

int main() {
	int n;
	scanf("%d", &n);
	vector<halfplane> H;
	for(int i = 0; i < n; i++) {
		point a = read();
		point b = read();
		point c = read();
		H.push_back(halfplane(a, b));
		H.push_back(halfplane(b, c));
		H.push_back(halfplane(c, a));
	}
	polygon ans = intersect(H, 1000);
	printf("%.20f\n", area(ans));
	return 0;
}*/

/*
 * ICPC Summer School 2019, Day 2, problem K
 */

bool check(const polygon &I, const polygon &O, double r) {
	int n = O.size(), m = I.size();
	vector<halfplane> H;
	for(int i = 0; i < n; i++) {
		point v = (O[(i+1)%n] - O[i]).normalized();
		point ref = O[i] + rotate(v, acos(-1.0)/2.0)*r;
		double trans = -INF;
		for(int j = 0; j < m; j++) {
			double cur = cross(v, I[j]- ref);
			trans = max(trans, -cur);
		}
		point a = rotate(v, acos(-1.0)/2.0)*trans;
		point b = a + v;
		H.push_back(halfplane(a, b));
	}
	return !intersect(H, 10000).empty();
}

point read() {
	double x, y;
	scanf("%lf %lf", &x, &y);
	return point(x, y);
}

int main() {
	int n;
	polygon O, I;
	scanf("%d", &n);
	for(int i = 0; i < n; i++) I.push_back(read());
	scanf("%d", &n);
	for(int i = 0; i < n; i++) O.push_back(read());
	if (area(O) < area(I)) I.swap(O);
	double lo = 0;
	double hi = 1000;
	for(int it = 0; it < 100; it++) {
		double mid = (hi + lo) / 2;
		if (check(I, O, mid)) lo = mid;
		else hi = mid;
	}
	printf("%.20f\n", lo/2);
	return 0;
}
\end{vncode}
\endgroup


\subsection{Điểm 3D}

Point/vector 3D.

\codepath{Computational Geometry/point3d.cpp}

\begingroup\scriptsize
\begin{vncode}
#include <cstdio>
#include <cmath>
#define EPS 1e-9

struct point {
	double x, y, z;
	point() { x = y = z = 0.0; }
	point(double _x, double _y, double _z) : x(_x), y(_y), z(_z) {}
	double norm() {
		return sqrt(x*x + y*y + z*z);
	}
	point normalized() {
		return point(x,y,z)*(1.0/norm());
	}
	bool operator < (point other) const {
		if (fabs(x - other.x) > EPS) return x < other.x;
		else if (fabs(y - other.y) > EPS) return y < other.y;
		else return z < other.z;
	}
	bool operator == (point other) const {
		return (fabs(x - other.x) < EPS && fabs(y - other.y) < EPS && fabs(z - other.z) < EPS);
	}
	point operator +(point other) const{
		return point(x + other.x, y + other.y, z + other.z);
	}
	point operator -(point other) const{
		return point(x - other.x, y - other.y, z - other.z);
	}
	point operator *(double k) const{
		return point(x*k, y*k, z*k);
	}
};

double dist(point p1, point p2) {
	return (p1-p2).norm();
}

double inner(point p1, point p2) {
	return p1.x*p2.x + p1.y*p2.y + p1.z*p2.z;
}

point cross(point p1, point p2) {
	point ans;
	ans.x = p1.y*p2.z - p1.z*p2.y;
	ans.y = p1.z*p2.x - p1.x*p2.z;
	ans.z = p1.x*p2.y - p1.y*p2.x;
	return ans;
}

bool collinear(point p, point q, point r) {
	return cross(p-q, r-p).norm() < EPS;
}

double angle(point a, point o, point b) {
	return acos(inner(a-o, b-o) / (dist(o,a)*dist(o,b)));
}

point proj(point u, point v) {
	return v*(inner(u,v)/inner(v,v));
}
\end{vncode}
\endgroup


\subsection{Tam giác 3D}

Triangle 3D.

\codepath{Computational Geometry/triangle3d.cpp}

\begingroup\scriptsize
\begin{vncode}
struct triangle{
	point a, b, c;
	triangle() { a = b = c = point(); }
	triangle(point _a, point _b, point _c) : a(_a), b(_b), c(_c) {}
	double perimeter() { return dist(a,b) + dist(b,c) + dist(c,a); }
	double semiPerimeter() { return perimeter()/2.0; }
	double area() {
		double s = semiPerimeter(), ab = dist(a,b),
			bc = dist(b,c), ca = dist(c,a);
		return sqrt(s*(s-ab)*(s-bc)*(s-ca));
	}
	double rInCircle() {
		return area()/semiPerimeter();
	}
	double rCircumCircle() {
		return dist(a,b)*dist(b,c)*dist(c,a)/(4.0*area());
	}
	point normalVector() {
		return cross(y-x, z-x).normalized();
	}
	int isInside(point p) {
		point n = normalVector();
		double u = proj(cross(b-a,p-a), n).normalized()*proj(cross(b-a,c-a), n).normalized();
		double v = proj(cross(c-b,p-b), n).normalized()*proj(cross(c-b,a-b), n).normalized();
		double w = proj(cross(a-c,p-c), n).normalized()*proj(cross(a-c,b-c), n).normalized();
		if (u > 0.0 && v > 0.0 && w > 0.0) return 0;
		else if (u < 0.0 || v < 0.0 || w < 0.0) return 2;
		else return 1;
	} //0 = inside/ 1 = border/ 2 = outside
	int isProjInside(point p) {
		return isInside(p + proj(a-p, normalVector()));
	} //0 = inside/ 1 = border/ 2 = outside
};

double rInCircle(point a, point b, point c) {
	return triangle(a,b,c).rInCircle();
}

double rCircumCircle(point a, point b, point c) {
	return triangle(a,b,c).rCircumCircle();
}

int isProjInsideTriangle(point a, point b, point c, point p) {
	return triangle(a,b,c).isProjInside(p);
} //0 = inside/ 1 = border/ 2 = outside
\end{vncode}
\endgroup


\subsection{Đường thẳng 3D}

Line 3D.

\codepath{Computational Geometry/line3d.cpp}

\begingroup\scriptsize
\begin{vncode}
#include <cstdio>
#include <cmath>
#define EPS 1e-9

struct point {
	double x, y, z;
	point() { x = y = z = 0.0; }
	point(double _x, double _y, double _z) : x(_x), y(_y), z(_z) {}
	double norm() {
		return hypot(x, y, z);
	}
	point normalized() {
		return point(x, y, z)*(1.0 / norm());
	}
	bool operator < (point other) const {
		if (fabs(x - other.x) > EPS) return x < other.x;
		else if (fabs(y - other.y) > EPS) return y < other.y;
		else return z < other.z;
	}
	bool operator == (point other) const {
		return (fabs(x - other.x) < EPS && fabs(y - other.y) < EPS && fabs(z - other.z) < EPS);
	}
	point operator +(point other) const {
		return point(x + other.x, y + other.y, z + other.z);
	}
	point operator -(point other) const {
		return point(x - other.x, y - other.y, z - other.z);
	}
	point operator *(double k) const {
		return point(x*k, y*k, z*k);
	}
};

double dist(point p1, point p2) {
	return (p1 - p2).norm();
}

double inner(point p1, point p2) {
	return p1.x*p2.x + p1.y*p2.y + p1.z*p2.z;
}

point cross(point p1, point p2) {
	point ans;
	ans.x = p1.y*p2.z - p1.z*p2.y;
	ans.y = p1.z*p2.x - p1.x*p2.z;
	ans.z = p1.x*p2.y - p1.y*p2.x;
	return ans;
}

bool collinear(point p, point q, point r) {
	return cross(p - q, r - p).norm() < EPS;
}

double angle(point a, point o, point b) {
	return acos(inner(a - o, b - o) / (dist(o, a)*dist(o, b)));
}

point proj(point u, point v) {
	return v*(inner(u, v) / inner(v, v));
}

struct line {
	point r;
	point v;
	line(point _r, point _v) {
		v = _v; r = _r;
	}
	bool operator == (line other) const {
		return fabs(cross(r - other.r, v).norm()) < EPS && fabs(cross(r - other.r, other.v).norm()) < EPS;
	}
};

point distVector(line l, point p) {
	point dr = p - l.r;
	return dr - proj(dr, l.v);
}

point distVectorBasePoint(line l, point p) {
	return proj(p - l.r, l.v) + l.r;
}

point distVectorEndPoint(line l, point p) {
	return p;
}

point distVector(line a, line b) {
	point dr = b.r - a.r;
	point n = cross(a.v, b.v);
	if (n.norm() < EPS) {
		return dr - proj(dr, a.v);
	}
	else return proj(dr, n);
}

double dist(line a, line b) {
	return distVector(a, b).norm();
}

point distVectorBasePoint(line a, line b) {
	if (cross(a.v, b.v).norm() < EPS) return a.r;
	point d = distVector(a, b);
	double lambda;
	if (fabs(b.v.x*a.v.y - a.v.x*b.v.y) > EPS)
		lambda = (b.v.x*(b.r.y - a.r.y - d.y) - b.v.y*(b.r.x - a.r.x - d.x)) / (b.v.x*a.v.y - a.v.x*b.v.y);
	else if (fabs(b.v.x*a.v.z - a.v.x*b.v.z) > EPS)
		lambda = (b.v.x*(b.r.z - a.r.z - d.z) - b.v.z*(b.r.x - a.r.x - d.x)) / (b.v.x*a.v.z - a.v.x*b.v.z);
	else if (fabs(b.v.z*a.v.y - a.v.z*b.v.y) > EPS)
		lambda = (b.v.z*(b.r.y - a.r.y - d.y) - b.v.y*(b.r.z - a.r.z - d.z)) / (b.v.z*a.v.y - a.v.z*b.v.y);
	return a.r + (a.v*lambda);
}

point distVectorEndPoint(line a, line b) {
	return distVectorBasePoint(a, b) + distVector(a, b);
}
\end{vncode}
\endgroup


\subsection{Hình học giải tích}

Công thức đường thẳng, mặt phẳng, khoảng cách và góc trong hình học giải tích.

\subsection{Tọa độ cực, trụ và cầu}

Đổi hệ tọa độ cực/trụ/cầu.

\subsection{Giải tích vector 2D}

Gradient, divergence, curl và các công thức vector 2D.

\subsection{Giải tích vector 3D}

Công thức vector 3D.

\subsection{Vấn đề độ chính xác, tổng ổn định và công thức Bhaskara}

Tránh trừ hai số gần nhau, dùng tổng ổn định và nghiệm phương trình bậc hai ổn định.

\codepath{Computational Geometry/precision.cpp}

\begingroup\scriptsize
\begin{vncode}
#include <cmath>
#include <algorithm>
using namespace std;
#define EPS 1e-9

/*
 * Auxiliary geometry objects
 */

struct point {
	double x, y;
	point() { x = y = 0.0; }
	point(double _x, double _y) : x(_x), y(_y) {}
	double norm() {
		return hypot(x, y);
	}
	point normalized() {
		return point(x,y)*(1.0/norm());
	}
	double angle() { return atan2(y, x);	}
	bool operator < (point other) const {
		if (fabs(x - other.x) > EPS) return x < other.x;
		else return y < other.y;
	}
	bool operator == (point other) const {
		return (fabs(x - other.x) < EPS && (fabs(y - other.y) < EPS));
	}
	point operator +(point other) const{
		return point(x + other.x, y + other.y);
	}
	point operator -(point other) const{
		return point(x - other.x, y - other.y);
	}
	point operator *(double k) const{
		return point(x*k, y*k);
	}
};

double dist(point p1, point p2) {
	return hypot(p1.x - p2.x, p1.y - p2.y);
}

double cross(point p1, point p2) {
	return p1.x*p2.y - p1.y*p2.x;
}

double inner(point p1, point p2) {
	return p1.x*p2.x + p1.y*p2.y;
}

point rotate(point p, double rad) {
	return point(p.x * cos(rad) - p.y * sin(rad),
	p.x * sin(rad) + p.y * cos(rad));
}

struct circle{
	point c;
	double r;
	circle() { c = point(); r = 0; }
	circle(point _c, double _r) : c(_c), r(_r) {}
	double chord(double rad) { return  2*r*sin(rad/2.0); }
	bool intersects(circle other) {
		return dist(c, other.c) < r + other.r;
	}
	bool contains(point p) { return dist(c, p) <= r + EPS; }
	pair<point, point> getTangentPoint(point p) {
		double d1 = dist(p, c), theta = asin(r/d1);
		point p1 = rotate(c-p,-theta);
		point p2 = rotate(c-p,theta);
		p1 = p1*(sqrt(d1*d1-r*r)/d1)+p;
		p2 = p2*(sqrt(d1*d1-r*r)/d1)+p;
		return make_pair(p1,p2);
	}
};

/*
 * Precision algorithms
 */
#include <vector>

struct StableSum {
	int cnt = 0;
	vector<double> v, pref(1, 0);
	void operator += (double a) { // a >= 0
		for (int s = ++cnt; s % 2 == 0; s >>= 1) {
			a += v.back();
			v.pop_back(), pref.pop_back();
		}
		v.push_back(a);
		pref.push_back(pref.back() + a);
	}
	double val() { return pref.back(); }
};

bool circleCircle(circle c1, circle c2, pair<point, point> & out) {
	double d = dist(c2.c, c1.c);
	double co = (d*d + c1.r*c1.r - c2.r*c2.r)/(2*d*c1.r);
	if (fabs(co) > 1.0) return false;
	double alpha = acos(co);
	point rad = (c2.c-c1.c)*(1.0/d*c1.r);
	out = {c1.c + rotate(rad, -alpha), c1.c + rotate(rad, alpha)};
	return true;
}

int quadRoots(double a, double b, double c, pair<double, double> &out) {
	if (fabs(a) < EPS) return 0;
	double delta = b*b - 4*a*c;
	if (delta < 0) return 0;
	double sum = (b >= 0) ? -b-sqrt(delta) : -b+sqrt(delta);
	out = {sum/(2*a), fabs(sum) > EPS ? 0 : (2*c)/sum};
	return 1 + (delta > 0);
}


/*
 * TEST MATRIX
 */
#include <cstdio>

int main() {
	double a, b, c;
	while(scanf("%lf %lf %lf", &a, &b, &c) != EOF) {
		pair<double, double> ans;
		int d = quadRoots(a, b, c, ans);
		printf("%d solutions for %.3fx^2 + %.3fx + %.3f:", d, a, b, c);
		if (d > 0) printf(" x = %.3f", ans.first);
		if (d > 1) printf(" x = %.3f", ans.second);
		printf("\n");
	}
	return 0;
}
\end{vncode}
\endgroup


\clearpage

\section{Linh tinh}

\subsection{Thuật toán Mo}

Mo's algorithm.

\codepath{Miscelaneous/mo.cpp}

\begingroup\scriptsize
\begin{vncode}
/*
 * Mo's Algorithm O(n^(3/2))
 */

#include <algorithm>
using namespace std;
#define SQ 500

int curans; // current ans
void add(int v) { /* add value */ }
void rem(int v) { /* remove value */ }

struct query {
	int i, l, r, ans;
};

bool qcomp(query a, query b) {
	return a.l/SQ == b.l/SQ ? a.r < b.r : a.l < b.l;
}
bool icomp(query a, query b) { return a.i < b.i; }

void mo(int v[], query qs[], int Q) {
	int l = 0, r = -1;
	sort(qs, qs + Q, qcomp);
	for (int i = 0; i < Q; i++) {
		query & q = qs[i];
		while (r < q.r) add(v[++r]);
		while (r > q.r) rem(v[r--]);
		while (l < q.l) rem(v[l++]);
		while (l > q.l) add(v[--l]);
		q.ans = curans;
	}
	sort(qs, qs + Q, icomp);
}

/*
 * COMPILATION TEST
 */

int main() {}
\end{vncode}
\endgroup


\subsection{Lịch Gregory}

Cấu trúc ngày tháng.

\codepath{Miscelaneous/date.cpp}

\begingroup\scriptsize
\begin{vncode}
/*
 * Date data structure
 */

int mnt[13] = {0, 31, 59, 90, 120, 151, 181, 212, 243, 273, 304, 334, 365};

struct Date {
	int d, m, y;
	Date() : d(1), m(1), y(1) {}
	Date(int d, int m, int y) : d(d), m(m), y(y) {}
	Date(int days) : d(1), m(1), y(1) { advance(days); }
	bool leap() { return (y%4 == 0 && y%100) || (y%400 == 0); }
	int mdays() { return mnt[m] - mnt[m-1] + (m == 2)*leap(); }
	int ydays() { return 365 + leap(); }
	int count() { // dist to 01/01/01
		return (d-1) + mnt[m-1] + (m > 2)*leap() +
			365*(y-1) + (y-1)/4 - (y-1)/100 + (y-1)/400;
	}
	int weekday() { return (count()+1) % 7; }
	void advance(int days) {
		days += count();
		d = m = 1, y = 1 + days/366;
		days -= count();
		while(days >= ydays()) days -= ydays(), y++;
		while(days >= mdays()) days -= mdays(), m++;
		d += days;
	}
};

/*
 * TEST MATRIX
 */

#include <cstdio>
#include <cassert>
#include <algorithm>
using namespace std;

Date advance1(Date date) {
	int d = date.d;
	int m = date.m;
	int y = date.y;
	int w = date.weekday();
	d++;
	int mcnt = mnt[m]- mnt[m-1];
	if (m == 2 && date.leap()) mcnt++;
	if (d > mcnt) {
		m++; d = 1;
	}
	if (m == 13) {
		y++; m = 1;
	}
	date.advance(1);
	assert(d == date.d);
	assert(m == date.m);
	assert(y == date.y);
	w = (w+1)%7;
	assert(w == date.weekday());
	return date;
}

void test() {
	Date date;
	for(int i = 0; i < 1000000; i++) {
		Date aux(i);
		assert(aux.d == date.d);
		assert(aux.m == date.m);
		assert(aux.y == date.y);
		date = advance1(date);
	}
	for(int t = 1; t <= 1000; t++) {
		int d1 = rand()%100000;
		int d2 = rand()%100000;
		if (d1 > d2) swap(d1, d2);
		Date date1(d1), date2(d2);
		Date aux = date1;
		for(int i = 0; i < d2-d1; i++) aux = advance1(aux);
		date1.advance(d2 - d1);
		assert(date1.d == date2.d);
		assert(date1.m == date2.m);
		assert(date1.y == date2.y);
		assert(aux.d == date2.d);
		assert(aux.m == date2.m);
		assert(aux.y == date2.y);
	}
}

int main() {
	test();
	int d, m, y;
	while (scanf("%d %d %d", &d, &m, &y) != EOF) {
		Date date(d, m, y);
		int lastw;
		for(int i = 0; i < 10; i++) {
			int w = date.weekday();
			printf("%02d/%02d/%04d %s\n", date.d, date.m, date.y, 
				w == 0 ? "Domingo" :
				w == 1 ? "Segunda-feira" :
				w == 2 ? "Terca-feira" :
				w == 3 ? "Quarta-feira" :
				w == 4 ? "Quinta-feira" :
				w == 5 ? "Sexta-feira" :
				w == 6 ? "Sabado" : "undefined"
			);
			if (i > 0) assert(w == (lastw+1)%7);
			lastw = w;
			date.advance(1);
		}
	}
	return 0;
}
\end{vncode}
\endgroup


\subsection{Simplex}

Quy hoạch tuyến tính bằng simplex.

\codepath{Miscelaneous/simplex.cpp}

\begingroup\scriptsize
\begin{vncode}
#include <bits/stdc++.h>

//#include <ext/pb_ds/assoc_container.hpp> // Common file
//#include <ext/pb_ds/tree_policy.hpp> // Including tree_order_statistics_node_update

#define pb(x) push_back(x)
#define sz(x) x.size()

#define eps 1e-9
#define Max 300005
#define INF 2000000000
#define mod 1000000007
#define pii pair<int,int>
#define _mp make_pair

typedef long long LL;
typedef unsigned long long ULL;

using namespace std;

#define MAXN 109
#define MAXM 10009
#define EPS 1e-9

#define MINIMIZE -1
#define MAXIMIZE +1
#define LESSEQ -1
#define EQUAL 0
#define GREATEQ 1
#define INFEASIBLE -1
#define UNBOUNDED 999

/***
Problem example: https://www.hackerrank.com/contests/zenhacks/challenges/circles-1

Expected complexity O(N*N*M)
Worst case exponential complexity
1.  Simplex Algorithm for Linear Programming
2.  Maximize or minimize f0*x0 + f1*x1 + f2*x2 + ... + fn-1*xn-1 subject to some constraints
3.  Constraints are of the form, c0x0 + c1x1 + c2x2 + ... + cn-1xn-1 (<= or >= or =) lim
4.  m is the number of constraints indexed from 1 to m, and n is the number of variables indexed from 0 to n-1
5.  ar[0] contains the objective function f, and ar[1] to ar[m] contains the constraints, ar[i][n] = lim_i
6.  It is assumed that all variables satisfies non-negativity constraint, i.e, xi >= 0
7.  If non-negativity constraint is not desired for a variable x, replace each occurrence
	by difference of two new variables r1 and r2 (where r1 >= 0 and r2 >= 0, handled automatically by simplex).
	That is, replace every x by r1 - r2 (Number of variables increases by one, -x, +r1, +r2)
8.  solution_flag = INFEASIBLE if no solution is possible and UNBOUNDED if no finite solution is possible
9.  Returns the maximum/minimum value of the linear equation satisfying all constraints otherwise
10. After successful completion, val[] contains the values of x0, x1 .... xn for the optimal value returned

*** If ABS(X) <= M in constraints, Replace with X <= M and -X <= M

*** Fractional LP:

	max/min
		3x1 + 2x2 + 4x3 + 6
		-------------------
		3x1 + 3x2 + 2x3 + 5

		s.t. 2x1 + 3x2 + 5x3 >= 23
		3x2 + 5x2 + 4x3 <= 30
		x1, x2, x3 ≥ 0

	Replace with:

	max/min
		3y1 + 2y2 + 4y3 + 6t

		s.t. 3y1 + 3y2 + 2y3 + 5t = 1
		2y1 + 3y2 + 53 - 23t >= 0
		3y1 + 5y2 + 4y3- 30t <= 0
		y1, y2, y3, t >= 0

***/

namespace lp {
double val[MAXN], ar[MAXM][MAXN];
int m, n, solution_flag, minmax_flag, basis[MAXM], index[MAXN];

void init(int nvars, double* f, int flag) {
	solution_flag = 0;
	ar[0][nvars] = 0.0;
	m = 0, n = nvars, minmax_flag = flag;
	for (int i = 0; i < n; i++)
		ar[0][i] = f[i] * minmax_flag;
}

void add_constraint(double* C, double lim, int cmp) {
	m++, cmp *= -1;
	if (cmp == EQUAL) {
		for (int i = 0; i < n; i++) ar[m][i] = C[i];
		ar[m++][n] = lim;
		for (int i = 0; i < n; i++) ar[m][i] = -C[i];
		ar[m][n] = -lim;
	}
	else {
		for (int i = 0; i < n; i++) ar[m][i] = C[i] * cmp;
		ar[m][n] = lim * cmp;
	}
}

void init() {
	for (int i = 0; i <= m; i++) basis[i] = -i;
	for (int j = 0; j <= n; j++)
		ar[0][j] = -ar[0][j], index[j] = j, val[j] = 0;
}

inline void pivot(int m, int n, int a, int b) {
	for (int i = 0; i <= m; i++)
		for (int j = 0; i != a && j <= n; j++)
			if (j != b) ar[i][j] -= (ar[i][b] * ar[a][j]) / ar[a][b];
	for (int j = 0; j <= n; j++)
		if (j != b) ar[a][j] /= ar[a][b];
	for (int i = 0; i <= m; i++)
		if (i != a) ar[i][b] = -ar[i][b] / ar[a][b];
	ar[a][b] = 1.0 / ar[a][b];
	swap(basis[a], index[b]);
}

inline double solve() {
	init();
	int i, j, k, l;
	while (true) {
		for (i = 1, k = 1; i <= m; i++)
			if ((ar[i][n] < ar[k][n]) || (ar[i][n] == ar[k][n] && basis[i] < basis[k] && (rand() & 1))) k = i;
		if (ar[k][n] >= -EPS) break;
		for (j = 0, l = 0; j < n; j++)
			if ((ar[k][j] < (ar[k][l] - EPS)) || (ar[k][j] < (ar[k][l] - EPS) && index[i] < index[j] && (rand() & 1))) l = j;
		if (ar[k][l] >= -EPS) {
			solution_flag = INFEASIBLE;
			return -1.0;
		}
		pivot(m, n, k, l);
	}
	while (true) {
		for (j = 0, l = 0; j < n; j++)
			if ((ar[0][j] < ar[0][l]) || (ar[0][j] == ar[0][l] && index[j] < index[l] && (rand() & 1))) l = j;
		if (ar[0][l] > -EPS) break;
		for (i = 1, k = 0; i <= m; i++)
			if (ar[i][l] > EPS && (!k || ar[i][n] / ar[i][l] < ar[k][n] / ar[k][l] - EPS || (ar[i][n] / ar[i][l] < ar[k][n] / ar[k][l] + EPS && basis[i] < basis[k]))) k = i;
		if (ar[k][l] <= EPS) {
			solution_flag = UNBOUNDED;
			return -999.0;
		}
		pivot(m, n, k, l);
	}
	for (i = 1; i <= m; i++)
		if (basis[i] >= 0) val[basis[i]] = ar[i][n];
	solution_flag = 1;
	return (ar[0][n] * minmax_flag);
}
}

double obj[109],cons[109];
int x[109], y[109];

int main()
{
	srand(time(NULL));
	int n;
	cin >> n;
	for (int i = 0; i < n; i++) {
		obj[i] = 1;
	}
	obj[n] = 0;
	lp::init(n,obj,MAXIMIZE);

	for (int i = 0; i < n; i++) {
		scanf("%d %d", &x[i], &y[i]);
		for (int j = 0; j < i; j++) {
			for (int k = 0; k < n; k++) {
				cons[k] = 0;
			}
			cons[i] = 1;
			cons[j] = 1;
			lp::add_constraint(cons, hypot(x[i]-x[j], y[i]-y[j]),LESSEQ);
		}
	}

	double ret = lp::solve();
   
	printf("%.8lf\n", ret*2*acos(-1.0));
	return 0;
}
\end{vncode}
\endgroup


\subsection{Duyệt polyomino}

Sinh/duyệt polyomino.

\codepath{Miscelaneous/polyomino\_it.cpp}

\begingroup\scriptsize
\begin{vncode}
#define MAXN 59
#include <cstdio>

/*
 * Polyomino iteration
 */

int N, M;
int num[MAXN][MAXN], qi[MAXN*MAXN], qj[MAXN*MAXN];
int field[MAXN][MAXN], par[MAXN][MAXN];
int di[4] = {0, 1, 0, -1};
int dj[4] = {-1, 0, 1, 0};
int cnt;

void assign(int i, int j, int k) {
	qi[k] = i; qj[k] = j; num[i][j] = k;
}

#define valid(i, j) (i >= 0 && i < N && j >= 0 && j < M)

int iterate(int k, int h) {
	if (h == 0) return 0;
	int i = qi[k], j = qj[k], ni, nj;
	for(int d = 0; d < 4; d++) {
		ni = i + di[d]; nj = j + dj[d];
		if (!valid(ni, nj)) continue;
		if (num[ni][nj] == 0) {
			par[ni][nj] = k;
			assign(ni, nj, ++cnt);
		}
	}
	int ans = 0, cur;
	for(int t = k+1; t <= cnt; t++) {
		cur = iterate(t, h-1);
		if (cur > ans) ans = cur;
	}
	for(int d = 0; d < 4; d++) {
		ni = i + di[d]; nj = j + dj[d];
		if (!valid(ni, nj)) continue;
		if (par[ni][nj] == k) {
			par[ni][nj] = 0; cnt--;
			assign(ni, nj, 0);
		}
	}
	return ans + field[i][j];
}

/*
 * URI 1712
 */

#include <cstdio>
#include <cstring>

int main() {
	int K;
	while(scanf("%d %d", &N, &K)!=EOF) {
		M = N;
		for (int i = 0; i < N; i++) {
			for (int j = 0; j < M; j++) {
				scanf("%d", &field[i][j]);
			}
		}
		int ans = 0;
		memset(&num, 0, sizeof num);
		memset(&par, 0, sizeof par);
		for (int i = 0; i < N; i++) {
			for (int j = 0; j < M; j++) {
				cnt = 0;
				assign(i, j, ++cnt);
				int cur = iterate(1, K);
				if (cur > ans) ans = cur;
				num[i][j] = -1;
			}
		}
		printf("%d\n", ans);
	}
	return 0;
}
\end{vncode}
\endgroup


\subsection{Ma phương bậc lẻ}

Sinh ma trận n x n, n lẻ, có tổng hàng/cột/chéo bằng nhau.

\codepath{Miscelaneous/magicsquare.cpp}

\begingroup\scriptsize
\begin{vncode}
#include <cstdio>
#include <cstring>
#define MAXN 1009

/*
 * Magic Square
 */

int mat[MAXN][MAXN], n;	//0-indexed

void magicsquare() {
	int i=n-1, j=n/2;
	memset(&mat, 0, sizeof mat);
	for(int k = 1; k <= n*n; k++) {
		mat[i][j] = k;
		if (mat[(i+1)%n][(j+1)%n] > 0) {
			i = (i-1+n)%n;
		}
		else {
			i = (i+1)%n;
			j = (j+1)%n;
		}
	}
}

/*
 * TEST MATRIX
 */

int main() {
	scanf("%d", &n);
	magicsquare();
	for(int i=0; i<n; i++) {
		for(int j=0; j<n; j++) {
			if (j > 0) printf(" ");
			printf("%d", mat[i][j]);
		}
		printf("\n");
	}
	return 0;
}
\end{vncode}
\endgroup


\subsection{Chu trình trong dãy: thuật toán Floyd}

Cycle finding.

\codepath{Math/floyd.cpp}

\begingroup\scriptsize
\begin{vncode}
#include <iostream>
using namespace std;

typedef pair<int, int> ii;

int f(int x) {
	return (x+1)%5;
}

ii floydCycleFinding(int x0) {
	// 1st part: finding k*start, hare’s speed is 2x tortoise’s
	int tortoise = f(x0), hare = f(f(x0)); // f(x0) is the node next to x0
	while (tortoise != hare) { tortoise = f(tortoise); hare = f(f(hare)); }
	// 2nd part: finding start, hare and tortoise move at the same speed
	int start = 0; hare = x0;
	while (tortoise != hare) { tortoise = f(tortoise); hare = f(hare); start++;}
	// 3rd part: finding period, hare moves, tortoise stays
	int period = 1; hare = f(tortoise);
	while (tortoise != hare) { hare = f(hare); period++; }
	return ii(start, period);
}
\end{vncode}
\endgroup


\subsection{Biểu thức ngoặc sang Ba Lan}

Chuyển biểu thức dạng ngoặc sang dạng Ba Lan.

\codepath{Miscelaneous/paren2polish.cpp}

\begingroup\scriptsize
\begin{vncode}
#include <cstdio>
#include <map>
#include <stack>
using namespace std;

/*
 * Parenthetic to polish expression conversion
 */

inline bool isOp(char c) {
	return c=='+' || c=='-' || c=='*' || c=='/' || c=='^';
}

inline bool isCarac(char c) {
	return (c>='a' && c<='z') || (c>='A' && c<='Z') || (c>='0' && c<='9');
}

int paren2polish(char* paren, char* polish) {
	map<char, int> prec;
	prec['('] = 0;
	prec['+'] = prec['-'] = 1;
	prec['*'] = prec['/'] = 2;
	prec['^'] = 3;
	int len = 0;
	stack<char> op;
	for (int i = 0; paren[i]; i++) {
		if (isOp(paren[i])) {
			while (!op.empty() && prec[op.top()] >= prec[paren[i]]) {
				polish[len++] = op.top(); op.pop();
			}
			op.push(paren[i]);
		}
		else if (paren[i]=='(') op.push('(');
		else if (paren[i]==')') {
			for (; op.top()!='('; op.pop())
				polish[len++] = op.top();
			op.pop();
		}
		else if (isCarac(paren[i]))
			polish[len++] = paren[i];
	}
	for(; !op.empty(); op.pop())
		polish[len++] = op.top();
	polish[len] = 0;
	return len;
}

/*
 * TEST MATRIX
 */

int main() {
	int N, len;
	char polish[400], paren[400];
	scanf("%d", &N);
	for (int j=0; j<N; j++) {
		scanf(" %s", paren);
		paren2polish(paren, polish);
		printf("%s\n", polish);
	}
	return 0;
}
\end{vncode}
\endgroup


\subsection{Bài toán Josephus}

Vòng n người, bỏ mỗi k người: f(n,k)=(f(n-1,k)+k)\%n, f(1,k)=0.

\subsection{Bài toán histogram}

Hình chữ nhật lớn nhất trong histogram trong O(n).

\codepath{Miscelaneous/histogram.cpp}

\begingroup\scriptsize
\begin{vncode}
#include <cstdio>
#include <stack>
using namespace std;
typedef long long ll;

#define MAXN 100009

/*
 * Solution to the histogram problem in O(n)
 */

ll histogram(ll vet[], int n) {
	stack<ll> s;
	ll ans = 0, tp, cur;
	int i = 0;
	while(i < n || !s.empty()) {
		if (i < n && (s.empty() || vet[s.top()] <= vet[i])) s.push(i++);
		else {
			tp = s.top();
			s.pop();
			cur = vet[tp] * (s.empty() ? i : i - s.top() - 1);
			if (ans < cur) ans = cur;
		}
	}
	return ans;
}

/*
 * URI 1683
 */

ll vet[MAXN];

int main() {
	int n;
	while (scanf("%d", &n), n) {
		for (int i = 0; i<n; i++) {
			scanf("%lld", vet + i);
		}
		printf("%lld\n", histogram(vet, n));
	}
	return 0;
}
\end{vncode}
\endgroup


\subsection{Bài toán hôn nhân ổn định}

Gale-Shapley O(n\textasciicircum{}2).

\codepath{Miscelaneous/stablemarriage.cpp}

\begingroup\scriptsize
\begin{vncode}
#include <cstdio>
#include <cstring>
#define MAXN 1009

int m, n;     //número de homens e de mulheres
int p[MAXN];
int L[MAXN][MAXN];    //L[i]: preferências do homem i (melhor primeiro)
int R[MAXN][MAXN];    //R[j][i]: nota dada pela mulher j ao homem i
int R2L[MAXN];  //R2L[i]==j: mulher i casada com homem j 
int L2R[MAXN];  //L2R[i]==j: homem i casado com mulher j 

void stableMarriage() {
    memset(R2L, -1, sizeof(R2L));
    memset(p, 0, sizeof(p));
    for(int i = 0, wom, hubby; i < m; i++) {
        for (int man = i; man >= 0;) {
            while(true) {
                wom = L[man][p[man]++];
                if( R2L[wom] < 0 || R[wom][man] > R[wom][R2L[wom]] ) break;
            }
            hubby = R2L[wom];
            R2L[L2R[man] = wom] = man;
            man = hubby;
        }
    }
}

int main() {
    int t, ma, wo, rank;
    scanf("%d", &t);
    while(t--) {
        scanf("%d", &n);
        m = n;
        for(int i=0; i<n; i++) {
            scanf("%d", &wo);
            rank = n;
            for(int j=0; j<n; j++, rank--) {
                scanf("%d", &ma);
                R[wo-1][ma-1] = rank;
            }
        }
        for(int i=0; i<n; i++) {
            scanf("%d", &ma);
            for(int j=0; j<n; j++) {
                scanf("%d", &L[ma-1][j]);
                L[ma-1][j]--;
            }
        }
        stableMarriage();
        for(int i=0; i<n; i++) {
            printf("%d %d\n", i+1, L2R[i]+1);
        }
    }
    return 0;
}
\end{vncode}
\endgroup


\subsection{Mã Huffman}

Chi phí cây Huffman.

\codepath{Trees/huffman.cpp}

\begingroup\scriptsize
\begin{vncode}
/*
 * Huffman Cost - O(nlogn)
 */

#include <queue>
using namespace std;

typedef long long ll;

ll huffman(ll* a, int n) {
	ll ans = 0, u, v;
	priority_queue<ll> pq;
	for(int i=0; i<n; i++) pq.push(-a[i]);
	while(pq.size() > 1) {
		u = -pq.top(); pq.pop();
		v = -pq.top(); pq.pop();
		pq.push(-u-v);
		ans += u + v;
	}
	return ans;
}

/*
 * Codeforces 101128C
 */
#include <cstdio>
#define MAXN 100009

int main() {
	int T, N;
	ll a[MAXN];
	scanf("%d", &T);
	while(T-->0) {
		scanf("%d", &N);
		for(int i=0; i<N; i++) scanf("%I64d", &a[i]);
		printf("%I64d\n", huffman(a, N));
	}
	return 0;
}
\end{vncode}
\endgroup


\subsection{Bài toán quân mã}

Tìm đường Hamilton của quân mã bằng backtracking và heuristic Warnsdorff.

\subsection{Giao matroid}

Tập độc lập lớn nhất trong giao của hai matroid; ví dụ cây khung lớn nhất với màu cạnh khác nhau.

\codepath{Graphs/colorfulspanningtree.cpp}

\begingroup\scriptsize
\begin{vncode}
#include <vector>
#include <set>
#include <queue>
using namespace std;
#define MAXN 109
#define MAXM 10009

typedef pair<int,int> ii;
vector<ii> adjList[MAXN];
vector<int> Q, A[MAXM];
ii edges[MAXM];
vector<bool> T, Fcolors;
set<int> F;
int c[MAXM], N, M, C;

class UnionFind {
private:
	vector<int> parent, rank;
public:
	UnionFind(int N) {
    	rank.assign(N+1, 0);
		parent.assign(N+1, 0);
		for (int i = 0; i <= N; i++) parent[i] = i;
	}
	int find(int i) {
		while(i != parent[i]) i = parent[i];
		return i;
	}
	bool isSameSet(int i, int j) {
		return find(i) == find(j);
	}
	void unionSet (int i, int j) {
		if (isSameSet(i, j)) return;
		int x = find(i), y = find(j);
		if (rank[x] > rank[y]) parent[y] = x;
		else {
			parent[x] = y;
			if (rank[x] == rank[y]) rank[y]++;
		}
	}
};

vector<int> findCycle(int s, int t) {
	vector<int> ans; vector<ii> prev;
	queue<int> q;
	prev.assign(N, ii(-1, -1));
	q.push(s); prev[s] = ii(s, -1);
	while (!q.empty()) {
		int u = q.front(); q.pop();
		if (u == t) break;
		for (int i = 0; i < (int)adjList[u].size(); i++) {
			int v = adjList[u][i].first;
			if (prev[v].first == -1) {
				prev[v] = ii(u, adjList[u][i].second);
				q.push(v);
			}
		}
	}
	while (t != s) {
		ans.push_back(prev[t].second);
		t = prev[t].first;
	}
	return ans;
}

void start() {
	T.assign(M, false); Q.clear();
	for (int i=0; i<N; i++) adjList[i].clear();
	for (int i=0; i<M; i++) A[i].clear();
}

void greedyAugmentation(UnionFind &UF) {
	int e, u, v;
	Fcolors.assign(C, false);
	for (set<int>::iterator it = F.begin(); it != F.end(); it++) {
		e = *it;
		u = edges[e].first, v = edges[e].second;
		UF.unionSet(u, v);
		adjList[u].push_back(ii(v, e));
		adjList[v].push_back(ii(u, e));
		Fcolors[c[e]] = true;
	}
	for (int i = 0; i < M; i++) {
		u = edges[i].first, v = edges[i].second;
		if (!UF.isSameSet(u, v) && !Fcolors[c[i]]) {
			UF.unionSet(u, v);
			adjList[u].push_back(ii(v, i));
			adjList[v].push_back(ii(u, i));
			Fcolors[c[i]] = true;
			F.insert(i);
		}
	}
}

void digraphConstruction(UnionFind &UF) {
	vector<int> cycle;
	for (int i = 0; i < M; i++) {
		if (F.count(i)) continue;
		if (!Fcolors[c[i]]) Q.push_back(i);
		else for (set<int>::iterator it = F.begin(); it != F.end(); it++) {
			if (c[*it] == c[i]) A[*it].push_back(i);
		}
		if (!UF.isSameSet(edges[i].first, edges[i].second)) T[i] = true;
		else {
			cycle = findCycle(edges[i].first, edges[i].second);
			for (int j = 0; j < (int)cycle.size(); j++)
            	A[i].push_back(cycle[j]);
		}
	}
}

bool augmentation() {
	queue<int> q;
	vector<int> prev;
	prev.assign(M, -1);
	for (int i = 0; i < (int)Q.size(); i++) {
		q.push(Q[i]);
		prev[Q[i]] = Q[i];
	}
	int u, v, found = -1;
	while (!q.empty()) {
		u = q.front(); q.pop();
		if (T[u]) {
			found = u; break;
		}
		for (int i = 0; i < (int)A[u].size(); i++) {
			v = A[u][i];
			if (prev[v] == -1) {
				prev[v] = u; q.push(v);
			}
		}
	}
	if (found == -1) return false;
	u = found;
	while (true) {
		if (F.count(u)) F.erase(u);
		else F.insert(u);
		if (u == prev[u]) break;
		u = prev[u];
	}
	return true;
}

void buildSpanningTree() {
	F.clear();
	while (true) {
		UnionFind UF(M);
		start();
		greedyAugmentation(UF);
		digraphConstruction(UF);
		if (!augmentation()) break;
	}
}

/*
 * URI 2128
 */

#include <cstdio>

int main() {
    int u,v,c1,t=1;
    while (scanf("%d%d%d",&N,&M,&C)!=EOF) {
        for (int i=0;i<M;i++) {
            scanf("%d%d%d",&u,&v,&c1);
            u--, v--, c1--;
            edges[i] = ii(u, v);
            c[i] = c1;
        }
        buildSpanningTree();
        if (F.size()!=N-1) {
            printf("Instancia %d\nnao\n\n",t++);
        }else{
            printf("Instancia %d\nsim\n\n",t++);
        }
    }
}
\end{vncode}
\endgroup


\end{document}
